\documentclass{article}
\usepackage[utf8]{inputenc}
\usepackage{amsmath,amssymb,amsthm}
\usepackage[margin=1in]{geometry}
\usepackage{longtable}
\usepackage{listings}
\lstset{basicstyle=\ttfamily\scriptsize,breaklines=true,columns=fullflexible,keepspaces=true,literate={§}{{\textrm{\S}}}1}

\newtheorem{proposition}{Proposition}
\newtheorem{theorem}[proposition]{Theorem}
\newtheorem{lemma}[proposition]{Lemma}
\newtheorem{corollary}[proposition]{Corollary}
\newtheorem{definition}[proposition]{Definition}
\newtheorem{inputstatement}[proposition]{Input}
\newtheorem{gate}[proposition]{Gate}
\theoremstyle{remark}
\newtheorem{remark}[proposition]{Remark}

\newcommand{\scopebox}[2]{%
  \par\medskip\noindent
  \fbox{\begin{minipage}{0.92\textwidth}
  \textbf{#1}\par\smallskip #2
  \end{minipage}}
  \par\medskip
}
\newcommand{\RL}{\mathfrak R_{\mathrm L}}
\newcommand{\RK}{\mathfrak R_{\mathrm K}}
\newcommand{\RM}{\mathfrak R_{\mathrm M}}
\newcommand{\Pc}{\mathcal P}
\newcommand{\JD}{\mathcal J_{\mathrm D}}
\newcommand{\sigeps}{\sigma_\varepsilon}
\newcommand{\chiI}{\chi_{\mathrm I}}
\newcommand{\Ccoh}{\mathcal C_{\mathrm{coh}}}
\newcommand{\Creg}{\mathcal C_{\mathrm{reg}}}
\newcommand{\Tcol}{\mathcal T_{\mathrm{collision}}}
\newcommand{\Coll}{\mathrm{Collapse}}

\title{Landmark XVI-M\\
\large From coherent graph renormalization to local counterterm functionals and symmetry-compatible effective quantum gravity\\
\large Draft: integration of work packages A--C (final verification); Gate~S-bv checked; M-C$_{\mathrm{curv}}$ proved obstructed; broad M-C not achieved ($\varpi$ unspecified); M-C$_{\mathrm{data}}$ satisfied}
\author{}
\date{}

\begin{document}
\setlength{\emergencystretch}{3em}
\maketitle

\begin{abstract}
XVI-M begins the passage from the coherent graph renormalization of XVI-L (a frozen dependency) toward local counterterm functionals.
(M.1, unconditional) The complete-source two-edge coefficients are assembled with independent residual-metric jets.
The general table specializes entry by entry to the XVI-L background table (all $1065$ grouped entries), and the formal label support is all $255$ admissible labels ($235$ on the XVI-L background).
(M.2) Gate~S-bv is certified (Theorem~\ref{thm:wpA-Sbv}): on $\mathbb R^4\setminus\{0\}$ the predecessor edge kernel is $S=h_0+iE_0^{\mathrm A}=-W_0+iE_0^{\mathrm A}=\frac1{4\pi^2}(Q-i0)^{-1}$, including across the nonzero null cone, so $c=1/(4\pi^2)$ with boundary side $Q-i0$. This formula is already printed in XVI-J (eq:rational, with $S=-d$ by eq:DF); XVI-M supplies a written real-time check of it, including a local boundary-value argument at nonzero null points (Lemma~\ref{lem:wpA-varreg}). The predecessor $S$ is globally defined (XVI-J); what is not established here is its position-space identification at the origin.
With a boundary-value calculus lemma this gives a two-direction equivalence between numerator identities and distributional relations, and hence, unconditionally on the punctured space, the graded presentation $d_N=(1,4,10,20,35)$, $\dim\mathcal V_{\mathrm{adm}}=70$, with $185$ certified relations (only $c\neq0$ and the single boundary side are used).
(M.3) The coherent contact space has dimension $3321$, against $14321$ raw slots.
The graph-contact map is injective in the formal universal-jet reading, while $1750$ coherent directions are invisible on the XVI-L background.
(M.4) Under a stated variational contract (Gate~V), the contact-collapse map is computed exactly.
It is injective in the formal reading, and its formally self-adjoint (variational) subspace is $26$-dimensional there and consists only of ultralocal potentials $\chiI\sigeps K(h,h)$; the formal universal-jet collapse image of this selected graph contains no nonzero variational tensor of total derivative degree $D\ge1$.
The same holds in the universal time-only reading $\sigeps=\sigma(t)$; if in addition $\chiI\equiv1$, there are $546$ further spatial-only, non-Lorentz variational directions.
In the double-plateau reading ($\sigeps\equiv1$, $\chiI\equiv1$) the variational image has dimension $1325$; it contains the linearized Euler tensors of $\int R^2$ and $\int R_{\mu\nu}R^{\mu\nu}$, with certified preimages, while the specified quadratic Einstein--Hilbert Euler tensor is not in that selected image.
These curvature-square matchings do not survive the time-only reading once $\sigeps'$, $\sigeps''$ are retained.
Stage~2: fixed-data criteria, universal region comparisons and the time-only classification are constructed; the transfer to the actual XII-C ramp on $(0,\varepsilon)\times\mathbb R^3$ is a proposition with explicit hypotheses; regions where the actual $\chiI$ is non-constant, and transitions, remain Unresolved.
M.5 records the named matching theorem \texttt{PlateauCurvatureSquareMatching} with explicit domain, order, pairing, and equivalence; ledgers~(3)--(4) remain Unresolved (EOM/redefs undefined).
(M.6) The $S^2$ scaling residue is $r=-2i\pi^2c^2=-i/(8\pi^2)$.
For the XVI-L prescription the scale derivative equals the dilation defect at every grade $N\le4$ ($A_\mu=\mathcal A_\mu$); top-order residues are $\varpi$-independent, while lower orders depend on the cutoff $\varpi$, which XVI-L does not fix.
In the double plateau, the cutoff-independent $D=4$ component is variational and Lorentz covariant, but not linearized-gauge invariant. The full scale tensor has cutoff-dependent lower-degree components and need not be variational. The $D=4$ component lies outside $\mathrm{span}(E_{R^2},E_{\mathrm{Ric}^2})$, certified by an exact separating functional ($\ell(v_{\mathrm{scale}})=-\tfrac1{16}$, $\ell(v_{R^2})=\ell(v_{\mathrm{Ric}^2})=0$). Hence the curvature-ledger scaling criterion M-C$_{\mathrm{curv}}$ is \emph{Proved obstructed} for this graph, while broad M-C (Definition~\ref{def:MCcurv}: matching plus an explicit cutoff and a functional representation of the entire collapsed scale variation) is \emph{not achieved} and remains Unresolved, because XVI-L leaves $\varpi$ unspecified (thm:L6); M-C$_{\mathrm{data}}$ is satisfied (not a claim about gravity generally).
M.7 (audited): with the contact distributions $\partial_\gamma\delta$ transported by the metric intertwiner, the Lorentz-equivariant contact assignments have dimension $14=5+9$; their exact intersection $V$ with the variational subspace has $\dim V=1$ (formal, $\chi$-plateau; $V\subset\Creg$) and $\dim V=12$ in the double plateau ($\dim V\cap\mathcal C_{\mathrm K}=3$, $\dim V\cap\Creg=7$); the curvature-square Euler tensors satisfy the linear divergence identity $\partial_\mu E^{\mu\nu}=0$ (this is not a Ward identity).
No Ward or BV statement is claimed; scale-variation coefficients are not physical beta functions.
\end{abstract}

\scopebox{Status (XVI-M, this draft).}{%
M.0--M.3: Constructed as stated (M.1 unconditional; M.2/M.3 use Input~(S-bv), now certified; Lemma~\ref{lem:bv}(c) by same-sign BV product).
Gate~S-bv: \textbf{certified} (Theorem~\ref{thm:wpA-Sbv}: $c=1/(4\pi^2)$, boundary side $Q-i0$, including across the nonzero null cone); a written check of XVI-J, eq:rational; the Euclidean motivation (Remark~\ref{rem:Sbv-motivation}) is not the certification; the predecessor $S$ is globally defined (XVI-J), and its position-space identification at the origin is not established here; not discharged by matching.
M.4: Constructed (Gate~V) in the formal, $\chi$-plateau, preparation-plateau, double-plateau and time-only readings (the formal, preparation-plateau and time-only readings are universal formal-jet readings). Stage~2: fixed-data criteria, universal region comparisons and the time-only classification constructed; actual-data transfer on $(0,\varepsilon)\times\mathbb R^3$ only as a proposition with explicit hypotheses; non-constant-$\chiI$ regions and transitions Unresolved.
M.5: Functional ledger with named theorem \texttt{PlateauCurvatureSquareMatching}; ledgers (1)--(2) filled; (3)--(4) Unresolved (EOM/redefs undefined).
M.6: dilation identity and $A_\mu=\mathcal A_\mu$ for all $N\le4$ Constructed; $S^2$ residue $r=-i/(8\pi^2)$ Constructed; top-order residues exact and $\varpi$-independent; lower-order values Unresolved ($\varpi$ not fixed by XVI-L). In the double plateau, the cutoff-independent $D=4$ component is variational and Lorentz covariant, but not linearized-gauge invariant. The full scale tensor has cutoff-dependent lower-degree components and need not be variational (Section~\ref{sec:wpB}).
M.7: \textbf{Constructed after audit} (Gate~V for the SA items; (S-bv) certified) for the exactly verified items: Lorentz action on $B_N$ and on $\partial_\gamma\delta$ with explicit $\eta$ intertwiner and boost check; equivariant dims $1,1,3,3,6$; exact $\mathrm{Equiv}\cap\mathrm{SA}$ ($1/1/12$) with projections and intersections; linear divergence identities. Obs$_L(R_L)$ Unresolved. See the audit log, Remark~\ref{rem:M7-audit}.
M.8--M.9: light stubs (Target II / BV Unresolved; M.8 should wait for XVI-N).
Release ladder: \textbf{M-A} Constructed (presentation unconditional on the punctured space); \textbf{M-B} Constructed in the formal reading (Gate~V convention); \textbf{M-C$_{\mathrm{curv}}$} (curvature-ledger scaling criterion, Definition~\ref{def:MCcurv}) Proved obstructed, graph-specific (Corollary~\ref{cor:wpB-MC}); broad \textbf{M-C} not achieved (Unresolved: $\varpi$ unspecified by XVI-L, thm:L6); \textbf{M-C$_{\mathrm{data}}$} Satisfied; \textbf{M-D/M-E} not claimed.
XVI-L is a frozen dependency.
No $\varepsilon\to0$, no loop-as-result, no UV claim, no interacting BRST/Ward identity as proved.
}

\scopebox{Programme arrow (domains/images/kernels/intersections/compatibility).}{%
\[
\boxed{\begin{aligned}
&\text{coherent numerical contact assignments}\\
&\quad\longrightarrow
\text{unrooted local collision tensors}\\
&\quad\longleftarrow
\text{variations of local counterterm functionals}\\
&\quad\longleftarrow
\text{symmetry-compatible EFT counterterms}.
\end{aligned}}
\]
The arrows need not be surjective or injective.
This draft computes the first arrow (M.3) and the intersection of its collapsed image with the variational tensors (M.4).
The symmetry arrow is reserved for M.7--M.9.
}

\noindent
Conventions: ambient $M=\mathbb R^4$, $\eta=\mathrm{diag}(1,-1,-1,-1)$, $\Pc r=r-\tfrac12\eta\,\mathrm{tr}_\eta r$, Gram $G=\mathrm{diag}(1,-2,-2,-2,$ $1,2,2,1,2,1)$, dual basis $B^i=(G^{-1})^{ii}B_i$; Hessian on field directions $B_i$; propagator/source contraction on dual slots $-\Pc B^i$.
The XI inverse $G_{\mathrm{gr}}^{\mathrm R}=-\Pc E_0^{\mathrm R}$ is not equated with the X-sheet $\Delta_R^{\mathrm{gr}}=-4\Pc E_0^{\mathrm R}$.
Series voice: Constructed / Proved obstructed / Unresolved / Non-claims.
Predecessor quotations are exact.

%=====================================================================
\section{M.0 Dependency, normalization, gates, and target ledger}
\label{sec:m0}

\subsection{Frozen XVI-L dependency}

XVI-L.tex and \texttt{XVI-L-artifacts/} are a frozen dependency: they are cited and read, not edited.
XVI-K-artifacts/\texttt{hessian\_full\_out.txt} is likewise frozen.
The following items are imported \emph{as stated in XVI-L}.
\begin{enumerate}
\item Printed cubic density $\ell_3$ (XVI-L, eq:L-ell3; XII-C, eq:ell3).
\item Complete weighted source $\JD=\sigeps J_3+2C_2\partial\partial\sigeps-2M_2\partial\sigeps$ (XVI-L, eq:L-JD).
\item Full cubic Hessian $L^{pq}_{ij}$ (XVI-K, prop:hess, eq:Lpq).
\item Metric-dual contraction on $-\Pc B^i$.
\item Ambient numerical span $\mathcal V_{\mathrm{adm}}$, $|\Lambda_{\mathrm{adm}}|=255$ (XVI-L, def:VL).
\item Selected extension $\RL$ with $\RL|_{\mathcal V_{\mathrm K}}=\RK$ (XVI-L, def:RL, prop:RL).
\item Coherent contact freedom (XVI-L, thm:L5); raw counts $14321/321/14000$ ``are not dimensions''.
\item Subtraction-scale family $\mathfrak R_{\mathrm L,\mu}$ and generator $A_\mu$ (XVI-L, thm:L6); explicit nonzero scaling residue Unresolved there.
\item Conditional root attachment under (R1)--(R3) (XVI-L, thm:L7).
\end{enumerate}

\paragraph{Preservation versus new results.}
XVI-M does not alter any release claim of XVI-L: L-A/L-C remain as constructed there, L-B remains as stated there, and L-D remains not claimed.
The new statements below are results \emph{about the objects of XVI-L} (its admissible span, its contact freedom, its graph coefficients).
They strengthen knowledge of the span without retroactively rewriting XVI-L.
In particular, the graded presentation of M.2 is a new result resting on an explicitly checked kernel identification (Gate~S-bv, now certified by Theorem~\ref{thm:wpA-Sbv}); it is not a correction of XVI-L, Theorem thm:L5.

\subsection{Two separate targets}

\begin{definition}[Target I and Target II]
\label{def:targets}
\textbf{Target I (unrooted counterterm matching).} Determine which coherent contact assignments of $\RL$ produce collision tensors that arise as variations of local metric or insertion functionals, without using the root hypotheses (R1)--(R3).
\textbf{Target II (gravitational perturbative-sector interpretation).} Interpret those matched functionals inside a specified gravitational action, insertion, gauge choice, ghost content, and diagram inventory.
\end{definition}
The M.4 results serve Target I.
Target II begins only in skeleton Section~\ref{sec:m8}.

\begin{definition}[Scaling criteria M-C, M-C$_{\mathrm{data}}$ and M-C$_{\mathrm{curv}}$ (final verification)]
\label{def:MCcurv}
\textbf{M-C (broad)} is the release-ladder \emph{matched-scale} criterion: it requires
(i)~at least one proved functional matching (Target~I), and
(ii)~an explicitly specified cutoff $\varpi$ together with a functional representation of the \emph{entire} collapsed scale variation in a stated reading.
\textbf{M-C$_{\mathrm{data}}$} is the weaker \emph{data} criterion: at least one proved functional matching together with explicitly specified, prescription-dependent scaling formulas (residues, $\nu_k[\varpi]$, top-order coefficients) for the collapsed scale variation; it does not require a fixed numerical cutoff.
\textbf{M-C$_{\mathrm{curv}}$} is the \emph{curvature-ledger scaling criterion}: the scale variation of the graph is matched inside the span of the M.5 curvature-ledger targets (the Euler tensors $E_{R^2},E_{\mathrm{Ric}^2}$ at $D=4$, and the lower-degree curvature targets of the ledger), with the matched-scale data of (ii).
M-C$_{\mathrm{curv}}$ implies broad M-C, not conversely; M-C$_{\mathrm{data}}$ is satisfied in this draft (Theorem~\ref{thm:wpB-B0}, Theorem~\ref{thm:wpB-dp}, Section~\ref{sec:wpB}), while broad M-C remains Unresolved because XVI-L leaves $\varpi$ unspecified (XVI-L, thm:L6; Remark rem:Lressretch).
Enlarging the target ledger by non-gauge-invariant quadratic representatives (for example $I^{\mathrm{top}}_{\mathrm{scale}}$ of Theorem~\ref{thm:wpB-dp}) changes the matching domain of M-C$_{\mathrm{curv}}$ but may remain within broad Target~I; gauge-fixing/ghost or further graph contributions are a different, independent enlargement and are not the only one.
\end{definition}

\subsection{Normalization table}
\label{sec:m0norm}

XVI-L determines the combinatorial graph factor $(1/2!)\cdot2=1$ and leaves coupling and loop weights undetermined (XVI-L, prop:Lnorm).
No number is invented below.

\begin{center}
\begin{longtable}{p{0.20\textwidth}p{0.44\textwidth}p{0.26\textwidth}}
\hline
Item & Value / source & Status\\
\hline
Metric fluctuation & $q_{\mu\nu}=\eta_{\mu\nu}+\kappa\gamma_{\mu\nu}$ (XII-C) & Imported\\
Residual $h$ vs $\gamma$ & Dummy of the same species; no rescaling (XVI-I, prop:Ft) & Imported\\
$\ell_3$ vs action & $\ell_3$ is the coefficient of $\kappa$ in the unintegrated signed EH density; $S_{\mathrm{EH}}=-\frac2{\kappa^2}\int\sqrt{-q}\,R$; $S_{\mathrm{grav,D}}=S_{\mathrm{EH}}+\kappa\int(\sigma-1)\ell_3$ (XII-C) & Imported\\
Graph vertex & $V(h)=\int\chiI\ell_3(h)$, $\chiI$ independent of $\sigeps,f,g$ (XVI-I, prop:second) & Imported; relation of $\chiI$ to $\kappa\int\sigma\ell_3$ Unresolved\\
Kinetic operator & $K=-\square\Pc=A+2D^*D$, $A=2G_{\mathrm{lin}}$ (XII-C; XVI-I) & Imported\\
Propagator & $D_F=H+iG^{\mathrm A}=-\Pc(h_0+iE_0^{\mathrm A})$, $KD_F=i\,\mathrm{id}$ (XVI-J, eq:DF, quoted in XVI-K, eq:J-DF); ``with $h_0=-W_0$ (XI-E.tex, eq:h)'' (XVI-K, K.0); $S=h_0+iE_0^{\mathrm A}=-\Delta_F^{\mathrm{FR}}$ (XVI-J, lem:identify) & Imported\\
Position-space form of $S$ on $\mathbb R^4\setminus\{0\}$ & $S=c\,(Q-i0)^{-1}$ with $c=\frac1{4\pi^2}$, boundary side $Q-i0$ (Theorem~\ref{thm:wpA-Sbv}); $W_0=-\frac1{4\pi^2}(Q-i0\,\mathrm{sgn}\,z^0)^{-1}$, $E_0^{\mathrm A}=\frac1{2\pi}\theta(-z^0)\delta(Q)$ & Certified (Gate~S-bv; written check of XVI-J, eq:rational); Euclidean motivation is not the certification; $S$ is global (XVI-J), its position-space identification at $z=0$ is not established\\
$S^2$ residue and scale units & $r=-2i\pi^2c^2=-\frac{i}{8\pi^2}$ (Theorem~\ref{thm:M6-S2}); top-order scale residues in units $i\pi^2c^2=\frac{i}{16\pi^2}$ (Section~\ref{sec:wpB}) & Constructed; not a beta function\\
Powers of $\kappa$ & Action cubic carries $\kappa$; $a_\lambda$ and the graph factor carry none & Overall order Unresolved\\
Powers of $\hbar$ & ``No power of $\hbar$ is assigned'' (XVI-I, prop:sector) & Unresolved\\
Factors of $i$ & $i$ in $D_F$ fixed by $KD_F=i\,\mathrm{id}$; vertex factors not assigned & Partial; remainder Unresolved\\
Insertion & $F_t(h,\phi)=\langle t_{\mathrm D},h\otimes h\otimes\phi\rangle$; $D_h^2F_t[a,c]=2\langle t_{\mathrm D},a\otimes c\otimes\phi\rangle$ (XVI-I) & Imported\\
Loop order & ``A count of one cycle is not a loop order of an interacting observable'' (XVI-I) & Unresolved\\
Variational contract & Gate~V (Section~\ref{sec:gates}) & Chosen Phase II convention; not certified by this table\\
\hline
\end{longtable}
\end{center}

\subsection{Gates}
\label{sec:gates}

\begin{gate}[Gate S-bv]
\label{gate:Sbv}
Input~(S-bv) (Input~\ref{inp:Sbv}) was introduced as an explicit input of XVI-M.
Certification had to establish the representation for the \emph{actual predecessor kernel} $S=h_0+iE_0^{\mathrm A}$ on the punctured space $\Omega$, including across the nonzero null cone---not only for a model kernel with the same wave equation, and not only off the cone.
(The alternative route, a direct certification of the relations of M.2 and of their completeness, is no longer needed.)
Earlier drafts proceeded conditionally on (S-bv); no downstream matching result was used to remove that condition.
The unconditional general-$h$ assembly (M.1) and the wave-operator bound (Proposition~\ref{prop:wave}) are stated separately and never depended on this gate.
\textbf{Gate~S-bv certified (Theorem~\ref{thm:wpA-Sbv}):} $S=\frac1{4\pi^2}(Q-i0)^{-1}$ in $\mathcal D'(\Omega)$, i.e.\ $c=1/(4\pi^2)$ with boundary side $Q-i0$ (Section~\ref{sec:wpA-Sbv}, Remark~\ref{rem:Sbv-status}).
\end{gate}

\begin{gate}[Gate V: variational contract for Phase II]
\label{gate:V}
\emph{Chosen Phase II convention; it is not certified by the normalization ledger.}
\begin{enumerate}
\item Field variable: the symmetric lowered residual field $h_{\mu\nu}$, with no rescaling and no trace reversal.
\item Variation: $h\mapsto h+t$ with $t$ symmetric and compactly supported.
\item Pairing: $\delta I[h](t)=\int d^4y\,t_{\mu\nu}E_I^{\mu\nu}[h]$ with the full symmetric-index contraction over all $16$ pairs $(\mu,\nu)$.
In the ten-component coordinates $h_{\mu\nu}$, $\mu\le\nu$:
\[
t_{\mu\nu}E^{\mu\nu}=\sum_{\mu\le\nu}w_{\mu\nu}\,t_{\mu\nu}E^{\mu\nu},\qquad w_{\mu\mu}=1,\quad w_{\mu\nu}=2\ (\mu<\nu),
\]
so the Euler tensor of a density $\mathcal L$ in the independent components is
\[
E_I^{\mu\nu}=\frac1{w_{\mu\nu}}\sum_\beta(-D)^\beta\frac{\partial\mathcal L}{\partial(\partial^\beta h_{\mu\nu})}.
\]
\item External data: $\sigeps$ and $\chiI$ are fixed under field variation, but their derivatives are kept in every integration by parts.
\item Order: $h$-linear collapsed tensors; quadratic functional ansatz modulo field-independent terms.
\item Helmholtz test: $\int r\cdot(DE_C[h]\cdot t)=\int t\cdot(DE_C[h]\cdot r)$; for homogeneous linear $E_C$ this is $\int\langle r,E_C[t]\rangle=\int\langle t,E_C[r]\rangle$.
\end{enumerate}
\end{gate}

\subsection{Verification boundary}
\label{sec:vb}

The following are \emph{reported computational results}, with exact certificates in the archive (Appendix~\ref{app:script}):
\begin{itemize}
\item the $1065$-entry specialization of M.1 and the support counts;
\item the numerator ranks $d_N$, reconstruction certificates and relation vectors of M.2;
\item the general-$h$ ranks $\dim U_N$ and the background/plateau kernels of M.3;
\item the collapse ranks, self-adjoint subspaces, Helmholtz certificates, and matching preimages of M.4;
\item the symbolic/numerical cross-checks of work package A (Lemma~\ref{lem:wpA-comp}); the top-order residues, coefficient files and scale-tensor collapse certificates of work package B; the time-only kernels, route~1/route~2 agreement and matching probes of work package C (Lemma~\ref{lem:to-comp});
\item the audit-3 checks: the separating functional on the archived target vectors, its pure-gauge sympy verification and the recollapse of $C_{\mathrm{top}}$ (Proposition~\ref{prop:wpB-sep}); the numerical variable-regulator pairing (Lemma~\ref{lem:wpA-comp}); the Beta-integral convolution of Lemma~\ref{lem:bv}(c).
\end{itemize}
The analytic statements (Lemma~\ref{lem:bv}, Lemma~\ref{lem:wpA-varreg}, Theorem~\ref{thm:wpA-Sbv}, Proposition~\ref{prop:wpB-sep}, Propositions~\ref{prop:relations}, \ref{prop:Gsplit}, \ref{prop:collinj}, Lemma~\ref{lem:helm}, Lemma~\ref{lem:wpB-dw}, Theorem~\ref{thm:wpB-B0}, Lemmas~\ref{lem:to-intertwine} and~\ref{lem:to-indep}, Corollaries~\ref{cor:wpB-timeonly} and~\ref{cor:wpB-MC}) are written proofs.
The outputs are reproduced verbatim in Appendix~\ref{app:out}.

\subsection{Release ladder}
\label{sec:ladder}

\begin{center}
\begin{longtable}{p{0.08\textwidth}p{0.32\textwidth}p{0.50\textwidth}}
\hline
Level & Target claim & Status in this draft\\
\hline
M-A & General-$h$ coefficients and an explicit presentation of the admissible span & \textbf{Constructed} (M.1 unconditional; the presentation depended only on (S-bv) and is now unconditional on the punctured space, Theorem~\ref{thm:wpA-Sbv})\\
M-B & Contact freedom $\to$ local collision tensors, with variational-integrability conditions & \textbf{Constructed in the formal reading} (M.3, M.4) under the Gate~V convention; no longer conditional on (S-bv). Time-only classification constructed (Section~\ref{sec:to-timeonly}); actual fixed-data dimensions unresolved (Stage~2)\\
M-C & Matched-scale criterion (Definition~\ref{def:MCcurv}): (i)~proved functional matching; (ii)~explicit $\varpi$ and a functional representation of the entire collapsed scale variation & \textbf{Not achieved (Unresolved).} Criterion~(i) holds (e.g.\ $I^{\mathrm{top}}_{\mathrm{scale}}$, Theorem~\ref{thm:wpB-dp}); criterion~(ii) fails because XVI-L leaves $\varpi$ unspecified (thm:L6). Their numerical values and the variationality of the resulting full tensor are not determined for the unspecified cutoff. \textbf{M-C$_{\mathrm{data}}$} (matching + explicit prescription-dependent scaling formulas) is \textbf{Satisfied}.\newline \textbf{M-C$_{\mathrm{curv}}$} (curvature-ledger scaling criterion, Definition~\ref{def:MCcurv}): \textbf{Proved obstructed (graph-specific).} For the specified two-edge graph, the XVI-L prescription and any admissible $\varpi$, the full collapsed scale variation does not lie in the span of the M.5 curvature targets, because its fixed $D=4$ component already lies outside (Theorem~\ref{thm:wpB-dp}, Proposition~\ref{prop:wpB-sep}, Corollary~\ref{cor:wpB-MC}); not a claim about gravity generally\\
M-D & Ward/BV compatibility for a diagrammatically closed sector & \textbf{Not claimed}\\
M-E & Full renormalization for a specified EFT truncation and order & \textbf{Not claimed}\\
\hline
\end{longtable}
\end{center}

%=====================================================================
\section{M.1 General-$h$ coefficient assembly}
\label{sec:m1}

XVI-L assembles all-sector coefficients on the background $h_{\mu\nu}=(w^1)^2\eta_{\mu\nu}$ (XVI-L, lem:L2comp: ``NOT a general-$h$ symbolic assembly'').
This section retains the independent residual-metric jets of $L^{pq}_{ij}$.

\begin{definition}[General-$h$ graph coefficients]
\label{def:aGeneral}
With $T_i^r:=(-\Pc B^i)^r$,
\[
\widetilde P_{A,ij}^{ab,\mu\nu}=\sum_{r,s}P_{A,rs}^{ab,\mu\nu}T_i^rT_j^s,\qquad
\hat a^{A\,\mu\nu}_{ab,pq}=(-1)^{|p|+|q|}\sum_{i,j}L^{pq}_{ij}(w;h)\widetilde P_{A,ij}^{ab,\mu\nu}(y),
\]
and $\hat a^{A\,\mu\nu}_\lambda$ is the accumulation over $\{a+p,b+q\}=\lambda$, exactly as in XVI-L, def:La / eq:L-apq, \emph{with $\chiI$ suppressed}.
The physical (switched) coefficient---XVI-L style---is the single product
\[
a^{A\,\mu\nu}_\lambda=\chiI(w)\,\hat a^{A\,\mu\nu}_\lambda.
\]
The hat marks the switch-free factor: $a_\lambda$ and $\hat a_\lambda$ are not the same object.
Here $L^{pq}_{ij}$ is kept as the symbolic linear forms of XVI-K in independent $h$-jet symbols $h_{ab}$, $\partial_ch_{ab}$, $\partial_c\partial_dh_{ab}$.
The recorded table \texttt{graph\_general\_table.txt} stores the switch-free $\hat a$ ($\chiI$ suppressed), as in XVI-L.
\end{definition}

\begin{proposition}[General-field coefficient construction and grading]
\label{prop:M1}
For every assembled switch-free coefficient $\hat a^{A\,\mu\nu}_\lambda(y,w;h)$:
\begin{enumerate}
\item it is linear in the residual-metric jet of order $2-v$, where $v=|p|+|q|$;
\item it is linear in the preparation jet of order $2-s$, where $s=|a|+|b|$ ($s=0,1,2$ for $C_2,M_2,J_3$);
\item $N=s+v$, and every monomial is one preparation-jet symbol times one $h$-jet symbol, with a rational coefficient;
\item it is symmetric in $(\mu,\nu)$;
\item it is consistent under edge interchange, $(A;a,b;p,q)\leftrightarrow(A;b,a;q,p)$;
\item no unintended generators occur, and the arithmetic is exact rational.
\end{enumerate}
The supports satisfy $\Lambda_{\mathrm{graph}}^{\mathrm{background}}\subseteq\Lambda_{\mathrm{graph}}^{\mathrm{formal}}\subseteq\Lambda_{\mathrm{adm}}$ with cardinalities $235\le255\le255$ (sectorwise $C_2$ $25/25/25$, $M_2$ $84/84/84$, $J_3$ $202/250/250$).
\end{proposition}

\begin{proof}
(1)--(3) and (6) are inherited from XVI-K, prop:hess, and XVI-L, prop:jet, through Definition~\ref{def:aGeneral}.
(4) holds because $P_A$ is symmetric in $(\mu,\nu)$.
(5) holds because the defining sum is symmetric under simultaneous interchange of the two edges.
The supports, and a coefficientwise confirmation of (1)--(6), are the reported computation of Lemma~\ref{lem:M1comp}.
\end{proof}

\begin{theorem}[Exact specialization to the XVI-L table]
\label{thm:M1spec}
After the specialization $h_{\mu\nu}=(w^1)^2\eta_{\mu\nu}$ (with $\chiI$ suppressed as in XVI-L), the general-$h$ table equals \texttt{two\_edge\_coeff\_table.txt} entry by entry.
All $1065$ nonzero grouped coefficients $(A,\lambda,(\mu,\nu))$ agree as polynomials in the preparation jet and $w^1$, and the nonzero key sets coincide.
\end{theorem}

\begin{lemma}[Reported computation for M.1]
\label{lem:M1comp}
\texttt{graph\_general.py} loads $P_A$ from the frozen XVI-L source table and $L^{pq}_{ij}$ from the frozen XVI-K Hessian list.
It assembles the switch-free $\hat a$ of Definition~\ref{def:aGeneral} in exact $\mathbb Q$ arithmetic and exits with status $0$, recording the following.
\begin{enumerate}
\item[(G1)] Proposition~\ref{prop:M1}(1)--(6). The edge-interchange check is \emph{exhaustive}: it runs over every enumerated ordered placement whose swap is distinct, which gives $1000$ unordered pairs.
\item[(G2)] The supports of Proposition~\ref{prop:M1}. Twenty labels lie in $\Lambda^{\mathrm{formal}}_{\mathrm{graph}}\setminus\Lambda^{\mathrm{background}}_{\mathrm{graph}}$.
\item[(G3)] Theorem~\ref{thm:M1spec} ($1065/1065$ entries).
\item[(G4)] An independent $C_2$ route via \texttt{xvi\_common.polar\_s} ($250$ comparisons).
\item[(G5)] $3282$ nonzero $(A,\lambda,(\mu,\nu))$ lines. The sector-summed table has $2322$ entries with $18646$ bilinear terms.
\end{enumerate}
\end{lemma}

%=====================================================================
\section{M.2 Numerical-kernel relations and Gate S-bv}
\label{sec:m2}

XVI-L does not compute the full relation space of its $255$ admissible labels (XVI-L, thm:L5).
This section constructs the presentation $\pi:\mathbb C^{\Lambda_{\mathrm{adm}}}\to\mathcal V_{\mathrm{adm}}$, $e_\lambda\mapsto q^\circ_\lambda$, degree by degree.

\begin{inputstatement}[Input (S-bv); certified by Theorem~\ref{thm:wpA-Sbv}]
\label{inp:Sbv}
On $\Omega:=\mathbb R^4\setminus\{0\}$,
\[
S=c\,(Q-i0)^{-1},\qquad Q(z)=\eta(z,z)=z_0^2-z_1^2-z_2^2-z_3^2,\qquad c\neq0,
\]
where $S=h_0+iE_0^{\mathrm A}$ is the predecessor edge kernel (XVI-K, eq:S) and $(Q-i0)^{-1}$ is the boundary-value distribution $T_1$ of Lemma~\ref{lem:bv}.
In its original formulation the constant $c$ was left unspecified; the graded results of M.2--M.4 and M.7 use only $c\neq0$ and the side $Q-i0$.
Theorem~\ref{thm:wpA-Sbv} certifies the input with $c=1/(4\pi^2)$.
\end{inputstatement}

\begin{remark}[Motivation for Input (S-bv)]
\label{rem:Sbv-motivation}
XVI-K quotes $D_F=-\Pc(h_0+iE_0^{\mathrm A})$ and states ``with $h_0=-W_0$ (XI-E.tex, eq:h)''.
XI-E, eq:h, reads $h_\mu=-W_\mu$, $h_\mu-h_\mu^{\mathsf T}=i\delta_\mu$.
Hence $S=h_0+iE_0^{\mathrm A}=-W_0+iE_0^{\mathrm A}$, and XVI-J, lem:identify, identifies this with $-\Delta_F^{\mathrm{FR}}$.
A Euclidean calculation motivates the candidate form.
Substituting $x^0-y^0=-i\tau_{\mathrm E}$ ($\tau_{\mathrm E}>0$) into the $W_0$ integral and using $\int_0^\infty e^{-p\tau_{\mathrm E}}\sin(pr)\,dp=r/(\tau_{\mathrm E}^2+r^2)$ (checked symbolically in \texttt{kernel\_relations.py}) gives $1/(4\pi^2(\tau_{\mathrm E}^2+r^2))$, which suggests the motivational value $c=1/(4\pi^2)$.
\emph{This Euclidean calculation motivates the candidate; it does not certify the real-time identification of $S$ on $\Omega$ as a single boundary value.}
The certification is the real-time argument of Theorem~\ref{thm:wpA-Sbv}, which uses the predecessor Fourier family of $W_0$ and the step representation of $E_0^{\mathrm A}$, not this Euclidean slice.
XVI-K states: ``No explicit position-space formula for \(S\) is used.''
\textbf{Boundary value versus $\varepsilon$.} The symbol $i0$ denotes a distributional boundary value.
It is unrelated to the fixed preparation parameter $\varepsilon$.
The series rule ``no $\varepsilon\to0$'' (XVI-I, prop:signs(5): ``There is no \(i\varepsilon\) prescription and no limit \(\varepsilon\to0\)'') does not forbid the $i0$ notation, and Input~(S-bv) introduces no prescription and changes no definition.
\end{remark}

\begin{remark}[Status of Gate S-bv]
\label{rem:Sbv-status}
\emph{What predecessors supply.} XVI-K, eq:S, defines $S:=h_0+iE_0^{\mathrm A}=-W_0+iE_0^{\mathrm A}=-\Delta_F^{\mathrm{FR}}$ through two globally defined distributions, with Fourier $z^0$-family formulae (XVI-K, lem:families / eq:W-fam, eq:E-fam) quoted from XVI-J and XI-E, eq:h ($h_0=-W_0$).
Homogeneity of degree $-2$ is proved (XVI-K, lem:homog); $\square S=0$ on $\Omega$ (XVI-K, lem:grading(3)); the H\"ormander product of two copies exists on $y\neq w$ (XVI-J / XVI-K, lem:prod).
XVI-K states explicitly: ``No explicit position-space formula for \(S\) is used.''
XI-E, eq:W and eq:h, likewise give the Fourier definition of $W_\mu$ and the normalization $h_\mu=-W_\mu$, with no $(Q-i0)^{-1}$ formula.
XVI-J, however, already prints the position-space formula: eq:DF sets $S=h_0+iE_0^{\mathrm A}=-d$ with $d=W_0-iE_0^{\mathrm A}$, and eq:rational states $d(z)=-\frac1{4\pi^2}(Q_{\mathrm{geom}}(z)-i0)^{-1}$ on $z\neq0$, including nonzero null points; XVI-J, prop:q, uses it (quotations in Section~\ref{sec:wpA-identify}).
Its printed justification is ``direct Fourier evaluation''.

\emph{What was missing.} A written real-time derivation of XVI-J, eq:rational, including the boundary-value argument at nonzero null points. The Euclidean Wick calculation of Remark~\ref{rem:Sbv-motivation} motivates the candidate $c=1/(4\pi^2)$; it is \emph{not} a real-time identification of $S$ with $c\,T_1$ on $\Omega$, and it does not settle equality across the nonzero null cone.
Lemma~\ref{lem:bv} establishes the model calculus for $T_m$; it does not identify $S$ with $c\,T_1$.
No matching probe, collapse rank, or Helmholtz certificate discharges (S-bv).

\emph{Certification (this cycle): a written check of XVI-J, eq:rational, not a new extension.} Theorem~\ref{thm:wpA-Sbv} establishes, from the predecessor Fourier families, all four required items on $\Omega$: (i) the distributional representation $S=c\,(Q-i0)^{-1}$; (ii) the boundary-value side, the same as for $T_1$ (namely $Q-i0$, not $Q+i0$); (iii) the normalization $c=1/(4\pi^2)\neq0$; (iv) equality across the nonzero null cone, not merely on $\{Q\neq0\}$.
The alternative route (certifying the $185$ relations and their completeness directly) is therefore not needed.
Consequently every graded M.2--M.4 (and M.7) conclusion that uses the presentation of Theorem~\ref{thm:M2} holds unconditionally on $\Omega$; only $c\neq0$ and the single side were ever used.
Statements that also use Gate~V remain relative to that chosen convention.
\end{remark}

\begin{lemma}[Boundary-value calculus on $\Omega$]
\label{lem:bv}
This lemma concerns the model family $T_m$ itself; it does not yet identify $T_m$ with derivatives of the predecessor kernel $S$.
For $m\ge1$ let $T_m:=(Q-i0)^{-m}:=\lim_{\tau\downarrow0}(Q-i\tau)^{-m}$ in $\mathcal D'(\Omega)$.
The limit exists, and:
\begin{enumerate}
\item[(a)] $Q\,T_{m+1}=T_m$;
\item[(b)] $\partial_\rho T_m=-m\,(\partial_\rho Q)\,T_{m+1}$;
\item[(c)] $T_mT_n=T_{m+n}$, where the product is the existing H\"ormander product on $\Omega$ (wavefront criterion).
\end{enumerate}
\end{lemma}

\begin{proof}
\emph{Points with $Q(z_0)\neq0$.} Near such a point $Q\neq0$ and $(Q-i\tau)^{-m}\to Q^{-m}$ in $C^\infty$, so $T_m=Q^{-m}$ there and (a)--(c) are identities of smooth functions.

\emph{Nonzero null points.} Let $Q(z_0)=0$ with $z_0\ne0$.
Then $dQ(z_0)=2\eta(z_0,\cdot)\neq0$, because $\eta$ is nondegenerate and $z_0\neq0$.
Hence $Q$ is a local coordinate: there is a neighbourhood $U$ and a diffeomorphism $\Phi:U\to(-\delta,\delta)\times V'$ with $\Phi_1=Q$.
On $(-\delta,\delta)$, $(s-i\tau)^{-m}\to(s-i0)^{-m}$ in $\mathcal D'$.
This is the boundary value from $\{\mathrm{Im}\,\zeta<0\}$ of the holomorphic function $\zeta^{-m}$, and $(s-i0)^{-m}=\frac{(-1)^{m-1}}{(m-1)!}\partial_s^{m-1}(s-i0)^{-1}$.
Pullback by the submersion $\Phi$ (H\"ormander, Thm.~6.1.2) is continuous, so $T_m|_U=\Phi^*((s-i0)^{-m}\otimes1)$ and the limit exists on $U$.

Its Fourier transform satisfies $\widehat{(s-i0)^{-m}}(\sigma)=\frac{2\pi i\,(-i\sigma)^{m-1}}{(m-1)!}\,\theta(-\sigma)$ for $\hat u(\sigma)=\int e^{-is\sigma}u(s)\,ds$.
This transform is supported in the single half-line $\sigma\le0$, the same for every $m$.
Hence $\mathrm{WF}((s-i0)^{-m})\subset\{(0,\sigma):\sigma<0\}$, and by the pullback theorem (H\"ormander, Thm.~8.2.4)
\[
\mathrm{WF}(T_m)\subset\Gamma_-:=\{(z,\sigma\,dQ(z)):Q(z)=0,\ z\ne0,\ \sigma<0\},
\]
with one sign for all $m$.

\emph{Compatibility.} All $T_m$ are defined by the same limit $\tau\downarrow0$ from the same side, globally on $\Omega$.
The local descriptions therefore agree on overlaps, and with $Q^{-m}$ off the cone.

(a) $Q(Q-i\tau)^{-m-1}=(Q-i\tau)^{-m}+i\tau(Q-i\tau)^{-m-1}$.
In the local coordinate, $\tau(s-i\tau)^{-m-1}$ is a constant times $\partial_s^{m}\bigl(\tau(s-i\tau)^{-1}\bigr)$, and $\tau(s-i\tau)^{-1}=\partial_s(\tau\log(s-i\tau))$ with $\tau\log(s-i\tau)\to0$ in $L^1_{\mathrm{loc}}$; so the remainder tends to $0$ in $\mathcal D'$.

(b) Derivatives commute with limits in $\mathcal D'$, and $\partial_\rho(Q-i\tau)^{-m}=-m(\partial_\rho Q)(Q-i\tau)^{-m-1}$.

(c) \emph{Same-sign boundary-value product.} Since $\mathrm{WF}(T_m),\mathrm{WF}(T_n)\subset\Gamma_-$ and $\Gamma_-$ contains no opposite covectors $(z,\xi),(z,-\xi)$, H\"ormander's criterion (Thm.~8.2.10) defines the product $T_mT_n$ in $\mathcal D'(\Omega)$.
It remains to identify that product with $T_{m+n}$.

All $T_m$ are boundary values from the \emph{same} lower half-plane.
Locally near a nonzero null point, in the coordinate $s=Q$ of the first paragraph,
\[
T_m|_U=\Phi^*\bigl((s-i0)^{-m}\otimes 1\bigr),\qquad
(s-i0)^{-m}=\mathrm{bv}_{\mathrm{Im}\,\zeta<0}(\zeta^{-m}).
\]
The one-variable distributions $(s-i0)^{-m}$ and $(s-i0)^{-n}$ therefore arise as boundary values of the holomorphic functions $\zeta\mapsto\zeta^{-m}$ and $\zeta\mapsto\zeta^{-n}$ from $\{\mathrm{Im}\,\zeta<0\}$.
We identify their H\"ormander product by an explicit Fourier computation.
Write $u_m:=(s-i0)^{-m}$ and use the one-variable convention $\hat u(\sigma)=\int e^{-is\sigma}u(s)\,\mathrm ds$, with inverse $u(s)=\frac1{2\pi}\int e^{is\sigma}\hat u(\sigma)\,\mathrm d\sigma$; then $\widehat{uv}=\frac1{2\pi}\,\hat u*\hat v$ (H\"ormander's $(2\pi)^{-n}$ with $n=1$).
As above, $\hat u_m(\sigma)=\frac{2\pi i\,(-i\sigma)^{m-1}}{(m-1)!}\theta(-\sigma)$ (for $m=1$: $\widehat{\mathrm{PV}\frac1s}=-i\pi\,\mathrm{sgn}\,\sigma$ and $\widehat{i\pi\delta}=i\pi$; then $\widehat{\partial_su}=i\sigma\hat u$).

\emph{Beta integral.} For $\sigma<0$ the convolution runs over the finite interval $[\sigma,0]$, where both factors are supported; with $\sigma'=\sigma x$, $x\in[0,1]$, $\mathrm d\sigma'$ over $[\sigma,0]$ becomes $|\sigma|\,\mathrm dx$, and
\[
\frac1{2\pi}(\hat u_m*\hat u_n)(\sigma)
=\frac{(2\pi i)^2}{2\pi\,(m-1)!(n-1)!}\int_\sigma^0\bigl(-i(\sigma-\sigma')\bigr)^{m-1}(-i\sigma')^{n-1}\,\mathrm d\sigma'
=\frac{2\pi i\cdot i|\sigma|\,(-i\sigma)^{m+n-2}}{(m-1)!(n-1)!}\,B(m,n).
\]
Since $B(m,n)=\frac{(m-1)!(n-1)!}{(m+n-1)!}$ and $i|\sigma|=-i\sigma$ for $\sigma<0$, this equals $\frac{2\pi i(-i\sigma)^{m+n-1}}{(m+n-1)!}=\hat u_{m+n}(\sigma)$; for $\sigma>0$ the integrand vanishes.
So $\frac1{2\pi}\hat u_m*\hat u_n=\hat u_{m+n}$, the integral being absolutely convergent.

\emph{Identification with the H\"ormander product.} For $\varphi\in C^\infty_c(\mathbb R)$ the H\"ormander product is characterized locally by $\widehat{\varphi^2uv}=\frac1{2\pi}\,\widehat{\varphi u}*\widehat{\varphi v}$, the integral converging absolutely under the wavefront condition (H\"ormander, proof of Thm.~8.2.10).
Here $\widehat{\varphi u_m}=\frac1{2\pi}\hat\varphi*\hat u_m$.
Because $\hat\varphi$ is Schwartz while $\hat u_m,\hat u_n$ are polynomially bounded and supported in $\sigma\le0$, the iterated integral $\int\!\!\int\!\!\int|\hat\varphi(a)\hat\varphi(b)\hat u_m(\sigma-a-b-c)\hat u_n(c)|$ runs in $c$ over an interval of length at most $|\sigma-a-b|$ and converges absolutely; Fubini allows regrouping:
\[
\tfrac1{2\pi}\widehat{\varphi u_m}*\widehat{\varphi u_n}
=\tfrac1{2\pi}\,\widehat{\varphi^2}*\Bigl(\tfrac1{2\pi}\hat u_m*\hat u_n\Bigr)
=\tfrac1{2\pi}\,\widehat{\varphi^2}*\hat u_{m+n}
=\widehat{\varphi^2u_{m+n}} .
\]
Hence $\varphi^2(u_mu_n)=\varphi^2u_{m+n}$ for every $\varphi$, i.e.
\[
(s-i0)^{-m}\,(s-i0)^{-n}=(s-i0)^{-(m+n)}
\quad\text{in }\mathcal D'(\mathbb R_s).
\]
\emph{Transfer to the chart.} With $p(s,y')=s$, $u_m\otimes1=p^*u_m$, and pullback is multiplicative on distributions satisfying the product condition (H\"ormander, Thms.~8.2.4 and~8.2.10: the product is the pullback of the tensor product by the diagonal map, and pullbacks compose); thus $p^*u_m\cdot p^*u_n=p^*u_{m+n}$, and, pulling back further by the diffeomorphism $\Phi$, $T_mT_n=T_{m+n}$ on $U$.
Off the cone the identity is the smooth identity $Q^{-m}Q^{-n}=Q^{-(m+n)}$.
On overlaps the two local descriptions agree by the global same-side limit $\tau\downarrow0$.
Thus $T_mT_n=T_{m+n}$ on $\Omega$.
\end{proof}

\begin{proposition}[Numerator representation]
\label{prop:num}
By Theorem~\ref{thm:wpA-Sbv}, Input~(S-bv) holds.
Then, using Lemma~\ref{lem:bv}(a),(b), $\partial^\alpha S=c\,p_\alpha\,T_{1+|\alpha|}$ on $\Omega$.
Here $p_0=1$ and $p_{\alpha+e_\rho}=Q\,\partial_\rho p_\alpha-(1+|\alpha|)\,p_\alpha\,\partial_\rho Q$, so $p_\alpha$ is homogeneous of degree $|\alpha|$.
For $\lambda=\{\alpha,\beta\}$ of order $N=|\alpha|+|\beta|$, Lemma~\ref{lem:bv}(c) gives
\[
q^\circ_\lambda=(\partial^\alpha S)(\partial^\beta S)=c^2\,n_\lambda\,T_{2+N},\qquad n_\lambda:=p_\alpha p_\beta\in\mathrm{Poly}_N(\mathbb R^4).
\]
\end{proposition}

\begin{proof}
By induction, $\partial_\rho(p_\alpha T_k)=(\partial_\rho p_\alpha)T_k-k\,p_\alpha(\partial_\rho Q)T_{k+1}$ by (b), and $T_k=QT_{k+1}$ by (a); this gives the recursion.
The product is the H\"ormander product used for $q^\circ$ on $y\neq w$ (XVI-J: ``the pointwise product of two copies of \(\Delta_F^{\mathrm{FR}}\) exists''), and (c) applies.
\texttt{kernel\_relations.py} computes $p_\alpha$ from the same recursion, realised as $\partial^\alpha(1/Q)=p_\alpha/Q^{1+|\alpha|}$, for all $|\alpha|\le4$.
\end{proof}

\begin{proposition}[Two-direction equivalence of relations]
\label{prop:relations}
Use Theorem~\ref{thm:wpA-Sbv} and fix $N$.
For constants $r_\lambda$:
\begin{enumerate}
\item if $\sum_\lambda r_\lambda n_\lambda=0$ as a polynomial, then $\sum_\lambda r_\lambda q^\circ_\lambda=0$ in $\mathcal D'(\Omega)$;
\item if $\sum_\lambda r_\lambda q^\circ_\lambda=0$ in $\mathcal D'(\Omega)$, then $\sum_\lambda r_\lambda n_\lambda=0$ as a polynomial.
\end{enumerate}
Relations mixing different $N$ split into homogeneous components, because $q^\circ_\lambda$ is homogeneous of degree $-4-N$.
Hence the nullspace of the numerator map is exactly the relation space of $\pi$.
\end{proposition}

\begin{proof}
(1) By Proposition~\ref{prop:num}, $\sum r_\lambda q^\circ_\lambda=c^2(\sum r_\lambda n_\lambda)T_{2+N}=0$.
(2) Restrict to the nonempty open set $\{Q<0\}\subset\Omega$ (it contains $(0,1,0,0)$).
There $T_{2+N}=Q^{-2-N}$ is smooth and nowhere zero, so $c^2(\sum r_\lambda n_\lambda)Q^{-2-N}=0$ pointwise.
Hence the polynomial $\sum r_\lambda n_\lambda$ vanishes on a nonempty open set and is therefore zero.
For mixed $N$, dilations $z\mapsto tz$ act on $\Omega$, and distinct homogeneity degrees separate.
\end{proof}

\begin{theorem}[Certified graded presentation]
\label{thm:M2}
With Input~(S-bv) certified (Theorem~\ref{thm:wpA-Sbv}):
\begin{center}
\begin{tabular}{ccccc}
\hline
$N$ & labels & $\dim\mathrm{Poly}_N$ & $d_N=\dim\mathcal H_{4+N}$ & $\dim\ker\pi_N$\\
\hline
$0$ & $1$ & $1$ & $1$ & $0$\\
$1$ & $4$ & $4$ & $4$ & $0$\\
$2$ & $20$ & $10$ & $10$ & $10$\\
$3$ & $60$ & $20$ & $20$ & $40$\\
$4$ & $170$ & $35$ & $35$ & $135$\\
\hline
Total & $255$ & & $70$ & $185$\\
\hline
\end{tabular}
\end{center}
The numerator map is onto $\mathrm{Poly}_N$.
A basis $\mathcal B_N$ extends a basis of $\mathcal V_{\mathrm K}=\mathcal H_4\oplus\mathcal H_5\oplus\mathcal H_6$ ($\dim15$).
Every label reconstructs as $q^\circ_\lambda=\sum_{b\in\mathcal B_N}R_{\lambda b}q^\circ_b$ with exact rational $R$.
\end{theorem}

\begin{proof}
Propositions~\ref{prop:num} and~\ref{prop:relations}, together with the exact $\mathbb Q$-linear algebra reported in Lemma~\ref{lem:M2comp}.
\end{proof}

\begin{proposition}[Unconditional wave-operator family]
\label{prop:wave}
Without using Input~(S-bv), $\sum_{\rho,\sigma}\eta^{\rho\sigma}q^\circ_{\{\alpha,\beta+e_\rho+e_\sigma\}}=0$ on $\Omega$, since $\square S=0$ there (XVI-K, lem:grading(3)).
The family spans a subspace of dimensions $(0,0,1,8,36)$, so $\dim\mathcal V_{\mathrm{adm}}\le210$ without any kernel identification.
With (S-bv) (certified), $140$ further independent relations appear.
\end{proposition}

\begin{remark}[A relation outside the wave family; sign check]
With $\eta=\mathrm{diag}(1,-1,-1,-1)$, $\partial_0Q=2z_0$ and $\partial_1Q=-2z_1$, so $p_{e_0}=-2z_0$, $p_{e_1}=2z_1$, $p_{e_0}p_{e_1}=-4z_0z_1$, and $p_{e_0+e_1}=-8z_0z_1$.
Hence $(\partial_0S)(\partial_1S)=\tfrac12\,S\,\partial_0\partial_1S$ on $\Omega$ (Theorem~\ref{thm:wpA-Sbv}); this is the relation \texttt{q\{1, 0\} - (1/2) q\{(), 01\} = 0}.
The script checks exactly this example (\texttt{require}: ``sign example'').
The relation is not generated by $\square S=0$.
\end{remark}

\begin{lemma}[Reported computation for M.2]
\label{lem:M2comp}
\texttt{kernel\_relations.py} exits with status $0$ and records:
\begin{itemize}
\item the table of Theorem~\ref{thm:M2};
\item reconstruction certificates for all labels (\texttt{kernel\_basis\_table.txt}, which also lists every numerator row $n_\lambda$);
\item $185$ certified relation vectors and $45$ wave-operator generators (\texttt{kernel\_relation\_matrix.txt});
\item the sign example.
\end{itemize}
\texttt{regression\_M.py} independently spot-checks every relation (on $(\partial^\alpha S)(\partial^\beta S)$ at three rational points).
\end{lemma}

\begin{remark}[Analytic boundary and dependency on (S-bv)]
\label{rem:M2boundary}
An identity of rational functions away from the null cone is not automatically a distributional identity on $\Omega$.
Lemma~\ref{lem:bv} is about the model family $T_m$ and does not by itself identify $S$ with $c\,T_1$.
The identification of the numerical kernels $q^\circ_\lambda$ with the corresponding products of $T_m$'s---and hence the numerical conclusions of M.2--M.4 that use the graded presentation---rests on Input~(S-bv), which is certified by Theorem~\ref{thm:wpA-Sbv}; these conclusions therefore hold unconditionally on $\Omega$.
Proposition~\ref{prop:wave} remains the part that needs no kernel identification.
The linear Helmholtz lemma (Lemma~\ref{lem:helm}) is pure variational calculus under Gate~V and never depended on (S-bv).
\end{remark}
%---------------------------------------------------------------------
\subsection{Kernel identification: Gate S-bv certified (work package A)}
\label{sec:wpA-Sbv}
\scopebox{Status (work package A).}{%
\textbf{Constructed (proved):} on $\Omega:=\mathbb R^4\setminus\{0\}$,
\[
S=h_0+iE_0^{\mathrm A}=-W_0+iE_0^{\mathrm A}
=\frac1{4\pi^2}\,(Q-i0)^{-1},
\]
with $Q(z)=\eta(z,z)$, $\eta=\mathrm{diag}(1,-1,-1,-1)$, including across the nonzero null cone (Theorem~\ref{thm:wpA-Sbv}).
\textbf{Constructed:} convention ledger (\texttt{wpA/kernel\_conventions.json}); symbolic and numerical cross-checks (\texttt{wpA/kernel\_identification\_checks.py}); local boundary-value lemma for the variable regulator (Lemma~\ref{lem:wpA-varreg}).
\textbf{Predecessor status:} the formula is printed in XVI-J, eq:rational (with $S=-d$, XVI-J, eq:DF); this subsection is a written check of it, not a new extension.\newline\textbf{Non-claims:} the predecessor $S$ is globally defined (XVI-J); what is not established here is its position-space identification at the origin; no $i\varepsilon$ Feynman contour is introduced as a definition; the massive case is not needed for this edge; nothing about M-C, Ward or BV.
Gate~S-bv / Input~(S-bv) is discharged by the theorem below, with the explicit constant $c=1/(4\pi^2)$ and boundary side $Q-i0$.
The Euclidean calculation of Remark~\ref{rem:Sbv-motivation} remains a motivation only; it is not the certification.
}

\paragraph{Convention ledger.}
\label{sec:wpA-conventions}
\begin{center}
\small
\begin{longtable}{p{0.20\textwidth}p{0.46\textwidth}p{0.25\textwidth}}
\hline
Item & Value (predecessor-authoritative) & Source\\
\hline
\endhead
Difference variable & $z=x-y$; $S$ is a translation-invariant distribution of $z$ & XVI-K, after eq:S\\
Caveat & XVI-K uses $z=y-w$ for the \emph{bubble extension} only; this is not a redefinition of the edge kernel $S$ & XVI-K (extension step)\\
Signature / $Q$ & $\eta=\mathrm{diag}(1,-1,-1,-1)$, $Q(z)=z_0^2-|\mathbf z|^2$ & XI-B; Input~\ref{inp:Sbv}\\
Wave operator & $\square=\partial_t^2-\Delta$; $\square E_0^{\mathrm{R/A}}=\mathrm{id}$ & XI-B, lem:Green; XVI-J, lem:signs\\
Spatial Fourier in $W_0$ & $W_0(x,y)=\int\frac{e^{-i|\mathbf p|(x^0-y^0)+i\mathbf p\cdot(\mathbf x-\mathbf y)}}{(2\pi)^3\,2|\mathbf p|}\,\mathrm d^3\mathbf p$ (massless) & XI-E, eq:W; XVI-K, eq:W-fam\\
Test-function FT & $\tilde\varphi(z^0,\mathbf p)=\int e^{i\mathbf p\cdot\mathbf z}\varphi\,\mathrm d^3\mathbf z$ & XVI-K, lem:families\\
$h_0$ & $h_0=-W_0$, $h_0-h_0^{\mathsf T}=i\delta_0$ & XI-E, eq:h\\
$E_0^{\mathrm A}$ (sine) & $(E_0^{\mathrm A}F)(t)=\int_t^\infty\frac{\sin((s-t)\Omega_0)}{\Omega_0}F(s)\,\mathrm ds$ & XVI-J, eq:sineA; XVI-K, eq:E-fam\\
$E_0^{\mathrm A}$ (step) & $E_0^{\mathrm A}=-\theta(y^0-x^0)\,\delta_0=-\theta(-z^0)\,\delta_0$ & XVI-J, eq:step\\
$S$ & $S:=h_0+iE_0^{\mathrm A}=-W_0+iE_0^{\mathrm A}=-\Delta_F^{\mathrm{FR}}$ & XVI-K, eq:S; XVI-J, lem:identify\\
Reconciliation & ``$S=h_0+iE_0^{\mathrm A}$'' and ``$S=-W_0+iE_0^{\mathrm A}$'' are equivalent via $h_0=-W_0$ & XI-E, eq:h; XVI-K, eq:S\\
Transposition & $\langle T^{\mathsf T},\phi(x,y)\rangle=\langle T,\phi(y,x)\rangle$; for $u(z)$, $u^{\mathsf T}(z)=u(-z)$ & XVI-J: $(E_0^{\mathrm A})^{\mathsf T}=E_0^{\mathrm R}$, $W_0^{\mathsf T}=W_0+i\delta_0$\\
$(Q-i0)^{-1}$ & $\lim_{\tau\downarrow0}(Q-i\tau)^{-1}$ in $\mathcal D'(\Omega)$; equivalently $\Phi^*((s-i0)^{-1}\otimes1)$ near nonzero null points & Lemma~\ref{lem:bv}\\
One-variable identity & $(s-i0)^{-1}=\mathrm{PV}(1/s)+i\pi\delta(s)$ (sign $+i\pi$) & classical; checked in \texttt{wpA/kernel\_\allowbreak identification\_\allowbreak checks.py}\\
Constant / side & $c=1/(4\pi^2)$; boundary side $Q-i0$ (the side of $T_1$) & Theorem~\ref{thm:wpA-Sbv}\\
\hline
\end{longtable}
\end{center}
No abstract $\mathbb R^4$ Fourier transform of $S$ is fixed by the predecessors---only the spatial Fourier formulae for $W_0$ and the sine kernel.
Signs are \emph{not} inferred from the word ``Feynman''; $\Delta_F^{\mathrm{FR}}$ enters only through XVI-J, lem:identify, as an identification of the already-defined combination $\theta_+W_0+\theta_-W_0^{\mathsf T}$.

\paragraph{Position-space identification.}
\label{sec:wpA-identify}
\emph{What the predecessors already print (audit~3).}
XI-E, eq:h (box copy \texttt{enter-companion/raw/tex/XI-E.tex}, lines 366--371, verified) only normalizes the Wightman factor and records the bisolution jump:
\begin{quote}\small
``To match the proposed block notation put
\(h_\mu=-W_\mu,\ h_\mu-h_\mu^{\mathsf T}=i\delta_\mu.\)''
\end{quote}
It gives no position-space boundary value.
XVI-J (version~1, \texttt{content/manuscripts/XVI-J.tex} of \texttt{cheungngo/enter-companion}), eq:W, eq:EA and eq:DF read (macros expanded):
\begin{quote}\small
``Define the massless positive-frequency distribution by
\(W_0(z)=\int_{\mathbb R^3}\frac{e^{-i|\mathbf p|z^0+i\mathbf p\cdot\mathbf z}}{(2\pi)^3\,2|\mathbf p|}\,\mathrm d^3\mathbf p.\)
The integral is distributional. Set \(h_0=-W_0\).''
\ldots\ ``\(\widehat E_0^{\mathrm A}(t,\mathbf p)=-\theta(-t)\frac{\sin(t|\mathbf p|)}{|\mathbf p|}\)'' \ldots\
``Put
\(d=W_0-iE_0^{\mathrm A},\quad S=h_0+iE_0^{\mathrm A}=-d,\quad D_F=-\mathcal PS=\mathcal Pd.\)''
\end{quote}
XVI-J, eq:rational:
\begin{quote}\small
``Let \(Q_{\mathrm{geom}}(z)=(z^0)^2-|\mathbf z|^2.\)
On \(z\ne0\), direct Fourier evaluation gives
\(d(z)=-\frac1{4\pi^2}\bigl(Q_{\mathrm{geom}}(z)-i0\bigr)^{-1}.\)
At a nonzero null point, \(Q_{\mathrm{geom}}\) is a submersion.
The boundary value in [eq:rational] has one, rather than both, conormal orientations.''
\end{quote}
XVI-J, prop:q:
\begin{quote}\small
``The distribution \(q^\circ(z)=S(z)^2=d(z)^2,\ z\ne0,\) exists on \(\mathbb R^4\setminus\{0\}\), including nonzero null points.
It is homogeneous of degree \(-4\), has scaling degree four at the origin, and is nonzero.''
\end{quote}
Thus $S=-d=\frac1{4\pi^2}(Q-i0)^{-1}$ on $z\neq0$, including nonzero null points, is the XVI-J formula ($Q=Q_{\mathrm{geom}}$), and XVI-J, prop:q, uses it.
Theorem~\ref{thm:wpA-Sbv} is a \emph{written check} of XVI-J, eq:rational, whose printed justification is ``direct Fourier evaluation''; it is not a new extension.
XVI-M's contribution is (i) the written real-time derivation $W_0=-\frac1{4\pi^2}\bigl(\mathrm{PV}\frac1Q+i\pi\,\mathrm{sgn}(z^0)\delta(Q)\bigr)$, $E_0^{\mathrm A}=\frac1{2\pi}\theta(-z^0)\delta(Q)$ and the cancellation giving $\frac1{4\pi^2}(Q-i0)^{-1}$; and (ii) the local justification at nonzero null points of the boundary value of the variable regulator $Q-2i\tau z^0-\tau^2$ (Lemma~\ref{lem:wpA-varreg}).
\emph{Predecessor provenance.} Each predecessor version used below is pinned by SHA-256 (local files) or by an archived corpus excerpt under \texttt{XVI-M-artifacts/\allowbreak provenance/} (regression re-checks those hashes). Citations are by label only.
\begin{itemize}\setlength{\itemsep}{0.15em}\small
\item \textbf{Local} \texttt{XVI-J.tex} (\texttt{dcba2602ab4a\ldots}): labels eq:W0, eq:sineA, eq:step, lem:signs, lem:identify. Role: supplementary (earlier box version).
\item \textbf{Repository fetch} of XVI-J (\texttt{d5558cd419c8\ldots}, \texttt{provenance/XVI-J\_github\_\allowbreak fetch.tex}): HTTP~404 stub. Role: failed retrieval; does not remove the corpus file.
\item \textbf{Corpus excerpt} of XVI-J (\texttt{eda87cb7400d\ldots}, \texttt{provenance/XVI-J\_corpus\_\allowbreak eqW\_DF\_\allowbreak rational\_propq.txt}): labels eq:W, eq:EA, eq:DF, eq:rational, prop:q. Role: \textbf{supplementary corroboration}; eq:rational is \emph{not} part of the frozen XVI-L dependency; Theorem~\ref{thm:wpA-Sbv} and Lemma~\ref{lem:wpA-varreg} stand independently.
\item \textbf{Frozen} \texttt{XVI-L.tex} (\texttt{db5f027ef08e\ldots}): thm:L5, thm:L6, def:RL,\ldots\ Role: \textbf{frozen dependency}. Corpus excerpt of thm:L6 (\texttt{1ef32e999bfc\ldots}).
\item \textbf{Local} \texttt{XVI-K.tex}/artifacts (\texttt{6264bc8ec685\ldots}): eq:S, prop:hess,\ldots\ Role: frozen interface. Corpus excerpt (\texttt{f0ebf560b5d6\ldots}): prop:hess, eq:Lpq.
\item \textbf{Corpus/local} XI-E (\texttt{ea40d58d5912\ldots}/\texttt{d686b6ec595e\ldots}): eq:W, eq:h. Role: supplementary corroboration.
\item \textbf{Corpus/local} XII-C (\texttt{79c1765ef7d8\ldots}/\texttt{83c6d0383b0a\ldots}): lem:weighted-curvature. Role: supplementary corroboration.
\end{itemize}
A failed repository retrieval does not make eq:rational a premise of XVI-L.

\begin{lemma}[Variable regulator near a nonzero null point (audit~3)]
\label{lem:wpA-varreg}
Let $z_\ast\in\Omega$ with $Q(z_\ast)=0$ and put $s_\ast:=\mathrm{sgn}(z_\ast^0)$.
Then $z^0_\ast\neq0$, and there are an open neighbourhood $U\Subset\Omega$ of $z_\ast$ and a chart $\Phi=(Q,y'):U\to(-\delta,\delta)\times V'$ such that, with $a_0:=|z^0_\ast|$, $s_\ast z^0\ge a_0/2>0$ on $U$.
For $0<\tau\le1$ put $g_\tau(z):=Q(z)-2i\tau z^0-\tau^2$.
Then $g_\tau$ has no zero on $\mathbb R^4$, and
\[
\lim_{\tau\downarrow0}\frac1{g_\tau}=\Phi^*\bigl((s-i0\,s_\ast)^{-1}\otimes1\bigr)
\quad\text{in }\mathcal D'(U),
\qquad
(s-i0\,s_\ast)^{-1}:=\lim_{\eta\downarrow0}(s-i\eta s_\ast)^{-1}=\mathrm{PV}\frac1s+i\pi s_\ast\,\delta(s).
\]
For $s_\ast=+1$ the right-hand side is $T_1|_U$ of Lemma~\ref{lem:bv}.
\end{lemma}

\begin{proof}
If $z^0_\ast=0$ then $Q(z_\ast)=-|\mathbf z_\ast|^2=0$ forces $z_\ast=0\notin\Omega$.
The chart with first coordinate $Q$ exists by the proof of Lemma~\ref{lem:bv}; shrinking it, we may assume $\overline U$ compact in $\Omega$ and $s_\ast z^0\ge a_0/2$ on $\overline U$, with $a_0:=|z^0_\ast|$.
A zero of $g_\tau$ needs $\mathrm{Im}\,g_\tau=-2\tau z^0=0$ and then $Q=-|\mathbf z|^2=\tau^2>0$, which is impossible.
In the chart, $g_\tau\circ\Phi^{-1}=w_\tau:=s-\tau^2-i\tau a$ with $a:=2z^0\circ\Phi^{-1}$ smooth, $s_\ast a\ge a_0$, $|a|\le a_1$, $|\partial_sa|\le a_2$ on the chart (after a further shrinking).
Pullback by the diffeomorphism $\Phi$ is continuous on $\mathcal D'$ and the Jacobian is absorbed in the test function, so it suffices to show, for $\psi\in C^\infty_c((-\delta,\delta)\times V')$,
$\int\psi\,w_\tau^{-1}\,\mathrm ds\,\mathrm dy'\to\langle(s-i0\,s_\ast)^{-1}\otimes1,\psi\rangle$.

Since $\mathrm{Im}\,w_\tau=-\tau a$ has the fixed sign $-s_\ast$, $w_\tau$ stays in one open half-plane, on which the principal logarithm $L_\tau:=\mathrm{Log}\,w_\tau$ is smooth, and $\partial_sL_\tau=(1-i\tau\partial_sa)/w_\tau$.
Hence
\[
\frac1{w_\tau}=\partial_sL_\tau+\frac{i\tau\,\partial_sa}{w_\tau}.
\]
(i) \emph{Remainder.} $|\tau\,\partial_sa/w_\tau|\le\tau a_2/|\mathrm{Im}\,w_\tau|\le a_2/a_0$, and for $s\neq0$ the term tends to $0$ (since $w_\tau\to s\ne0$).
By dominated convergence on the compact support of $\psi$, $\int\psi\,i\tau\partial_sa/w_\tau\to0$.

(ii) \emph{Main term.} Integrating by parts in $s$, $\int\psi\,\partial_sL_\tau=-\int(\partial_s\psi)L_\tau$.
For $s\ne0$, $L_\tau\to\mathrm{Log}_{s_\ast}(s):=\log|s|-i\pi s_\ast\theta(-s)$ (for $s<0$ the point $w_\tau$ approaches the negative axis from the half-plane $\mathrm{sgn}\,\mathrm{Im}=-s_\ast$).
Moreover $|\mathrm{Im}\,L_\tau|\le\pi$ and $|s-\tau^2|\le|w_\tau|\le C:=\delta+1+a_1$, so $|\mathrm{Re}\,L_\tau|\le|\log C|+\bigl|\log|s-\tau^2|\bigr|$.
The family $\{\log|s-\tau^2|\}_{0<\tau\le1}$ is bounded in $L^2(-\delta,\delta)$ (each is a translate of $\log|u|$ restricted to $(-\delta-1,\delta+1)$), hence uniformly integrable on $\mathrm{supp}\,\psi$; by Vitali's theorem $\int(\partial_s\psi)L_\tau\to\int(\partial_s\psi)\,\mathrm{Log}_{s_\ast}(s)$.

(iii) \emph{Identification.} The case $a\equiv2s_\ast$ without the shift $\tau^2$ (i.e.\ $w=s-i\eta s_\ast$, remainder $\equiv0$) gives by the same argument $(s-i\eta s_\ast)^{-1}\to\partial_s\mathrm{Log}_{s_\ast}(s)$ in $\mathcal D'$; hence both limits equal $\partial_s\mathrm{Log}_{s_\ast}$.
Finally $\partial_s\log|s|=\mathrm{PV}\frac1s$ and $\partial_s\bigl(-i\pi s_\ast\theta(-s)\bigr)=i\pi s_\ast\delta(s)$.
For $s_\ast=+1$, $(s-i\eta)^{-1}\otimes1$ is $(Q-i\eta)^{-1}$ in the chart, whose limit defines $T_1|_U$.
\end{proof}

\begin{theorem}[Certified identification for Gate S-bv]
\label{thm:wpA-Sbv}
Let $W_0$, $h_0=-W_0$ and $E_0^{\mathrm A}$ be as in XVI-K, eq:S / lem:families (quoting XI-E, eq:W, eq:h and XVI-J, eq:sineA, eq:step), massless case.
Set $S:=h_0+iE_0^{\mathrm A}$ and $Q(z)=\eta(z,z)$.
Then, in $\mathcal D'(\Omega)$,
\begin{equation}
\label{eq:wpA-id}
S=\frac1{4\pi^2}\,(Q-i0)^{-1}
=\frac1{4\pi^2}\Bigl(\mathrm{PV}\frac1Q+i\pi\,\delta(Q)\Bigr),
\end{equation}
where $(Q-i0)^{-1}$ is the boundary-value distribution $T_1$ of Lemma~\ref{lem:bv}, and $\delta(Q)$ denotes the pullback surface distribution $\Phi^*(\delta(s)\otimes1)$ in any local chart with first coordinate $Q$.
In particular Input~(S-bv) holds with
\[
c=\frac1{4\pi^2}\neq0
\]
and with the same boundary side $Q-i0$ as $T_1$.
Since $S=-d$ (XVI-J, eq:DF), this is XVI-J, eq:rational; the theorem is a written check of it, not a new extension.
\end{theorem}

\begin{proof}
\emph{(1) Position-space formulae from the Fourier definitions.}
\textit{Wightman factor.}
From XVI-K, eq:W-fam / XI-E, eq:W, with $z=x-y$ and $\omega_0=|\mathbf p|$,
\[
W_0(z)
=\int_{\mathbb R^3}\frac{e^{-i|\mathbf p|z^0+i\mathbf p\cdot\mathbf z}}{(2\pi)^3\,2|\mathbf p|}\,\mathrm d^3\mathbf p .
\]
Angular integration yields the radial formula
\[
W_0(z)=\frac1{4\pi^2\,r}\int_0^\infty\sin(pr)\,e^{-ipz^0}\,\mathrm dp,
\qquad r=|\mathbf z|.
\]
Insert the positive-frequency damping $z^0\mapsto z^0-i\tau$ ($\tau\downarrow0$), i.e.\ $\alpha=\tau+iz^0$ with $\mathrm{Re}\,\alpha>0$, and use
$\int_0^\infty e^{-\alpha p}\sin(pr)\,\mathrm dp=r/(\alpha^2+r^2)$.
(The damping parameter is called $\tau$ throughout; it is unrelated to the preparation width $\varepsilon$. The factor $e^{-\tau|\mathbf p|}$ only realises the oscillatory integral as a limit in $\mathcal D'(\mathbb R^4)$: for a test function $\varphi$, $\langle W_0^{(\tau)},\varphi\rangle=\int\mathrm dz^0\int\frac{e^{-(\tau+iz^0)|\mathbf p|}}{(2\pi)^32|\mathbf p|}\tilde\varphi(z^0,-\mathbf p)\,\mathrm d^3\mathbf p$, and the integrand is dominated by $|\tilde\varphi(z^0,-\mathbf p)|/((2\pi)^32|\mathbf p|)$, which is integrable; it changes no definition and introduces no $i\varepsilon$ prescription.)
One obtains the smooth function
\[
W_0^{(\tau)}(z)=\frac1{4\pi^2(\alpha^2+r^2)}
=-\frac1{4\pi^2\bigl(Q-2i\tau z^0-\tau^2\bigr)}=-\frac1{4\pi^2\,g_\tau(z)} .
\]
Off the cone, $g_\tau^{-1}\to Q^{-1}$ in $C^\infty$; near each nonzero null point, where $z^0$ has a fixed sign, Lemma~\ref{lem:wpA-varreg} gives $g_\tau^{-1}\to(Q-i0\,\mathrm{sgn}(z^0))^{-1}$.
Passing to $\tau\downarrow0$ therefore gives the distributional identity on $\Omega$
\begin{equation}
\label{eq:wpA-W}
W_0(z)
=-\frac1{4\pi^2}\,\frac1{Q-i0\,\mathrm{sgn}(z^0)}
=-\frac1{4\pi^2}\Bigl(\mathrm{PV}\frac1Q+i\pi\,\mathrm{sgn}(z^0)\,\delta(Q)\Bigr).
\end{equation}
(The Euclidean slice $z^0=-i\tau_{\mathrm E}$, $\tau_{\mathrm E}>0$, recovers $W_0=1/(4\pi^2(\tau_{\mathrm E}^2+r^2))$, matching Remark~\ref{rem:Sbv-motivation} and the check in \texttt{kernel\_relations.py}.)

\textit{Causal factor.}
From $W_0-W_0^{\mathsf T}=-i\delta_0$ (XI-E, after eq:W; XVI-J) and \eqref{eq:wpA-W},
$\delta_0=\frac1{2\pi}\,\mathrm{sgn}(z^0)\,\delta(Q)$ on $\Omega$.
XVI-J, eq:step, gives $E_0^{\mathrm A}=-\theta(-z^0)\,\delta_0$, hence
\begin{equation}
\label{eq:wpA-EA}
E_0^{\mathrm A}(z)=\frac1{2\pi}\,\theta(-z^0)\,\delta(Q)
\quad\text{on }\Omega .
\end{equation}
(Equivalently, $E_0^{\mathrm R}=\theta(z^0)\,\delta_0=\frac1{2\pi}\theta(z^0)\delta(Q)$ is the standard massless retarded kernel for $\square E=\delta$.)

\emph{(2) Combination: cancellation of $\mathrm{sgn}(z^0)$.}
Using $h_0=-W_0$,
\[
-W_0=\frac1{4\pi^2}\mathrm{PV}\frac1Q+\frac{i}{4\pi}\,\mathrm{sgn}(z^0)\,\delta(Q),
\qquad
iE_0^{\mathrm A}=\frac{i}{2\pi}\,\theta(-z^0)\,\delta(Q).
\]
On $\{z^0>0\}$ one has $\mathrm{sgn}=+1$, $\theta(-z^0)=0$; on $\{z^0<0\}$, $\mathrm{sgn}=-1$, $\theta(-z^0)=1$.
The coefficient of $\delta(Q)$ becomes
\[
\frac{i}{4\pi}\mathrm{sgn}(z^0)+\frac{i}{2\pi}\theta(-z^0)
=\frac{i}{4\pi}\bigl(\theta(z^0)+\theta(-z^0)\bigr)
=\frac{i}{4\pi}
=\frac1{4\pi^2}\cdot i\pi
\]
on both sheets.
(The identity $\theta(z^0)+\theta(-z^0)=1$ holds on $\Omega\cap\{Q=0\}$, because $Q=0$ and $z^0=0$ force $\mathbf z=0$, hence $z=0\notin\Omega$.)
Therefore
\[
S=-W_0+iE_0^{\mathrm A}
=\frac1{4\pi^2}\Bigl(\mathrm{PV}\frac1Q+i\pi\,\delta(Q)\Bigr)
=\frac1{4\pi^2}\,(Q-i0)^{-1},
\]
where the last step is the pullback of $(s-i0)^{-1}=\mathrm{PV}(1/s)+i\pi\delta(s)$.
The $\mathrm{sgn}(z^0)$-dependence of $W_0$ is exactly cancelled by $iE_0^{\mathrm A}$, which yields the \emph{single} boundary value $Q-i0$ (not $Q+i0$).

\emph{(3) Normalization off the cone.}
On the open sets $\{Q>0\}$ and $\{Q<0\}$, $\delta(Q)=0$ and $E_0^{\mathrm A}=0$, so $S=c/Q$ with $c=1/(4\pi^2)$; e.g.\ $S=-c$ at $(0,1,0,0)$ and $S=c$ at $(1,0,0,0)$.

\emph{(4) Equality across the nonzero null cone.}
Let $z_\ast\in\Omega$ with $Q(z_\ast)=0$.
Then $\mathrm dQ(z_\ast)=2\eta(z_\ast,\cdot)\neq0$, so $Q$ is a local coordinate: there is a chart $\Phi$ with $\Phi_1=Q$ (proof of Lemma~\ref{lem:bv}).
In that chart $(Q-i0)^{-1}=\Phi^*((s-i0)^{-1}\otimes1)$, and step~(2) shows that $-W_0+iE_0^{\mathrm A}$ equals $c$ times this same pullback on both time sheets.
The surface distribution $\delta(Q)$ is well defined on the cone minus the origin by the same submersion, and the sheet-wise boundary values $(Q-i0\,\mathrm{sgn}\,z^0)^{-1}$ used in~\eqref{eq:wpA-W} are those of Lemma~\ref{lem:wpA-varreg}.
Thus \eqref{eq:wpA-id} holds near every nonzero null point, and therefore on all of $\Omega$.

\emph{(5) Cross-checks.}
(a) $\square W_0=0$ and $\square E_0^{\mathrm A}=\delta$ (XI-B, lem:Green), so $\square S=i\delta$ on $\mathbb R^4$ and $\square S=0$ on $\Omega$, matching XVI-K, lem:grading(3); off the cone $\square(c/Q)=0$ as a rational identity in four dimensions.
(b) Both sides are homogeneous of degree $-2$ (XVI-K, lem:homog).
(c) $W_0-W_0^{\mathsf T}=-i\delta_0$ and $E_0^{\mathrm R}-E_0^{\mathrm A}=\delta_0$ imply $S^{\mathsf T}=S$, consistently with $D_F=-\Pc S$ (XVI-J, lem:signs).
(d) $KD_F=i\,\mathrm{id}$ with $K=-\square\Pc$ and $D_F=-\Pc S$ is equivalent to $\square S=i\delta$.
\end{proof}

\begin{remark}[Downstream impact]
\label{rem:wpA-M6}
\emph{M.2--M.4, M.7.} The numerator relations, the graded presentation and everything built on it use only $c\neq0$ and the single side $T_1$; they are now unconditional on $\Omega$ (statements that also use Gate~V stay relative to that convention).
\emph{M.6.} Theorem~\ref{thm:M6-S2} gives $r=-2i\pi^2c^2$; with $c=1/(4\pi^2)$,
$r=-2i\pi^2\cdot\frac1{16\pi^4}=-\frac{i}{8\pi^2}$.
The boundary side is $Q-i0$, so the angular integral $\int_{S^3}(Q-i0)^{-2}\,\mathrm d\sigma=-2i\pi^2$ of \texttt{scaling\_anomaly.py} is unchanged in sign; the counterfactual $Q+i0$ would conjugate that integral and flip $r$, and it is excluded by Theorem~\ref{thm:wpA-Sbv}.
The top-order scale-residue unit of Section~\ref{sec:wpB} is $i\pi^2c^2=i/(16\pi^2)$.
\end{remark}

\begin{remark}[What remains open after work package A]
\label{rem:wpA-open}
The predecessor $S$ is globally defined (XVI-J); what is not established here is its position-space identification at the origin. (Gate~S-bv concerns $\Omega$ only.)
The massive formulae ($\mu>0$ in XI-E, eq:W) are not needed for the massless edge kernel of XVI-K/J.
Quoted predecessor displays (XVI-K, eq:S and eq:W-fam; XVI-J, eq:step; XI-E, eq:h) are reproduced in \texttt{wpA/kernel\_identification.tex} and \texttt{wpA/kernel\_conventions.json}.
\end{remark}

\begin{lemma}[Reported computation for work package A]
\label{lem:wpA-comp}
\texttt{wpA/kernel\_identification\_checks.py} exits with status $0$ (about $2$\,s) and prints \texttt{ALL CHECKS PASSED}.
It checks: the radial Laplace integral and the damped closed form of $W_0$ (symbolic, with numerical spot checks); the sign of $(s\mp i0)^{-1}$; the cancellation of $\mathrm{sgn}(z^0)$ and the coefficients $c$ and $c\,i\pi$; $\delta_0$ and $E_0^{\mathrm A}$ on $\Omega$; off-cone normalization at $12$ points; numerical spatial Fourier integrals against the closed form; the null-cone chart $s=Q$ on both sheets; the variable regulator of Lemma~\ref{lem:wpA-varreg} (numerical pairing for $a=2+\sin s$ and $a=-(1.5+0.5\cos s)$ at $\tau=10^{-2},10^{-3},10^{-4}$, error decreasing linearly in $\tau$; audit~3); $\square(c/Q)=0$, homogeneity, transposition, propagator normalization; and $r=-i/(8\pi^2)$.
The output is deterministic (fixed seeds) and is re-generated and compared byte-for-byte by \texttt{regression\_M.py}.
\end{lemma}

%=====================================================================
\section{M.3 Contact space, graph maps, and invisible directions}
\label{sec:m3}

\begin{definition}[Coherent contact space]
\label{def:Ccoh}
Let $\mathcal D_{\le N}=\mathrm{span}\{\partial^\gamma\delta:|\gamma|\le N\}$.
After Theorem~\ref{thm:M2}, and using XVI-L, thm:L5,
\[
\Ccoh=\bigoplus_{N=0}^4\mathrm{Hom}\bigl(\mathcal H_{4+N},\mathcal D_{\le N}\bigr),\qquad\dim\Ccoh=3321
\]
(using Theorem~\ref{thm:M2}, hence Theorem~\ref{thm:wpA-Sbv}; raw labelled slots $14321$).
Coordinates are $C_{N,b,\gamma}$ with $b\in\mathcal B_N$ and $|\gamma|\le N$.
The regression-preserving subspace $\Creg=\{\Delta\RL|_{\mathcal V_{\mathrm K}}=0\}$ consists of the coordinates with $N\ge3$; it has dimension $3150$.
\end{definition}

\begin{definition}[Three graph maps and two named variants]
\label{def:Gmaps}
Write $\bar a_b:=\sum_\lambda R_{\lambda b}a_\lambda$ for the reduced coefficient of a basis kernel $b$, using the physical (switched) $a_\lambda=\chiI\hat a_\lambda$ of Definition~\ref{def:aGeneral}.
Equivalently $\bar a_b=\chiI(w)\,\hat{\bar a}_b$ with the switch-free factor $\hat{\bar a}_b:=\sum_\lambda R_{\lambda b}\hat a_\lambda$ (built from the table $\hat a$).
For $C\in\Ccoh$ put
\[
\mathcal G(C)=\sum_{N,b,\gamma}C_{N,b,\gamma}\,\bar a_b^{\mu\nu}(y,w)\,(\partial^\gamma\delta)(y-w)
\]
(equivalently $\sum C_{N,b,\gamma}\,\chiI(w)\,\hat{\bar a}_b^{\mu\nu}(y,w)\,(\partial^\gamma\delta)(y-w)$; the preferred form matches $E_C$ below).
\begin{enumerate}
\item $\mathcal G_{\mathrm{formal}}$, the \emph{formal universal-jet reading}.
External data $\sigma$ (preparation), $\chi$ (switch) and the field $h$ are arbitrary smooth functions.
$\Tcol^{\mathrm{formal}}$ is the space of collision tensors $\sum_\gamma F_\gamma(y,w)\,\partial^\gamma\delta(y-w)$ whose coefficients are polynomials in the jet symbols: $F_\gamma\in\mathbb Q[\sigma\text{-jet}(y)]\otimes\mathbb Q[h\text{-jet}(w),\chi(w)]\otimes\mathbb Q^{10}$ (output components).
$\mathcal G_{\mathrm{formal}}(C)=0$ means vanishing for \emph{all} smooth data.
Since every finite jet at a point is realised by a polynomial function, this is equivalent to vanishing as a polynomial identity in independent jet symbols.
\item $\mathcal G_{\mathrm{background}}$ (the XVI-L specialization): $h=(w^1)^2\eta$, with $\sigma,\chi$ arbitrary; $\Tcol^{\mathrm{background}}$ has coefficients in $\mathbb Q[\sigma\text{-jet}(y)]\otimes\mathbb Q[w^1]\,\chi(w)\otimes\mathbb Q^{10}$.
\item $\mathcal G_{\mathrm{fixed}}$: the actual fixed preparation $\sigeps$ and switch $\chiI$, with $h$ arbitrary; $\Tcol^{\mathrm{fixed}}$ is the space of $h$-linear maps into $\mathcal D'(\mathbb R^4\times\mathbb R^4)$.
\end{enumerate}
Named variants of the formal reading are \emph{time-only}, where $\sigma$ depends on time only and only $\sigma,\partial_0\sigma,\partial_0^2\sigma$ survive, and \emph{plateau}, on $\operatorname{supp}\chiI\cap\{\sigeps=1\}$, where only the $J_3$ sector survives.
\end{definition}

\paragraph{Reduced coefficient expansion.}
Because $\bar a_b=\chiI\,\hat{\bar a}_b$, expand the switch-free factor $\hat{\bar a}_b^{\mu\nu}=\sum_m c_{m,b}\,m$ over coefficient monomials $m=(\mu\nu)\otimes(\sigma\text{-symbol})\otimes(h\text{-symbol})$; the output component is part of the monomial.
(The overall factor $\chiI(w)$ in $\bar a_b$ factors out of the linear algebra on coefficients.)
Then
\[
\sum_{b\in\mathcal B_N}\hat{\bar a}_b\otimes b=\sum_m m\otimes u_m,\qquad u_m:=\sum_bc_{m,b}\,b\in\mathcal H_{4+N},\qquad U_N:=\mathrm{span}\{u_m\}.
\]
\texttt{graph\_contact\_map.txt} lists a basis of each $U_N$ as vectors $u_m$, labelled by their pivot monomials $m$.

\begin{proposition}[Reduction of $\mathcal G$]
\label{prop:Gsplit}
In the formal universal-jet reading of Definition~\ref{def:Gmaps} (with the presentation of Theorem~\ref{thm:M2}):
Then $\mathcal G_{\mathrm{formal}}(C)=0$ if and only if $C_N|_{U_N}=0$ for every $N$.
Hence
\[
\ker\mathcal G_{\mathrm{formal}}=\bigoplus_N\mathrm{Hom}(\mathcal H_{4+N}/U_N,\mathcal D_{\le N}),\qquad
\dim\operatorname{im}\mathcal G_{\mathrm{formal}}=\sum_N\dim U_N\cdot\dim\mathcal D_{\le N}.
\]
\end{proposition}

\begin{proof}
Coherence gives $\mathcal G(C)=\sum_\gamma F_\gamma(y,w)\,\partial^\gamma\delta(y-w)$ with $F_\gamma=\sum_{N,b}C_{N,b,\gamma}\bar a_b=\chiI(w)\sum_{N,b}C_{N,b,\gamma}\hat{\bar a}_b$.
Taylor expansion of $F_{\gamma'}(w+z,w)$ in $z$ gives the unique canonical form $\sum_\gamma g_\gamma(w)\partial^\gamma\delta(y-w)$, with
$g_\gamma=\sum_{\gamma'\ge\gamma}(-1)^{|\gamma'-\gamma|}\binom{\gamma'}{\gamma'-\gamma}\partial_y^{\gamma'-\gamma}F_{\gamma'}|_{y=w}$.
Proceed by descending induction on $|\gamma|$: the top coefficient is $F_\gamma(w,w)$.
In the formal reading the $\sigma$-jet symbols (at $y$) and the $h$-jet and $\chi$ symbols (at $w$) are independent, and evaluation on the diagonal does not identify any two of them.
So $F_\gamma(w,w)=0$ for all data if and only if $\sum_bC_{N,b,\gamma}\hat{\bar a}_b\equiv0$ as a polynomial (the factor $\chiI(w)$ in $\bar a_b$ is independent), and then all its $y$-derivatives vanish too.
Thus $\mathcal G(C)=0$ if and only if $\sum_bC_{N,b,\gamma}\hat{\bar a}_b\equiv0$ for every $\gamma$.
By the reduced expansion this means $\sum_bC_{N,b,\gamma}c_{m,b}=0$ for all $m$, i.e.\ $C_N(\cdot)_\gamma$ annihilates $U_N$.
Monomials of different $N$ are disjoint (preparation order $2-s$, $h$-order $2-N+s$, checked in \texttt{contact\_space.py}), so the conditions split by $N$.
\end{proof}

\begin{remark}[Faithfulness for the specializations]
For $\mathcal G_{\mathrm{background}}$ the coefficient class is $\mathbb Q[\sigma\text{-jet}(y)]\otimes\mathbb Q[w^1]\chi(w)$.
Diagonal evaluation $F\mapsto F(w,w)$ is injective on this class: at a point $w$ with $\chi(w)\ne0$, the $\sigma$-jet at $w$ ranges over all values independently of $w^1$, because $\sigma$ is arbitrary.
The induction above therefore applies with $U_N^{\mathrm{background}}$ computed from the specialized expansion.
The same holds for the time-only class, whose remaining symbols are independent, and for the plateau class, which is a sector restriction.
For $\mathcal G_{\mathrm{fixed}}$ the coefficients are fixed functions (the imported $\sigeps$, $\chiI$).
In general one only has $\ker\mathcal G_{\mathrm{formal}}\subseteq\ker\mathcal G_{\mathrm{fixed}}$.
Equality with the formal kernel is not automatic from formal $D\ge1$ vanishing; it requires a separate injectivity analysis.
Stage~2 records sufficient criteria (Proposition~\ref{prop:fixed-crit}) and a region ledger (Proposition~\ref{prop:fixed-regions}); the time-only graph map being injective does not by itself settle the time-only collapse/Helmholtz problem; that problem is settled in Section~\ref{sec:to-timeonly}, the transfer to the concrete XII-C ramp is Proposition~\ref{prop:to-actual} (explicit hypotheses), and regions where the actual $\chiI$ is non-constant remain Unresolved.
\end{remark}

\begin{theorem}[Ambient freedom versus graph-visible collision freedom]
\label{thm:M3}
These are reported exact ranks, cross-checked mod $p$ (Lemma~\ref{lem:M3comp}):
\begin{center}
\begin{tabular}{lcccc}
\hline
Map / variant & $\dim U_N$ & $\dim\operatorname{im}$ & $\dim\ker$ & on $\Creg$: im / ker\\
\hline
$\mathcal G_{\mathrm{formal}}$ & $(1,4,10,20,35)$ & $3321$ & $0$ & $3150/0$\\
time-only (formal) & $(1,4,10,20,35)$ & $3321$ & $0$ & $3150/0$\\
$\mathcal G_{\mathrm{background}}$ & $(1,4,10,20,10)$ & $1571$ & $1750$ & $1400/1750$\\
time-only + background & $(1,4,10,16,10)$ & $1431$ & $1890$ & $1260/1890$\\
plateau ($J_3$ only) & $(0,0,10,20,35)$ & $3300$ & $21$ & $3150/0$\\
$\mathcal G_{\mathrm{fixed}}$ & --- & --- & --- & Criteria only; Unresolved (Stage~2)\\
\hline
\end{tabular}
\end{center}
Injectivity is claimed only for $\mathcal G_{\mathrm{formal}}$ (and its time-only variant).
\end{theorem}

\begin{lemma}[Reported computation for M.3]
\label{lem:M3comp}
\texttt{contact\_space.py} exits with status $0$.
It records the table of Theorem~\ref{thm:M3}, the disjointness of monomials across $N$, the vectors $u_m$ spanning $U_N$, and certified annihilators: for every annihilator vector $x$ of $U_N$, $\sum_bx_b\hat{\bar a}_b=0$ exactly (\texttt{graph\_contact\_map.txt}).
\end{lemma}

%=====================================================================
\section{M.4 Contact collapse and Helmholtz conditions}
\label{sec:m4}

\begin{definition}[Contact collapse]
\label{def:collapse}
For $C\in\Ccoh$, with the physical reduced coefficient $\bar a_b$ of Definition~\ref{def:Gmaps},
\[
E_C^{\mu\nu}(y)=\sum_{N,b,\gamma}C_{N,b,\gamma}\,\partial_w^\gamma\bigl(\bar a_b^{\mu\nu}(y,w;h)\bigr)\Big|_{w=y}
=\int dw\;\mathcal G(C)^{\mu\nu}(y,w).
\]
There is no second $\chiI$ outside $\bar a_b$: the single switch factor already sits inside $\bar a_b=\chiI\,\hat{\bar a}_b$.
The second equality holds because $\int dw\,F(y,w)(\partial^\gamma\delta)(y-w)=(\partial_w^\gamma F)(y,y)$ and $\mathcal G(C)=\sum C\,\bar a\,\partial^\gamma\delta$.
$\partial_w^\gamma$ distributes by Leibniz over the factors of $\bar a_b=\chiI(w)\,\hat{\bar a}_b(y,w;h)$ (the $h$-jet and $\chiI$); the $\sigeps$-jet sits at $y$ and is not hit (XVI-L, prop:Ljet(2)).
$\Coll:C\mapsto E_C$ maps into local tensors that are linear in the $h$-jet, with coefficients that are monomials in the external jets.
Four computational readings are used, together with the Stage~2 fixed-data region ledger:
\begin{itemize}
\item \emph{formal}: $\sigeps$- and $\chiI$-jets are independent symbols;
\item \emph{$\chi$-plateau}: $\chiI\equiv1$ locally and $\sigeps$ formal;
\item \emph{preparation plateau}: $\sigeps\equiv1$ locally and $\chiI$ formal (variable switch; only $J_3$ survives);
\item \emph{double plateau}: $\chiI\equiv1$ and $\sigeps\equiv1$ locally, so only $J_3$ terms survive and the coefficients are constant.
\end{itemize}
\end{definition}

\begin{remark}[Single switch; preferred presentation]
\label{rem:singleswitch}
The preferred three-equation form keeps the physical coefficient $a_\lambda=\chiI\hat a_\lambda$ (XVI-L style) and the reduced $\bar a_b=\sum_\lambda R_{\lambda b}a_\lambda$, and collapses by $\partial^\gamma\bar a_b$ with no further $\chiI$ outside $\bar a_b$.
Equivalently one may write everything in switch-free hats with a single explicit $\chiI$ in $\mathcal G$ and in $E_C=\partial(\chiI\hat{\bar a})$; both presentations are the same if and only if exactly one overall switch factor is used.
Earlier drafts that wrote $a_\lambda=\chiI\hat a_\lambda$ \emph{and} a second $\chiI$ in the collapse $\partial(\chiI\,\bar a)$ were double-counting in prose.
The recorded computation uses a \emph{single} switch factor: \texttt{graph\_general\_table.txt} stores $\hat a$ with $\chiI$ suppressed, and \texttt{m4\_common.collapse\_column} applies Leibniz once to $\partial^\gamma(\chiI\hat{\bar a})=\partial^\gamma\bar a$.
Certificates A/B/B2 already used one switch factor; the numerical collapse matrices are unchanged by this notation alignment.
\end{remark}

\begin{lemma}[Joint derivative budget]
\label{lem:joint}
If $r_h$ of the $|\gamma|$ collapse derivatives hit the metric coefficient and $r_\chi$ hit the switch, then $(2-v+r_h)+(2-s)+r_\chi\le4-N+|\gamma|\le4$.
Consequently $D:=4-N+|\gamma|$, the total derivative count on $(h,\sigeps,\chiI)$, grades $\Coll$.
\end{lemma}

\begin{proof}
$r_h+r_\chi=|\gamma|\le N$; add $4-s-v=4-N$.
\end{proof}

\begin{proposition}[Collapse rank]
\label{prop:collinj}
\begin{enumerate}
\item In the formal reading, $\Coll$ is injective on $\Ccoh$.
\item In the $\chi$-plateau reading, $\Coll$ is injective (reported computation).
\item In the double-plateau reading, $\dim\ker\Coll=22$ (with $4$ at $D=3$ and $18$ at $D=4$), $\dim\operatorname{im}\Coll=3299$, and $\ker\Coll\cap\Creg=0$ (reported computation).
\end{enumerate}
Injectivity of $\mathcal G$ does not by itself imply injectivity after collapse.
In the plateau reading, $\ker\mathcal G\subseteq\ker\Coll$ grows from $21$ to $22$.
\end{proposition}

\begin{proof}
(1) Factor $\bar a_b=\chiI\,\hat{\bar a}_b$ and expand the Leibniz rule for $\partial_w^\gamma\bar a_b$:
write $E_C=\sum_\rho\chi_\rho F_\rho$ with $\chi_\rho=\partial^\rho\chiI$ and
$F_\rho=\sum_{\gamma\ge\rho}\binom{\gamma}{\rho}\sum_{N,b}C_{N,b,\gamma}\,\partial_h^{\gamma-\rho}\hat{\bar a}_b$, where $\partial_h$ acts on the $h$-jet only (the displayed $F_\rho$ use the switch-free $\hat{\bar a}_b$).
In the formal reading the $\chi_\rho$ are independent, so $E_C=0$ forces $F_\rho=0$ for all $\rho$.
Proceed by descending induction on $k$.
If $C_{\cdot,\cdot,\gamma}=0$ for $|\gamma|>k$, then for $|\rho|=k$ one has $F_\rho=\sum_{N,b}C_{N,b,\rho}\hat{\bar a}_b=0$.
Since $U_N=\mathcal H_{4+N}$ (Theorem~\ref{thm:M3}) and monomials of different $N$ are disjoint, $C_{\cdot,\cdot,\rho}=0$.
(2), (3): Lemma~\ref{lem:M4comp}.
\end{proof}

\begin{definition}[Formal adjoint and self-adjoint subspace]
\label{def:adjoint}
Write $E^{\mu\nu}[h]=\sum_{a\le b,\beta}K^{\mu\nu}_{ab,\beta}\,\partial^\beta h_{ab}$, where $K$ depends on the external jets according to the reading.
Under Gate~V the formal adjoint is
\[
E^{\dagger\,ab}[t]=\frac1{w_{ab}}\sum_{\mu\le\nu,\beta}(-\partial)^\beta\bigl(w_{\mu\nu}K^{\mu\nu}_{ab,\beta}\,t_{\mu\nu}\bigr),
\]
with derivatives hitting the external coefficients retained (Leibniz on $\sigeps$- and $\chiI$-jets, unless these are constant in the reading).
Then $\int\langle t,E[h]\rangle=\int\langle E^\dagger[t],h\rangle$ for compactly supported $t$.
The \emph{self-adjoint subspace} is $\mathrm{SA}:=\{C\in\Ccoh:E_C=E_C^\dagger\}$, as an identity in formal jets.
By Lemma~\ref{lem:joint}, $\mathrm{SA}=\bigoplus_D\mathrm{SA}_D$.
\end{definition}

\begin{lemma}[Helmholtz lemma for linear operators]
\label{lem:helm}
Independent of Input~(S-bv).
Let $E$ be a linear differential operator $h\mapsto E[h]$ with fixed smooth (or formal-jet) coefficients.
Let $\mathcal L_E:=\tfrac12\langle h,E[h]\rangle$ with Gate~V weights, and let $\mathrm{EL}$ denote the Gate~V Euler operator.
Then:
\begin{enumerate}
\item $\mathrm{EL}(\mathcal L_E)=\tfrac12(E+E^\dagger)$;
\item $E$ is the Euler derivative of \emph{some} local Lagrangian if and only if $E=E^\dagger$, and then $E=\mathrm{EL}(\mathcal L_E)$, i.e.\ $E=\delta I_E/\delta h$ with $I_E[h]=\tfrac12\int\langle h,E[h]\rangle$, modulo divergences and field-independent terms.
\end{enumerate}
\end{lemma}

\begin{proof}
(1) For compactly supported $t$: $\frac{d}{d\epsilon}\big|_0\int\mathcal L_E[h+\epsilon t]=\tfrac12\int\langle t,E[h]\rangle+\tfrac12\int\langle h,E[t]\rangle$.
Integrating by parts, with the boundary terms vanishing because $t$ has compact support and derivatives hitting the external coefficients retained, the second term equals $\tfrac12\int\langle t,E^\dagger[h]\rangle$.
Hence $\int\langle t,\mathrm{EL}(\mathcal L_E)\rangle=\int\langle t,\tfrac12(E+E^\dagger)[h]\rangle$ for all $t$.
(2) If $E=E^\dagger$, then (1) gives $E=\mathrm{EL}(\mathcal L_E)$.
Conversely, if $E=\mathrm{EL}(\mathcal L)$ for some local $\mathcal L$, the Fr\'echet derivative $D_E$ is formally self-adjoint (Helmholtz conditions; Olver, Thm.~5.92).
Since $E$ is linear, $D_E=E$, so $E=E^\dagger$.
\end{proof}

Thus $\mathrm{SA}$ is the preimage of the variational tensors, and $\Coll(\mathrm{SA})=\mathcal T_{\mathrm{variational}}\cap\operatorname{im}\Coll$ at this order.

\begin{theorem}[Variational subspace of the contact collapse]
\label{thm:M4}
Assume Gate~V. These are reported exact computations (Lemma~\ref{lem:M4comp}):
\begin{center}
\small
\begin{tabular}{lccccc}
\hline
Reading & $\ker\Coll$ & $\mathrm{SA}$ & $\mathrm{SA}_D$, $D=0..4$ & var.\ image$^\ast$ & $\mathrm{SA}\cap\Creg$\\
\hline
formal & $0$ & $26$ & $(26,0,0,0,0)$ & $26$ & $26$\\
$\chi$-plateau & $0$ & $26$ & $(26,0,0,0,0)$ & $26$ & $26$\\
prep.\ plateau$^\dagger$ & $21$ & $47$ & $(26,0,0,4,17)$ & $26$ $=(26,0,0,0,0)$ & (D=0: $26$)\\
double plateau & $22$ & $1347$ & $(26,0,271,12,1038)$ & $1325$ $=(26,0,271,8,1020)$ & $1208$\\
time-only$^\ddagger$ & $0$ & $26$ & $(26,0,0,0,0)$ & $26$ & $26$\\
time-only, $\chiI\equiv1$$^\ddagger$ & $0$ & $572$ & $(26,0,156,0,390)$ & $572$ & $572$\\
fixed data & \multicolumn{5}{c}{Criteria only; actual dimensions Unresolved (Stage~2; Prop.~\ref{prop:fixed-regions})}\\
\hline
\end{tabular}
\normalsize

$^\ast$var.\ image $=\dim(\mathcal T_{\mathrm{variational}}\cap\operatorname{im}\Coll)=\dim\mathrm{SA}-\dim\ker\Coll$; all entries are dimensions.
$^\dagger$universal reading with formal $\chiI$-jets (exact for that reading; not a statement about the actual fixed $\chiI$).
$^\ddagger$universal time-only readings, $\sigeps=\sigma(t)$ with formal $\sigma^{(k)}$ (SA $=\ker H_t$; Section~\ref{sec:to-timeonly}).
\end{center}
In particular:
\begin{enumerate}
\item[(i)] \textbf{Formal-reading vanishing (not a global obstruction).} Under Gate~V, the formal universal-jet collapse image of this selected graph contains no nonzero variational tensor of total derivative degree $D\ge1$.
Every variational collapsed tensor in that image is an ultralocal multiplication operator $E^{\mu\nu}=\chiI\sigeps\sum_{a\le b}K^{\mu\nu}_{ab}h_{ab}$ with the Gate~V weighted symmetry
$w_{\mu\nu}K^{\mu\nu}_{ab}=w_{ab}K^{ab}_{\mu\nu}$
(a $26$-dimensional family, all in $\Creg$, all from $N=4$, $\gamma=0$).
(Certificate~A checks $w_{00}K^{00}_{02}=-56/3=w_{02}K^{02}_{00}$; the index-reversed form fails.)
The matched functionals are $I_C=\tfrac12\int\chiI\sigeps\langle h,Kh\rangle$.
This is a statement about the selected two-edge graph in the formal reading; it does not claim absence of derivative counterterms in a completed gravitational sector, failure for every fixed preparation, or impossibility after adding other diagrams.
\item[(ii)] The $\chi$-plateau reading has the same $26$-dimensional SA basis.
\item[(iii)] In the double-plateau reading, derivative counterterms do occur: $271$ second-order, $8$ third-order, and $1020$ fourth-order variational directions in the image (variational per-$D$ dimensions $(26,0,271,8,1020)$).
\end{enumerate}
\end{theorem}

\begin{remark}[Eight third-order double-plateau variational directions]
\label{rem:thirdorder}
In the double-plateau reading the SA subspace has dimension $12$ at total derivative degree $D=3$, of which $4$ lie in $\ker\Coll$ (hence contribute $0$ to the variational image) and $8$ survive in $\mathcal T_{\mathrm{variational}}\cap\operatorname{im}\Coll$.
Each of the eight image directions collapses to a pure third-order operator on the residual field $h$ (every monomial in the collapsed tensor has field multi-index of order exactly $3$; no undifferentiated or second-order $h$ terms).
They arise from $N=2$ coherent contact assignments with $|\gamma|=1$ (and linear combinations thereof), not from $N=4$ ultralocal potentials.
\textbf{They are not gravitational scalar counterterms.}
In particular they are not Euler tensors of $\int R$, $\int R^2$, $\int\mathrm{Ric}^2$, or of any other Lorentz-scalar curvature polynomial at this order: those live in the $D=0$ ultralocal block or the $D=2$/$D=4$ matched curvature-squared sector of Proposition~\ref{prop:match}.
The eight directions are variational collision tensors of total derivative degree $3$ inside the selected double-plateau image; their functional interpretation (if any) is deferred to the M.5 ledger and is not claimed here.
Exact sparse SA supports are recorded in \texttt{helmholtz\_certificates.txt} (double\_plateau, $D=3$).
\end{remark}

\begin{remark}[Explicit matched functional; certificates]
\label{rem:cert}
Certificate A, in the formal reading, is the single coordinate $C[N=4,\,b=\{(),0222\},\,\gamma=()]=1$. It gives
\[
E^{00}=-\tfrac{56}3\chiI\sigeps h_{02},\quad E^{02}=-\tfrac{28}3\chiI\sigeps(h_{00}+h_{22}),\quad E^{22}=-\tfrac{56}3\chiI\sigeps h_{02},
\]
with all other components zero, and
\[
I_C=-\tfrac{56}3\int d^4y\;\chiI\sigeps\,h_{02}\,(h_{00}+h_{22}).
\]
The Gate~V Euler derivative of $I_C$ is computed by an independent sympy total-derivative route (\texttt{jet\_euler.py}) and equals $E_C$ exactly.
Certificates B ($D=4$) and B2 ($D=2$), in the double plateau, are $C[N=4,\{(),0222\},\gamma=0002]$ and $C[\ldots,\gamma=02]$.
For example, B gives $E^{00}=-\tfrac{56}3\partial_0^3\partial_2h_{02}$, $E^{02}=-\tfrac{28}3\partial_0^3\partial_2(h_{00}+h_{22})$, $E^{22}=-\tfrac{56}3\partial_0^3\partial_2h_{02}$, and $I_C=-\tfrac{28}3\int(h_{00}\partial_0^3\partial_2h_{02}+\partial_0^3\partial_2h_{00}\,h_{02}+h_{02}\partial_0^3\partial_2h_{22}+\partial_0^3\partial_2h_{02}\,h_{22})$.
For a non-self-adjoint unit vector (negative control C) the Euler derivative of $I_C$ differs from $E_C$.
All certificates are in \texttt{helmholtz\_certificates.txt}.
\end{remark}

\begin{proposition}[Matching probes]
\label{prop:match}
Assume Gate~V.
Membership in a linear space does not depend on normalization, so the target normalizations below only fix the preimage coefficients, which are not Wilson coefficients.
Targets, at quadratic order on the flat background, with $R_{\mathrm{lin}}=\partial^\mu\partial^\nu h_{\mu\nu}-\square h$ and $R^{\mathrm{lin}}_{\mu\nu}=\tfrac12(\partial^\rho\partial_\mu h_{\rho\nu}+\partial^\rho\partial_\nu h_{\rho\mu}-\square h_{\mu\nu}-\partial_\mu\partial_\nu h)$:
\[
P_1=\int w\,h_{\mu\nu}h^{\mu\nu},\ \ P_2=\int w\,h^2,\ \ \mathrm{EH}=\int h^{\mu\nu}G^{\mathrm{lin}}_{\mu\nu},\ \ R2=\int R_{\mathrm{lin}}^2,\ \ RC2=\int R^{\mathrm{lin}}_{\mu\nu}R_{\mathrm{lin}}^{\mu\nu}.
\]
Here the scalar weight is
$w=\chiI\sigeps$ in the formal reading,
$w=\sigeps$ in the $\chi$-plateau reading, and
$w=1$ in the double-plateau reading.
$\int\mathrm{Riem}_{\mathrm{lin}}^2$ is not separate: at quadratic order $\int(\mathrm{Riem}^2-4\mathrm{Ric}^2+R^2)_{\mathrm{lin}}$ is a total derivative.
\begin{enumerate}
\item In all three readings at $D=0$: $\mathrm{span}\{E_{P_1},E_{P_2}\}\cap\operatorname{im}\Coll$ is one-dimensional, spanned by the Euler tensor of $P_1-\tfrac{89}{140}P_2$.
The preimage is the single coordinate $C[N=4,\{33,33\},()]=3/70\in\Creg$.
Neither $P_1$ nor $P_2$ alone, nor the Fierz--Pauli combination $P_1-P_2$, is an image.
\item \textbf{Non-membership (double plateau).} Under Gate~V, the specified quadratic Einstein--Hilbert Euler tensor $E_{\mathrm{EH}}\propto G^{\mathrm{lin}}$ is not in the selected double-plateau collapse image ($D=2$), either in $\Ccoh$ or in $\Creg$.
\item \textbf{Constructed (double plateau).} $E_{R2}$ and $E_{RC2}$ are each collapse images ($D=4$), with explicit exact preimages in $\Ccoh$ ($286$ nonzero coefficients each).
In $\Creg$ only the single combination $R2-\tfrac{188813}{88456}RC2$ is an image.
Matching $R^2$ or $\mathrm{Ric}^2$ individually therefore requires coherent changes on the $\mathcal V_{\mathrm K}$ block, which do not preserve regression.
\item In the formal reading no target with $D\ge1$ is an image of a variational collapse (Theorem~\ref{thm:M4}(i)).
The double-plateau matchings hold locally on the open set where $\sigeps\equiv1$ and $\chiI\equiv1$; they fail once switch and preparation derivatives are retained.
This is a separation of plateau matching from full switched matching.
In the time-only reading the weighted $R^2$ and $\mathrm{Ric}^2$ Euler tensors are not images either, even with $\chiI\equiv1$ (Proposition~\ref{prop:to-match}(3)): the double-plateau matching does not survive the time-only reading.
\end{enumerate}
\end{proposition}

\begin{proof}
For each target $T$: $T$ is self-adjoint (checked).
Membership is decided by the exact verified kernel of the matrix whose columns are the $\Coll$ columns of the $(D,\text{parity }0)$ block, augmented by $T$.
Each preimage is reassembled and compared with $T$ exactly; \texttt{regression\_M.py} re-verifies the $R2$ preimage with an independently coded $R_{\mathrm{lin}}$.
Targets are Lorentz scalars, hence of reflection parity $0$, so only parity-$0$ blocks contribute.
The target construction is checked by the linearized Bianchi identity in jets.
Audit (Proposition~\ref{prop:M7-id}): $E_{R2}$ and $E_{RC2}$ are rebuilt independently from the metric linearization and coincide exactly with the targets; both targets and the certified preimages (after $\Coll$) satisfy $\partial_\mu E^{\mu\nu}=0$.
\end{proof}

\begin{lemma}[Reported computation for M.4]
\label{lem:M4comp}
\texttt{contact\_collapse.py}, \texttt{helmholtz\_check.py} and \texttt{matching\_probe.py} exit with status $0$.
Kernels are computed exactly over $\mathbb Q$ in each block of $(D,\text{reflection parity})$; there are $40$ blocks, of maximal size $292$.
The procedure has three steps: rows independent mod $p$ are selected; the exact nullspace of the selected rows is computed; and every candidate kernel vector is verified exactly against \emph{all} rows.
Since $\ker(\text{all rows})\subseteq\ker(\text{selected})$, verified equality is exact.
The adjoint implementation is unit-tested against Euler derivatives of random quadratic Lagrangians computed by the independent route.
Archives: \texttt{collapse\_kernel\_certificates.txt}, \texttt{helmholtz\_certificates.txt}, \texttt{matching\_certificates.txt}.
The full collapse matrices are regenerated deterministically by the scripts and are not stored; they have about $6.7\times10^5$ (formal) and $1.5\times10^6$ (adjoint-defect) nonzero entries.
\end{lemma}

\paragraph{M.4 status.}
Constructed in the formal, $\chi$-plateau, preparation-plateau, double-plateau and time-only readings (Gate~S-bv certified): the collapse ranks, the SA subspaces, the variational images, the Helmholtz certificates, and the matching probes.
Under Gate~V: the formal collapse image has no variational $D\ge1$ tensors; the preparation-plateau variational image is likewise only the $26$-dimensional ultralocal family; the specified quadratic Einstein--Hilbert Euler tensor is not in the selected double-plateau collapse image.
Fixed-data: criteria, universal region comparisons and the time-only classification in Stage~2 (Propositions~\ref{prop:fixed-crit},~\ref{prop:fixed-regions}, Theorem~\ref{thm:to-timeonly}); actual fixed-data dimensions unresolved where $\chiI$ is non-constant; not assumed settled by formal $D\ge1$ vanishing.
Unresolved: physical normalization ($\kappa$, $\hbar$, $i$); actual fixed-data dimensions where $\chiI$ is non-constant; whether transition-region contributions near plateau edges can be absorbed by insertion or switch counterterms.

%=====================================================================
\subsection{Stage~2: fixed-data and region-wise readings}
\label{sec:m4-fixed}

Gate~V is unchanged.
The imported fixed external data are those of the XVI-L Scope box: the preparation $\sigeps$ is the fixed time-only ramp of XII-C (Section~sec:freeze), never varied and with no $\varepsilon\to0$; the interaction switch $\chiI\in C_c^\infty(M)$ is independent and compact.
The closed form of $\sigeps$ is the XII-C profile recalled in Lemma~\ref{lem:to-indep}, with $\varepsilon$ fixed but not numerically specified; the predecessors give no closed form of $\chiI$, and none is selected inside XVI-M.

\begin{proposition}[Sufficient criteria for fixed data]
\label{prop:fixed-crit}
Assume Gate~V.
Always $\ker\mathcal G_{\mathrm{formal}}\subseteq\ker\mathcal G_{\mathrm{fixed}}$ and $\ker\Coll_{\mathrm{formal}}\subseteq\ker\Coll_{\mathrm{fixed}}$ (specialization).
Let $U$ be an open set meeting the relevant support, and let $\mathrm{ev}_U$ be evaluation of formal external-jet monomials (the $\sigeps$- and $\chiI$-jets occurring) on the actual fixed data over $U$.
Consider the finite-dimensional coefficient spaces spanned by the external-jet coefficients occurring in
\begin{enumerate}
\item[(i)] the canonical graph map $\mathcal G$;
\item[(ii)] the collapse map $\Coll$;
\item[(iii)] the adjoint-defect map $C\mapsto E_C^\dagger-E_C$ (whose coefficients include the external-jet derivatives produced by integration by parts).
\end{enumerate}
If $\mathrm{ev}_U$ is injective on the space in (i), then $\ker\mathcal G_{\mathrm{fixed}}|_U=\ker\mathcal G_{\mathrm{formal}}=0$; if it is injective on the space in (ii), then $\ker\Coll_{\mathrm{fixed}}|_U=\ker\Coll_{\mathrm{formal}}=0$; if it is injective on the space in (iii), then $\mathrm{SA}_{\mathrm{fixed}}|_U=\mathrm{SA}_{\mathrm{formal}}$, the $26$-dimensional ultralocal family.
Condition (iii) is needed for the SA statement: injectivity in (i) or (ii) alone does not control the defect map.
\begin{itemize}
\item \emph{Time-only caveat.} The actual $\sigeps$ is a time-only ramp, so every spatial $\sigeps$-derivative vanishes identically and the full spacetime preparation jets are not independent; conditions (i)--(iii) cannot hold for the formal coefficient spaces.
A time-only comparison class (spatial $\sigeps$-jets set to zero, $\chiI$ formal) needs its own Helmholtz analysis, with its own collapse and adjoint-defect maps; injectivity of the time-only graph map alone is insufficient.
That analysis is carried out in Section~\ref{sec:to-timeonly} (Theorem~\ref{thm:to-timeonly}, Proposition~\ref{prop:to-chiplateau}); its transfer to the actual ramp is Proposition~\ref{prop:to-actual}, which has explicit hypotheses.
\item \emph{Universal versus actual.} Readings with formal $\chiI$-jets (formal; preparation plateau) are universal variable-switch readings: their dimensions are exact for that universal reading.
They coincide with the actual fixed $\chiI$ only under (i)--(iii) for the actual $\chiI$ on $U$, which are not verified here.
\end{itemize}
These are sufficient criteria, not a certification of the concrete XII-C ramp.
Formal universal-jet vanishing at $D\ge1$ does \emph{not} by itself settle any fixed-data case.
\end{proposition}

\begin{proposition}[Region ledger for fixed preparation/switch]
\label{prop:fixed-regions}
Assume Gate~V. Distinguish four regions.
\begin{enumerate}
\item \textbf{Full switched/prepared region.} Comparison classes: formal universal reading, and the time-only class, whose Helmholtz analysis is now exact (Theorem~\ref{thm:to-timeonly}). Status: sufficient criteria (Proposition~\ref{prop:fixed-crit}); on $(0,\varepsilon)\times\mathbb R^3$ with a universal switch the time-only table applies under the explicit hypotheses of Proposition~\ref{prop:to-actual}; actual dimensions where $\chiI$ is non-constant Unresolved.
\item \textbf{Preparation plateau with variable switch, universal reading.} On $\{\sigeps\equiv1\}$ with formal $\chiI$-jets (reading \texttt{prep\_plateau}): support kills $C_2$/$M_2$; $\dim\ker\Coll=21$ exactly; $\mathrm{SA}=(26,0,0,4,17)$ by $D$; the $D\ge1$ SA directions lie in $\ker\Coll$, so the variational image is exactly the $26$-dimensional ultralocal family. Exact for the universal reading; for the actual fixed $\chiI$ it holds under (ii)--(iii) only.
\item \textbf{Double plateau.} On $\{\sigeps\equiv1,\chiI\equiv1\}$: recover Theorem~\ref{thm:M4} exactly ($\dim\ker\Coll=22$, variational image $1325$).
\item \textbf{Transition regions} near plateau boundaries: exact dimensions Unresolved. Only the domain-size bound $0\le\dim\ker\Coll_{\mathrm{fixed}}\le3321$ is asserted. (The previously stated upper bound $22$ was unsupported: rank is only lower semicontinuous, so special points of the fixed data could have larger kernels. No sharper bound is proved.)
\end{enumerate}
Vanishing by support, additional relations among fixed coefficient functions, and genuine Gate~V self-adjointness are kept distinct: no concrete ramp identity enlarging $\ker\mathcal G_{\mathrm{fixed}}$ is certified.
\textbf{Status: fixed-data criteria, universal region comparisons and the time-only classification constructed; actual fixed-data dimensions unresolved where $\chiI$ is non-constant} (status table at the end of Section~\ref{sec:to-timeonly}).
Archives: \texttt{fixed\_data\_reading.py}, \texttt{fixed\_data\_reading.json}, \texttt{prep\_plateau\_sa.json}; \texttt{wpC/}.
\end{proposition}

%---------------------------------------------------------------------
\subsection{Time-only universal Helmholtz classification (work package C)}
\label{sec:to-timeonly}

Everything in this subsection is relative to the current coherent presentation $\Ccoh$ of M.2 ($\dim\Ccoh=3321$, $\dim\Creg=3150$, K block $\mathcal C_{\mathrm K}=\{N\le2\}$ of dimension $171$); since Gate~S-bv is certified (Theorem~\ref{thm:wpA-Sbv}), no (S-bv) condition remains.
Gate~V is unchanged: $\eta=\mathrm{diag}(1,-1,-1,-1)$; $\int t_{\mu\nu}E^{\mu\nu}=\int\sum_{\mu\le\nu}w_{\mu\nu}t_{\mu\nu}E^{\mu\nu}$ with $w=1$ on the diagonal and $w=2$ off it; symmetric lowered $h$; $\sigeps$ and $\chiI$ are fixed under the variation, and their derivatives are retained in every integration by parts.
The collapse is the single-switch collapse of Definition~\ref{def:collapse}, $E_C=\partial_w^\gamma\bar a|_{w=y}$ with $a_\lambda=\chiI\hat a_\lambda$, and there is no second $\chiI$.
This subsection supplies the time-only analysis that Proposition~\ref{prop:fixed-crit} (time-only caveat) had left Unresolved.

\begin{definition}[Time-only formal reading and the specialization $S_t$]
\label{def:to-timeonly}
In the \emph{time-only formal} reading, $\sigeps=\sigma(t)$, $t=x^0$, and the symbols $\sigma,\sigma',\sigma'',\dots$ are independent.
The $\chiI$-jets $\chi_\rho$ are independent spacetime jets, and $h$ is a general spacetime field with independent jets.
The time-only total derivatives are as follows: $D_0^{(t)}$ acts by Leibniz on $h$- and $\chi$-jets and maps $\sigma^{(k)}\mapsto\sigma^{(k+1)}$; $D_i^{(t)}$ ($i=1,2,3$) acts by Leibniz on $h$- and $\chi$-jets and annihilates every $\sigma^{(k)}$.
The specialization $S_t$ from the formal reading sends $\sigma_\alpha\mapsto\sigma^{(\alpha_0)}$ if $\alpha=(\alpha_0,0,0,0)$, and $\sigma_\alpha\mapsto0$ if $\alpha$ has a spatial index; it leaves the $h$- and $\chi$-jets unchanged.
Put $\Coll_t:=S_t\circ\Coll_{\mathrm{formal}}$ and $H_t(C):=S_t(E_C^\dagger-E_C)$.
The supplementary reading \emph{time-only $\chi$-plateau} additionally sets $\chiI\equiv1$, so that all $\chi$-derivatives vanish.
\end{definition}

\begin{lemma}[Intertwining]
\label{lem:to-intertwine}
Independent of Input~(S-bv).
For every polynomial $F$ in the formal jets, $S_t(D_\mu F)=D_\mu^{(t)}S_t(F)$ for $\mu=0,\dots,3$.
Hence $S_t\circ\dagger_{\mathrm{formal}}=\dagger_t\circ S_t$, where $\dagger_t$ is the Gate~V formal adjoint built from $D^{(t)}$, and $H_t(C)=E_{C,t}^{\dagger_t}-E_{C,t}$ with $E_{C,t}=\Coll_t(C)$.
\end{lemma}

\begin{proof}
Both sides are derivations that agree on generators.
On $h$- and $\chi$-jets $S_t$ is the identity.
For $\sigma_\alpha$: $D_\mu\sigma_\alpha=\sigma_{\alpha+e_\mu}$.
If $\mu=0$, then $\alpha+e_0$ is time-only exactly when $\alpha$ is, and both sides equal $\sigma^{(\alpha_0+1)}$ or $0$.
If $\mu=i\ge1$, then $\alpha+e_i$ is spatial, so the left side is $0$, and the right side is $D_i^{(t)}S_t(\sigma_\alpha)=0$.
The adjoint is the composite of multiplications and iterated total derivatives, so it commutes with the ring homomorphism $S_t$.
\end{proof}

\paragraph{Tests (\texttt{wpC/time\_only\_specialization\_tests.txt}).}
The intertwining is tested exactly on $400$ random rational polynomials (seed $20261006$) in all four directions, and for random iterated multi-indices $|\beta|\le4$.
It is tested by an independent sympy route with $\sigma=\sigma(x^0)$, $\chi(x)$ and $h(x)$ as functions, and on $960$ actual collapse and formal adjoint-defect expressions ($120$ contact variables $\times$ $4$ directions).
A negative control, a $D_1^{(t)}$ that differentiates $\sigma$, violates the identity.
\textbf{Route~1 versus route~2.}
Route~1 is $S_t\circ\Coll_{\mathrm{formal}}$, and $S_t$ applied to the formal adjoint defect.
Route~2 applies $S_t$ to the graph table and then builds the collapse by iterated single $\partial_w$ (Leibniz on $\chiI$ and $h$ only) and the adjoint by iterated $D^{(t)}$ on (coefficient $\times$ test field), with no binomial closed forms.
The two routes agree exactly over $\mathbb Q$ on all $3321$ columns of both maps: $446292$ nonzero entries for $\Coll_t$ and $573582$ for $H_t$.
In the $\chi$-plateau variant they also agree, with $91568$ and $130110$ entries.

\begin{theorem}[Time-only classification]
\label{thm:to-timeonly}
Assume Gate~V.
In the time-only formal reading the following dimensions are exact (Lemma~\ref{lem:to-comp}):
\begin{center}
\small
\begin{tabular}{c|c|cc|cc|cc|cc|c}
\hline
$D$ & size & $k_C$ & $i_C$ & $k_H$ & $v$ & $k_C^{\mathrm{reg}}$ & $k_C^{\mathrm K}$ & $k_H^{\mathrm{reg}}$ & $k_H^{\mathrm K}$ & $v^{\mathrm{reg}}$\\
\hline
0 & 35 & 0 & 35 & 26 & 26 & 0 & 0 & 26 & 0 & 26\\
1 & 160 & 0 & 160 & 0 & 0 & 0 & 0 & 0 & 0 & 0\\
2 & 440 & 0 & 440 & 0 & 0 & 0 & 0 & 0 & 0 & 0\\
3 & 944 & 0 & 944 & 0 & 0 & 0 & 0 & 0 & 0 & 0\\
4 & 1742 & 0 & 1742 & 0 & 0 & 0 & 0 & 0 & 0 & 0\\
\hline
$\Sigma$ & 3321 & 0 & 3321 & 26 & 26 & 0 & 0 & 26 & 0 & 26\\
\hline
\end{tabular}
\normalsize
\end{center}
Here $k_C=\dim\ker\Coll_t$, $i_C=\dim\operatorname{im}\Coll_t$, $k_H=\dim\ker H_t$, $v=\dim\Coll_t(\ker H_t)$ (variational image), $k^{\mathrm{reg}}_\bullet$, $k^{\mathrm K}_\bullet$ are the dimensions of the intersections with $\Creg$ and $\mathcal C_{\mathrm K}$, and $v^{\mathrm{reg}}=\dim\Coll_t(\ker H_t\cap\Creg)$ (regression-preserving variational image).
Thus $\Coll_t$ is injective for every $D$, and the time-only variational image is exactly the $26$-dimensional ultralocal family $E^{\mu\nu}=\chiI\sigeps\sum_{a\le b}K^{\mu\nu}_{ab}h_{ab}$ with $w_{\mu\nu}K^{\mu\nu}_{ab}=w_{ab}K^{ab}_{\mu\nu}$.
The family is supported on $N=4$, $\gamma=0$, lies entirely in $\Creg$ (so it preserves regression), and coincides with the formal SA family: the $26$ formal SA vectors lie in $\ker H_t$, and $\dim\ker H_t=26$.
There is no time-only variational direction of total derivative degree $D\ge1$ for a variable (formal) switch.
\end{theorem}

\begin{proposition}[Time-only $\chi$-plateau: new spatial variational directions]
\label{prop:to-chiplateau}
Assume Gate~V.
With $\chiI\equiv1$ and $\sigma=\sigma(t)$ formal, $\Coll_t$ is injective, $\ker H_t$ has dimensions $(26,0,156,0,390)$ by $D$ ($572$ in total), all of it lies in $\Creg$, $\ker H_t\cap\mathcal C_{\mathrm K}=0$, and $\Coll_t(\ker H_t)$ has the same dimensions.
Exactly,
\[
\ker H_t\big|_{D}=
\begin{cases}
\operatorname{span}\Bigl\{\displaystyle\sum_b K_b\,C[N{=}4,b,\gamma]:\
K\in\mathcal K_{\mathrm{SA}},\ \gamma\text{ purely spatial},\ |\gamma|=D\Bigr\},
& D=0,2,4,\\[0.6em]
0,& D=1,3,
\end{cases}
\]
so that for $D=0,2,4$ one has $E_C=\sigma(t)\,K\,\partial^\gamma h$ with $\gamma$ purely spatial; here $26\times6=156$ and $26\times15=390$, and the dimensions $(26,0,156,0,390)$ are unchanged.
Reason: for $E[h]=\sigma(t)K\,\partial_1 h$ one has $E^\dagger=-E$, since $\partial_1\sigma=0$ (odd spatial order is anti-self-adjoint); likewise for every odd-order purely spatial multi-index. Even-order purely spatial derivatives are formally self-adjoint and commute with $\sigma(t)$, so the even-$D$ directions are self-adjoint; they are variational only by virtue of the time-only structure.
The $D=0$ ultralocal family is the formal self-adjoint family already recorded above; the sentence that none is self-adjoint in the spacetime-formal $\chi$-plateau reading refers to the \emph{extra} $D=2,4$ directions (checked), so those are time-only-specific.
All even-$D$ directions remain self-adjoint in the double plateau.
For all $546$ directions with $D=2$ or $D=4$, $\mathrm{EL}(\tfrac12\langle h,E_C h\rangle)=E_C$ is verified by an independent sympy time-only jet route, and for samples also by \texttt{euler\_equations} on functions.
They are spatial-only and not Lorentz covariant, and they are not Euler tensors of $\int\sigma R$, $\int\sigma R^2$ or $\int\sigma\mathrm{Ric}^2$ (Proposition~\ref{prop:to-match}).
\end{proposition}

\begin{proposition}[Time-only matching probes]
\label{prop:to-match}
Assume Gate~V.
Take the weight $w=\chiI(x)\sigeps(t)$, or $w=\sigeps(t)$ in the $\chi$-plateau variant, with all preparation and switch derivatives retained in the Euler tensors.
\begin{enumerate}
\item $D=0$: $\mathrm{span}\{E_{wP_1},E_{wP_2}\}\cap\operatorname{im}\Coll_t$ is one-dimensional, spanned by $P_1-\tfrac{89}{140}P_2$, with preimage $C[N=4,\{33,33\},()]=3/70\in\Creg$ (regression-preserving), in both $\Ccoh$ and $\Creg$.
Neither $P_1$ nor $P_2$ is an image by itself.
The Helmholtz check $\mathrm{EL}(\tfrac12\langle h,E_C h\rangle)=E_C$ is verified by two independent routes (time-only jets; \texttt{euler\_equations} with $\sigma(x^0)$, $\chi(x)$, $h(x)$).
\item $D=2$: the weighted Einstein--Hilbert Euler tensor is not in $\operatorname{im}\Coll_t$, in either $\Ccoh$ or $\Creg$.
\item $D=4$: $\mathrm{span}\{E_{wR^2},E_{w\mathrm{Ric}^2}\}\cap\operatorname{im}\Coll_t=0$, in both $\Ccoh$ and $\Creg$.
\textbf{The double-plateau matchings of Proposition~\ref{prop:match}(3) and Theorem~\ref{thm:M5-plateau-R2} therefore do not survive the time-only reading:} once $\sigeps'(t)$, $\sigeps''(t),\dots$ are retained, the weighted $R^2$ and $\mathrm{Ric}^2$ Euler tensors are not images, even with $\chiI\equiv1$.
\item Divergence identities (upper-index $E$, no metric factor; time-only jets): $\partial_\mu E^{\mu\nu}=0$ for $\int wR_{\mathrm{lin}}^2$ and $\int w\mathrm{Ric}_{\mathrm{lin}}^2$; $\partial_\mu E_{\mathrm{EH}}^{\mu\nu}=(\partial_\mu w)\,G_{\mathrm{lin}}^{\mu\nu}$; nonzero for the ultralocal $P_1$.
\end{enumerate}
Each weighted target is checked to be self-adjoint under $\dagger_t$.
Membership is decided by the exact verified kernel of the $(D,\text{parity }0)$ block augmented by the targets.
\end{proposition}

\begin{lemma}[Linear independence of the actual ramp derivatives]
\label{lem:to-indep}
Independent of Input~(S-bv) and of Gate~V (analytic argument).
The actual $\sigeps$ is explicit: XII-C, Section~sec:freeze (imported by XVI-L, Scope), defines $\sigeps(t)=\rho(t/\varepsilon)$, where
$\rho(s)=0$ for $s\le0$, $\rho(s)=e^{-1/s}/(e^{-1/s}+e^{-1/(1-s)})$ for $0<s<1$, and $\rho(s)=1$ for $s\ge1$; here $0<\varepsilon<T$ is fixed but not numerically specified.
For every open interval $J\subset(0,\varepsilon)$ and every $m$, the functions $\sigeps,\sigeps',\dots,\sigeps^{(m)}$ are linearly independent over $\mathbb R$ on $J$.
The actual $\chiI$ is \emph{not} explicit: XVI-J/K/L give only $\chiI\in C_c^\infty(M)$, an independent compact switch.
\end{lemma}

\begin{proof}
Since $\sigeps^{(k)}(t)=\varepsilon^{-k}\rho^{(k)}(t/\varepsilon)$, it suffices to treat $\rho$ on $J'\subset(0,1)$.
Suppose $\sum_{k\le m}c_k\rho^{(k)}=0$ on $J'$ with $c_m\ne0$.
If $m=0$, then $\rho=0$ on $J'$, which is impossible because $\rho>0$ on $(0,1)$.
Otherwise, $\rho$ is real-analytic on $(0,1)$ (a quotient of analytic functions with positive denominator), so the identity holds on all of $(0,1)$.
Then $\rho|_{(0,1)}$ solves a constant-coefficient linear ODE of order $m$.
Since $\rho$ is smooth on $\mathbb R$ and vanishes for $s\le0$, $\rho^{(j)}(s)\to0$ as $s\to0^+$ for all $j$, so the solution has zero Cauchy data at $0$ and vanishes identically, a contradiction.
\end{proof}
A supplementary floating-point (60-digit) evaluation of the Wronskian of $\rho,\dots,\rho^{(4)}$ at $s=\tfrac12,\tfrac13,\tfrac17$ is nonzero ($\approx8.8\times10^{13}$, $2.6\times10^{15}$, $4.1\times10^{17}$); it is not used in the proof and is re-evaluated independently by \texttt{regression\_M.py}.

\begin{proposition}[Actual preparation versus universal switch; explicit hypotheses]
\label{prop:to-actual}
\emph{Status: a proposition with explicit hypotheses; it is not counted as a Constructed fixed-data result.}
Assume Gate~V and:
\begin{enumerate}
\item[(H1)] $\sigeps$ is the actual XII-C ramp of Lemma~\ref{lem:to-indep}, and $U$ is a nonempty open set with $U\subset(0,\varepsilon)\times\mathbb R^3$;
\item[(H2)] either the switch is treated \emph{universally} on $U$ (the identities $E_C=0$, resp.\ $E_C=E_C^\dagger$, are required for all $\chiI$), or $U\subset\{\chiI=1\}^\circ$ for the actual $\chiI$;
\item[(H3)] on $U$ every coefficient of $\Coll_t(C)$ and of $H_t(C)$ is a finite sum $\sum c_{k,\rho,\beta}\,\sigma^{(k)}(t)\chi_\rho(x)\,\partial^\beta h$ with $k\le4$, linear in the $\sigma$-jet (each graph coefficient carries exactly one $\sigma$-jet factor, Proposition~\ref{prop:M1}(3), and collapse and adjoint preserve this; Lemma~\ref{lem:joint}).
\end{enumerate}
Then, with a universal switch, $\Coll$ and $\mathrm{SA}$ on $U$ are exactly those of Theorem~\ref{thm:to-timeonly}; and if $U\subset\{\chiI=1\}^\circ\cap(0,\varepsilon)\times\mathbb R^3$, with actual $\sigeps$ and actual $\chiI$, they are those of Proposition~\ref{prop:to-chiplateau} (vacuous if that set is empty).
\end{proposition}

\begin{proof}
By (H3), vanishing of a coefficient for all $h$, and (in the universal-switch case) for all $\chiI$, forces $\sum_kc_{k,\rho,\beta}\sigma^{(k)}=0$ on the time projection of $U$, hence $c=0$ by Lemma~\ref{lem:to-indep}.
In the plateau case the $\chi_\rho$ with $\rho\ne0$ vanish identically on $U$, and the same argument applies.
\end{proof}

\begin{lemma}[Reported computation for the time-only classification]
\label{lem:to-comp}
\texttt{wpC/time\_only\_helmholtz.py} exits with status $0$ ($98$ checks; total runtime about $155$\,s on the box).
Kernels are computed exactly over $\mathbb Q$ in each of the $40$ $(D,\text{reflection parity})$ blocks, for $\Ccoh$, $\Creg$ and $\mathcal C_{\mathrm K}$.
The grading is checked: every key has total derivative count $D$ and the block parity, so the blocks have disjoint rows.
Rows independent mod $p=2^{61}-1$ are selected, their exact nullspace is computed, and every kernel vector is verified exactly against \emph{all} rows and re-verified by route~2.
\emph{Completeness:} the rank mod $p$ of the selected integer rows equals $n-\#\ker$, which gives $\operatorname{rank}_{\mathbb Q}\ge n-\#\ker$; the verified independent kernel vectors give the reverse inequality.
Ranks mod $10^9+7$ and $998244353$ of all rows agree in every block; they are supplementary only.
$\ker\Coll_t\subseteq\ker H_t$ is checked, and $\dim\Coll_t(\ker H_t)$ is certified as a mod-$p$ lower bound meeting the upper bound $\dim\ker H_t-\dim\ker\Coll_t$.
Archives: \texttt{time\_only\_collapse\_certificates.json}, \texttt{time\_only\_sa\_certificates.json}, \texttt{time\_only\_matching\_certificates.json}, \texttt{time\_only\_summary.json}, \texttt{time\_only\_specialization\_tests.txt}, \texttt{time\_only\_helmholtz\_out.txt} (all in \texttt{wpC/}); their status strings ``CONDITIONAL on (S-bv)'' predate Theorem~\ref{thm:wpA-Sbv} and are superseded by it.
\end{lemma}

\paragraph{Status of the time-only analysis.}
\begin{itemize}
\item \textbf{Constructed} (Gate~V; on the current coherent presentation):
the intertwining lemma and its tests; exact route~1 / route~2 agreement;
the time-only classification (Theorem~\ref{thm:to-timeonly}): $\Coll_t$ injective, variational image $=$ the $26$ ultralocal directions, all regression-preserving, no $D\ge1$ direction;
the time-only $\chi$-plateau classification $(26,0,156,0,390)$ with its exact spatial-operator identification;
the matching probes and divergence identities (Proposition~\ref{prop:to-match});
the linear independence of $\sigeps^{(k)}$ on open subintervals of $(0,\varepsilon)$ (Lemma~\ref{lem:to-indep}).
\item \textbf{Proposition with explicit hypotheses} (not Constructed): the transfer to the actual XII-C ramp with a universal switch on $(0,\varepsilon)\times\mathbb R^3$, and to open parts of $\{\chiI=1\}^\circ$ in the slab $(0,\varepsilon)\times\mathbb R^3$ (Proposition~\ref{prop:to-actual}, hypotheses (H1)--(H3)).
\item \textbf{Unresolved}:
the actual fixed pair $(\sigeps,\chiI)$ wherever $\chiI$ is non-constant, because $\chiI$ is not given explicitly;
transition neighbourhoods of $\partial\{\chiI=1\}$ and of $t\in\{0,\varepsilon\}$.
\emph{Sufficient test (recorded):} on an open $U$, a full-rank exact evaluation matrix, on $U$, of the external-data coefficient functions $\sigeps^{(k)}\partial^\rho\chiI$ that occur in (i) the graph map, (ii) $\Coll_t$ and (iii) the adjoint-defect map $H_t$.
Full rank for (iii) gives $\mathrm{SA}_{\mathrm{fixed}}|_U=\ker H_t$.
The condition concerns the \emph{joint} family $\{\sigeps^{(k)}(t)\,\partial^\rho\chiI(x)\}$; it is not verified for any concrete $\chiI$, and it does not follow from separate linear independence of the $\sigeps^{(k)}$ (Lemma~\ref{lem:to-indep}) and of the $\partial^\rho\chiI$ (Remark~\ref{rem:to-joint}).
\item \textbf{Non-claims}:
the $546$ spatial $\chi$-plateau directions are not gravitational (Lorentz-scalar) counterterms and are not matched to any curvature functional;
no statement is made after other diagrams are added, and none about a completed gravitational sector;
no $\varepsilon\to0$ limit is taken;
preimage coefficients are not Wilson coefficients;
the double-plateau $R^2/\mathrm{Ric}^2$ matchings (Theorem~\ref{thm:M5-plateau-R2}) are unchanged as double-plateau statements, but they do not admit the stated universal time-only completion (Proposition~\ref{prop:to-match}(3)); the actual-ramp conclusion holds under Proposition~\ref{prop:to-actual}.
\end{itemize}

\begin{remark}[Separate independence does not give joint independence (audit~3)]
\label{rem:to-joint}
Linear independence of $\{\sigeps^{(k)}\}_k$ on an interval and of $\{\partial^\rho\chiI\}_\rho$ on an open set does not imply linear independence of the products $\{\sigeps^{(k)}(t)\,\partial^\rho\chiI(x)\}_{k,\rho}$ on $U$.
Counterexample: if locally in the slab $\chiI(x)=\sigeps(t)$ (a switch that is a function of time on $U$), then
\[
\sigeps'\,\chiI-\sigeps\,\partial_0\chiI=\sigeps'\sigeps-\sigeps\sigeps'=0
\]
on $U$: a nontrivial linear relation between the products $\sigeps^{(1)}\cdot\chiI$ ($k=1$, $\rho=0$) and $\sigeps^{(0)}\cdot\partial_0\chiI$ ($k=0$, $\rho=e_0$), although $\{\sigeps,\sigeps'\}$ is independent (Lemma~\ref{lem:to-indep}) and so is $\{\chiI,\partial_0\chiI\}$.
The sufficient test above is therefore stated for the joint family.
The actual-ramp transfer of Proposition~\ref{prop:to-actual} is unaffected: with a universal switch the jets of $\chiI$ at a point can be prescribed independently of $\sigeps$, and on $\{\chiI=1\}^\circ$ only $\rho=0$ occurs.
\end{remark}

\paragraph{Stage~2 / fixed-data status after work package C.}
\begin{center}
\small
\begin{longtable}{p{0.30\textwidth}p{0.35\textwidth}p{0.26\textwidth}}
\hline
Region / reading & Result & Status\\
\hline
\endhead
Formal universal reading & $\ker\Coll=0$; $\mathrm{SA}=(26,0,0,0,0)$ & Constructed (Theorem~\ref{thm:M4})\\
Time-only universal reading ($\sigma(t)$ formal, $\chiI$ formal) & $\ker\Coll_t=0$ for all $D$; $\ker H_t=(26,0,0,0,0)$; variational image $=26$ ultralocals, regression-preserving & Constructed (Theorem~\ref{thm:to-timeonly})\\
Time-only, $\chiI\equiv1$ & $\ker H_t=(26,0,156,0,390)$; spatial-only, non-Lorentz, variational only by the time-only structure & Constructed (Proposition~\ref{prop:to-chiplateau})\\
Preparation plateau, universal $\chiI$ & $\ker\Coll=21$; $\mathrm{SA}=(26,0,0,4,17)$; variational image $26$ & Constructed for the universal reading (Proposition~\ref{prop:fixed-regions})\\
Double plateau & $\ker\Coll=22$; variational image $1325$ & Constructed (Theorem~\ref{thm:M4})\\
Actual $\sigeps$ on $(0,\varepsilon)\times\mathbb R^3$, universal switch & time-only table exact & Proposition~\ref{prop:to-actual} under (H1)--(H3); not counted as Constructed\\
Actual $\sigeps$ and $\chiI$ on open parts of $\{\chiI=1\}^\circ$ in the slab & time-only $\chi$-plateau table exact (vacuous if empty) & Proposition~\ref{prop:to-actual} under (H1)--(H3)\\
Actual $\chiI$ where non-constant & --- & Unresolved ($\chiI$ not explicit); sufficient test recorded\\
Preparation plateau with actual $\chiI$ & universal numbers under criteria (ii)--(iii) only & Unresolved for the actual $\chiI$\\
Transitions ($\partial\{\chiI=1\}$; $t\in\{0,\varepsilon\}$) & only $0\le\dim\ker\Coll_{\mathrm{fixed}}\le3321$ & Unresolved\\
\hline
\end{longtable}
\end{center}


\section{M.5 Counterterm-functionals ledger}
\label{sec:m5}

All statements in this section use Gate~V; Input~(S-bv) enters through the presentation of M.2 and is certified (Theorem~\ref{thm:wpA-Sbv}).
Four ledgers are maintained separately:
\begin{enumerate}
\item literal local representatives;
\item equality modulo total derivatives;
\item EOM reduction;
\item field-redefinition equivalence.
\end{enumerate}
\textbf{EOM and permitted field redefinitions are not defined in Target~I.}
Consequently ledgers~(3) and~(4) are left empty/Unresolved; work proceeds in~(1)--(2).
(Defining EOM/redefs is a Target~II / M.8 obligation.)

\paragraph{Ultralocal family and Lorentz scalars.}
The $26$-dimensional ultralocal SA family ($N=4$, $\gamma=()$, all in $\Creg$) consists of Gate~V-symmetric multiplication operators
$E^{\mu\nu}=\chiI\sigeps\sum K^{\mu\nu}_{ab}h_{ab}$, with matched functionals $I=\tfrac12\int\chiI\sigeps\langle h,Kh\rangle$.
The SA basis realises $26$ distinct nonzero $K$-patterns (\texttt{counterterm\_functionals.py}).
Inside $\mathrm{span}\{E_{P_1},E_{P_2}\}\cap\operatorname{im}\Coll$ there is a single Lorentz-scalar line $P_1-\tfrac{89}{140}P_2$, with preimage $C[N=4,\{33,33\},()]=3/70$.

\paragraph{$R^2$/$\mathrm{Ric}^2$ preimages and K-block changes.}
In the double plateau, $E_{R2}$ and $E_{RC2}$ are collapse images in $\Ccoh$ ($286$ nonzero coefficients each).
Their supports use $N=2$ slots, hence change the $\mathcal V_{\mathrm K}$ block and are \emph{not} regression-preserving.
In $\Creg$ only the combination $R2-\tfrac{188813}{88456}RC2$ is an image.
At quadratic order $\int(\mathrm{Riem}^2-4\mathrm{Ric}^2+R^2)_{\mathrm{lin}}$ is a total derivative, so $\int\mathrm{Riem}_{\mathrm{lin}}^2$ is not a separate probe (ledger~(2)).

\paragraph{Eight third-order directions.}
The eight double-plateau variational directions at $D=3$ (Remark~\ref{rem:thirdorder}) are recorded in ledger~(1) as Gate~V variational collision tensors.
They are \textbf{not} gravitational scalar counterterms.

\paragraph{Plateau completion gate.}
\begin{proposition}[Completion obstruction inside the selected formal/$\chi$-plateau image]
\label{prop:completion}
Under Gate~V: the formal and $\chi$-plateau variational images have $\mathrm{SA}_D=0$ for $D\ge1$ (Theorem~\ref{thm:M4}).
The double-plateau Euler tensors of $\int R_{\mathrm{lin}}^2$ and $\int\mathrm{Ric}_{\mathrm{lin}}^2$ have total degree $D=4$, hence cannot lie in those images.
Therefore the plateau matchings do \emph{not} complete inside the already-stated selected formal/$\chi$-plateau image.
Any enlarged ansatz (retained switch/preparation derivatives, or additional diagrams) is a \textbf{new domain}, not a completion established by M.4.
\end{proposition}

\begin{theorem}[Plateau curvature-square matching]
\label{thm:M5-plateau-R2}
\textbf{(Named theorem \texttt{PlateauCurvatureSquareMatching}.)}
Assume Gate~V.
\begin{itemize}
\item \textbf{Domain:} double-plateau reading of the selected XVI-L complete-source two-edge graph (open set where $\sigeps\equiv1$ and $\chiI\equiv1$ locally); coherent contact space $\Ccoh$ and regression-preserving subspace $\Creg$.
\item \textbf{Order:} quadratic in the residual field $h$ on the flat background; total derivative degree $D=4$.
\item \textbf{Pairing:} Gate~V.
\item \textbf{Equivalence:} literal local Euler-tensor representatives (ledger~(1)), with the optional quadratic Gauss--Bonnet identification of $\int\mathrm{Riem}_{\mathrm{lin}}^2$ in ledger~(2); EOM and field-redefinition equivalences are not used.
\end{itemize}
Then the Gate~V Euler tensors of $\int R_{\mathrm{lin}}^2$ and of $\int\mathrm{Ric}_{\mathrm{lin}}^2$ each lie in $\Coll(\Ccoh)$, with explicit exact preimages.
In $\Coll(\Creg)$ only the single combination $R2-\tfrac{188813}{88456}RC2$ is an image.
The specified quadratic Einstein--Hilbert Euler tensor is not an image.
These matchings are local to the double-plateau domain and do not complete inside the formal or $\chi$-plateau selected images (Proposition~\ref{prop:completion}).
\emph{Audit status.} The theorem still holds after the M.7 audit: the matching code uses neither the corrected divergence routine nor the equivariant vectors, the coefficient $188813/88456$ is unchanged on rerun, the targets agree with an independent metric-based construction, and the targets and their certified preimages satisfy the linear divergence identity $\partial_\mu E^{\mu\nu}=0$ (Proposition~\ref{prop:M7-id}).
\end{theorem}

\begin{remark}[Scope of \texttt{PlateauCurvatureSquareMatching}: time-only reading and scale tensor]
\label{rem:M5-timeonly}
Theorem~\ref{thm:M5-plateau-R2} is a double-plateau statement and remains true as such.
\textbf{It does not survive the time-only reading:} with weight $w=\chiI\sigeps(t)$, or $w=\sigeps(t)$ when $\chiI\equiv1$, and with $\sigeps'$, $\sigeps''$ retained, $\mathrm{span}\{E_{wR^2},E_{w\mathrm{Ric}^2}\}\cap\operatorname{im}\Coll_t=0$ in $\Ccoh$ and in $\Creg$, and $E_{w\mathrm{EH}}$ is not an image; only the ultralocal line $P_1-\tfrac{89}{140}P_2$ ($D=0$) survives (Proposition~\ref{prop:to-match}).
Moreover the controlled scale tensor of Section~\ref{sec:wpB} is not in the span of the M.5 curvature targets in any reading, because its fixed $D=4$ component already lies outside (Theorem~\ref{thm:wpB-dp}, Corollary~\ref{cor:wpB-MC}); it is recorded as an unmatched entry relative to the curvature ledger, not projected.
\end{remark}

\paragraph{M.5 status.}
Ledgers~(1)--(2) constructed from the M.4 probes; named theorem recorded (double-plateau domain only; not surviving the time-only reading, Remark~\ref{rem:M5-timeonly}); completion-gate obstruction recorded; ledgers~(3)--(4) Unresolved (definitions missing).
Archives: \texttt{counterterm\_functionals.py}, \texttt{matching\_matrix.*}, \texttt{matching\_certificates\_M5.txt}, \texttt{ledger\_*.json}, \texttt{ledger\_tables.txt}.


%=====================================================================
\section{M.6 Dilation identity, residues, and matched scale variation}
\label{sec:m6}

This section imports the XVI-L scale family $\varpi_\mu$, $\mathfrak R_{\mathrm L,\mu}$, $A_\mu=\mu\,\mathrm d/\mathrm d\mu\,\mathfrak R_{\mathrm L,\mu}$, and the scale-variation coefficients of XVI-L,~L.6, without modifying XVI-L.
The term \emph{scaling residue} means the anomaly form $\mathcal A(u)=(E+4+N)\mathfrak R(u)$ (or its contact coefficients), as reserved in XVI-L,~rem:Lressretch.
Scale-variation coefficients are \textbf{not} physical gravitational beta functions (XVI-L,~rem:Lscale; Target~II unfinished).

\paragraph{Dilation identity for the XVI-L prescription.}
\begin{proposition}[Dilation identity for $W^{(\mu)}_\rho$ / zero-moment extensions]
\label{prop:M6-dil}
Let $\varpi_\mu(z)=\varpi((\mu/\mu_0)z)$ and let $W^{(\mu)}_\rho$, $w^{(\rho,\mu)}_\alpha$, $\mathfrak R_{\mathrm L,\mu}$ be as in XVI-L, Definition~def:RLmu.
Write $E=z^\rho\partial_\rho$.
\begin{enumerate}
\item For every $\rho$, $\mu\,\partial_\mu\varpi_\mu=E\varpi_\mu$.
\item For $\rho=0$ (so $N=0$, $w^{(0,\mu)}_0=\varpi_\mu$):
\[
\bigl\langle A_\mu(u),\varphi\bigr\rangle
=
-\varphi(0)\,\bigl\langle u,\,E\varpi_\mu\bigr\rangle,
\qquad
A_\mu(u)=-\bigl\langle u,\,E\varpi_\mu\bigr\rangle\,\delta.
\]
\item For $u\in\mathcal H_4$ homogeneous of degree $-4$ on $\mathbb R^4\setminus\{0\}$, the dilation defect satisfies
\[
\mathcal A(u):=(E+4)\mathfrak R_{\mathrm L,\mu}(u)=A_\mu(u).
\]
\end{enumerate}
\end{proposition}

\begin{proof}
(1) Chain rule: $\partial_\mu\varpi((\mu/\mu_0)z)=(z/\mu_0)\cdot(\nabla\varpi)((\mu/\mu_0)z)$, hence $\mu\partial_\mu\varpi_\mu=E\varpi_\mu$.
(2) For $\rho=0$, $\langle\mathfrak R_{\mathrm L,\mu}(u),\varphi\rangle=\langle u,\varphi-\varphi(0)\varpi_\mu\rangle$; differentiate in $\mu$ and use (1).
(3) $\langle(E+4)\mathfrak R,\varphi\rangle=-\langle\mathfrak R,E\varphi\rangle$.
Since $(E\varphi)(0)=0$, $W_0(E\varphi)=E\varphi$, so $\langle\mathfrak R,E\varphi\rangle=\langle u,E\varphi\rangle$.
Write $\varphi=W_0\varphi+\varphi(0)\varpi_\mu$ and $E\varphi=E(W_0\varphi)+\varphi(0)E\varpi_\mu$.
The continuous extension of $u$ to $\{\psi:\psi(0)=0\}$ obeys $\langle u,E(W_0\varphi)\rangle=0$ (homogeneity of degree $-4$; polar identity on the Euclidean sphere), whence $\langle\mathcal A(u),\varphi\rangle=-\varphi(0)\langle u,E\varpi_\mu\rangle=\langle A_\mu(u),\varphi\rangle$.
\end{proof}

\begin{remark}[General $N$]
\label{rem:M6-N}
For $u\in\mathcal H_{4+N}$ and $\rho=N$, differentiating the zero-moment formula yields the contact expansion
\[
A_\mu(u)=\sum_{|\alpha|\le N}a_\mu(u,\alpha)\,\partial^\alpha\delta,
\qquad
a_\mu(u,\alpha)=(-1)^{|\alpha|+1}\bigl\langle u,\,\mu\partial_\mu w^{(\rho,\mu)}_\alpha\bigr\rangle
\]
(with the pairing understood via $W^{(\mu)}_\rho$).
Equality $A_\mu=(E+4+N)\mathfrak R_{\mathrm L,\mu}$ for $N\ge1$, previously Unresolved here, is now \textbf{resolved}: the $\mu$-dependence of $M^{(\rho,\mu)}$ is accounted for by $\mu\,\mathrm dM^{-1}=DM^{-1}-M^{-1}D$, and $B_\mu=0$ for all $N\le4$ (Lemma~\ref{lem:wpB-dw}, Theorem~\ref{thm:wpB-B0}).
\end{remark}

\paragraph{Explicit $S^2$ residue.}
\begin{theorem}[$S^2$ scaling residue]
\label{thm:M6-S2}
By Theorem~\ref{thm:wpA-Sbv}, $S=c(Q-i0)^{-1}$ on $\mathbb R^4\setminus\{0\}$ with $c=1/(4\pi^2)$ and $Q=\eta(z,z)$.
Let $u=S^2=c^2(Q-i0)^{-2}\in\mathcal H_4$.
Then, for every $\mu>0$ and every admissible $\varpi$ with $\varpi(0)=1$,
\[
A_\mu(S^2)=r\,\delta,
\qquad
r=-2i\pi^2 c^2=-\frac{i}{8\pi^2}\neq0,
\]
and $\mathcal A(S^2)=(E+4)\mathfrak R_{\mathrm L,\mu}(S^2)=r\delta$.
The value of $r$ is independent of $\mu$ and of the radial profile of $\varpi$.
\end{theorem}

\begin{proof}
By Proposition~\ref{prop:M6-dil}, $r=-\langle S^2,E\varpi_\mu\rangle$.
In Euclidean polar coordinates $z=r\omega$ on the unit sphere $S^3$, the obstruction identity for a homogeneous density of degree $-4$ reduces $r$ to the spherical integral of the angular part of $S^2$.
By Theorem~\ref{thm:wpA-Sbv} that angular part is $c^2(Q-i0)^{-2}$.
Parametrizing $S^3$ by $\omega_0=\cos\chi$ and spatial radius $\sin\chi$, one has $Q=\cos(2\chi)$ and
\[
\int_{S^3}(Q-i0)^{-2}\,\mathrm d\sigma
=
4\pi\int_0^\pi\frac{\sin^2\chi}{(\cos(2\chi)-i0)^2}\,\mathrm d\chi
=
\pi\int_0^{2\pi}\frac{1-\cos u}{(\cos u-i0)^2}\,\mathrm d u.
\]
The $u$-integral equals $-2i\pi$ by the residue at the interior double pole $z_m=i(e-\sqrt{e^2+1})$ of the unit-circle integrand $2i(z-1)^2/(z^2-2iez+1)^2$ after $e\downarrow0$ (certificate: \texttt{scaling\_anomaly.py}).
Hence $\int_{S^3}(Q-i0)^{-2}\,\mathrm d\sigma=-2i\pi^2$ and $r=c^2(-2i\pi^2)$.
Nonvanishing uses $c\neq0$; with $c=1/(4\pi^2)$, $r=-2i\pi^2/(16\pi^4)=-i/(8\pi^2)$.
The side $Q-i0$ fixes the sign of the angular integral; the counterfactual $Q+i0$ would conjugate it and flip $r$ (Remark~\ref{rem:wpA-M6}).
Independence of the radial profile of $\varpi$ is the polar telescope $\int_0^\infty\partial_r\varpi_\mu(r\omega)\,\mathrm d r=-\varpi_\mu(0)=-1$.
\end{proof}

\begin{remark}[Status of the residue]
\label{rem:M6-Sbv}
Theorem~\ref{thm:M6-S2} uses Input~(S-bv), which is certified by Theorem~\ref{thm:wpA-Sbv}; the logic does not run the other way (the residue does not certify the gate), and matching played no role in the certification (Remark~\ref{rem:Sbv-status}).
The explicit value $r=-i/(8\pi^2)$ fills, as a result about XVI-L objects, the Unresolved stretch of XVI-L,~rem:Lressretch, for the $S^2$ kernel; XVI-L itself is not edited.
Higher graded residues ($N\ge1$) are computed in Section~\ref{sec:wpB}: exact and $\varpi$-independent at top order, $\varpi$-dependent at lower orders.
\end{remark}

\paragraph{The $N=0$ anomaly component in the coherent basis.}
\begin{definition}[$N=0$ anomaly component]
\label{def:M6-Canom}
The \emph{$N=0$ anomaly component} is the coherent contact vector $C_{\mathrm{anom}}\in\Ccoh$ with
\[
C_{\mathrm{anom}}\bigl[N=0,\,\{(),()\},\,\gamma=()\bigr]=r=-2i\pi^2 c^2=-\frac{i}{8\pi^2}
\]
($\dim\mathcal H_4=1$), and with all higher-grade coordinates ($N\ge1$) \emph{explicitly set to zero for this probe}.
Coherence of the $N=0$ coordinate follows from linearity of $A_\mu$ and uniqueness of the $N=0$ basis element.
$C_{\mathrm{anom}}$ is not the full graph scale variation; the full $C_{\mathrm{scale}}$, including all higher-grade coordinates, is treated in Section~\ref{sec:wpB}.
\end{definition}

\paragraph{Collapse and variational compatibility with the M.5 ledger.}
\begin{theorem}[Collapse of $C_{\mathrm{anom}}$; unmatched remainder retained]
\label{thm:M6-collapse}
Assume Gate~V. Let $C_{\mathrm{anom}}$ be as in Definition~\ref{def:M6-Canom} (unit $r=1$ for the linear probe; overall factor $r$ restores the physical coefficient).
\begin{enumerate}
\item \textbf{Formal / $\chi$-plateau.} $\Coll(C_{\mathrm{anom}})$ is nonzero (744 nonzero keys) and \emph{not} self-adjoint ($\mathrm{SA}$-defect nonempty). It is therefore \textbf{not} a nonzero variational tensor.
\item \textbf{Preparation / double plateau.} $\Coll(C_{\mathrm{anom}})=0$.
Structural reason: the $S^2$ coefficient $\hat a$ is supported only on $|\partial\sigeps|=2$ ($C_2$ sector); both plateaus set $\sigeps\equiv1$ and kill those derivatives.
\item \textbf{Matched projection (double plateau).} The orthogonal projection of $\Coll(C_{\mathrm{anom}})$ onto $\mathrm{span}\{E_{R2},E_{RC2}\}$ (the M.5 matched Euler tensors) is $(0,0)$, since the vector itself is $0$.
\item \textbf{Unmatched remainder.} In the double plateau the remainder is $0$.
In the formal reading the \emph{entire} non-self-adjoint $\Coll(C_{\mathrm{anom}})$ is retained as unmatched remainder.
No unmatched component is discarded and renamed ``the anomaly.''
\end{enumerate}
\end{theorem}

\begin{proof}
Exact collapse columns and SA defects: \texttt{scaling\_anomaly.py} (reads the same $\bar a$ and Collapse as M.4).
Linear independence of $E_{R2},E_{RC2}$ and vanishing of $\Coll(C_{\mathrm{anom}})$ in the double plateau give the projection claim.
\end{proof}

\begin{corollary}[Scope of the $N=0$ component relative to the M.5 ledger]
\label{cor:M6-compat}
The $N=0$ anomaly component does \emph{not} furnish a nonzero Gate~V-variational matched scale variation inside the double-plateau domain of Theorem~\ref{thm:M5-plateau-R2}: the $N=0$ component alone contributes no nonzero plateau matching.
In the formal/$\chi$-plateau readings it collapses to a non-variational tensor.
The full graph scale variation, including all higher-grade residues, is treated in Section~\ref{sec:wpB}; it leads to the graph-specific obstruction of M-C$_{\mathrm{curv}}$, while broad M-C is not achieved (Corollary~\ref{cor:wpB-MC}).
\end{corollary}

%---------------------------------------------------------------------
\subsection{Scale derivative versus dilation defect at every grade (work package B)}
\label{sec:wpB}

\scopebox{Status (work package B).}{%
\textbf{Constructed} (with the certified $c=1/(4\pi^2)$ of Theorem~\ref{thm:wpA-Sbv}; Gate~V for (iv)):
(i) the complete $\mu$-derivative of the XVI-L subtraction formula, including $\mu\,\mathrm d M^{-1}=-M^{-1}(\mu\,\mathrm dM)M^{-1}=DM^{-1}-M^{-1}D$, at every order $\rho=N\le4$;
(ii) the identity $A_\mu=\mathcal A_\mu$, i.e.\ $B_\mu=0$, on every $b\in\mathcal B_N$, $N=0,\dots,4$, for every admissible $\varpi$;
(iii) the exact, $\mu$- and $\varpi$-independent top-order ($|\gamma|=N$) coefficients of $A_\mu(b)=\mathcal A_\mu(b)$, consistent with all $185$ relations and all labels;
(iv) the collapse, full adjoint defect, variationality and ledger membership of the assembled scale tensor in the formal, $\chi$-plateau, preparation-plateau and double-plateau readings, and (Corollary~\ref{cor:wpB-timeonly}) its non-variationality in the time-only readings.
In the double plateau, the cutoff-independent $D=4$ component is variational and Lorentz covariant, but not linearized-gauge invariant. The full scale tensor has cutoff-dependent lower-degree components and need not be variational.
\textbf{Proved obstructed:} M-C$_{\mathrm{curv}}$ (Definition~\ref{def:MCcurv}), graph-specific (Corollary~\ref{cor:wpB-MC}). \textbf{Not achieved (Unresolved):} broad M-C ($\varpi$ unspecified by XVI-L, thm:L6). \textbf{Satisfied:} M-C$_{\mathrm{data}}$.
\textbf{Unresolved:} the numerical values of the lower-order ($|\gamma|<N$) coefficients; they depend on the cutoff $\varpi$ through the functionals $\nu_k[\varpi]$, $k=1,\dots,4$, of \eqref{eq:wpB-nu}, and XVI-L does not fix $\varpi$ (the radial comparison is a separately named specialization).
\textbf{Non-claims:} no physical beta function; no claim about gravity generally; no orthogonal projection of an unmatched tensor; $B_\mu=0$ does not remove prescription dependence.}

\paragraph{Prescription (imported unchanged).}
XVI-L, Definition~def:RLmu: $\varpi_\mu(z)=\varpi(\lambda z)$, $\lambda=\mu/\mu_0$, $M^{(\rho,\mu)}_{\beta\alpha}=\partial^\beta(z^\alpha\varpi_\mu/\alpha!)(0)$,
$w^{(\rho,\mu)}_\alpha=\sum_\gamma((M^{(\rho,\mu)})^{-1})_{\gamma\alpha}z^\gamma\varpi_\mu/\gamma!$, $\langle\mathfrak R_{\mathrm L,\mu}(u),\varphi\rangle=\langle u,W^{(\mu)}_\rho\varphi\rangle$, $\rho=N$.
$A_\mu(u)=\mu\,\mathrm d_\mu\mathfrak R_{\mathrm L,\mu}(u)$; $\mathcal A_\mu(u)=(E+4+N)\mathfrak R_{\mathrm L,\mu}(u)$ with $E=z^\rho\partial_\rho$ and $\langle ET,\varphi\rangle=-\langle T,(E+4)\varphi\rangle$, so $\langle\mathcal A_\mu(u),\varphi\rangle=\langle\mathfrak R_{\mathrm L,\mu}(u),(N-E)\varphi\rangle$; $B_\mu=A_\mu-\mathcal A_\mu$.
Contacts: $T=\sum_\gamma t_\gamma\partial_\gamma\delta$, $\langle\partial_\gamma\delta,\varphi\rangle=(-1)^{|\gamma|}\partial^\gamma\varphi(0)$, $t_\gamma=(-1)^{|\gamma|}\langle T,w_\gamma\rangle$; these are the coordinates $C_{N,b,\gamma}$ of $\Ccoh$.

\begin{lemma}[Complete scale derivative of the subtraction functions]
\label{lem:wpB-dw}
For every $\rho\le4$, every $\varpi$ with $\varpi(0)=1$, and $D=\mathrm{diag}(|\beta|)$:
\[
\mu\,\partial_\mu M^{(\rho,\mu)}=[D,M^{(\rho,\mu)}],\qquad
\mu\,\partial_\mu (M^{(\rho,\mu)})^{-1}=-M^{-1}(\mu\,\partial_\mu M)M^{-1}=D M^{-1}-M^{-1}D,
\]
\[
\mu\,\partial_\mu w^{(\rho,\mu)}_\alpha=(E-|\alpha|)\,w^{(\rho,\mu)}_\alpha,\qquad
w^{(\rho,\mu)}_\alpha=\frac{z^\alpha}{\alpha!}\,T_{\rho-|\alpha|}\!\bigl[1/\varpi_\mu\bigr]\,\varpi_\mu ,
\]
with $T_k$ the Taylor polynomial of order $k$ at $0$. Equivalently $w^{(\rho,\mu)}_\alpha(z)=\lambda^{-|\alpha|}w^{(\rho)}_\alpha(\lambda z)$.
\end{lemma}
\begin{proof}
$M_{\beta\alpha}=\binom\beta\alpha\lambda^{|\beta|-|\alpha|}\partial^{\beta-\alpha}\varpi(0)$ (Leibniz), so $\mu\partial_\mu M=[D,M]$; differentiate $MM^{-1}=1$.
For the basis functions $f_\gamma:=z^\gamma\varpi_\mu/\gamma!$, $\mu\partial_\mu\varpi_\mu=E\varpi_\mu$ (Proposition~\ref{prop:M6-dil}(1)) and $Ez^\gamma=|\gamma|z^\gamma$ give $\mu\partial_\mu f_\gamma=(z^\gamma/\gamma!)E\varpi_\mu=(E-|\gamma|)f_\gamma$.
With $\mu\partial_\mu(M^{-1})_{\gamma\alpha}=(|\gamma|-|\alpha|)(M^{-1})_{\gamma\alpha}$ (the entries of $DM^{-1}-M^{-1}D$),
\[
\mu\partial_\mu w_\alpha=\sum_\gamma(M^{-1})_{\gamma\alpha}\bigl[(|\gamma|-|\alpha|)+(E-|\gamma|)\bigr]f_\gamma=\sum_\gamma(M^{-1})_{\gamma\alpha}(E-|\alpha|)f_\gamma=(E-|\alpha|)w_\alpha .
\]
The closed form holds because $z^\alpha T_k[1/\varpi_\mu]\varpi_\mu/\alpha!=z^\alpha/\alpha!+O(|z|^{\rho+1})$ and $M$ is invertible.
Exact certificates (\texttt{wpB/higher\_grade\_scaling.py}, Section~A): symbolic jets for $\rho\le3$, two random rational jets for $\rho=3,4$, and an explicit non-radial $\varpi=e^{q}$ as functions, all $\rho\le4$; \texttt{regression\_M.py} re-checks the matrix identities independently for $\rho\le2$ (symbolic) and $\rho=3$ (rational jets).
\end{proof}

\begin{theorem}[Scale derivative equals dilation defect at every grade]
\label{thm:wpB-B0}
For every $N\in\{0,\dots,4\}$, $b\in\mathcal B_N$, $\mu>0$ and admissible $\varpi$,
\[
A_\mu(b)=\mathcal A_\mu(b)=\sum_{|\gamma|\le N}a_{N,b,\gamma}(\mu)\,\partial_\gamma\delta,\qquad B_\mu(b)=0,
\qquad
a_{N,b,\gamma}(\mu)=(-1)^{|\gamma|+1}\bigl\langle b,(E-|\gamma|)w^{(N,\mu)}_\gamma\bigr\rangle .
\]
\begin{enumerate}
\item \emph{Top order} ($|\gamma|=N$): $a_{N,b,\gamma}=\frac{(-1)^N}{\gamma!}\int_{S^3}b\,\omega^\gamma\,\mathrm d\sigma$, independent of $\mu$ and $\varpi$; it is the intrinsic residue of $b$ (common to all extensions of degree $4+N$).
\item \emph{Lower orders} ($k=N-|\gamma|\ge1$):
\begin{equation}
a_{N,b,\gamma}(\mu)=\frac{(-1)^{|\gamma|+1}}{\gamma!}\Bigl(\frac\mu{\mu_0}\Bigr)^{k}\bigl\langle b|_{S^3}\,\omega^\gamma,\nu_k[\varpi]\bigr\rangle_{S^3},
\qquad
\nu_k[\varpi](\omega)=k\!\int_0^\infty\! r^{-k-1}\bigl(T_k[1/\varpi]\varpi\,(r\omega)-1\bigr)\mathrm dr .
\label{eq:wpB-nu}
\end{equation}
\end{enumerate}
The identity $B_\mu=0$ is pure prescription calculus; the explicit top-order values below use Theorem~\ref{thm:wpA-Sbv}.
\end{theorem}
\begin{proof}
$A_\mu$: differentiate $\langle b,W^{(\mu)}_N\varphi\rangle$; $\mu\partial_\mu w_\gamma$ vanishes to order $N$ at $0$ (since $\partial^\beta w_\gamma(0)=\delta^\beta_\gamma$ for all $\mu$), so the pairing is the unique extension; Lemma~\ref{lem:wpB-dw}.
$\mathcal A_\mu$: by homogeneity $\mathcal A_\mu(b)$ is supported at $0$ of order $\le N$; its coefficients are $(-1)^{|\gamma|}\langle\mathfrak R_{\mathrm L,\mu}b,(N-E)w_\gamma\rangle=(-1)^{|\gamma|}\langle\mathfrak R_{\mathrm L,\mu}b,(N-|\gamma|)w_\gamma-(E-|\gamma|)w_\gamma\rangle$; the first term is a zero moment, the second vanishes to order $N$. Hence the coefficients agree; $B_\mu=0$.
Polar coordinates, $r^{-|\gamma|}w_\gamma(r\omega)=\frac{\omega^\gamma}{\gamma!}T_k[1/\varpi_\mu]\varpi_\mu(r\omega)$ and one integration by parts give (1)--(2); for $k=0$ the radial integral telescopes to $-\varpi(0)=-1$.
Certificates: top order by two independent exact routes (contour residues; Hadamard finite part plus $i\pi\delta^{(m-1)}$ in $x=Q|_{S^3}$), equal on every slot; finite exact end-to-end tests with an explicit compactly supported non-radial $\varpi$ for $N=0,1,2$ (Euclidean toy kernels, Section~B of the script).
\end{proof}

\paragraph{Exact top-order data.}
All top-order values are $i\pi^2c^2\times\mathbb Q$, and with the certified $c$ the unit is $i\pi^2c^2=i/(16\pi^2)$ (the real, finite-part contributions vanish identically).
$N=0$: $a=-2i\pi^2c^2=-i/(8\pi^2)$, the M.6 residue $r$ of Theorem~\ref{thm:M6-S2}. $N=1$: $a_{1,\{(),\mu\},\mu}=-i\pi^2c^2=-i/(16\pi^2)$, others $0$.
Nonzero top slots per grade: $1,4,22,68,194$ of $1,16,100,400,1225$ ($289$ in total).
In the audited equivariant coordinates of M.7 ($C_{N,b,\gamma}=(-1)^{\gamma_1+\gamma_2+\gamma_3}[Q^k\square^k n_b]_\gamma$), the residue map is, in units $i\pi^2c^2=i/(16\pi^2)$,
$-2\phi_{00}$; $\tfrac12\phi_{00}$; $-\tfrac1{12}\phi_{00}-\tfrac1{48}\phi_{11}$; $\tfrac1{96}\phi_{00}+\tfrac1{384}\phi_{11}$; $-\tfrac1{960}\phi_{00}-\tfrac1{3840}\phi_{11}-\tfrac1{61440}\phi_{22}$ for $N=0,\dots,4$.
Tests where the hypotheses hold: Lorentz equivariance on all six generators (pull-back action, intertwiner $z^\mu\leftrightarrow\eta^{\mu\nu}\partial_\nu\delta$); $\mathrm{Res}(\partial_\alpha u)=\partial_\alpha\mathrm{Res}(u)$ for all $340$ label/direction pairs with $N\le3$; every one of the $185$ relation vectors annihilates the top-order, the radial-specialization and the general-$\varpi$ coefficient vectors; all $255$ labels reconstructed and checked against direct evaluation.
The derivative and Lorentz identities are \emph{not} used for, and do not hold for, the lower-order coefficients ($\varpi$ is neither covariant nor Leibniz-compatible).

\paragraph{Prescription dependence.}
$B_\mu=0$ does not remove prescription dependence: it resides in the lower-order coefficients, common to $A_\mu$ and $\mathcal A_\mu$, which scale as $(\mu/\mu_0)^{N-|\gamma|}$ and depend on $\nu_k[\varpi]$, $k=1,\dots,4$.
\emph{Missing datum:} XVI-L fixes $\varpi\in C_c^\infty$, $\varpi(0)=1$, but does not specify it; the values of $\nu_k[\varpi]$ (equivalently the functionals $\Lambda_k(P)=\langle P(Q-i0)^{-6}|_{S^3},\nu_k\rangle$ on $\mathrm{Poly}_{8-k}$) are what is missing.
For a reflection-even $\varpi$ the odd-$k$ coefficients vanish.
The named specialization $\mathfrak R_{\mathrm L,\mu}[\varpi_{\mathrm{rad}}]$ (Euclidean-radial $\varpi$; \emph{not} asserted by XVI-L) has constant $\nu_k^{\mathrm{rad}}$, only $k=2,4$ survive, and its contact difference from $\mathfrak R_{\mathrm L,\mu}[\varpi]$ has coefficients $-(-1)^{|\gamma|}\langle b,w_\gamma[\varpi]-w_\gamma[\varpi_{\mathrm{rad}}]\rangle$, i.e.\ $\frac{(-1)^{|\gamma|+1}}{k\,\gamma!}\langle b\,\omega^\gamma,\nu_k[\varpi]-\nu_k^{\mathrm{rad}}\rangle$ for $k\ge1$ and $\frac{(-1)^{N+1}}{\gamma!}\langle b\,\omega^\gamma,\int_0^\infty r^{-1}(\varpi-\varpi_{\mathrm{rad}})(r\omega)\,\mathrm dr\rangle$ for $k=0$.

\paragraph{Collapse of the full scale tensor.}
$C_{\mathrm{scale}}(\mu)=(a_{N,b,\gamma}(\mu))$ for all $N,b,\gamma$.
Collapse and adjoint defect are homogeneous in the total derivative degree $D=4-N+|\gamma|$, so the top part ($D=4$) and the lower parts ($D=4-k$) cannot cancel each other; cross-grade cancellations at fixed $D$ are included.
\begin{center}
\small
\begin{tabular}{lcccc}
\hline
reading & $\Coll$ nnz ($D=4$) & adjoint-defect nnz & variational & per-$N$ nnz ($N=0..4$)\\
\hline
formal & 14112 & 16608 & no & 744, 4440, 10278, 10040, 4604\\
$\chi$-plateau & 2522 & 4090 & no & 744, 1728, 2386, 1550, 514\\
preparation plateau & 4580 & 4090 & no & 0, 0, 2274, 4072, 4604\\
double plateau & 490 & 0 & \textbf{yes} & 0, 0, 474, 514, 514\\
\hline
\end{tabular}
\normalsize
\end{center}
Lower-order envelope (all $\varpi$): $\Coll$ is injective on each $\Lambda_k$-block; self-adjoint parameter subspaces have dimensions $0,0,0,26$ ($k=1,2,3,4$) in the formal, $\chi$- and preparation readings and $0,49,0,26$ in the double plateau; there $E_{\mathrm{EH}}\notin$ image ($k=2$) and $\mathrm{span}(E_{P_1},E_{P_2})\cap$ image is one-dimensional ($k=4$).
The radial specialization is non-self-adjoint at $k=2$ and $k=4$ in every reading.
Time-only readings: Corollary~\ref{cor:wpB-timeonly}.

\begin{theorem}[Double-plateau scale tensor: top component and full tensor]
\label{thm:wpB-dp}
Assume Gate~V. In the double plateau, the cutoff-independent $D=4$ component is variational and Lorentz covariant, but not linearized-gauge invariant. The full scale tensor has cutoff-dependent lower-degree components and need not be variational. In detail:
\begin{enumerate}
\item[(i)] \emph{Top component.} The $D=4$ component $\Coll(C_{\mathrm{top}})$ of $\Coll(C_{\mathrm{scale}})$ (the slots with $|\gamma|=N$; independent of $\mu$ and $\varpi$) is nonzero and variational, with
\[
I^{\mathrm{top}}_{\mathrm{scale}}=\frac{i}{16\pi^2}\!\int\Bigl(-\tfrac{71}{120}h_{\mu\nu}\square^2h^{\mu\nu}-\tfrac{27}{80}h\square^2h-\tfrac{11}{10}(\partial^\mu h_{\mu\nu})\square(\partial_\rho h^{\rho\nu})+\tfrac{19}{30}h\square\partial^\mu\partial^\nu h_{\mu\nu}-\tfrac{13}{15}(\partial^\mu\partial^\nu h_{\mu\nu})^2\Bigr)
\]
(mod total derivatives; the prefactor is $i\pi^2c^2$ with $c=1/(4\pi^2)$).
It is Lorentz covariant (a combination of Lorentz scalars), but $\partial_\mu E^{\mu\nu}\neq0$ (not linearized-gauge invariant), and it is \emph{not} in $\mathrm{span}(E_{R^2},E_{\mathrm{Ric}^2})$: in the basis above $R_{\mathrm{lin}}^2\mapsto(0,1,0,-2,1)$ and $\mathrm{Ric}_{\mathrm{lin}}^2\mapsto(\frac14,\frac14,\frac12,-\frac12,\frac12)$, the rank goes from $2$ to $3$ when the scale vector is adjoined, and Proposition~\ref{prop:wpB-sep} gives an exact separating functional.
In particular it is not a multiple of $E_{R^2-(188813/88456)\mathrm{Ric}^2}$, of $E_{\mathrm{EH}}$, or of an ultralocal.
\item[(ii)] \emph{Full tensor.} The lower-degree components ($D=4-k$, $k=1,\dots,4$) depend on the cutoff through $\Lambda_k$, i.e.\ through $\nu_k[\varpi]$ of~\eqref{eq:wpB-nu}.
The full tensor is variational iff $\Lambda_1=\Lambda_3=0$, $\Lambda_2$ lies in the $49$-dimensional and $\Lambda_4$ in the $26$-dimensional self-adjoint subspaces; so it need not be variational.
Lorentz covariance of the full tensor is \emph{not} established: the lower-order terms depend on the possibly noncovariant $\varpi$, and the Lorentz tests apply to the top order only.
\item[(iii)] \emph{Curvature span.} For \emph{every} $\mu$ and every admissible $\varpi$ the full collapsed scale tensor lies outside the span of the M.5 curvature targets: collapse is homogeneous in $D$, and its fixed $D=4$ component already lies outside $\mathrm{span}(E_{R^2},E_{\mathrm{Ric}^2})$, while the targets of degree $D=0,2$ lie in other blocks.
\end{enumerate}
The whole tensor is retained as unmatched relative to the curvature ledger.
\end{theorem}
\begin{proof}
(i) Exact double-plateau collapse and adjoint defect of $C_{\mathrm{top}}$ (\texttt{wpB/higher\_grade\_scaling.py}, Section~D; \texttt{full\_scale\_collapse\_certificates.json}, \texttt{readings.double\_plateau.top\_D4}): adjoint defect $0$, so Lemma~\ref{lem:helm} gives the potential; its coordinates in the five-term Lorentz-scalar basis are the displayed coefficients; the divergence is nonzero exactly; the rank statement is exact over $\mathbb Q$.
\texttt{regression\_M.py} independently recollapses $C_{\mathrm{top}}$ from \texttt{\_wpB\_state.json} and checks the pure-gauge coefficient of Proposition~\ref{prop:wpB-sep}.
(ii) Theorem~\ref{thm:wpB-B0}(2) and the lower-order envelope above; the Lorentz and derivative identities hold for the top order only (paragraph ``Exact top-order data'').
(iii) Homogeneity in $D$ (Lemma~\ref{lem:to-comp} and the paragraph above the table) and (i).
\end{proof}

\begin{proposition}[Exact separating functional for the displayed scale functional (audit~3)]
\label{prop:wpB-sep}
In the ordered basis $(L_1,\dots,L_5)=\bigl(h_{\mu\nu}\square^2h^{\mu\nu},\ h\square^2h,\ (\partial^\mu h_{\mu\nu})\square(\partial_\rho h^{\rho\nu}),\ h\square\partial^\mu\partial^\nu h_{\mu\nu},\ (\partial^\mu\partial^\nu h_{\mu\nu})^2\bigr)$ of Theorem~\ref{thm:wpB-dp} (the order of \texttt{COV4} in \texttt{higher\_grade\_scaling.py} and of \texttt{conventions.json}), let
\[
\ell(a,b,c,d,e):=a+b-c+d+e .
\]
Then $\ell(v_{R^2})=0+1-0-2+1=0$, $\ell(v_{\mathrm{Ric}^2})=\frac14+\frac14-\frac12-\frac12+\frac12=0$, and
\[
\ell(v_{\mathrm{scale}})=-\tfrac{71}{120}-\tfrac{27}{80}+\tfrac{11}{10}+\tfrac{19}{30}-\tfrac{13}{15}=-\tfrac1{16}\neq0,
\]
so $v_{\mathrm{scale}}\notin\mathrm{span}(v_{R^2},v_{\mathrm{Ric}^2})$ exactly.
\emph{Interpretation.} On the pure-gauge configuration $h_{00}=f(t)$, all other components zero ($h_{\mu\nu}=\partial_\mu\xi_\nu+\partial_\nu\xi_\mu$ with $\xi_0=\frac12\int f$), one has $h=h^{00}=f$, $\partial^\mu h_{\mu0}=f'$, $\partial^\mu\partial^\nu h_{\mu\nu}=f''$, $\square=\partial_t^2$, so modulo total derivatives
\[
(L_1,L_2,L_3,L_4,L_5)\equiv(1,1,-1,1,1)\,f f'''' ,
\]
($L_3=f'f'''\equiv-ff''''$, $L_5=(f'')^2\equiv ff''''$). Hence $\int\sum_ja_jL_j\equiv\ell(a)\int ff''''$: the curvature squares vanish on pure gauge, as they must, while the scale functional leaves $-\frac1{16}\,ff''''$ times the unit $i/(16\pi^2)$; its Euler--Lagrange expression in $f$ is $2\ell(a)f''''$.
This is verified symbolically with \texttt{sympy} in \texttt{regression\_M.py} (audit-3 section).
\emph{Scope.} The proposition certifies the displayed functional only; the identification of the graph's collapsed top scale tensor with $E[I^{\mathrm{top}}_{\mathrm{scale}}]$ (the graph$\to$functional obligation) is carried by the collapse certificate of Theorem~\ref{thm:wpB-dp}(i); \texttt{regression\_M.py} re-verifies it exactly (the recollapse of the archived $C_{\mathrm{top}}$ equals the Euler tensor of the displayed $I^{\mathrm{top}}_{\mathrm{scale}}$ entrywise, $490$ entries, and its $f''''$ coefficient in $E^{00}[f\,e_{00}]$ is $-\frac18=2\ell(v_{\mathrm{scale}})$).
\end{proposition}

\begin{remark}[Top contact versus full tensor; relation to XVI-L]
\label{rem:wpB-topfull}
The split between the cutoff-independent top component and the cutoff-dependent lower-degree components matches the residue left open by XVI-L, thm:L6: ``No explicit nonzero scaling residue (anomaly form) is claimed (Remark rem:Lressretch)'', where rem:Lressretch notes that ``a nonzero homogeneous distribution of degree \(-4\) may still have vanishing spherical average against the radial derivative of \(\varpi_\mu\)''.
The value $r=-i/(8\pi^2)$ (Theorem~\ref{thm:M6-S2}) discharges only the top contact (the nonvanishing spherical average at $N=0$); the lower-degree, $\mu\,\mathrm d\varpi$-dependent contacts at $N\ge1$ are not discharged.
\end{remark}

\begin{corollary}[Time-only readings of the scale tensor]
\label{cor:wpB-timeonly}
Assume Gate~V.
In the time-only reading and in the time-only $\chiI\equiv1$ reading of Section~\ref{sec:to-timeonly}, for every $\mu>0$ and every admissible $\varpi$, the collapsed scale tensor $\Coll_t(C_{\mathrm{scale}}(\mu))$ is not variational.
(The proof combines the archived top-order coordinates with the exact kernels of Theorem~\ref{thm:to-timeonly} and Proposition~\ref{prop:to-chiplateau}; it is complete relative to the reported computations of Lemma~\ref{lem:to-comp}; audit~3 review.)
\end{corollary}
\begin{proof}
The top-order part $C_{\mathrm{top}}$ of $C_{\mathrm{scale}}$ is a nonzero vector of the $D=4$ block (all its slots have $|\gamma|=N$, hence $D=4-N+|\gamma|=4$), and its $N=0$ coordinate is $-2i\pi^2c^2\neq0$ (archived in \texttt{wpB/\_wpB\_state.json}; re-checked by \texttt{regression\_M.py}).
$H_t$ is homogeneous in $D$ (Lemma~\ref{lem:to-comp}: the blocks have disjoint rows), so the $D=4$ component of $H_t(C_{\mathrm{scale}})$ is $H_t(C_{\mathrm{top}})$, whatever the $\varpi$-dependent lower-order parts are.
In the time-only reading $\ker H_t$ is $0$ at $D=4$ (Theorem~\ref{thm:to-timeonly}), so $H_t(C_{\mathrm{top}})\neq0$.
In the time-only $\chiI\equiv1$ reading $\ker H_t$ at $D=4$ is spanned by vectors supported on $N=4$ (Proposition~\ref{prop:to-chiplateau}); $C_{\mathrm{top}}$ has a nonzero $N=0$ coordinate, so again $H_t(C_{\mathrm{top}})\neq0$.
By Lemma~\ref{lem:helm}, applied with the time-only adjoint of Lemma~\ref{lem:to-intertwine}, a collapsed tensor with nonzero adjoint defect is not an Euler tensor.
\end{proof}

\begin{corollary}[M-C$_{\mathrm{curv}}$ obstructed; broad M-C not achieved]
\label{cor:wpB-MC}
Assume Gate~V. For the specified two-edge graph and the XVI-L prescription:
\begin{enumerate}
\item \textbf{M-C$_{\mathrm{curv}}$} (Definition~\ref{def:MCcurv}) is \textbf{Proved obstructed} (graph-specific): for every $\mu$ and every admissible $\varpi$ the full collapsed scale variation does not lie in the span of the M.5 curvature targets, because its fixed $D=4$ component already lies outside.
\item \textbf{Broad M-C} is \textbf{not achieved} and remains \textbf{Unresolved}: criterion~(i) of Definition~\ref{def:MCcurv} holds (the top scale tensor has its own variational representative $I^{\mathrm{top}}_{\mathrm{scale}}$, Theorem~\ref{thm:wpB-dp}(i)), but criterion~(ii) fails because XVI-L leaves $\varpi$ unspecified (thm:L6). Their numerical values and the variationality of the resulting full tensor are not determined for the unspecified cutoff. \textbf{M-C$_{\mathrm{data}}$} is satisfied.
\end{enumerate}
This is not a claim about gravity generally. Adding non-gauge-invariant quadratic representatives to the ledger changes the target ledger of M-C$_{\mathrm{curv}}$ but may remain within broad Target~I; gauge-fixing/ghost or further graph contributions are another, independent enlargement.
\end{corollary}
\begin{proof}
(1) M-C$_{\mathrm{curv}}$ requires a matching in the curvature-target span together with the matched-scale data of Definition~\ref{def:MCcurv}.
The scale derivative equals the dilation defect at every grade ($B_\mu=0$; exact $\varpi$-independent top residues; Theorem~\ref{thm:wpB-B0}); the numerical values of the lower-degree components and the variationality of the resulting full tensor are not determined for the unspecified cutoff (XVI-L, thm:L6).
But for every $\mu$ and every admissible $\varpi$ the collapsed scale tensor is not in the span of the ledger targets ($E_{P_1}$, $E_{P_2}$, $E_{\mathrm{EH}}$, $E_{R^2}$, $E_{\mathrm{Ric}^2}$, with the reading's weight): in the formal, $\chi$-plateau and preparation-plateau readings and in both time-only readings its $D=4$ component is not variational (table above; Corollary~\ref{cor:wpB-timeonly}), whereas every target is self-adjoint; in the double plateau its $D=4$ component is nonzero and outside $\mathrm{span}(E_{R^2},E_{\mathrm{Ric}^2})$ (Theorem~\ref{thm:wpB-dp}, Proposition~\ref{prop:wpB-sep}), and targets of degree $D=0,2$ cannot cancel it because collapse is homogeneous in $D$.
The lower-order, $\varpi$-dependent part cannot repair this, for the same reason.
(2) Theorem~\ref{thm:wpB-dp}(i),(ii): the $D=4$ component is the Euler tensor of $I^{\mathrm{top}}_{\mathrm{scale}}$, which lies outside the curvature ledger but is a local quadratic functional, so criterion~(i) of broad M-C holds; criterion~(ii) requires an explicitly specified $\varpi$ and a functional representation of the entire collapsed scale variation, which XVI-L, thm:L6, does not supply. M-C$_{\mathrm{data}}$ (matching + explicit prescription-dependent scaling formulas) is satisfied.
\end{proof}

\paragraph{Archives.}
\begin{sloppypar}\noindent
\texttt{wpB/higher\_grade\_scaling.py} ($185$ checks; output \texttt{higher\_grade\_scaling\_out.txt}), \texttt{scale\_derivative\_coefficients.json}, \texttt{dilation\_defect\_coefficients.json}, \texttt{scale\_defect\_difference.json}, \texttt{full\_scale\_collapse\_certificates.json}, \texttt{matched\_scale\_ledger.txt}, \texttt{\_wpB\_state.json}, \texttt{\_wpB\_sectionD.json}.
The script keeps $c^2$ symbolic; its status strings ``CONDITIONAL on Input (S-bv)'' predate Theorem~\ref{thm:wpA-Sbv} and are superseded by it.
\end{sloppypar}


\paragraph{M.6 status.}
\begin{itemize}
\item Dilation identity and $A_\mu=\mathcal A$ for $N=0$: \textbf{Constructed} (Proposition~\ref{prop:M6-dil}).
\item $A_\mu=\mathcal A_\mu$ ($B_\mu=0$) for all $N=0,\dots,4$ and every admissible $\varpi$: \textbf{Constructed} (Theorem~\ref{thm:wpB-B0}).
\item $S^2$ residue $r=-2i\pi^2 c^2=-i/(8\pi^2)\neq0$: \textbf{Constructed} (Theorem~\ref{thm:M6-S2}, with Theorem~\ref{thm:wpA-Sbv}).
\item Top-order residues, all grades: \textbf{Constructed}, exact and $\varpi$-independent (units $i\pi^2c^2=i/(16\pi^2)$); lower-order values: \textbf{Unresolved} (depend on $\nu_k[\varpi]$, $k=1,\dots,4$; $\varpi$ not fixed by XVI-L; the radial comparison is a separately named specialization).
\item $N=0$ anomaly component $C_{\mathrm{anom}}$ (higher-grade coordinates set to zero for this probe): \textbf{Constructed} (Definition~\ref{def:M6-Canom}).
\item Collapse of the $N=0$ component: \textbf{Constructed} (Theorem~\ref{thm:M6-collapse}); it contributes no nonzero plateau matching (Corollary~\ref{cor:M6-compat}).
\item Collapse of the full scale tensor: \textbf{Constructed} in all readings. In the double plateau, the cutoff-independent $D=4$ component is variational and Lorentz covariant, but not linearized-gauge invariant; it lies outside $\mathrm{span}(E_{R^2},E_{\mathrm{Ric}^2})$ (Theorem~\ref{thm:wpB-dp}, Proposition~\ref{prop:wpB-sep}). The full scale tensor has cutoff-dependent lower-degree components and need not be variational; its Lorentz covariance is not established. Time-only readings non-variational (Corollary~\ref{cor:wpB-timeonly}).
\item Release: \textbf{M-C$_{\mathrm{curv}}$ Proved obstructed} (graph-specific); \textbf{broad M-C not achieved} (Unresolved; $\varpi$ unspecified by XVI-L, thm:L6); \textbf{M-C$_{\mathrm{data}}$ Satisfied} (Corollary~\ref{cor:wpB-MC}); not a claim about gravity generally.
\end{itemize}
\begin{sloppypar}\noindent
Archives: \texttt{scaling\_anomaly.py}, \texttt{scaling\_anomaly\_out.txt}, \texttt{scaling\_anomaly\_cert.json}, \texttt{scaling\_anomaly\_cert.txt}; \texttt{wpB/}.
\end{sloppypar}


\section{M.7 Covariance and Ward-compatibility constraints}
\label{sec:m7}

This section treats \emph{prescription} covariance of an extension $\RM=\RL+\Delta\RL$ on the flat background.
It is kept distinct from interacting-sector Ward/BV symmetry (M.8--M.9; Target~II).
Gate~S-bv is certified (Theorem~\ref{thm:wpA-Sbv}), so the (S-bv) dependence of this section is discharged; M-C$_{\mathrm{curv}}$ is Proved obstructed (graph-specific), broad M-C is not achieved (Corollary~\ref{cor:wpB-MC}), and M-D is not claimed.
The audited M.7 revision repaired three defects of the previous M.7 computation (Remark~\ref{rem:M7-audit}).

\paragraph{Flat-background Lorentz covariance conditions (LC0--LC3).}
\begin{enumerate}
\item[(LC0)] Lie algebra $\mathfrak{so}(1,3)\subset\mathfrak{gl}(4)$, $X=(X^\mu{}_\nu)$ acting on coordinates $z^\mu$, with $X^{\mathsf T}\eta+\eta X=0$; generators $\{J_{12},J_{13},J_{23},K_{01},K_{02},K_{03}\}$ ($K_{0i}$ symmetric, $J_{ij}$ antisymmetric as matrices).
\item[(LC1)] \emph{Coordinate polynomials.} On functions of the difference variable, $(X\cdot f)(z)=-(Xz)^\mu\partial_\mu f(z)$.
By Theorem~\ref{thm:wpA-Sbv}, $\mathcal H_{4+N}\cong\mathrm{Poly}_N$ via numerators ($Q$ is invariant), inducing matrices $\rho_N(X)$ on the selected basis $B_N$.
\item[(LC2)] \emph{Contact distributions (repaired).} $W_N=\mathrm{span}\{\partial_\gamma\delta:|\gamma|\le N\}$, where $\partial_\gamma=\partial_0^{\gamma_0}\cdots\partial_3^{\gamma_3}$ are lower-index coordinate derivatives in $z$ (the same $\partial^\gamma$ as in $\mathcal G(C)=\sum C\,\bar a\,\partial^\gamma\delta$ and in the collapse $\partial_w^\gamma$).
$\mathfrak{so}(1,3)$ acts by pull-back, $\langle X\cdot T,f\rangle:=-\langle T,X\cdot f\rangle=\langle T,(Xz)^\mu\partial_\mu f\rangle$; for $T=\sum_\nu c_\nu\partial_\nu\delta$ this gives $X\cdot\partial_\nu\delta=X^\mu{}_\nu\,\partial_\mu\delta$, i.e.\ $c\mapsto Xc$.
On coordinate polynomials $p=\sum_\nu c_\nu z^\nu$ the action (LC1) is $c\mapsto-X^{\mathsf T}c$.
These differ for boosts.
The intertwiner is the metric identification $z^\mu\leftrightarrow\partial^\mu\delta=\eta^{\mu\nu}\partial_\nu\delta$ (index raising by $\eta$, using $\eta X\eta=-X^{\mathsf T}$), extended multiplicatively: a coordinate polynomial $q(z)$ corresponds to $q(\partial^\cdot)\delta$, so the coefficient of $\partial_\gamma\delta$ equals $(-1)^{\gamma_1+\gamma_2+\gamma_3}$ times the coefficient of $z^\gamma$.
Writing $P=\mathrm{diag}((-1)^{\gamma_1+\gamma_2+\gamma_3})$, the action on $W_N$ in the basis $\partial_\gamma\delta$ is $\rho^\delta_N(X)=P\,\rho^{\mathrm{poly}}_N(X)\,P$.
\item[(LC3)] A contact change $\Delta\RL(u)=\sum_{|\gamma|\le N}C_\gamma(u)\,\partial_\gamma\delta$ is infinitesimally covariant iff $C\circ\rho_N(X)=\rho^\delta_N(X)\circ C$ for all $X$.
Homogeneous solutions are the Lorentz-equivariant contact assignments.
The inhomogeneous correction of a concrete obstruction $\mathrm{Obs}_L(X,u)=X_*\RL(u)-\RL(X\cdot u)$ requires an explicit $\varpi$ and is \textbf{Unresolved} (a Euclidean-radial bump is not Lorentz-invariant; a boost obstruction is expected).
\end{enumerate}
A linear divergence identity of an Euler tensor is \textbf{not} a Ward identity.

\paragraph{Lorentz action on $B_N$ and on contact distributions.}
\begin{proposition}[so(1,3) action on $B_N$ and on $W_N$]
\label{prop:M7-rho}
For each $N=0,\ldots,4$ the maps $\rho_N(X)$ on $B_N$ and $\rho^\delta_N(X)$ on $\{\partial_\gamma\delta\}$ define Lie-algebra homomorphisms (in particular $[\rho(K_{01}),\rho(K_{02})]=\rho(J_{12})$ for both).
The $N=0$ action on $B_0$ is trivial.
\end{proposition}

\begin{proposition}[Boost verification on $\partial_\mu\delta$; exact]
\label{prop:M7-boost}
In the basis $(\partial_0\delta,\partial_1\delta,\partial_2\delta,\partial_3\delta)$ the pull-back action of every generator is the matrix $X$ itself; for $K_{01}$,
\[
\rho^\delta(K_{01})\big|_{|\gamma|=1}=\begin{pmatrix}0&1&0&0\\1&0&0&0\\0&0&0&0\\0&0&0&0\end{pmatrix},
\qquad
\rho^{\mathrm{poly}}(K_{01})\big|_{|\gamma|=1}=-\begin{pmatrix}0&1&0&0\\1&0&0&0\\0&0&0&0\\0&0&0&0\end{pmatrix}.
\]
The second matrix is the coordinate matrix used by the previous revision for the SA intersection (it identified $\partial_\gamma\delta$ with $z^\gamma$).
For the three boosts the two differ by a sign; for the three rotations they coincide; in all cases $\rho^\delta=\eta\,\rho^{\mathrm{poly}}\eta$ at first order and $\rho^\delta_N=P\rho^{\mathrm{poly}}_NP$ for all $N\le4$ (checked exactly for all six generators).
A symbolic test-function computation gives $\langle K_{01}\cdot\partial_1\delta,f\rangle=-\partial_0f(0)=\langle\partial_0\delta,f\rangle$, confirming $K_{01}\cdot\partial_1\delta=+\partial_0\delta$.
\end{proposition}

\begin{proof}
$z^\rho\partial_\nu\delta=-\delta^\rho_\nu\delta$ and $z^\rho\partial_\mu\partial_\nu\delta=-\delta^\rho_\nu\partial_\mu\delta-\delta^\rho_\mu\partial_\nu\delta$ give $-(Xz)^\mu\partial_\mu\partial_\nu\delta=X^\mu{}_\nu\partial_\mu\delta+(\mathrm{tr}X)\partial_\nu\delta$, and $\mathrm{tr}X=0$.
The higher-order statement follows because $q\mapsto q(\partial^\cdot)\delta$ is multiplicative and $\delta$ is invariant; it is also checked exactly by pairing with $z^\gamma/\gamma!$ for all $|\gamma|\le4$ (\texttt{covariance\_constraints.py}).
\end{proof}

Ultralocal Lorentz scalars $P_1=\int\eta^{\mu\alpha}\eta^{\nu\beta}h_{\mu\nu}h_{\alpha\beta}$ and $P_2=\int(\mathrm{tr}_\eta h)^2$, and the matched line $P_1-\tfrac{89}{140}P_2$, are annihilated by all six generators.

\paragraph{Compatibility matrix: solutions, constraint rank, cokernel.}
\begin{theorem}[Equivariant contact assignments]
\label{thm:M7-compat}
Write $\mathrm{Hom}_G(\mathcal H_{4+N},W_N)$ for $\mathfrak{so}(1,3)$-equivariant contact maps at grade $N$.
\begin{enumerate}
\item \emph{Abstract count.} $\mathrm{Poly}_N=\bigoplus_{k=0}^{\lfloor N/2\rfloor}H_{N-2k}$ with $\dim H_m=(m+1)^2$, $W_N\cong\mathrm{Poly}_{\le N}$ as representations (via the intertwiner of (LC2)), and
\[
\dim\mathrm{Hom}_G(\mathcal H_{4+N},W_N)=\frac{(\lfloor N/2\rfloor+1)(\lfloor N/2\rfloor+2)}{2}.
\]
Explicit basis: $C_{k,j}=P\circ\varphi_{k,j}$ with $\varphi_{k,j}:p\mapsto Q^j\square^k p$, $k=0,\ldots,\lfloor N/2\rfloor$, $j=0,\ldots,k$; in words, $Q^j\square^kp$ read as a polynomial in $\partial^\mu$ acting on $\delta$.
\item \emph{Coordinate verification.} Each $C_{k,j}$ satisfies $C\rho_N(X)=\rho^\delta_N(X)C$ exactly for all six generators and all $N\le4$.
The uncorrected maps $\varphi_{k,j}$ used previously fail exactly the three boost equations whenever $P$ acts nontrivially on their image, i.e.\ for $(N,k,j)\in\{(1,0,0),\allowbreak (2,0,0),\allowbreak (3,0,0),\allowbreak (3,1,0),\allowbreak (3,1,1),\allowbreak (4,0,0),\allowbreak (4,1,0),\allowbreak (4,1,1)\}$; the six maps whose image lies at $|\gamma|$ even with $P=+1$ ($N=0$; $(2,1,\cdot)$; $(4,2,\cdot)$) were unaffected.
For $N\le2$ the full constraint map was additionally solved by an exact verified kernel computation.
\item Graded compatibility matrix. Slots $=\dim\mathcal H_{4+N}\cdot\dim W_N$; the constraint map is $C\mapsto(C\rho_N(X)-\rho^\delta_N(X)C)_{X}$ with codomain $\bigoplus_{6\text{ generators}}\mathrm{Hom}(\mathcal H_{4+N},W_N)$ of dimension $6\cdot$slots;\allowbreak\ rank $=$ slots $-\dim$ sol.\ $=$ codimension of the solution space; coker $=6\cdot$slots $-$ rank.
\begin{center}
\begin{tabular}{crrrrrrl}
$N$ & $\dim H$ & $\dim W$ & slots & rank (=codim) & $\dim$ codomain & coker & $\dim$ sol./block\\
\hline
0 & 1 & 1 & 1 & 0 & 6 & 6 & 1 / K\\
1 & 4 & 5 & 20 & 19 & 120 & 101 & 1 / K\\
2 & 10 & 15 & 150 & 147 & 900 & 753 & 3 / K\\
3 & 20 & 35 & 700 & 697 & 4200 & 3503 & 3 / $\Creg$\\
4 & 35 & 70 & 2450 & 2444 & 14700 & 12256 & 6 / $\Creg$
\end{tabular}
\end{center}
The cokernel is dominated by the redundancy of the six generator equations and is not, by itself, an obstruction count.
Total equivariant freedom: $14=5$ (K-block, $N\le2$)$+9$ ($\Creg$, $N\ge3$) (unchanged by the repair, since $P$ is an isomorphism).
\item Separation: (a) $\Creg$ solutions ($N\ge3$) are strict $K$-regression-preserving ($\Delta\RL|_{\mathcal V_{\mathrm K}}=0$); (b) K-block solutions ($N\le2$) modify $\RL|_{\mathcal V_{\mathrm K}}$ and break XVI-K regression if used.
\item The inhomogeneous $\mathrm{Obs}_L(\RL)$ is \textbf{Unresolved}.
\end{enumerate}
\end{theorem}

\begin{proof}
Schur count from the harmonic decomposition; linear independence of $\{C_{k,j}\}$; equivariance against $\rho^\delta$ on all six generators; $N\le2$ exact kernel of the full constraint map (\texttt{m4\_common.block\_kernel}, verified on all rows).
For $N=3,4$ the rank follows from rank--nullity with the abstract kernel dimension and the explicit verified basis.
Certificates: \texttt{covariance\_matrix.txt}, \texttt{covariance\_certificates.txt}, \texttt{covariance\_summary.json}, \texttt{covariance\_exact\_certificates.json}.
\end{proof}

\paragraph{Intersection with the variational subspace.}
Let $\mathrm{Equiv}\subset\Ccoh$ be the $14$-dimensional span of the corrected $C_{k,j}$ and $V:=\mathrm{Equiv}\cap\mathrm{SA}$ (reading-dependent).
Write $\Ccoh=\mathcal C_{\mathrm K}\oplus\Creg$ by grade ($N\le2$ / $N\ge3$), with projections $\pi_K,\pi_{\mathrm{reg}}$.
Four distinct objects are reported: the projections $\pi_K(V)$, $\pi_{\mathrm{reg}}(V)$ (equivariant, but \emph{not} self-adjoint in general) and the intersections $V\cap\mathcal C_{\mathrm K}$, $V\cap\Creg$ (genuine SA assignments).
Always $\dim(V\cap\mathcal C_{\mathrm K})=\dim V-\dim\pi_{\mathrm{reg}}(V)$, $\dim(V\cap\Creg)=\dim V-\dim\pi_K(V)$, hence $\dim(V\cap\mathcal C_{\mathrm K})+\dim(V\cap\Creg)\le\dim V$.

\begin{proposition}[Equivariant $\cap$ SA, exact; Gate~V]
\label{prop:M7-inter}
$V$ is the exact rational kernel of the adjoint-defect map $x\mapsto E_{C(x)}^\dagger-E_{C(x)}$ on the $14$ equivariant coordinates:
\begin{center}
\begin{tabular}{lccccccc}
\hline
Reading & $\dim V$ & $\dim\pi_K(V)$ & $\dim V\cap\mathcal C_{\mathrm K}$ & $\dim\pi_{\mathrm{reg}}(V)$ & $\dim V\cap\Creg$ & $\dim\Coll(V)$\\
\hline
formal & $1$ & $0$ & $0$ & $1$ & $1$ & $1$\\
$\chi$-plateau & $1$ & $0$ & $0$ & $1$ & $1$ & $1$\\
double plateau & $12$ & $5$ & $3$ & $9$ & $7$ & $9$\\
\hline
\end{tabular}
\end{center}
\begin{enumerate}
\item Formal and $\chi$-plateau: $V$ is the line of $C_{2,0}$ at $N=4$ ($\square^2$, $\gamma=0$), contained in $\Creg$; $\Coll(V)$ is spanned by the Euler tensor of $P_1-\tfrac{89}{140}P_2$ (ratio $17920$ in the formal reading), i.e.\ the matched ultralocal line of Proposition~\ref{prop:match}.
\item Double plateau: $\dim V=12$, of which $3$ dimensions lie in $\ker\Coll$.
$\pi_K(V)$ is all of $\mathrm{Equiv}_K$ ($5$), but only $V\cap\mathcal C_{\mathrm K}$ ($3$) consists of SA assignments; the remaining projected directions have nonzero adjoint defect and are \emph{not} SA assignments.
Likewise $\pi_{\mathrm{reg}}(V)=\mathrm{Equiv}_{\mathrm{reg}}$ ($9$) while $V\cap\Creg$ has dimension $7$.
Check: $3+7=10\le12$.
\end{enumerate}
Every basis vector of $V$, $V\cap\mathcal C_{\mathrm K}$ and $V\cap\Creg$ is re-verified from scratch against all constraints: equivariance (six generators, all grades) and zero adjoint defect; \texttt{regression\_M.py} re-verifies them independently, including the Helmholtz route $\mathrm{EL}(\tfrac12\langle h,E_Ch\rangle)=E_C$.
Mod-$p$ ranks at three primes agree and are supplementary only.
Prescription-covariant variational contacts exist but do not by themselves produce a Ward identity.
\end{proposition}

\paragraph{Reconciliation with previously stated numbers.}
The equivariant freedom $14=5+9$ and the formal/$\chi$-plateau value $\dim(\mathrm{Equiv}\cap\mathrm{SA})=1$ are confirmed (the formal SA vector sits at $|\gamma|=0$, where $P=1$).
The previous double-plateau values ``$5$ globally, $5$ on K, $3$ on $\Creg$'' were wrong on three counts: the coordinate action was wrong for boosts (Proposition~\ref{prop:M7-boost}); the grade-restricted numbers were $\dim(\mathrm{Equiv}_G\cap\pi_G(\mathrm{SA}))$ (intersection with a \emph{projection} of SA), not $V\cap\mathcal C_G$; and the modular reduction \texttt{int(Fraction(c) \% p)} truncated the $14639$ non-integer rational entries of the SA basis.
With the old vectors and a correct reduction the old routine gives $5/4/6$; it reproduces $5/3/5$ only with the truncating reduction (\texttt{m7\_legacy\_reconstruction.py}).

\paragraph{Linearized divergence identities (external data retained).}
\begin{proposition}[Linear divergence identities; not Ward identities]
\label{prop:M7-id}
Let $E^{\mu\nu}$ denote Gate~V Euler tensors; they carry \emph{upper} indices because they are paired with the lowered $h_{\mu\nu}$, and $(\operatorname{div}E)^\nu:=\partial_\mu E^{\mu\nu}$ involves no metric factor.
\begin{enumerate}
\item[(A)] Contracted Bianchi: $\eta^{\mu\rho}\partial_\rho G^{\mathrm{lin}}_{\mu\nu}=0$ identically in jets ($G$ with lower indices).
\item[(B$'$)] \emph{Curvature squares.} For $I_w=\int w\,(\alpha R_{\mathrm{lin}}^2+\beta R^{\mathrm{lin}}_{\mu\nu}R_{\mathrm{lin}}^{\mu\nu})$ with constant $\alpha,\beta$ and \emph{any} fixed weight $w$ (in particular $w=1$ in the double plateau, $w=\sigeps$ in the $\chi$-plateau reading, $w=\chiI\sigeps$ formally),
\[
\partial_\mu E_{I_w}^{\mu\nu}=0\quad\text{identically.}
\]
Reason: $R_{\mathrm{lin}}$ and $R^{\mathrm{lin}}_{\mu\nu}$ are invariant \emph{pointwise} under $\delta h_{\mu\nu}=\partial_\mu\xi_\nu+\partial_\nu\xi_\mu$, so $0=\delta I_w=2\int E^{\mu\nu}\partial_\mu\xi_\nu=-2\int\xi_\nu\partial_\mu E^{\mu\nu}$; weight-derivative terms arise in the Euler tensor and cancel in its divergence.
(For the weighted Einstein--Hilbert action the analogous weight-derivative terms are $\sigma G^{\mu\nu}+q^{\mu\nu}\square\sigma-\nabla^\mu\nabla^\nu\sigma$ of XII-C, lem:weighted-curvature, eq:weighted-full-variation; no landmark constructs the weighted $R^2/\mathrm{Ric}^2$ Euler tensors, which are internal to XVI-M.)
This holds separately for $R^2$, $\mathrm{Ric}^2$ and for $R^2-\tfrac{188813}{88456}\mathrm{Ric}^2$ (exact assertions in all three readings).
\item[(EH)] For $I_w=\int w\,h^{\mu\nu}G^{\mathrm{lin}}_{\mu\nu}$: $\partial_\mu E_{I_w}^{\mu\nu}=(\partial_\mu w)\,G_{\mathrm{lin}}^{\mu\nu}[h]$ (indices raised with $\eta$), which vanishes identically for arbitrary $h$ iff $w$ is locally constant (double plateau).
Reason: $\delta(h\cdot G)=2\partial_\mu\xi_\nu G^{\mu\nu}$ and $\partial_\mu G^{\mu\nu}=0$.
\item[(B)] \emph{Ultralocal.} For the matched line $\int w(h_{\mu\nu}h^{\mu\nu}-\tfrac{89}{140}h^2)$: $\partial_\mu E^{\mu\nu}=2\partial_\mu[w(h^{\mu\nu}-\tfrac{89}{140}\eta^{\mu\nu}h)]\ne0$ in every reading; this nonvanishing is genuine, because the ultralocal term is not invariant under $\delta h=\partial\xi+\partial\xi$.
\end{enumerate}
The Euler tensors of $\int R_{\mathrm{lin}}^2$ and $\int\mathrm{Ric}_{\mathrm{lin}}^2$ are also built independently, from $g=\eta+\epsilon h$ by first-order expansion of the Christoffel symbols and Riemann tensor (no hand-written linearized formula), with \texttt{sympy}'s \texttt{euler\_equations} and the Gate~V weights; they coincide exactly with the targets used in M.4/M.5 (and $R_{\mathrm{lin}}$, $R^{\mathrm{lin}}_{\mu\nu}$ agree with the formulas of Proposition~\ref{prop:match} with sign $+1$).
The certified preimages of $E_{R2}$, $E_{RC2}$ (in $\Ccoh$) and of $E_{R2}-\tfrac{188813}{88456}E_{RC2}$ (in $\Creg$) collapse exactly to these independent targets and satisfy the identity.
These are linear (free-field gauge) identities; none of (A), (B), (B$'$), (EH) is an interacting Ward identity (no ghost/gauge-fixed interacting sector is present; Target~II unfinished).
\end{proposition}

\begin{remark}[Seeking $\RM=\RL+\Delta\RL$]
\label{rem:M7-RM}
Within the homogeneous constraint (LC3), any $\Delta\RL$ with contact map in $\mathrm{Hom}_G$ yields a prescription-covariant change.
Correcting a nonzero $\mathrm{Obs}_L(\RL)$ is Unresolved pending an explicit $\varpi$.
No claim is made that $\RL$ itself is Lorentz covariant (cf.\ XVI-L, non-properties of $\RL$).
\end{remark}

\begin{remark}[Audit log (M.7 audited revision; retained)]
\label{rem:M7-audit}
\begin{enumerate}
\item \emph{Curvature-square divergence (bug).} Old: \texttt{covariance\_constraints.py} computed $\sum_\mu\eta^{\mu\mu}\partial_\mu E^{\mu\nu}$, inserting a metric sign into the divergence of an upper-index Euler tensor, and recorded ``does not vanish'' through a check that always passed\linebreak[1] (\texttt{require(True,}\allowbreak\texttt{\ldots)}). The Euler tensor itself (\texttt{jet\_euler.euler}, factor $1/w$) and the off-diagonal multiplicities were correct. Fix: $\partial_\mu E^{\mu\nu}$ with no metric factor; hard assertions $=0$ for $R^2$, $\mathrm{Ric}^2$ and the combination in all readings; independent metric-based construction; targets and preimages checked. The matched coefficient $188813/88456$ is unchanged (the matching code never used the divergence routine).
\item \emph{Intersection dimensions (bug).} Old: grade-restricted SA \emph{projections} intersected with $\mathrm{Equiv}_G$, and a truncating modular reduction; reported $5/5/3$. Fix: exact kernel of the adjoint-defect map on $\mathrm{Equiv}$, four objects per reading (Proposition~\ref{prop:M7-inter}); legacy reconstruction archived.
\item \emph{Stage~2 fixed data (overstatement).} Old: criterion without the adjoint-defect map; transition bound $0\le\dim\ker\Coll_{\mathrm{fixed}}\le22$ unsupported. Fix: criteria (i)--(iii), time-only caveat, universal-versus-actual separation, domain-size bound $3321$ (Section~\ref{sec:m4-fixed}).
\item \emph{Contact representation (bug).} Old: $\partial_\gamma\delta$ identified with $z^\gamma$ ``by the same formula'', giving the wrong boost action; ``homogeneous cokernel $=$ codimension'' was a misnomer. Fix: explicit $\eta$ intertwiner, exact boost verification (Proposition~\ref{prop:M7-boost}), corrected equivariant vectors; rank/codimension and cokernel (with codomain) reported separately.
\item \emph{M.6 scope (overstatement).} $C_{\mathrm{anom}}$ renamed the $N=0$ anomaly component; the M-C sentence is replaced (Section~\ref{sec:m6}).
\item \emph{Status.} M.7 downgraded to ``Constructed after audit'' for exactly verified items; the former ``ALL CHECKS PASSED'' covering wrong expected identities is replaced by assertions of the corrected identities in \texttt{covariance\_constraints.py} and \texttt{regression\_M.py}.
\end{enumerate}
\end{remark}

\paragraph{M.7 status.}
\textbf{Constructed after audit} (Gate~V for the SA items; (S-bv) certified), exactly verified: the actions $\rho_N$, $\rho^\delta_N$ and the $\eta$ intertwiner; the boost verification; the abstract equivariant count $14=5+9$ and its coordinate verification; the exact $V=\mathrm{Equiv}\cap\mathrm{SA}$ with its four derived objects in the formal, $\chi$-plateau and double-plateau readings; the linear divergence identities (A), (B), (B$'$), (EH).
\textbf{Unresolved}: $\mathrm{Obs}_L(\RL)$ for a concrete $\varpi$; any interacting Ward/BV statement (M.8--M.9).


\section{M.8 Action, insertion, ghosts, and diagrammatic closure (stub)}
\label{sec:m8}
\textbf{Unresolved.} Target~II requires the action, gauge fixing, ghosts and antifields, the composite insertion, BRST/BV transformations, and a diagram inventory beyond the XVI-L two-edge graph.
No finite fully specified BV slice is available from predecessors; building a full BV complex here would be speculative.
\textbf{Recommendation:} M.8 should wait for XVI-N (or an explicitly finite BV slice).
Completion ledger stub: \texttt{completion\_ledger\_M89.json}.
\textbf{M-D and M-E are not claimed.} Prescription covariance (M.7) is not interacting-sector symmetry.

\section{M.9 Conditional symmetry-restoration result (stub)}
\label{sec:m9}
\textbf{Unresolved.} Depends on M.8 (closed sector).
Possible outcomes once M.8 is specified: Constructed / Conditionally constructed / Obstructed / Unresolved.
M.4 already shows that not every collision term can be made into a bulk gravitational counterterm: in the formal reading, no derivative term can.
M.7 supplies prescription-covariant variational contacts of low dimension; it does not restore interacting Ward identities.

%=====================================================================
\section{Status}
\label{sec:status}

\begin{remark}[Revision log (this cycle)]
\label{rem:revlog}
This cycle integrates three completed work packages; XVI-A--XVI-L and their artifacts are not modified.
\begin{enumerate}
\item \textbf{Work package A} (Section~\ref{sec:wpA-Sbv}): Gate~S-bv certified (Theorem~\ref{thm:wpA-Sbv}); $c=1/(4\pi^2)$, boundary side $Q-i0$, including across the nonzero null cone; convention ledger; position-space forms of $W_0$ and $E_0^{\mathrm A}$.
Consequences applied: the M.2 presentation, the M.3/M.4 ranks and the M.7 constructions are no longer conditional on (S-bv) (only $c\neq0$ and the single side were used); $r=-i/(8\pi^2)$ (Theorem~\ref{thm:M6-S2}); the identification at $z=0$ is not established ($S$ itself is global, XVI-J).
Every former ``Gate~S-bv conditional'' statement in the body was replaced by a reference to Theorem~\ref{thm:wpA-Sbv}; the Euclidean motivation (Remark~\ref{rem:Sbv-motivation}) is kept and is explicitly not the certification.
\item \textbf{Work package B} (Section~\ref{sec:wpB}): $A_\mu=\mathcal A_\mu$ ($B_\mu=0$) at every grade $N\le4$; top-order residues $\varpi$-independent; lower orders through $\nu_k[\varpi]$; double-plateau scale tensor variational but not matched (superseded in audit~3: only the top $D=4$ component is variational; Theorem~\ref{thm:wpB-dp}); M-C \emph{Proved obstructed} in the graph-specific matching domain (superseded in audit~3 by: M-C$_{\mathrm{curv}}$ Proved obstructed, broad M-C not achieved; Corollary~\ref{cor:wpB-MC}).
The former M.6 remark ``$A_\mu=\mathcal A$ unresolved for $N\ge1$'' is marked resolved.
\item \textbf{Work package C} (Section~\ref{sec:to-timeonly}): universal time-only classification; $\chiI\equiv1$ time-only reading; route~1 $=$ route~2; the double-plateau curvature-square matching does not survive the time-only reading (Remark~\ref{rem:M5-timeonly}).
Downgrade applied: the actual-data consequence is Proposition~\ref{prop:to-actual} with hypotheses (H1)--(H3), not a Constructed claim; the linear independence of the $\sigeps^{(k)}$ is a separately proved Lemma~\ref{lem:to-indep}.
\item \textbf{Integration consistency checks.}
(a) \emph{Normalization of $c$.} The fragments agree: wpA certifies $c=1/(4\pi^2)$ with side $Q-i0$; the motivational value of Remark~\ref{rem:Sbv-motivation} coincides; wpB works with symbolic $c^2$, and substituting gives the unit $i\pi^2c^2=i/(16\pi^2)$, with $N=0$ top coefficient $-2\,i\pi^2c^2=r$; wpC uses only $c\neq0$. No sign or $2\pi$ discrepancy was found; \texttt{conventions.json} records $c$ and the side.
(b) \emph{Labels.} No clashes between the fragments and the manuscript; the wpC labels \texttt{sec:time-only}, \texttt{def:time-only}, \texttt{thm:time-only} were renamed into the \texttt{to-} prefix (\texttt{sec:to-timeonly}, \texttt{def:to-timeonly}, \texttt{thm:to-timeonly}); wpA/wpB use the prefixes \texttt{wpA-}, \texttt{wpB-}.
(c) \emph{Fragment typos.} A malformed longtable column specification and a retarded/advanced slip in the wpA fragment were corrected; plain-text cross-references were turned into \verb|\ref|s.
(d) \emph{Superseded statements.} The fragments' own ``conditional on (S-bv)'' or ``not discharged'' clauses (written before integration) are superseded by Theorem~\ref{thm:wpA-Sbv}; wpC keeps ``on the current coherent presentation''.
(e) \emph{Time-only scale tensor.} wpB left open how the scale tensor behaves in the time-only readings; Corollary~\ref{cor:wpB-timeonly} (new, derived at integration from the archived top-order coordinate and the wpC kernels) closes this: it is non-variational there.
(f) \emph{Historical outputs.} Earlier-cycle script outputs and certificates keep their historical status strings (e.g.\ ``S-bv CONDITIONAL'', ``M-C \dots unfinished''); they are reproduced verbatim and are superseded by this revision.
\end{enumerate}
\end{remark}

\begin{remark}[Revision log (audit 3)]
\label{rem:revlog3}
Integration and proof audit following expert review; no new work package, no expansion; XVI-A--XVI-L, their artifacts and the GitHub repository are not modified.
\begin{enumerate}
\item \emph{Gate~S-bv provenance.} XVI-J already prints the formula (eq:W, eq:DF: $h_0=-W_0$, $d=W_0-iE_0^{\mathrm A}$, $S=-d$; eq:rational: $d=-\frac1{4\pi^2}(Q_{\mathrm{geom}}-i0)^{-1}$ on $z\neq0$ including nonzero null points; prop:q uses it); XI-E, eq:h, only sets $h=-W$ and the jump. These are quoted in Section~\ref{sec:wpA-identify}, and Theorem~\ref{thm:wpA-Sbv} is reframed as a written check of XVI-J, eq:rational (corroboration, not a frozen XVI-L premise; the direct Fourier proof and Lemma~\ref{lem:wpA-varreg} stand independently). New Lemma~\ref{lem:wpA-varreg} proves the boundary value of the variable regulator $Q-2i\tau z^0-\tau^2$ near nonzero null points. The damping parameter is renamed $\varepsilon\to\tau$ (and the Euclidean slice variable $\tau\to\tau_{\mathrm E}$) in the manuscript and the wpA fragment and script. The scope sentence on ``extension across $z=0$'' is replaced by: the predecessor $S$ is globally defined (XVI-J); what is not established here is its position-space identification at the origin.
\item \emph{Top versus full scale tensor.} In the double plateau, the cutoff-independent $D=4$ component is variational and Lorentz covariant, but not linearized-gauge invariant. The full scale tensor has cutoff-dependent lower-degree components and need not be variational; its Lorentz covariance is not established (Theorem~\ref{thm:wpB-dp}). The relation to XVI-L, thm:L6 / rem:Lressretch, is recorded (Remark~\ref{rem:wpB-topfull}); $r=-i/(8\pi^2)$ discharges only the top contact.
\item \emph{Separating functional.} Proposition~\ref{prop:wpB-sep}: $\ell=a+b-c+d+e$ annihilates $v_{R^2},v_{\mathrm{Ric}^2}$ and gives $\ell(v_{\mathrm{scale}})=-\frac1{16}$; basis order and signs verified against \texttt{full\_scale\_collapse\_certificates.json}, \texttt{COV4} and \texttt{conventions.json}; pure-gauge interpretation verified with \texttt{sympy}; the $f''''$ coefficient $-\frac18$ of $E^{00}[f\,e_{00}]$ is recomputed from the archived $C_{\mathrm{top}}$.
\item \emph{M-C split.} Definition~\ref{def:MCcurv}: broad M-C is the matched-scale criterion (proved matching + explicit $\varpi$ and a functional representation of the entire collapsed scale variation); M-C$_{\mathrm{data}}$ (matching + explicit prescription-dependent scaling formulas) is satisfied; M-C$_{\mathrm{curv}}$ Proved obstructed (graph-specific); broad M-C not achieved (Unresolved) because XVI-L leaves $\varpi$ unspecified (thm:L6). The implication that ghosts/diagrams are the only enlargement is removed (Corollary~\ref{cor:wpB-MC}; release ladder; abstract).
\item \emph{Time-only.} The step ``By Lemma~\ref{lem:to-indep} this reduces to linear independence of the occurring switch jets'' is removed; the sufficient condition is kept for the joint family $\{\sigeps^{(k)}\partial^\rho\chiI\}$, with the counterexample Remark~\ref{rem:to-joint}. ``Do not extend to any region where $\sigeps$ varies'' is replaced by ``do not admit the stated universal time-only completion; the actual-ramp conclusion holds under Proposition~\ref{prop:to-actual}''. Corollary~\ref{cor:wpB-timeonly} was reviewed; its proof is complete relative to the reported computations and it is kept as a corollary.
\item \emph{Proof corrections.} (a) Lemma~\ref{lem:wpB-dw}: the basis derivative is $\mu\partial_\mu f_\gamma=(E-|\gamma|)f_\gamma$, combined with $(|\gamma|-|\alpha|)$ from $M^{-1}$. (b) Lemma~\ref{lem:bv}(c): explicit $\widehat{uv}=\frac1{2\pi}\hat u*\hat v$ and Beta integral, with identification of the H\"ormander product by the local Fourier formula; the ``sequential continuity'' shortcut is removed. (c) Proposition~\ref{prop:M7-id}(B$'$): weight-derivative terms arise in the Euler tensor and cancel in its divergence (cf.\ XII-C, eq:weighted-full-variation). (d) Shell annotations in the run listing and README are \texttt{\#} comments. (e) \texttt{regression\_M.py} labels current-status versus historical-record checks and has an audit-3 section.
\end{enumerate}
\end{remark}

\begin{remark}[Revision log (final verification)]
\label{rem:revlog-final}
Final verification edits; no expansion (no M.8--M.9); XVI-A--XVI-L, their artifacts and the GitHub repository are not modified.
\begin{enumerate}
\item \emph{Time-only $\chi$-plateau kernel.} Proposition~\ref{prop:to-chiplateau}: $\ker H_t|_D$ is displayed by cases ($D=0,2,4$ versus $D=1,3$); odd spatial order is anti-self-adjoint ($E[h]=\sigma(t)K\partial_1 h$ gives $E^\dagger=-E$ since $\partial_1\sigma=0$); the spacetime-formal non-self-adjointness remark refers to the extra $D=2,4$ directions; dimensions $(26,0,156,0,390)$ unchanged; regression asserts the odd-$D$ anti-self-adjointness exactly.
\item \emph{Broad M-C matched-scale criterion.} Definition~\ref{def:MCcurv}: broad M-C requires (i)~a proved functional matching and (ii)~an explicit cutoff $\varpi$ with a functional representation of the entire collapsed scale variation; M-C$_{\mathrm{data}}$ (matching + explicit prescription-dependent scaling formulas) is satisfied; broad M-C Unresolved because XVI-L leaves $\varpi$ unspecified (thm:L6). Replaced ``uncontrolled'' by: numerical values and variationality of the full tensor are not determined for the unspecified cutoff.
\item \emph{Predecessor provenance.} Pinned local/frozen files and archived corpus excerpts under \texttt{XVI-M-artifacts/\allowbreak provenance/}; XVI-J eq:rational marked as corroboration, not a frozen XVI-L premise; ``resting on a new input'' $\to$ ``an explicitly checked kernel identification''.
\item \emph{Wording.} Lemma~\ref{lem:M2comp}: independently spot-checks every relation on $(\partial^\alpha S)(\partial^\beta S)$ at three rational points. Proposition~\ref{prop:M7-id}(EH): vanishes identically for arbitrary $h$ iff $w$ is locally constant.
\item \emph{Full scratch rerun.} \texttt{XVIM\_FULL\_WP\allowbreak =1} regression recorded as \texttt{regression\_full\_\allowbreak out.txt} (artifacts and package); appendix listing updated.
\end{enumerate}
\end{remark}

\subsection{Constructed}
\begin{enumerate}
\item M.0: frozen dependency ledger; Targets I/II; normalization table (now with the certified $S$ row); Gates S-bv (certified) and V; verification boundary; release ladder.
\item M.1 (unconditional): general-$h$ assembly; exact specialization to XVI-L ($1065/1065$); supports $255/255/235$.
\item Gate~S-bv: \textbf{certified} (Theorem~\ref{thm:wpA-Sbv}); $S=\frac1{4\pi^2}(Q-i0)^{-1}$ on $\Omega$, across the nonzero null cone; $W_0$ and $E_0^{\mathrm A}$ in position space.
\item M.2 (unconditional on $\Omega$): boundary-value calculus lemma; two-direction relation equivalence; $d_N=(1,4,10,20,35)$; $\dim\mathcal V_{\mathrm{adm}}=70$; $185$ certified relations; wave-operator bound $\le210$.
\item M.3: $\dim\Ccoh=3321$; $\mathcal G_{\mathrm{formal}}$ injective; background kernel $1750$; plateau kernel $21$.
\item M.4 (Gate~V): collapse injective (formal, $\chi$-plateau, time-only), kernel $21$ (prep.\ plateau) / $22$ (double plateau); SA dimensions $26/26/47/1347$ (time-only: $26$; time-only with $\chiI\equiv1$: $572$); variational images $26/26/26/1325$ (time-only: $26$ and $572$); Helmholtz lemma and certificates; $R^2$ and $\mathrm{Ric}^2$ matchings in the double plateau.
\item Stage~2: fixed-data criteria (i)--(iii), universal region comparisons, and the universal time-only classification (Theorem~\ref{thm:to-timeonly}, Proposition~\ref{prop:to-chiplateau}); explicit $\sigeps$ and the linear independence of its derivatives (Lemma~\ref{lem:to-indep}).
\item M.5: named theorem~\ref{thm:M5-plateau-R2} (double-plateau domain); ledgers~(1)--(2); completion-gate obstruction; ledgers~(3)--(4) Unresolved.
\item M.6: dilation identity; $A_\mu=\mathcal A_\mu$ for all $N\le4$ (Theorem~\ref{thm:wpB-B0}); $S^2$ residue $r=-i/(8\pi^2)$; top-order residues ($\varpi$-independent); $N=0$ component collapse; full scale-tensor collapse; in the double plateau, the cutoff-independent $D=4$ component is variational and Lorentz covariant, but not linearized-gauge invariant, and lies outside the curvature span (exact separating functional, Proposition~\ref{prop:wpB-sep}); the full scale tensor has cutoff-dependent lower-degree components and need not be variational (Theorem~\ref{thm:wpB-dp}).
\item M.7 (\textbf{Constructed after audit}; Gate~V for SA items; exactly verified items only): so(1,3) action on $B_N$ and pull-back action on $\partial_\gamma\delta$ with explicit $\eta$ intertwiner and exact boost verification; abstract equivariant count $1,1,3,3,6$ (K $5$ + $\Creg$ $9$) with coordinate verification; exact $V=\mathrm{Equiv}\cap\mathrm{SA}$ ($\dim V=1/1/12$) with $\pi_K(V)$, $V\cap\mathcal C_{\mathrm K}$, $\pi_{\mathrm{reg}}(V)$, $V\cap\Creg$ ($0,0,1,1$ / $0,0,1,1$ / $5,3,9,7$); linear divergence identities (A), (B), (B$'$), (EH). $\mathrm{Obs}_L(\RL)$ Unresolved.
\item M.8--M.9: light stubs; Target~II / BV Unresolved; M.8 recommended to wait for XVI-N.
\end{enumerate}
Release: \textbf{M-A} Constructed (presentation unconditional on $\Omega$). \textbf{M-B} Constructed in the formal reading (Gate~V convention). \textbf{M-C$_{\mathrm{curv}}$} Proved obstructed (graph-specific); broad \textbf{M-C} not achieved (Unresolved; $\varpi$ unspecified by XVI-L, thm:L6); \textbf{M-C$_{\mathrm{data}}$} Satisfied. \textbf{M-D/M-E} not claimed.

\subsection{Proved obstructed}
\begin{enumerate}
\item \textbf{M-C$_{\mathrm{curv}}$ (graph-specific; Definition~\ref{def:MCcurv}).} For the specified two-edge graph, the XVI-L prescription and any admissible $\varpi$, the full collapsed scale variation does not lie in the span of the M.5 curvature targets, because its fixed $D=4$ component already lies outside (Theorem~\ref{thm:wpB-dp}, Proposition~\ref{prop:wpB-sep}, Corollaries~\ref{cor:wpB-timeonly},~\ref{cor:wpB-MC}). Not a claim about gravity generally, and not an obstruction of broad M-C.
\end{enumerate}

\subsection{Selected-image statements (Gate~V; formal-versus-fixed kept)}
\begin{enumerate}
\item Formal reading: the formal universal-jet collapse image of this selected graph contains no nonzero variational tensor of total derivative degree $D\ge1$ ($\mathrm{SA}_D=0$ for $D\ge1$).
This is a vanishing inside one selected image under one reading; it is not a global obstruction to derivative counterterms.
\item Double plateau: the specified quadratic Einstein--Hilbert Euler tensor is not in the selected double-plateau collapse image.
Membership outside that selected image, or after other diagrams, is not addressed.
\item Time-only readings: the variational image is the $26$ ultralocals ($\chiI$ formal) or $572$ spatial-only, non-Lorentz directions ($\chiI\equiv1$); only $P_1-\tfrac{89}{140}P_2$ among the probed targets is an image; the weighted EH, $R^2$ and $\mathrm{Ric}^2$ Euler tensors are not (Proposition~\ref{prop:to-match}).
\end{enumerate}
These do not claim absence of derivative counterterms in a completed gravitational sector, failure for every fixed preparation, or impossibility after adding other diagrams.

\subsection{Unresolved}
\begin{enumerate}
\item Broad M-C (Definition~\ref{def:MCcurv}: proved functional matching plus an explicitly specified cutoff $\varpi$ and a functional representation of the entire collapsed scale variation): not achieved, because XVI-L leaves $\varpi$ unspecified (thm:L6). The top scale tensor has its own variational representative $I^{\mathrm{top}}_{\mathrm{scale}}$, so the M-C$_{\mathrm{curv}}$ obstruction does not decide broad M-C; M-C$_{\mathrm{data}}$ is satisfied. Their numerical values and the variationality of the resulting full tensor are not determined for the unspecified cutoff (Corollary~\ref{cor:wpB-MC}).
\item Lorentz covariance and variationality of the full (not only top) double-plateau scale tensor for the actual, unspecified $\varpi$.
\item The position-space identification of the predecessor $S$ at the origin $z=0$ ($S$ itself is globally defined, XVI-J; not needed for the punctured-space results).
\item Actual fixed-data dimensions where the actual $\chiI$ is non-constant, and transition regions (only $0\le\dim\ker\Coll_{\mathrm{fixed}}\le3321$); whether the actual $\chiI$ satisfies the hypotheses of Proposition~\ref{prop:to-actual} on given regions. Sufficient test (recorded in Section~\ref{sec:to-timeonly}): on an open $U$, a full-rank exact evaluation matrix of the joint family of external-data coefficient functions $\sigeps^{(k)}(t)\partial^\rho\chiI(x)$ occurring in the graph map, in $\Coll_t$ and in the adjoint-defect map $H_t$; not verified for any concrete $\chiI$ (separate independence of the two factors does not suffice, Remark~\ref{rem:to-joint}).
\item Lower-order scale residues: their values depend on $\nu_k[\varpi]$, $k=1,\dots,4$; XVI-L does not fix $\varpi$; the radial comparison is a separately named specialization, not the XVI-L prescription.
\item Inhomogeneous $\mathrm{Obs}_L(\RL)$ for a concrete $\varpi$ (M.7); the closed sector (M.8); restoration (M.9).
\item EOM and permitted field redefinitions for ledgers~(3)--(4).
\item Powers of $\kappa$ and $\hbar$, the remaining factors of $i$, loop order, and the relation of $\chiI$ to the action weighting.
\item Physical (normalized) values of the matched coefficients.
\end{enumerate}

\subsection{Non-claims}
\begin{enumerate}
\item M-D and M-E are not claimed. The M-C$_{\mathrm{curv}}$ obstruction is graph-specific: it concerns the specified two-edge graph, the XVI-L prescription and the M.5 curvature targets; it is not a claim about gravity generally and not an obstruction of broad M-C. Enlarging the ledger (by non-gauge-invariant quadratic representatives, which may remain within broad Target~I, or by gauge-fixing/ghost or further graph contributions) gives a different matching domain, not covered here.
\item The certification of Gate~S-bv is a written check of XVI-J, eq:rational, on $\mathbb R^4\setminus\{0\}$ only: the predecessor $S$ is globally defined (XVI-J), and its position-space identification at the origin is not established here; no $i\varepsilon$ prescription of the predecessors is modified, and the Euclidean motivation is not used as a proof.
\item The $546$ extra variational directions of the time-only $\chiI\equiv1$ reading are spatial-only and non-Lorentz; they are not proposed as counterterms.
\item Proposition~\ref{prop:to-actual} is a conditional transfer, not a certification of actual fixed-data dimensions.
\item $B_\mu=0$ does not remove prescription dependence: lower-order residues still depend on $\varpi$.
\item The linear divergence identities of M.7 are free-field gauge identities, not interacting Ward identities; prescription covariance is not interacting-sector symmetry.
\item The double-plateau intersection numbers of the previous M.7 draft ($5/5/3$) and its statement that the curvature-square Euler tensor has nonzero divergence are withdrawn (Remark~\ref{rem:M7-audit}).
\item No $\varepsilon\to0$; no UV completion; no interacting BRST/Ward identity as proved; no all-orders renormalization.
\item The XI inverse $G_{\mathrm{gr}}^{\mathrm R}=-\Pc E_0^{\mathrm R}$ is not equated with the X-sheet $\Delta_R^{\mathrm{gr}}=-4\Pc E_0^{\mathrm R}$.
\item Preimage coefficients and collision dimensions are not physical Wilson coefficients.
\item Scale-dependent subtraction and scale-variation coefficients are not physical beta functions.
\item No release claim of XVI-L is altered.
\item Computational certificates support the stated identities; they do not replace the written proofs.
\end{enumerate}

\scopebox{Release.}{%
\textbf{XVI-M draft:}
\begin{itemize}
\item General-$h$ coefficients: Constructed; they specialize exactly to XVI-L.
\item Gate~S-bv: certified (written check of XVI-J, eq:rational), $S=\frac1{4\pi^2}(Q-i0)^{-1}$ on $\mathbb R^4\setminus\{0\}$; graded presentation, contact maps, collapse and variational subspaces: Constructed, unconditional on (S-bv).
\item Formal, (universal) preparation-plateau and time-only readings: only ultralocal $\chiI\sigeps K(h,h)$ potentials are variational ($546$ further spatial-only directions if $\chiI\equiv1$ in the time-only reading).
\item Double plateau: linearized $R^2$ and $\mathrm{Ric}^2$ are matched (Theorem~\ref{thm:M5-plateau-R2}); the specified quadratic Einstein--Hilbert Euler tensor is not in the selected collapse image; the matching does not survive the time-only reading.
\item Fixed data: criteria, region comparisons and time-only classification constructed; actual-ramp transfer as a proposition with explicit hypotheses; non-constant $\chiI$ Unresolved.
\item M.6: $r=-i/(8\pi^2)$ (top contact only); $A_\mu=\mathcal A_\mu$ at all grades. In the double plateau, the cutoff-independent $D=4$ component is variational and Lorentz covariant, but not linearized-gauge invariant, and lies outside $\mathrm{span}(E_{R^2},E_{\mathrm{Ric}^2})$ ($\ell(v_{\mathrm{scale}})=-\frac1{16}$). The full scale tensor has cutoff-dependent lower-degree components and need not be variational.
\item M.7 (Constructed after audit): prescription Lorentz compatibility matrix (equivariant $1,1,3,3,6$) with corrected contact representation; exact $\mathrm{Equiv}\cap\mathrm{SA}$ ($1/1/12$); linear divergence identities; Obs$_L(\RL)$ Unresolved. M.8 waits for XVI-N.
\end{itemize}
M-A Constructed. M-B Constructed in the formal reading. M-C$_{\mathrm{curv}}$ Proved obstructed (graph-specific; not a claim about gravity generally). Broad M-C not achieved (Unresolved; $\varpi$ unspecified by XVI-L, thm:L6). M-C$_{\mathrm{data}}$ Satisfied. M-D/M-E not claimed.
}


\begin{thebibliography}{9}
\bibitem{BF} R.~Brunetti, K.~Fredenhagen, \emph{Microlocal analysis and interacting quantum field theories}, Commun.\ Math.\ Phys.\ 208 (2000) 623, arXiv:math-ph/9903028.
\bibitem{FR} K.~Fredenhagen, K.~Rejzner, \emph{Perturbative algebraic quantum field theory}, arXiv:1208.1428.
\bibitem{Hor} L.~H\"ormander, \emph{The Analysis of Linear Partial Differential Operators I}, 2nd ed., Springer, 1990 (Thms.~6.1.2, 8.1.6, 8.2.4, 8.2.10; theorem numbers as in that edition).
\bibitem{GS} I.~M.~Gel'fand, G.~E.~Shilov, \emph{Generalized Functions, Vol.~1}, Academic Press.
\bibitem{Olver} P.~J.~Olver, \emph{Applications of Lie Groups to Differential Equations}, 2nd ed., Springer (Thm.~5.92).
\bibitem{series} Series files XI-E, XII-C, XVI-I, XVI-J, XVI-K, XVI-L (quoted by label).
\end{thebibliography}

%=====================================================================
\appendix
\section{Computational artifacts}
\label{app:script}

Inventory of \texttt{XVI-M-artifacts/}:
\begin{sloppypar}
\begin{itemize}
\item Shared code and copies: \texttt{xvi\_common.py} (byte-identical to XVI-L); \texttt{m4\_common.py}; \texttt{jet\_euler.py}.
\item M.1: \texttt{graph\_general.py}, with \texttt{graph\_general\_out.txt} and \texttt{graph\_general\_table.txt}.
\item M.2: \texttt{kernel\_relations.py}, with \texttt{kernel\_relations\_out.txt}, \texttt{kernel\_basis\_table.txt} (reconstruction certificates and numerator rows), \texttt{kernel\_relation\_matrix.txt} (nullspace vectors) and \texttt{kernel\_basis.json}.
\item M.3: \texttt{contact\_space.py}, with \texttt{contact\_space\_out.txt} and \texttt{graph\_contact\_map.txt} ($u_m$, annihilators).
\item M.4: \texttt{contact\_collapse.py}, with \texttt{contact\_collapse\_out.txt}, \texttt{collapse\_kernel.json} and \texttt{collapse\_kernel\_certificates.txt}; \texttt{helmholtz\_check.py}, with \texttt{helmholtz\_check\_out.txt}, \texttt{helmholtz\_certificates.txt} and \texttt{helmholtz\_summary.json}; \texttt{matching\_probe.py}, with \texttt{matching\_probe\_out.txt} and \texttt{matching\_certificates.txt}.
\item Stage~2: \texttt{fixed\_data\_reading.py}, with \texttt{fixed\_data\_reading\_out.txt}, \texttt{fixed\_data\_reading.json}, \texttt{prep\_plateau\_sa.json}.
\item M.5: \texttt{counterterm\_functionals.py}, with \texttt{counterterm\_functionals\_out.txt}, \texttt{matching\_matrix.json}/\texttt{.txt}, \texttt{matching\_certificates\_M5.txt}, \texttt{ledger\_*.json}, \texttt{ledger\_tables.txt}, \texttt{m5\_ledger\_summary.json}.
\item M.6: \texttt{scaling\_anomaly.py}, with \texttt{scaling\_anomaly\_out.txt}, \texttt{scaling\_anomaly\_cert.json}, \texttt{scaling\_anomaly\_cert.txt}.
\item M.7 (audited): \texttt{covariance\_constraints.py} (imports \texttt{curvature\_independent.py}, the independent metric-linearization route), with \texttt{covariance\_constraints\_out.txt}, \texttt{covariance\_matrix.txt}, \texttt{covariance\_certificates.txt}, \texttt{covariance\_summary.json}, \texttt{covariance\_exact\_certificates.json} (exact rational bases of $V$, $\pi_K(V)$, $\pi_{\mathrm{reg}}(V)$, $V\cap\mathcal C_{\mathrm K}$, $V\cap\Creg$ and the corrected equivariant vectors); \texttt{m7\_legacy\_reconstruction.py} with \texttt{m7\_legacy\_reconstruction\_out.txt} and \texttt{.json} (reconstruction of the withdrawn numbers); completion stub \texttt{completion\_ledger\_M89.json}.
\item Work package A (\texttt{wpA/}, Gate~S-bv): \texttt{kernel\_identification.tex} (source fragment of Section~\ref{sec:wpA-Sbv}), \texttt{kernel\_identification\_checks.py} (symbolic/numerical cross-checks; writes \texttt{kernel\_identification\_out.txt}), \texttt{kernel\_conventions.json} (convention ledger, $c$, boundary side).
\item Work package B (\texttt{wpB/}, scaling): \texttt{scaling\_prescription\_identity.tex} (source fragment of Section~\ref{sec:wpB}), \texttt{higher\_grade\_scaling.py} with \texttt{higher\_grade\_scaling\_out.txt}, \texttt{scale\_derivative\_coefficients.json}, \texttt{dilation\_defect\_coefficients.json}, \texttt{scale\_defect\_difference.json}, \texttt{full\_scale\_collapse\_certificates.json}, \texttt{matched\_scale\_ledger.txt}, and the state files \texttt{\_wpB\_state.json}, \texttt{\_wpB\_sectionD.json}.
\item Work package C (\texttt{wpC/}, time-only): \texttt{time\_only\_section.tex} (source fragment of Section~\ref{sec:to-timeonly}), \texttt{time\_only\_helmholtz.py} with \texttt{time\_only\_helmholtz\_out.txt}, \texttt{time\_only\_collapse\_certificates.json}, \texttt{time\_only\_sa\_certificates.json}, \texttt{time\_only\_matching\_certificates.json}, \texttt{time\_only\_summary.json}, \texttt{time\_only\_specialization\_tests.txt}.
The fragments are archived as delivered, except the audit-3 renaming of the damping parameter $\varepsilon\to\tau$ (and of the Euclidean slice variable to $\tau_{\mathrm E}$) in \texttt{wpA/kernel\_identification.tex} and \texttt{wpA/kernel\_identification\_checks.py}, whose new check section~7b is also audit-3; the manuscript sections are edited versions (Remarks~\ref{rem:revlog} and~\ref{rem:revlog3}).
\item Bookkeeping: \texttt{regression\_M.py}, \texttt{regression\_M\_out.txt} and \texttt{regression\_full\_\allowbreak out.txt} (full scratch rerun with \texttt{XVIM\_FULL\_WP\allowbreak =1}); \texttt{make\_manifest.py} and \texttt{provenance\_manifest.txt}; \texttt{provenance/} (predecessor pins and corpus excerpts); \texttt{conventions.json}; \texttt{README.txt}.
\end{itemize}
\end{sloppypar}
Run from the archive parent directory, \texttt{/workspace}:
\begin{lstlisting}
PY=/workspace/tmp/venv/bin/python ; cd /workspace ; A=XVI-M-artifacts
$PY $A/graph_general.py           > $A/graph_general_out.txt
$PY $A/kernel_relations.py        > $A/kernel_relations_out.txt
$PY $A/contact_space.py           > $A/contact_space_out.txt
$PY $A/contact_collapse.py        > $A/contact_collapse_out.txt
$PY $A/helmholtz_check.py         > $A/helmholtz_check_out.txt
$PY $A/matching_probe.py          > $A/matching_probe_out.txt
$PY $A/fixed_data_reading.py      > $A/fixed_data_reading_stdout.txt   # writes fixed_data_reading_out.txt
$PY $A/counterterm_functionals.py > $A/counterterm_functionals_out.txt
$PY $A/scaling_anomaly.py         > $A/scaling_anomaly_out.txt
$PY $A/covariance_constraints.py  > $A/covariance_constraints_out.txt  # imports curvature_independent.py
$PY $A/m7_legacy_reconstruction.py > $A/m7_legacy_reconstruction_out.txt
$PY $A/wpA/kernel_identification_checks.py   # writes wpA/kernel_identification_out.txt
$PY $A/wpB/higher_grade_scaling.py > $A/wpB/higher_grade_scaling_out.txt   # ~5.5 min; writes wpB/*.json
$PY $A/wpC/time_only_helmholtz.py   # ~2.5 min; writes wpC/time_only_helmholtz_out.txt, wpC/*.json
$PY $A/make_manifest.py
$PY $A/regression_M.py            > $A/regression_M_out.txt
# full scratch rerun of wpB/wpC (about 8 min); recorded output:
XVIM_FULL_WP=1 $PY $A/regression_M.py > $A/regression_full_out.txt
\end{lstlisting}
Exit status $0$ together with \texttt{ALL CHECKS PASSED} means every check passed.
XVI-L and XVI-K artifacts are not modified; \texttt{regression\_M.py} verifies their SHA-256.
\texttt{regression\_M.py} reruns the wpA script and compares its output byte for byte, checks the archived wpB/wpC results and recomputes the key ones independently; setting \texttt{XVIM\_FULL\_WP\allowbreak =1} also reruns the wpB and wpC scripts in a scratch copy and compares outputs, ignoring timing lines.
\emph{Historical output strings.} Outputs and certificates from earlier cycles (M.2--M.7, and the wp scripts written before integration) carry status strings such as ``CONDITIONAL on (S-bv)'', ``Gate S-bv not discharged'' or ``M-C \dots unfinished''. They are reproduced verbatim below as records of those runs and are superseded by this revision (Theorem~\ref{thm:wpA-Sbv}, Corollary~\ref{cor:wpB-MC}); the numbers they certify are unchanged. In \texttt{regression\_M.py}, checks are labelled \texttt{[current status]} (assertions about the present manuscript and ledger) or \texttt{[historical record]} (checks that an archived output still carries its historical status string).

\section{Script outputs (verbatim copies of the current recorded outputs)}
\label{app:out}
\paragraph{M.1 (\texttt{graph\_general\_out.txt}).}
\begin{lstlisting}
# graph_general.py output   (python 3.13.5, sympy 1.14.0, exact QQ arithmetic via fractions.Fraction)
# M.1 GENERAL-h assembly of the XVI-L complete-source two-edge graph (independent h-jets)
CHECK frozen input present: XVI-L-artifacts/source_full_table.txt: True
# sha256(XVI-L-artifacts/source_full_table.txt) = e13ce0ae29c769ac3b154b338a7dba2eaa75c2f12fa3427c93bdbe1cfb08430f
CHECK frozen input present: XVI-K-artifacts/hessian_full_out.txt: True
# sha256(XVI-K-artifacts/hessian_full_out.txt) = ce840f99f6df6ba31f6d415d78ec050a5ba09bb0c3153192d94c829f585ce490
CHECK frozen input present: XVI-L-artifacts/two_edge_coeff_table.txt: True
# sha256(XVI-L-artifacts/two_edge_coeff_table.txt) = cf53538133e87da76fa126e884223bab2871ac3a177f740005f66df4d3e0b1c7
CHECK frozen input present: XVI-L-artifacts/xvi_common.py: True
# sha256(XVI-L-artifacts/xvi_common.py) = 5b7b0c483a31489548efc901e3388afae7a86efc376073ab25416b34367b6f3d
CHECK xvi_common.py is byte-identical to XVI-L-artifacts/xvi_common.py: True
CHECK P_A table has 5056 entries (XVI-L L.1): True
CHECK source placements 1/8/36 for C_2/M_2/J_3: True
CHECK P_C2 uses only sigma 2-jet symbols: True
CHECK P_M2 uses only sigma 1-jet symbols: True
CHECK P_J3 uses only sigma: True
CHECK L coefficient count 2988 (XVI-K): True
CHECK slot symmetry L_ij^{pq} = L_ji^{qp} (entrywise, symbolic): True
CHECK each L^{pq}_{ij} is linear in h-jets of order exactly 2-|p|-|q| (XVI-K prop:hess(1)): True
# distinct residual-metric jet symbols occurring in L: 122
CHECK 45 ordered vertex placements (symbolic L): True
# T_i^r = (-P B^i)^r ; nnz = 22/100
# === GENERAL-h ASSEMBLY (all sectors) ===
# sector C_2 assembled in 0.1s; ordered nonzero (placement,mn) contributions 450
# sector M_2 assembled in 0.5s; ordered nonzero (placement,mn) contributions 3600
# sector J_3 assembled in 0.5s; ordered nonzero (placement,mn) contributions 10644
# assembly wall time 1.1s
CHECK edge interchange: ordered contribution (A;a,b;p,q) equals (A;b,a;q,p) for EVERY ordered placement with distinct swap (exhaustive over all enumerated placements: 1000 unordered pairs): True
CHECK Lambda_adm has 255 labels (XVI-L L.3): True
CHECK admissible per sector 25/84/250 (XVI-L): True
CHECK every monomial is exactly (one sigma-jet symbol) x (one h-jet symbol): bilinear in prep jet and residual-metric jet: True
CHECK sector grading: prep-jet order = 2-s, h-jet order = 2-v, N = s+v (s=0,1,2 for C2,M2,J3): True
CHECK no unintended generators (only sector prep symbols and XVI-K h-jet symbols): True
CHECK all coefficients exact rationals: True
CHECK output tensor stored as symmetric (mu<=nu) components (P_A symmetric in (mu,nu), XVI-L L.1 (C3)): True
CHECK C_2 formal support lies in its admissible block: True
CHECK M_2 formal support lies in its admissible block: True
CHECK J_3 formal support lies in its admissible block: True
CHECK XVI-L two_edge_coeff_table.txt parsed: 1065 nonzero entries: True
CHECK general-h table specialized to h=(w^1)^2 eta has exactly 1065 nonzero (A,lambda,mn) entries: True
CHECK specialized nonzero key set equals XVI-L key set (sector, label, component): True
CHECK EXACT SPECIALIZATION: every grouped coefficient equals the XVI-L background coefficient (all 1065 entries, coefficientwise in prep jet and w^1): True
CHECK background support recomputed: C2 25, M2 84, J3 202, union 235 (XVI-L): True
# === SUPPORT DISTINCTIONS ===
# C_2: |Lambda_adm|=25 {0: 1, 1: 4, 2: 20} ; |Lambda_graph^formal|=25 {0: 1, 1: 4, 2: 20} ; |Lambda_graph^background|=25 {0: 1, 1: 4, 2: 20}
# M_2: |Lambda_adm|=84 {1: 4, 2: 20, 3: 60} ; |Lambda_graph^formal|=84 {1: 4, 2: 20, 3: 60} ; |Lambda_graph^background|=84 {1: 4, 2: 20, 3: 60}
# J_3: |Lambda_adm|=250 {2: 20, 3: 60, 4: 170} ; |Lambda_graph^formal|=250 {2: 20, 3: 60, 4: 170} ; |Lambda_graph^background|=202 {2: 20, 3: 32, 4: 150}
# complete: |Lambda_adm|=255 {0: 1, 1: 4, 2: 20, 3: 60, 4: 170} ; |Lambda_graph^formal|=255 {0: 1, 1: 4, 2: 20, 3: 60, 4: 170} ; |Lambda_graph^background|=235 {0: 1, 1: 4, 2: 20, 3: 60, 4: 150}
CHECK Lambda_graph^background subset Lambda_graph^formal subset Lambda_adm: True
# labels in Lambda_adm but NOT in Lambda_graph^formal (0): 
# labels in Lambda_graph^formal but vanishing on the XVI-L background (20): {(), 0123} {3, 012} {33, 12} {33, 02} {33, 01} {2, 013} {23, 11} {23, 01} {23, 00} {22, 13} {22, 03} {22, 01} {1, 023} {13, 02} {13, 00} {12, 03} {12, 00} {123, 0} {11, 03} {11, 02}
# nonzero (lambda,mn) entries by sector: {'C_2': 250, 'M_2': 840, 'J_3': 2192}  bilinear terms: {'C_2': 4008, 'M_2': 8092, 'J_3': 6546}
# sector-summed a_lambda: nonzero (lambda,mn) entries 2322 ; bilinear terms 18646
# sample: a^{C_2,(0,0)}_{{0,0}} has 78 bilinear terms (sigma_2jet x d^2h)
#   first terms: 5*sig_11*h_00_11 ; 1*sig_11*h_00_22 ; 1*sig_11*h_00_33 ; -10*sig_11*h_01_01 ; -2*sig_11*h_02_02 ; -2*sig_11*h_03_03
# === INDEPENDENT C_2 ROUTE, GENERAL h (xvi_common.polar_s, all 25 labels x 10 mn) ===
CHECK GENERAL-h C_2 regression: all 25 labels x 10 components equal the independent polar_s route, coefficientwise in (sigma 2-jet) x (h-jet) (250 comparisons, 23.2s): True
# wrote graph_general_table.txt : 3282 nonzero (A,lambda,mn) lines
# sha256(graph_general_table.txt) = 4cd29d3c9f691340b314aecbfe7d6aef3590b9cee21dcdb486fe9fbcdd679e49
# total runtime 27.1s
ALL CHECKS PASSED
\end{lstlisting}

\paragraph{M.2 (\texttt{kernel\_relations\_out.txt}).}
\begin{lstlisting}
# kernel_relations.py output   (python 3.13.5, sympy 1.14.0, exact QQ arithmetic)
# M.2 numerical-kernel relations: presentation C^{Lambda_adm} -> V_adm, degree by degree
# --- finite symbolic check behind the MOTIVATION for (S-bv) (Euclidean calculation; does not certify (S-bv)) ---
CHECK int_0^oo e^{-p tau} sin(p r) dp = r/(tau^2+r^2): True
CHECK XVI-K eq:W-fam integrand at z0=-i*tau gives 1/(4 pi^2 (tau^2+r^2)) (radial integration): True
CHECK d^alpha (1/Q) = p_alpha / Q^{1+|alpha|} with p_alpha homogeneous of degree |alpha| (all 70 |alpha|<=4): True
CHECK p_0 = 1: True
CHECK p_{e_0} = -2 z0  (coordinate derivative d_0 Q = 2 z0): True
CHECK p_{e_1} = +2 z1  (coordinate derivative d_1 Q = -2 z1): True
CHECK sign example: p_{e0} p_{e1} = -4 z0 z1: True
CHECK sign example: p_{e0+e1} = -8 z0 z1: True
CHECK sign example: (d0 S)(d1 S) = (1/2) S d0 d1 S  (numerators -4 z0 z1 = (1/2)(-8 z0 z1)): True
# labelled generators per N: [1, 4, 20, 60, 170]  total 255
CHECK label counts 1,4,20,60,170 (=Lambda_adm, XVI-L): True
CHECK qo_lambda = c^2 n_lambda (Q-i0)^{-2-N}, n_lambda=p_alpha p_beta homogeneous of degree N, for all 255 labels: True
CHECK N=0: greedy basis has rank-many elements (1): True
CHECK N=0: exact reconstruction n_lambda = sum_b R_{lambda b} n_b for all 1 labels: True
CHECK N=0: every exported relation vector annihilates the numerator map exactly (0 relations): True
CHECK N=0: dim ker pi_N = #labels - rank = 0: True
CHECK N=1: greedy basis has rank-many elements (4): True
CHECK N=1: exact reconstruction n_lambda = sum_b R_{lambda b} n_b for all 4 labels: True
CHECK N=1: every exported relation vector annihilates the numerator map exactly (0 relations): True
CHECK N=1: dim ker pi_N = #labels - rank = 0: True
CHECK N=2: greedy basis has rank-many elements (10): True
CHECK N=2: exact reconstruction n_lambda = sum_b R_{lambda b} n_b for all 20 labels: True
CHECK N=2: every exported relation vector annihilates the numerator map exactly (10 relations): True
CHECK N=2: dim ker pi_N = #labels - rank = 10: True
CHECK N=3: greedy basis has rank-many elements (20): True
CHECK N=3: exact reconstruction n_lambda = sum_b R_{lambda b} n_b for all 60 labels: True
CHECK N=3: every exported relation vector annihilates the numerator map exactly (40 relations): True
CHECK N=3: dim ker pi_N = #labels - rank = 40: True
CHECK N=4: greedy basis has rank-many elements (35): True
CHECK N=4: exact reconstruction n_lambda = sum_b R_{lambda b} n_b for all 170 labels: True
CHECK N=4: every exported relation vector annihilates the numerator map exactly (135 relations): True
CHECK N=4: dim ker pi_N = #labels - rank = 135: True
# === DEGREE-BY-DEGREE PRESENTATION (conditional on (S-bv)) ===
#  N | labelled generators | dim Poly_N | d_N = dim H_{4+N} | dim ker pi_N
#  0 |   1 |  1 |  1 |   0
#  1 |   4 |  4 |  4 |   0
#  2 |  20 | 10 | 10 |  10
#  3 |  60 | 20 | 20 |  40
#  4 | 170 | 35 | 35 | 135
CHECK d_N = dim Poly_N = C(N+3,3) = 1,4,10,20,35 : numerator map is ONTO homogeneous polynomials of degree N: True
# dim V_adm = 70 (vs 255 labels); dim ker pi = 185
CHECK dim V_adm = 70 (conditional on (S-bv)): True
CHECK every label with N<=2 is an XVI-K label (V_K block = C_2 block): True
# V_K = H_4 + H_5 + H_6 ; dim V_K = 15 = {0: 1, 1: 4, 2: 10} ; V_K cap H_7 = V_K cap H_8 = 0
CHECK dim V_K = 1+4+10 = 15 (25 XVI-K generators, 10 relations, all at N=2): True
CHECK selected basis of H_4+H_5+H_6 consists of XVI-K labels (basis extends a basis of V_K): True
# === UNCONDITIONAL STRUCTURAL RELATIONS (square S = 0 on R^4\{0}; no (S-bv)) ===
#  N=0: wave-operator relations generated 0, span dim 0 ; dim ker pi_N = 0 ; unconditional upper bound d_N <= 1
#  N=1: wave-operator relations generated 0, span dim 0 ; dim ker pi_N = 0 ; unconditional upper bound d_N <= 4
CHECK N=2: every wave-operator relation lies in ker pi_N (consistency with (S-bv)): True
#  N=2: wave-operator relations generated 1, span dim 1 ; dim ker pi_N = 10 ; unconditional upper bound d_N <= 19
CHECK N=3: every wave-operator relation lies in ker pi_N (consistency with (S-bv)): True
#  N=3: wave-operator relations generated 8, span dim 8 ; dim ker pi_N = 40 ; unconditional upper bound d_N <= 52
CHECK N=4: every wave-operator relation lies in ker pi_N (consistency with (S-bv)): True
#  N=4: wave-operator relations generated 36, span dim 36 ; dim ker pi_N = 135 ; unconditional upper bound d_N <= 134
# unconditional upper bound: dim V_adm <= 210 ; (S-bv) value: 70
CHECK wave-operator span is contained in the full kernel (dimension check): True
# relations NOT generated by the wave-operator family (need (S-bv)): {0: 0, 1: 0, 2: 9, 3: 32, 4: 99}
# sample N=2 relations (reconstruction form  e_lambda - sum_b R e_b = 0):
#   (-1)*q{(), 33} + (-1)*q{(), 22} + (-1)*q{(), 11} + (1)*q{(), 00}
#   (-1)*q{(), 33} + (-1/2)*q{(), 22} + (-1/2)*q{(), 11} + (1)*q{3, 3} + (1)*q{0, 0}
#   (-1/2)*q{(), 01} + (1)*q{1, 0}
#   (-1/2)*q{(), 02} + (1)*q{2, 0}
# dim D_{<=N} = {0: 1, 1: 5, 2: 15, 3: 35, 4: 70}
# dim C_coh = sum_N d_N dim D_{<=N} = 3321   (raw labelled slots 14321, XVI-L count 14321)
# dim C_reg (Delta R|_{V_K}=0) = d_3*35 + d_4*70 = 3150
CHECK raw slot count reproduces XVI-L 14321: True
CHECK dim C_coh = 3321, dim C_reg = 3150 (conditional on (S-bv)): True
# wrote kernel_basis_table.txt  sha256=267344b69a488ec7f99fa3a86180fe1cf25431c60a31beeec3d777da579d88e8
# wrote kernel_relation_matrix.txt  sha256=f25684da36b2f19709e344ab3cb599c093bf8d4e9b236ba3ef069ac2fc8dcbec
# wrote kernel_basis.json  sha256=fcbe2876abd2a6b410813f732356f829c970c22c3a333f7e4ec34909c0e0793b
# total runtime 2.9s
ALL CHECKS PASSED
\end{lstlisting}

\paragraph{M.3 (\texttt{contact\_space\_out.txt}).}
\begin{lstlisting}
# contact_space.py output   (python 3.13.5, exact QQ arithmetic via fractions.Fraction)
# M.3 coherent contact space, graph-contact map G, kernel/image (conditional on (S-bv) through M.2)
CHECK input present: graph_general_table.txt: True
# sha256(graph_general_table.txt) = 4cd29d3c9f691340b314aecbfe7d6aef3590b9cee21dcdb486fe9fbcdd679e49
CHECK input present: kernel_basis.json: True
# sha256(kernel_basis.json) = fcbe2876abd2a6b410813f732356f829c970c22c3a333f7e4ec34909c0e0793b
CHECK graph_general_table.txt: 3282 sector-resolved nonzero lines (M.1): True
CHECK d_N = 1,4,10,20,35 (M.2, conditional on (S-bv)): True
CHECK reconstruction available for all 255 labels: True
# dim D_{<=N} = {0: 1, 1: 5, 2: 15, 3: 35, 4: 70}  ; dim C_coh = 3321  ; dim C_reg = 3150  ; raw labelled slots = 14321
CHECK FORMAL N=0: exact rank 1 equals independent mod-p rank: True
CHECK FORMAL N=1: exact rank 4 equals independent mod-p rank: True
CHECK FORMAL N=2: exact rank 10 equals independent mod-p rank: True
CHECK FORMAL N=3: exact rank 20 equals independent mod-p rank: True
CHECK FORMAL N=4: exact rank 35 equals independent mod-p rank: True
# FORMAL                     dim U_N = [1, 4, 10, 20, 35] ; dim im G = 3321 ; dim ker G = 0 ; on C_reg: dim im = 3150, dim ker = 0  (0.1s)
CHECK TIME-ONLY N=0: exact rank 1 equals independent mod-p rank: True
CHECK TIME-ONLY N=1: exact rank 4 equals independent mod-p rank: True
CHECK TIME-ONLY N=2: exact rank 10 equals independent mod-p rank: True
CHECK TIME-ONLY N=3: exact rank 20 equals independent mod-p rank: True
CHECK TIME-ONLY N=4: exact rank 35 equals independent mod-p rank: True
# TIME-ONLY                  dim U_N = [1, 4, 10, 20, 35] ; dim im G = 3321 ; dim ker G = 0 ; on C_reg: dim im = 3150, dim ker = 0  (0.1s)
CHECK BACKGROUND N=0: exact rank 1 equals independent mod-p rank: True
CHECK BACKGROUND N=1: exact rank 4 equals independent mod-p rank: True
CHECK BACKGROUND N=2: exact rank 10 equals independent mod-p rank: True
CHECK BACKGROUND N=3: exact rank 20 equals independent mod-p rank: True
CHECK BACKGROUND N=4: exact rank 10 equals independent mod-p rank: True
# BACKGROUND                 dim U_N = [1, 4, 10, 20, 10] ; dim im G = 1571 ; dim ker G = 1750 ; on C_reg: dim im = 1400, dim ker = 1750  (0.0s)
CHECK TIME+BG N=0: exact rank 1 equals independent mod-p rank: True
CHECK TIME+BG N=1: exact rank 4 equals independent mod-p rank: True
CHECK TIME+BG N=2: exact rank 10 equals independent mod-p rank: True
CHECK TIME+BG N=3: exact rank 16 equals independent mod-p rank: True
CHECK TIME+BG N=4: exact rank 10 equals independent mod-p rank: True
# TIME+BG                    dim U_N = [1, 4, 10, 16, 10] ; dim im G = 1431 ; dim ker G = 1890 ; on C_reg: dim im = 1260, dim ker = 1890  (0.0s)
CHECK PLATEAU(J_3 only, formal) N=0: exact rank 0 equals independent mod-p rank: True
CHECK PLATEAU(J_3 only, formal) N=1: exact rank 0 equals independent mod-p rank: True
CHECK PLATEAU(J_3 only, formal) N=2: exact rank 10 equals independent mod-p rank: True
CHECK PLATEAU(J_3 only, formal) N=3: exact rank 20 equals independent mod-p rank: True
CHECK PLATEAU(J_3 only, formal) N=4: exact rank 35 equals independent mod-p rank: True
# PLATEAU(J_3 only, formal)  dim U_N = [0, 0, 10, 20, 35] ; dim im G = 3300 ; dim ker G = 21 ; on C_reg: dim im = 3150, dim ker = 0  (0.0s)
CHECK C_2 only, formal N=0: exact rank 1 equals independent mod-p rank: True
CHECK C_2 only, formal N=1: exact rank 4 equals independent mod-p rank: True
CHECK C_2 only, formal N=2: exact rank 10 equals independent mod-p rank: True
CHECK C_2 only, formal N=3: exact rank 0 equals independent mod-p rank: True
CHECK C_2 only, formal N=4: exact rank 0 equals independent mod-p rank: True
# C_2 only, formal           dim U_N = [1, 4, 10, 0, 0] ; dim im G = 171 ; dim ker G = 3150 ; on C_reg: dim im = 0, dim ker = 3150  (0.0s)
CHECK M_2 only, formal N=0: exact rank 0 equals independent mod-p rank: True
CHECK M_2 only, formal N=1: exact rank 4 equals independent mod-p rank: True
CHECK M_2 only, formal N=2: exact rank 10 equals independent mod-p rank: True
CHECK M_2 only, formal N=3: exact rank 20 equals independent mod-p rank: True
CHECK M_2 only, formal N=4: exact rank 0 equals independent mod-p rank: True
# M_2 only, formal           dim U_N = [0, 4, 10, 20, 0] ; dim im G = 870 ; dim ker G = 2451 ; on C_reg: dim im = 700, dim ker = 2450  (0.0s)
CHECK specialized visible ranks are bounded by the formal ones (specialization cannot raise rank): True
CHECK N=0: S^2 direction is graph-visible (U_0 = H_4): True
CHECK coefficient monomials of different N are disjoint (grading prep order 2-s, h order 2-N+s): True
CHECK every reduced coefficient monomial has prep order + h order = 4 - N: True
CHECK N=0: annihilator of U_N has dimension d_N - dim U_N = 0: True
CHECK N=1: annihilator of U_N has dimension d_N - dim U_N = 0: True
CHECK N=2: annihilator of U_N has dimension d_N - dim U_N = 0: True
CHECK N=3: annihilator of U_N has dimension d_N - dim U_N = 0: True
CHECK N=4: annihilator of U_N has dimension d_N - dim U_N = 0: True
CHECK kernel certificate: for every annihilator vector x of U_N, sum_b x_b abar_b = 0 exactly (graph-invisible numerical direction): True
# wrote graph_contact_map.txt sha256=49b59c00dc788efcb48681f4bb711fa11350f251404977aba5d251df6018628e
# === COMPARISON: ambient prescription freedom vs graph-visible collision freedom (FORMAL) ===
#   raw labelled contact slots (XVI-L count)          : 14321
#   coherent assignments dim C_coh (M.2, (S-bv))       : 3321
#   graph-visible collision freedom dim im G (formal)  : 3321
#   graph-invisible coherent directions dim ker G      : 0
#   regression-preserving: dim C_reg 3150 -> dim im 3150, ker 0
# total runtime 0.3s
ALL CHECKS PASSED
\end{lstlisting}

\paragraph{M.4 collapse (\texttt{contact\_collapse\_out.txt}).}
\begin{lstlisting}
# contact_collapse.py output   (exact QQ arithmetic; conditional on (S-bv) through M.2)
# M.4 contact-collapse map: rank, kernel; injectivity of G does NOT imply injectivity after collapse
CHECK graph_general_table loaded (3282 lines; field indices a<=b): True
CHECK d_N = 1,4,10,20,35 (M.2): True
CHECK dim C_coh = 3321 contact variables: True
# sha256(graph_general_table.txt) = 4cd29d3c9f691340b314aecbfe7d6aef3590b9cee21dcdb486fe9fbcdd679e49
# sha256(kernel_basis.json) = fcbe2876abd2a6b410813f732356f829c970c22c3a333f7e4ec34909c0e0793b
# number of (D, parity) blocks: 40 ; max block size: 292
CHECK formal: rank-nullity: rank 3321 + ker 0 = 3321: True
# formal           dim im Collapse = 3321 ; dim ker Collapse = 0 ; nnz entries = 666507  (5.5s)
CHECK chi_plateau: rank-nullity: rank 3321 + ker 0 = 3321: True
# chi_plateau      dim im Collapse = 3321 ; dim ker Collapse = 0 ; nnz entries = 157549  (3.2s)
CHECK double_plateau: rank-nullity: rank 3299 + ker 22 = 3321: True
# double_plateau   dim im Collapse = 3299 ; dim ker Collapse = 22 ; nnz entries = 73320  (91.0s)
CHECK FORMAL: Collapse injective on C_coh (dim ker = 0): True
CHECK chi_plateau: Collapse injective: True
CHECK double_plateau: dim ker Collapse = 22: True
# NOTE: injectivity of G_formal (M.3) does not by itself imply injectivity of Collapse;
#       the formal reading happens to make Collapse injective as well (verified above).
# dim C_reg = 3150 (expected 3150)
CHECK dim C_reg = 3150: True
# formal           dim (ker Collapse cap C_reg) = 0 ; dim Collapse(C_reg) = 3150
# chi_plateau      dim (ker Collapse cap C_reg) = 0 ; dim Collapse(C_reg) = 3150
# double_plateau   dim (ker Collapse cap C_reg) = 0 ; dim Collapse(C_reg) = 3150
# wrote collapse_kernel.json collapse_kernel_certificates.txt
# === SUMMARY ===
# formal           dim im = 3321 ; dim ker = 0 ; ker_cap_C_reg = 0 ; perD (size,ker) = {0: [35, 0], 1: [160, 0], 2: [440, 0], 3: [944, 0], 4: [1742, 0]}
# chi_plateau      dim im = 3321 ; dim ker = 0 ; ker_cap_C_reg = 0 ; perD (size,ker) = {0: [35, 0], 1: [160, 0], 2: [440, 0], 3: [944, 0], 4: [1742, 0]}
# double_plateau   dim im = 3299 ; dim ker = 22 ; ker_cap_C_reg = 0 ; perD (size,ker) = {0: [35, 0], 1: [160, 0], 2: [440, 0], 3: [944, 4], 4: [1742, 18]}
# total runtime 105.4s
ALL CHECKS PASSED
\end{lstlisting}

\paragraph{M.4 Helmholtz (\texttt{helmholtz\_check\_out.txt}).}
\begin{lstlisting}
# helmholtz_check.py output   (exact QQ; conditional on (S-bv) through M.2)
# M.4 formal self-adjoint (variational) subspace under Gate V; Helmholtz certificates
CHECK dim C_coh = 3321: True
CHECK unit test: Euler derivatives of 30 random quadratic Lagrangians (independent sympy route) are self-adjoint under m4_common.adjoint, all readings, D=0..4: True
CHECK negative control: E^{00}=d_0 h_{01} is NOT self-adjoint in every reading: True
CHECK formal: every exported SA vector satisfies E_C = E_C^dag exactly: True
# formal           dim SA = 26 ; per D = {0: 26, 1: 0, 2: 0, 3: 0, 4: 0} ; defect-map nnz = 1482596 (37.2s)
CHECK chi_plateau: every exported SA vector satisfies E_C = E_C^dag exactly: True
# chi_plateau      dim SA = 26 ; per D = {0: 26, 1: 0, 2: 0, 3: 0, 4: 0} ; defect-map nnz = 645026 (16.6s)
CHECK double_plateau: every exported SA vector satisfies E_C = E_C^dag exactly: True
# double_plateau   dim SA = 1347 ; per D = {0: 26, 1: 0, 2: 271, 3: 12, 4: 1038} ; defect-map nnz = 60016 (28.2s)
CHECK formal: dim Collapse(SA) = dim SA - dim ker Collapse = 26 - 0 = 26 (mod-p rank lower bound meets the upper bound): True
# formal           dim (T_variational cap im Collapse) = 26 ; per D = {0: 26, 1: 0, 2: 0, 3: 0, 4: 0}
CHECK chi_plateau: dim Collapse(SA) = dim SA - dim ker Collapse = 26 - 0 = 26 (mod-p rank lower bound meets the upper bound): True
# chi_plateau      dim (T_variational cap im Collapse) = 26 ; per D = {0: 26, 1: 0, 2: 0, 3: 0, 4: 0}
CHECK double_plateau: dim Collapse(SA) = dim SA - dim ker Collapse = 1347 - 22 = 1325 (mod-p rank lower bound meets the upper bound): True
# double_plateau   dim (T_variational cap im Collapse) = 1325 ; per D = {0: 26, 1: 0, 2: 271, 3: 8, 4: 1020}
CHECK FORMAL: SA = 26-dimensional, supported only at D=0 (no variational collapse with any derivative): True
CHECK chi_plateau: same as formal: True
# formal           dim (SA cap C_reg) = 26
# chi_plateau      dim (SA cap C_reg) = 26
# double_plateau   dim (SA cap C_reg) = 1208
CHECK Helmholtz certificate A (formal, D=0): Euler derivative of I_C = (1/2) int <h,E_C[h]> equals E_C: True
CHECK Certificate A K-matrix satisfies CORRECT Gate-V identity w_mn K^{mn}_ab = w_ab K^{ab}_mn: True
CHECK Certificate A K-matrix FAILS the reversed identity w_ab K^{mn}_ab = w_mn K^{ab}_mn (as required): True
CHECK Certificate A numeric check: w00 K00_02 = -56/3 = w02 K02_00: True
CHECK Helmholtz certificate B (double plateau, D=4, fourth-order operator): Euler(I_C) = E_C: True
CHECK Helmholtz certificate B2 (double plateau, D=2, second-order operator): Euler(I_C) = E_C: True
CHECK negative control: for a non-self-adjoint collapse E_C, Euler(I_C) != E_C: True
# wrote helmholtz_certificates.txt
# --- certificate A (formal, D=0): explicit matched quadratic functional ---
#   C[N=4, b={(),0222}, gamma=()] = 1
#   E^{00} = -56*h_02_*s_*x_/3
#   E^{02} = -28*h_00_*s_*x_/3 - 28*h_22_*s_*x_/3
#   E^{22} = -56*h_02_*s_*x_/3
#   I_C = int d^4y L ,  L = -56*h_00_*h_02_*s_*x_/3 - 56*h_02_*h_22_*s_*x_/3
# --- certificate B (double plateau, D=4): 4 E-terms ; B2 (D=2): 4 E-terms (see certificates file)
# === SUMMARY ===
# formal           {'dim_SA': 26, 'perD': {0: 26, 1: 0, 2: 0, 3: 0, 4: 0}, 'dim_SA_cap_Creg': 26, 'dim_variational_image': 26, 'variational_perD': {0: 26, 1: 0, 2: 0, 3: 0, 4: 0}}
# chi_plateau      {'dim_SA': 26, 'perD': {0: 26, 1: 0, 2: 0, 3: 0, 4: 0}, 'dim_SA_cap_Creg': 26, 'dim_variational_image': 26, 'variational_perD': {0: 26, 1: 0, 2: 0, 3: 0, 4: 0}}
# double_plateau   {'dim_SA': 1347, 'perD': {0: 26, 1: 0, 2: 271, 3: 12, 4: 1038}, 'dim_SA_cap_Creg': 1208, 'dim_variational_image': 1325, 'variational_perD': {0: 26, 1: 0, 2: 271, 3: 8, 4: 1020}}
# total runtime 124.6s
ALL CHECKS PASSED
\end{lstlisting}

\paragraph{M.4/M.5 matching probes (\texttt{matching\_probe\_out.txt}).}
\begin{lstlisting}
# matching_probe.py output   (exact QQ; conditional on (S-bv) through M.2)
CHECK target construction: linearized Bianchi identity d^mu G_lin_{mu nu} = 0 in jets: True
CHECK formal/P1 target has Lorentz-reflection parity 0: True
CHECK formal/P1 target is self-adjoint under Gate V (Euler derivative): True
CHECK formal/P2 target has Lorentz-reflection parity 0: True
CHECK formal/P2 target is self-adjoint under Gate V (Euler derivative): True
CHECK formal/C_coh: explicit preimage of the intersection direction ['1', '-89/140'] reproduces it exactly: True
# formal          C_coh D=0 block(parity 0) size=10 targets=['P1', 'P2'] : dim(span cap im) = 1 ; members: {'P1': False, 'P2': False} ; direction(s) in target coords: [['1', '-89/140']] (block wall 0.0s; probe elapsed 0.0s)
CHECK formal/C_reg: explicit preimage of the intersection direction ['1', '-89/140'] reproduces it exactly: True
# formal          C_reg D=0 block(parity 0) size=10 targets=['P1', 'P2'] : dim(span cap im) = 1 ; members: {'P1': False, 'P2': False} ; direction(s) in target coords: [['1', '-89/140']] (block wall 0.0s; probe elapsed 0.0s)
CHECK chi_plateau/P1 target has Lorentz-reflection parity 0: True
CHECK chi_plateau/P1 target is self-adjoint under Gate V (Euler derivative): True
CHECK chi_plateau/P2 target has Lorentz-reflection parity 0: True
CHECK chi_plateau/P2 target is self-adjoint under Gate V (Euler derivative): True
CHECK chi_plateau/C_coh: explicit preimage of the intersection direction ['1', '-89/140'] reproduces it exactly: True
# chi_plateau     C_coh D=0 block(parity 0) size=10 targets=['P1', 'P2'] : dim(span cap im) = 1 ; members: {'P1': False, 'P2': False} ; direction(s) in target coords: [['1', '-89/140']] (block wall 0.0s; probe elapsed 0.0s)
CHECK chi_plateau/C_reg: explicit preimage of the intersection direction ['1', '-89/140'] reproduces it exactly: True
# chi_plateau     C_reg D=0 block(parity 0) size=10 targets=['P1', 'P2'] : dim(span cap im) = 1 ; members: {'P1': False, 'P2': False} ; direction(s) in target coords: [['1', '-89/140']] (block wall 0.0s; probe elapsed 0.0s)
CHECK double_plateau/P1 target has Lorentz-reflection parity 0: True
CHECK double_plateau/P1 target is self-adjoint under Gate V (Euler derivative): True
CHECK double_plateau/P2 target has Lorentz-reflection parity 0: True
CHECK double_plateau/P2 target is self-adjoint under Gate V (Euler derivative): True
CHECK double_plateau/C_coh: explicit preimage of the intersection direction ['1', '-89/140'] reproduces it exactly: True
# double_plateau  C_coh D=0 block(parity 0) size=10 targets=['P1', 'P2'] : dim(span cap im) = 1 ; members: {'P1': False, 'P2': False} ; direction(s) in target coords: [['1', '-89/140']] (block wall 0.0s; probe elapsed 0.0s)
CHECK double_plateau/C_reg: explicit preimage of the intersection direction ['1', '-89/140'] reproduces it exactly: True
# double_plateau  C_reg D=0 block(parity 0) size=10 targets=['P1', 'P2'] : dim(span cap im) = 1 ; members: {'P1': False, 'P2': False} ; direction(s) in target coords: [['1', '-89/140']] (block wall 0.0s; probe elapsed 0.0s)
CHECK double_plateau/EH target has Lorentz-reflection parity 0: True
CHECK double_plateau/EH target is self-adjoint under Gate V (Euler derivative): True
# double_plateau  C_coh D=2 block(parity 0) size=84 targets=['EH'] : dim(span cap im) = 0 ; members: {'EH': False} ; direction(s) in target coords: None (block wall 0.0s; probe elapsed 0.2s)
# double_plateau  C_reg D=2 block(parity 0) size=80 targets=['EH'] : dim(span cap im) = 0 ; members: {'EH': False} ; direction(s) in target coords: None (block wall 0.0s; probe elapsed 0.3s)
CHECK double_plateau/R2 target has Lorentz-reflection parity 0: True
CHECK double_plateau/R2 target is self-adjoint under Gate V (Euler derivative): True
CHECK double_plateau/RC2 target has Lorentz-reflection parity 0: True
CHECK double_plateau/RC2 target is self-adjoint under Gate V (Euler derivative): True
CHECK double_plateau/C_coh/R2: explicit preimage C reproduces the target Euler tensor exactly: True
CHECK double_plateau/C_coh/RC2: explicit preimage C reproduces the target Euler tensor exactly: True
CHECK double_plateau/C_coh: explicit preimage of the intersection direction ['1', '0'] reproduces it exactly: True
# double_plateau  C_coh D=4 block(parity 0) size=292 targets=['R2', 'RC2'] : dim(span cap im) = 2 ; members: {'R2': True, 'RC2': True} ; direction(s) in target coords: [['1', '0'], ['0', '1']] (block wall 38.6s; probe elapsed 40.0s)
CHECK double_plateau/C_reg: explicit preimage of the intersection direction ['1', '-188813/88456'] reproduces it exactly: True
# double_plateau  C_reg D=4 block(parity 0) size=265 targets=['R2', 'RC2'] : dim(span cap im) = 1 ; members: {'R2': False, 'RC2': False} ; direction(s) in target coords: [['1', '-188813/88456']] (block wall 5.9s; probe elapsed 45.9s)
# wrote matching_certificates.txt
# === SUMMARY (conditional on (S-bv)) ===
# formal          C_coh  P1+P2    D=0 : intersection dim 1 ; member: {'P1': False, 'P2': False} ; direction [['1', '-89/140']]
# formal          C_reg  P1+P2    D=0 : intersection dim 1 ; member: {'P1': False, 'P2': False} ; direction [['1', '-89/140']]
# chi_plateau     C_coh  P1+P2    D=0 : intersection dim 1 ; member: {'P1': False, 'P2': False} ; direction [['1', '-89/140']]
# chi_plateau     C_reg  P1+P2    D=0 : intersection dim 1 ; member: {'P1': False, 'P2': False} ; direction [['1', '-89/140']]
# double_plateau  C_coh  P1+P2    D=0 : intersection dim 1 ; member: {'P1': False, 'P2': False} ; direction [['1', '-89/140']]
# double_plateau  C_reg  P1+P2    D=0 : intersection dim 1 ; member: {'P1': False, 'P2': False} ; direction [['1', '-89/140']]
# double_plateau  C_coh  EH       D=2 : intersection dim 0 ; member: {'EH': False} ; direction None
# double_plateau  C_reg  EH       D=2 : intersection dim 0 ; member: {'EH': False} ; direction None
# double_plateau  C_coh  R2+RC2   D=4 : intersection dim 2 ; member: {'R2': True, 'RC2': True} ; direction [['1', '0'], ['0', '1']]
# double_plateau  C_reg  R2+RC2   D=4 : intersection dim 1 ; member: {'R2': False, 'RC2': False} ; direction [['1', '-188813/88456']]
CHECK double plateau: linearized R^2 and Ric^2 Euler tensors are collapse images (explicit preimages certified): True
CHECK double plateau: linearized Einstein tensor is NOT a collapse image: True
# total wall-clock runtime 46.4s (sum of sequential block walls; not parallel)
ALL CHECKS PASSED
\end{lstlisting}

\paragraph{Stage 2 fixed-data (\texttt{fixed\_data\_reading\_out.txt}).}
\begin{lstlisting}
# fixed_data_reading_out.txt -- Stage 2 summary (conditional on S-bv)
# Imported: sigma_eps = XII-C/XVI-L time-only ramp; chi_I = compact switch; Gate V unchanged
# S-bv: CONDITIONAL (not discharged)

STATUS: fixed-data criteria and universal region comparisons constructed; actual fixed-data dimensions unresolved
REGION 1 (full switched/prepared): status = sufficient criteria (i) G, (ii) Coll, (iii) adjoint-defect injectivity of ev_U
  ker G_formal subseteq ker G_fixed; (iii) needed for SA equality; time-only ramp: own Helmholtz analysis needed (Unresolved)
  comparison class formal: dim ker Coll = 0
REGION 2 (prep plateau, UNIVERSAL formal chi-jets): status = EXACT for the universal reading (actual chi_I: Unresolved)
  dim ker Coll = 21 ; SA_D0 = 26
  full SA dim = 47 ; variational image = 26 ; perD SA = {'0': 26, '1': 0, '2': 0, '3': 4, '4': 17}
REGION 3 (double plateau): status = EXACT (M.4)
  dim ker Coll = 22 ; SA = 1347 ; var image = 1325
REGION 4 (transitions): status = Unresolved (domain-size bound 0 <= dim ker Coll_fixed <= 3321)

DISTINCTIONS:
  support vanishing: C_2/M_2 die on {sigma=1}; chi-Leibniz dies on {chi=1}
  additional relations: possible for concrete ramp; none certified
  genuine SA: Gate V self-adjointness of collapsed operator

ULTRALOCAL: 26-dim family; 26 K-patterns in SA basis; Lorentz-scalar matched line: P1-(89/140)P2
ACCEPTANCE: criteria / exact-universal / exact / unresolved for regions 1..4
total runtime 120.3s
ALL CHECKS PASSED
\end{lstlisting}

\paragraph{M.5 ledger (\texttt{counterterm\_functionals\_out.txt}).}
\begin{lstlisting}
# counterterm_functionals.py output   (M.5 ledger; conditional on (S-bv))
CHECK matching_certificates.txt present: True
CHECK expected matching preimages from M.4 probes: True
CHECK 26-dim ultralocal SA basis: True
CHECK P1/P2 matching row present: True
CHECK R2 C_coh preimage recorded as non-regression-preserving (K-block): True
CHECK C_reg R2/RC2 combination is regression-preserving: True
CHECK completion gate obstruction recorded: True
CHECK EOM ledger empty because EOM undefined: True
CHECK field-redef ledger empty because undefined: True
# wrote matching_matrix.json/.txt matching_certificates_M5.txt ledger_*.json ledger_tables.txt
# wrote counterterm_functionals_out.txt m5_ledger_summary.json
# Named theorem: PlateauCurvatureSquareMatching
# total runtime 0.1s
ALL CHECKS PASSED
\end{lstlisting}

\paragraph{M.6 scaling anomaly (\texttt{scaling\_anomaly\_out.txt}).}
\begin{lstlisting}
# scaling_anomaly.py output
# Stage 4 / M.6: dilation identity, S^2 residue, anomaly assignment, collapse probe
# CONDITIONAL: residue evaluation uses Input (S-bv); Gate S-bv not discharged
CHECK dilation identity: mu d/dmu varpi_mu = E varpi_mu (rho=0): True
# identity (N=0): A_mu(u) = (E+4) R_{L,mu}(u) = - <u, E varpi_mu> delta
CHECK N=0 structural identity A_mu = (E+4) R_{L,mu} recorded (analytic; see M.6): True
CHECK contour residue at z_m tends to -1: True
CHECK angular u-integral Iu0 = -2 i pi: True
CHECK sphere integral of (Q-i0)^{-2} = -2 i pi^2: True
# CONDITIONAL ON (S-bv): int_{S^3} (Q-i0)^{-2} d sigma = -2 i pi^2
CHECK S^2 residue coefficient r = -2 i pi^2 c^2: True
CHECK S^2 residue nonzero for every c != 0 (S-bv): True
# CONDITIONAL ON (S-bv): A_mu(S^2) = r delta with r = -2 i pi^2 c^2 != 0
# Residue independent of mu and of radial profile of varpi (varpi(0)=1)
# Higher-N residues (N>=1): Unresolved in this dispatch (formula recorded in M.6)
CHECK N=0 coherent basis is the single S^2 label {(),()}: True
CHECK anomaly assignment sits at contact variable index 0 (N=0,gamma=()): True
# N=0 anomaly component C_anom: C[0] = r; higher-grade (N>=1) coordinates explicitly set to ZERO for this probe
# (higher-N residues Unresolved: C_anom is NOT the full graph scale variation)
# collapse[formal]: nnz=744 sa_defect_nnz=4602 variational_nonzero=False
# collapse[chi_plateau]: nnz=744 sa_defect_nnz=1938 variational_nonzero=False
# collapse[prep_plateau]: nnz=0 sa_defect_nnz=0 variational_nonzero=False
# collapse[double_plateau]: nnz=0 sa_defect_nnz=0 variational_nonzero=False
CHECK formal: S^2 anomaly collapse nonzero: True
CHECK formal: S^2 anomaly collapse NOT self-adjoint: True
CHECK formal: S^2 anomaly is NOT a nonzero variational tensor: True
CHECK chi-plateau: S^2 anomaly collapse nonzero: True
CHECK chi-plateau: S^2 anomaly NOT self-adjoint: True
CHECK prep-plateau: S^2 anomaly collapses to 0 (sigma-derivatives): True
CHECK double-plateau: S^2 anomaly collapses to 0 (matching domain): True
CHECK S^2 coefficient ahatbar supported only on |partial sigma|=2 (C_2 sector): True
CHECK double-plateau R2 and RC2 Euler tensors are linearly independent: True
CHECK double-plateau Collapse(C_anom)=0, so matched projection = 0: True
# Matched projection onto span(E_R2, E_RC2) in double plateau: (0,0)
# Unmatched remainder in double plateau: 0 (collapse itself vanishes)
# Unmatched remainder in formal reading: full Collapse(C_anom) retained (non-variational)
CHECK M-C not claimable: the full controlled matched scaling calculation is unfinished: True
# M-C remains incomplete because the full controlled matched scaling calculation is unfinished.
# The N=0 component alone contributes no nonzero plateau matching.
# wrote scaling_anomaly_cert.json
# wrote scaling_anomaly_cert.txt
ALL CHECKS PASSED
\end{lstlisting}

\paragraph{M.7 covariance (\texttt{covariance\_constraints\_out.txt}, \texttt{covariance\_matrix.txt}).}
\begin{lstlisting}
# covariance_constraints.py output  (AUDITED REVISION)
# Stage 5 / M.7: Lorentz action, equivariant contacts, exact SA intersections, linear divergence identities
# CONDITIONAL on Input (S-bv); Gate S-bv not discharged; M-C/M-D not claimed
# Prescription covariance != interacting Ward/BV; linear divergence identity != Ward identity
CHECK d_N=(1,4,10,20,35) under (S-bv): True
CHECK dim C_coh = 3321: True
CHECK W dims and C_coh: True
CHECK generator J12 is eta-skew: True
CHECK generator J13 is eta-skew: True
CHECK generator J23 is eta-skew: True
CHECK generator K01 is eta-skew: True
CHECK generator K02 is eta-skew: True
CHECK generator K03 is eta-skew: True
CHECK p_0=1: True
CHECK B_N basis of Poly_N, N=0: True
CHECK B_N basis of Poly_N, N=1: True
CHECK B_N basis of Poly_N, N=2: True
CHECK B_N basis of Poly_N, N=3: True
CHECK B_N basis of Poly_N, N=4: True
CHECK N=0: [rho(K01),rho(K02)]=rho(J12) on B_N: True
CHECK N=1: [rho(K01),rho(K02)]=rho(J12) on B_N: True
CHECK N=2: [rho(K01),rho(K02)]=rho(J12) on B_N: True
CHECK N=3: [rho(K01),rho(K02)]=rho(J12) on B_N: True
CHECK N=4: [rho(K01),rho(K02)]=rho(J12) on B_N: True
CHECK N=0 trivial action: True
CHECK N=0 J12: pull-back action on d_gamma delta = P rho^poly P (eta intertwiner): True
CHECK N=0 J13: pull-back action on d_gamma delta = P rho^poly P (eta intertwiner): True
CHECK N=0 J23: pull-back action on d_gamma delta = P rho^poly P (eta intertwiner): True
CHECK N=0 K01: pull-back action on d_gamma delta = P rho^poly P (eta intertwiner): True
CHECK N=0 K02: pull-back action on d_gamma delta = P rho^poly P (eta intertwiner): True
CHECK N=0 K03: pull-back action on d_gamma delta = P rho^poly P (eta intertwiner): True
CHECK N=0: [rhoD(K01),rhoD(K02)]=rhoD(J12) on d_gamma delta: True
CHECK N=1 J12: pull-back action on d_gamma delta = P rho^poly P (eta intertwiner): True
CHECK N=1 J13: pull-back action on d_gamma delta = P rho^poly P (eta intertwiner): True
CHECK N=1 J23: pull-back action on d_gamma delta = P rho^poly P (eta intertwiner): True
CHECK N=1 K01: pull-back action on d_gamma delta = P rho^poly P (eta intertwiner): True
CHECK N=1 K02: pull-back action on d_gamma delta = P rho^poly P (eta intertwiner): True
CHECK N=1 K03: pull-back action on d_gamma delta = P rho^poly P (eta intertwiner): True
CHECK N=1: [rhoD(K01),rhoD(K02)]=rhoD(J12) on d_gamma delta: True
CHECK N=2 J12: pull-back action on d_gamma delta = P rho^poly P (eta intertwiner): True
CHECK N=2 J13: pull-back action on d_gamma delta = P rho^poly P (eta intertwiner): True
CHECK N=2 J23: pull-back action on d_gamma delta = P rho^poly P (eta intertwiner): True
CHECK N=2 K01: pull-back action on d_gamma delta = P rho^poly P (eta intertwiner): True
CHECK N=2 K02: pull-back action on d_gamma delta = P rho^poly P (eta intertwiner): True
CHECK N=2 K03: pull-back action on d_gamma delta = P rho^poly P (eta intertwiner): True
CHECK N=2: [rhoD(K01),rhoD(K02)]=rhoD(J12) on d_gamma delta: True
CHECK N=3 J12: pull-back action on d_gamma delta = P rho^poly P (eta intertwiner): True
CHECK N=3 J13: pull-back action on d_gamma delta = P rho^poly P (eta intertwiner): True
CHECK N=3 J23: pull-back action on d_gamma delta = P rho^poly P (eta intertwiner): True
CHECK N=3 K01: pull-back action on d_gamma delta = P rho^poly P (eta intertwiner): True
CHECK N=3 K02: pull-back action on d_gamma delta = P rho^poly P (eta intertwiner): True
CHECK N=3 K03: pull-back action on d_gamma delta = P rho^poly P (eta intertwiner): True
CHECK N=3: [rhoD(K01),rhoD(K02)]=rhoD(J12) on d_gamma delta: True
CHECK N=4 J12: pull-back action on d_gamma delta = P rho^poly P (eta intertwiner): True
CHECK N=4 J13: pull-back action on d_gamma delta = P rho^poly P (eta intertwiner): True
CHECK N=4 J23: pull-back action on d_gamma delta = P rho^poly P (eta intertwiner): True
CHECK N=4 K01: pull-back action on d_gamma delta = P rho^poly P (eta intertwiner): True
CHECK N=4 K02: pull-back action on d_gamma delta = P rho^poly P (eta intertwiner): True
CHECK N=4 K03: pull-back action on d_gamma delta = P rho^poly P (eta intertwiner): True
CHECK N=4: [rhoD(K01),rhoD(K02)]=rhoD(J12) on d_gamma delta: True
CHECK first-order basis order d_0..d_3: True
CHECK first order J12: X.(d_nu delta) = X^mu_nu d_mu delta (abstract action matrix = X): True
CHECK first order J12: old coordinate matrix = -X^T: True
CHECK first order J12 (rotation): abstract action = old coordinate matrix: True
CHECK first order J12: abstract = eta (old) eta: True
CHECK first order J13: X.(d_nu delta) = X^mu_nu d_mu delta (abstract action matrix = X): True
CHECK first order J13: old coordinate matrix = -X^T: True
CHECK first order J13 (rotation): abstract action = old coordinate matrix: True
CHECK first order J13: abstract = eta (old) eta: True
CHECK first order J23: X.(d_nu delta) = X^mu_nu d_mu delta (abstract action matrix = X): True
CHECK first order J23: old coordinate matrix = -X^T: True
CHECK first order J23 (rotation): abstract action = old coordinate matrix: True
CHECK first order J23: abstract = eta (old) eta: True
CHECK first order K01: X.(d_nu delta) = X^mu_nu d_mu delta (abstract action matrix = X): True
CHECK first order K01: old coordinate matrix = -X^T: True
CHECK first order K01 (boost): abstract action = -(old coordinate matrix): OLD MATRIX WRONG FOR BOOSTS: True
CHECK first order K01: abstract = eta (old) eta: True
CHECK first order K02: X.(d_nu delta) = X^mu_nu d_mu delta (abstract action matrix = X): True
CHECK first order K02: old coordinate matrix = -X^T: True
CHECK first order K02 (boost): abstract action = -(old coordinate matrix): OLD MATRIX WRONG FOR BOOSTS: True
CHECK first order K02: abstract = eta (old) eta: True
CHECK first order K03: X.(d_nu delta) = X^mu_nu d_mu delta (abstract action matrix = X): True
CHECK first order K03: old coordinate matrix = -X^T: True
CHECK first order K03 (boost): abstract action = -(old coordinate matrix): OLD MATRIX WRONG FOR BOOSTS: True
CHECK first order K03: abstract = eta (old) eta: True
# first-order boost K01 on (d0 delta,d1 delta,d2 delta,d3 delta): abstract = [['0', '1', '0', '0'], ['1', '0', '0', '0'], ['0', '0', '0', '0'], ['0', '0', '0', '0']]
#                                                 old coordinate matrix = [['0', '-1', '0', '0'], ['-1', '0', '0', '0'], ['0', '0', '0', '0'], ['0', '0', '0', '0']]
CHECK symbolic test function: K01.(d_1 delta) = + d_0 delta (old matrix gave - d_0 delta): True
CHECK Poly_N = sum H_{N-2k} dims, N=0: True
CHECK Poly_N = sum H_{N-2k} dims, N=1: True
CHECK Poly_N = sum H_{N-2k} dims, N=2: True
CHECK Poly_N = sum H_{N-2k} dims, N=3: True
CHECK Poly_N = sum H_{N-2k} dims, N=4: True
CHECK N=0: number of phi_{k,j} = Schur count 1: True
CHECK N=0: corrected equivariant maps linearly independent: True
# N=0: corrected C = P o phi_{k,j} equivariant on all 6 generators (exact); maps [(0, 0)]
CHECK N=1: number of phi_{k,j} = Schur count 1: True
CHECK N=1: corrected equivariant maps linearly independent: True
# N=1: corrected C = P o phi_{k,j} equivariant on all 6 generators (exact); maps [(0, 0)]
CHECK N=2: number of phi_{k,j} = Schur count 3: True
CHECK N=2: corrected equivariant maps linearly independent: True
# N=2: corrected C = P o phi_{k,j} equivariant on all 6 generators (exact); maps [(0, 0), (1, 0), (1, 1)]
CHECK N=3: number of phi_{k,j} = Schur count 3: True
CHECK N=3: corrected equivariant maps linearly independent: True
# N=3: corrected C = P o phi_{k,j} equivariant on all 6 generators (exact); maps [(0, 0), (1, 0), (1, 1)]
CHECK N=4: number of phi_{k,j} = Schur count 6: True
CHECK N=4: corrected equivariant maps linearly independent: True
# N=4: corrected C = P o phi_{k,j} equivariant on all 6 generators (exact); maps [(0, 0), (1, 0), (1, 1), (2, 0), (2, 1), (2, 2)]
CHECK N=0 old vector trivially equivariant: True
CHECK OLD (uncorrected) vectors: those changed by P fail exactly the three boost equations; unchanged ones (image only at even spatial order) pass: True
# old vectors failing boost equivariance: ['(1, 0, 0)', '(2, 0, 0)', '(3, 0, 0)', '(3, 1, 0)', '(3, 1, 1)', '(4, 0, 0)', '(4, 1, 0)', '(4, 1, 1)'] ; unchanged by P: ['(0, 0, 0)', '(2, 1, 0)', '(2, 1, 1)', '(4, 2, 0)', '(4, 2, 1)', '(4, 2, 2)']
CHECK equivariant dims 1,1,3,3,6: True
CHECK N=0: exact kernel of the equivariance constraint map (pull-back action) has dim 1, rank 0: True
CHECK N=1: exact kernel of the equivariance constraint map (pull-back action) has dim 1, rank 19: True
CHECK N=2: exact kernel of the equivariance constraint map (pull-back action) has dim 3, rank 147: True
# cokernel = dim codomain - rank, codomain = (+)_6 Hom(H_{4+N},W_N): {0: 6, 1: 101, 2: 753, 3: 3503, 4: 12256}
CHECK equivariant freedom 14 = 5 (K, N<=2) + 9 (C_reg, N>=3): True
CHECK P1, P2, P1-(89/140)P2 are so(1,3)-invariant: True
# --- exact V = Equiv cap SA (14 equivariant coordinates; adjoint defect map; exact QQ kernel) ---
CHECK formal: all 1 basis vectors of V re-verified from scratch: zero adjoint defect AND equivariance (6 generators, all N): True
CHECK formal: dim(V cap K)=dimV-dim pi_reg(V), dim(V cap C_reg)=dimV-dim pi_K(V): True
CHECK formal: dim(V cap K)+dim(V cap C_reg)=1 <= dim V=1: True
CHECK formal: V cap K <= pi_K(V), V cap C_reg <= pi_reg(V): True
CHECK formal: pi_K(V) cap SA = V cap K and pi_reg(V) cap SA = V cap C_reg (projections NOT relabelled as SA): True
# formal          dim V=1 | pi_K(V)=0  V cap K=0 | pi_reg(V)=1  V cap C_reg=1 | dim Coll(V)=1 (2.6s)
CHECK chi_plateau: all 1 basis vectors of V re-verified from scratch: zero adjoint defect AND equivariance (6 generators, all N): True
CHECK chi_plateau: dim(V cap K)=dimV-dim pi_reg(V), dim(V cap C_reg)=dimV-dim pi_K(V): True
CHECK chi_plateau: dim(V cap K)+dim(V cap C_reg)=1 <= dim V=1: True
CHECK chi_plateau: V cap K <= pi_K(V), V cap C_reg <= pi_reg(V): True
CHECK chi_plateau: pi_K(V) cap SA = V cap K and pi_reg(V) cap SA = V cap C_reg (projections NOT relabelled as SA): True
# chi_plateau     dim V=1 | pi_K(V)=0  V cap K=0 | pi_reg(V)=1  V cap C_reg=1 | dim Coll(V)=1 (0.7s)
CHECK double_plateau: all 12 basis vectors of V re-verified from scratch: zero adjoint defect AND equivariance (6 generators, all N): True
CHECK double_plateau: dim(V cap K)=dimV-dim pi_reg(V), dim(V cap C_reg)=dimV-dim pi_K(V): True
CHECK double_plateau: dim(V cap K)+dim(V cap C_reg)=10 <= dim V=12: True
CHECK double_plateau: V cap K <= pi_K(V), V cap C_reg <= pi_reg(V): True
CHECK double_plateau: pi_K(V) cap SA = V cap K and pi_reg(V) cap SA = V cap C_reg (projections NOT relabelled as SA): True
# double_plateau  dim V=12 | pi_K(V)=5  V cap K=3 | pi_reg(V)=9  V cap C_reg=7 | dim Coll(V)=9 (0.9s)
CHECK formal and chi-plateau: dim(Equiv cap SA) = 1: True
CHECK formal: the SA equivariant line lies in C_reg (N=4): True
CHECK formal: Coll(V) = span of E[P1-(89/140)P2] (ratio 17920); V = phi_{2,0} at N=4 (Box^2, gamma=0): True
CHECK formal: supplementary mod-p kernel dims [1, 1, 1] agree with the exact dim V (correct rational reduction): True
CHECK chi_plateau: supplementary mod-p kernel dims [1, 1, 1] agree with the exact dim V (correct rational reduction): True
CHECK double_plateau: supplementary mod-p kernel dims [12, 12, 12] agree with the exact dim V (correct rational reduction): True
# --- linear (gauge) divergence identities: d_mu E^{mu nu}, E with UPPER indices, no metric factor ---
CHECK Identity A: eta^{mu rho} d_rho G^lin_{mu nu} = 0 (G LOWER indices, metric factor correct here): True
CHECK double_plateau: d_mu E^{mu nu}[R2] == 0 identically (weight w=1): True
CHECK double_plateau: d_mu E^{mu nu}[RC2] == 0 identically (weight w=1): True
CHECK double_plateau: d_mu E^{mu nu}[R2-(188813/88456)RC2] == 0 identically (weight w=1): True
CHECK chi_plateau: d_mu E^{mu nu}[R2] == 0 identically (weight w=s_): True
CHECK chi_plateau: d_mu E^{mu nu}[RC2] == 0 identically (weight w=s_): True
CHECK chi_plateau: d_mu E^{mu nu}[R2-(188813/88456)RC2] == 0 identically (weight w=s_): True
CHECK formal: d_mu E^{mu nu}[R2] == 0 identically (weight w=s_*x_): True
CHECK formal: d_mu E^{mu nu}[RC2] == 0 identically (weight w=s_*x_): True
CHECK formal: d_mu E^{mu nu}[R2-(188813/88456)RC2] == 0 identically (weight w=s_*x_): True
CHECK diagnosis: the OLD formula sum_mu eta^{mu mu} d_mu E^{mu nu} is nonzero (this produced the false 'does not vanish'): True
CHECK independent metric linearization: R_lin(metric) = +1 x R_lin(manuscript formula): True
CHECK independent metric linearization: Ric_lin(metric) = +1 x Ric_lin(manuscript formula): True
CHECK independent route: d_mu E^{mu nu}[R2] == 0 (Functions, euler_equations): True
CHECK independent route: old eta^{mu mu}-weighted formula nonzero for R2: True
CHECK independent Euler tensor of int R2 equals the matching_probe/jet_euler target exactly: True
CHECK independent route: d_mu E^{mu nu}[RC2] == 0 (Functions, euler_equations): True
CHECK independent route: old eta^{mu mu}-weighted formula nonzero for RC2: True
CHECK independent Euler tensor of int RC2 equals the matching_probe/jet_euler target exactly: True
# independent curvature route done (13.8s); sign(R)=+1 sign(Ric)=+1 (squares are sign-insensitive)
CHECK preimage [double_plateau C_coh R2]: Coll(C) equals the INDEPENDENTLY built target exactly: True
CHECK preimage [double_plateau C_coh R2]: d_mu E_C^{mu nu} == 0: True
CHECK preimage [double_plateau C_coh R2]: support meets K-block (N<=2): True
CHECK preimage [double_plateau C_coh RC2]: Coll(C) equals the INDEPENDENTLY built target exactly: True
CHECK preimage [double_plateau C_coh RC2]: d_mu E_C^{mu nu} == 0: True
CHECK preimage [double_plateau C_coh RC2]: support meets K-block (N<=2): True
CHECK preimage [double_plateau C_reg (1)R2+(-188813/88456)RC2]: Coll(C) equals the INDEPENDENTLY built target exactly: True
CHECK preimage [double_plateau C_reg (1)R2+(-188813/88456)RC2]: d_mu E_C^{mu nu} == 0: True
CHECK preimage [double_plateau C_reg (1)R2+(-188813/88456)RC2]: support inside C_reg (N>=3): True
CHECK double_plateau: EH identity d_mu E^{mu nu} = (d_mu w) G_lin^{mu nu}[h]  (zero iff w constant): True
CHECK chi_plateau: EH identity d_mu E^{mu nu} = (d_mu w) G_lin^{mu nu}[h]  (zero iff w constant): True
CHECK formal: EH identity d_mu E^{mu nu} = (d_mu w) G_lin^{mu nu}[h]  (zero iff w constant): True
CHECK double_plateau: ultralocal P1-(89/140)P2: d_mu E^{mu nu} = 2 d_mu[w(h^{mu nu}-(89/140)eta^{mu nu} h)] != 0 (not gauge invariant): True
CHECK chi_plateau: ultralocal P1-(89/140)P2: d_mu E^{mu nu} = 2 d_mu[w(h^{mu nu}-(89/140)eta^{mu nu} h)] != 0 (not gauge invariant): True
CHECK formal: ultralocal P1-(89/140)P2: d_mu E^{mu nu} = 2 d_mu[w(h^{mu nu}-(89/140)eta^{mu nu} h)] != 0 (not gauge invariant): True
# wrote covariance_matrix.txt covariance_certificates.txt covariance_summary.json covariance_exact_certificates.json completion_ledger_M89.json
CHECK M-C/M-D not claimed: True
# total runtime 241.8s
ALL CHECKS PASSED
# ---------------- covariance_matrix.txt ----------------
# covariance_matrix.txt -- XVI-M M.7 compatibility matrix (CONDITIONAL on (S-bv)), AUDITED REVISION
# Contact maps C: H_{4+N} -> W_N = span{d_gamma delta : |gamma|<=N}, pull-back action (eta intertwiner)
# rank = rank of the constraint map C -> (C rho_V(X) - rho_delta(X) C)_X = codim of the solution space
# codomain = (+)_{6 generators} Hom(H_{4+N},W_N), dim 6*slots ; coker = dim codomain - rank
#
N  dimH  dimW  slots  rank(=codim)  dim_codomain  dim_coker  dim_sol  block
0     1     1      1             0             6          6        1  K
1     4     5     20            19           120        101        1  K
2    10    15    150           147           900        753        3  K
3    20    35    700           697          4200       3503        3  C_reg
4    35    70   2450          2444         14700      12256        6  C_reg
#
K_block_dim_sol = 5
C_reg_dim_sol = 9
total_equivariant = 14
#
# Exact V = Equiv cap SA and the four objects:
V formal dim_V=1 pi_K(V)=0 V_cap_K=0 pi_reg(V)=1 V_cap_Creg=1 dim_Coll(V)=1
V chi_plateau dim_V=1 pi_K(V)=0 V_cap_K=0 pi_reg(V)=1 V_cap_Creg=1 dim_Coll(V)=1
V double_plateau dim_V=12 pi_K(V)=5 V_cap_K=3 pi_reg(V)=9 V_cap_Creg=7 dim_Coll(V)=9
#
# Linear divergence identities (d_mu E^{mu nu}, E upper indices):
ID_A eta^{mu rho} d_rho G^lin_{mu nu} = 0 : HOLDS (not Ward)
ID_B ultralocal P1-(89/140)P2 : d_mu E^{mu nu} = 2 d_mu[w(h^{mu nu}-(89/140)eta^{mu nu}h)] != 0 (not gauge invariant)
ID_Bprime R2, RC2, R2-(188813/88456)RC2 : d_mu E^{mu nu} = 0 in formal, chi-plateau, double plateau (linear gauge identity; not Ward)
ID_EH : d_mu E^{mu nu} = (d_mu w) G_lin^{mu nu} ; zero in double plateau
#
# Gate S-bv CONDITIONAL. M-C not claimed. M-D not claimed.
\end{lstlisting}
\paragraph{M.7 audit legacy reconstruction (\texttt{m7\_legacy\_reconstruction\_out.txt}).}
\begin{lstlisting}
# m7_legacy_reconstruction.py output (supplementary; modular ranks)
CHECK double_plateau SA basis 1347 loaded: True
CHECK double_plateau SA basis has 14639 non-integer rational entries (truncation-sensitive): True
# double_plateau  vectors=old       reduction=buggy_int_truncation G=K     dim(Equiv_G cap pi_G(SA)) at 3 primes = [5, 5, 5]
# double_plateau  vectors=old       reduction=buggy_int_truncation G=C_reg dim(Equiv_G cap pi_G(SA)) at 3 primes = [3, 3, 3]
# double_plateau  vectors=old       reduction=buggy_int_truncation G=all   dim(Equiv_G cap pi_G(SA)) at 3 primes = [5, 5, 5]
# double_plateau  vectors=old       reduction=correct_reduction    G=K     dim(Equiv_G cap pi_G(SA)) at 3 primes = [5, 5, 5]
# double_plateau  vectors=old       reduction=correct_reduction    G=C_reg dim(Equiv_G cap pi_G(SA)) at 3 primes = [4, 4, 4]
# double_plateau  vectors=old       reduction=correct_reduction    G=all   dim(Equiv_G cap pi_G(SA)) at 3 primes = [6, 6, 6]
# double_plateau  vectors=corrected reduction=buggy_int_truncation G=K     dim(Equiv_G cap pi_G(SA)) at 3 primes = [5, 5, 5]
# double_plateau  vectors=corrected reduction=buggy_int_truncation G=C_reg dim(Equiv_G cap pi_G(SA)) at 3 primes = [3, 3, 3]
# double_plateau  vectors=corrected reduction=buggy_int_truncation G=all   dim(Equiv_G cap pi_G(SA)) at 3 primes = [5, 5, 5]
# double_plateau  vectors=corrected reduction=correct_reduction    G=K     dim(Equiv_G cap pi_G(SA)) at 3 primes = [5, 5, 5]
# double_plateau  vectors=corrected reduction=correct_reduction    G=C_reg dim(Equiv_G cap pi_G(SA)) at 3 primes = [9, 9, 9]
# double_plateau  vectors=corrected reduction=correct_reduction    G=all   dim(Equiv_G cap pi_G(SA)) at 3 primes = [12, 12, 12]
CHECK verbatim old computation reproduces the previously reported 5 / 3 / 5: True
CHECK old vectors with correct reduction: dim(Equiv_old cap SA) = 6 (old report 5 was a truncation artefact): True
CHECK corrected vectors with correct reduction: dim(Equiv cap SA) = 12 (agrees with exact): True
# wrote m7_legacy_reconstruction.json
ALL CHECKS PASSED
\end{lstlisting}
\paragraph{Work package A (\texttt{wpA/kernel\_identification\_out.txt}).}
\begin{lstlisting}
# kernel_identification_checks.py output
# xvi-M wpA: Gate S-bv identification of S with c (Q-i0)^{-1}
# python 3.13.5, sympy 1.14.0, mpmath 1.3.0, numpy 2.5.3
# signature eta=diag(1,-1,-1,-1); massless; z=x-y

# --- 1. Radial Fourier integral for W_0 (symbolic) ---
CHECK int_0^oo e^{-tau p} sin(p r) dp = r/(tau^2+r^2) for tau>0: True
CHECK analytic continuation spot (tau,t,r)=(0.2,0.5,1.0), err=1.88e-05: True
CHECK analytic continuation spot (tau,t,r)=(0.1,-1.2,0.7), err=5.25e-05: True
CHECK analytic continuation spot (tau,t,r)=(0.05,2.0,1.5), err=4.50e-04: True
CHECK damped radial W equals -1/(4 pi^2 (Q - 2 i tau t - tau^2)): True
CHECK Euclidean slice z0=-i tE: W=1/(4 pi^2 (tE^2+r^2)) (matches rem:Sbv-motivation): True
# => W_0 = -1/(4 pi^2) * 1/(Q - i0 sgn(z^0)) as the boundary value of the damped integral

# --- 2. Sign check for (s-i0)^{-1} = PV(1/s) + i pi delta(s) ---
CHECK 1D Im~pi f(0) Gaussian#0, err=4.02e-04: True
CHECK 1D Re~PV Gaussian#0, err=0.00e+00: True
CHECK 1D Im~pi f(0) Gaussian#1, err=2.80e-04: True
CHECK 1D Re~PV Gaussian#1, err=3.38e-04: True
CHECK 1D Im~pi f(0) Gaussian#2, err=2.57e-04: True
CHECK 1D Re~PV Gaussian#2, err=2.87e-04: True
CHECK 1D Im~pi f(0) Gaussian#3, err=4.16e-05: True
CHECK 1D Re~PV Gaussian#3, err=7.97e-05: True
CHECK opposite side Im~-pi f(0): True
# confirmed: (s-i0)^{-1}=PV(1/s)+i pi delta(s); (s+i0)^{-1}=PV-i pi delta

# --- 3. Algebraic combination of singular terms (symbolic coefficients) ---
CHECK sgn(z0) dependence cancels: delta coeff same on both time sheets: True
CHECK common delta coeff equals c * i pi = i/(4 pi): True
CHECK PV coeff equals c = 1/(4 pi^2): True
# S = c (PV(1/Q) + i pi delta(Q)) = c (Q - i0)^{-1},  c = 1/(4 pi^2)

# --- 4. Causal propagator delta_0 from W - W^T ---
CHECK delta_0 = (1/(2 pi)) sgn(z0) delta(Q): True
CHECK E_0^A = (1/(2 pi)) theta(-z0) delta(Q) on Omega: True

# --- 5. Off-cone pointwise normalization (spacelike & timelike) ---
CHECK off-cone point#0 has Q!=0 (Q=-1.0): True
CHECK off-cone S=-W=c/Q at point#0 z=(0.0, 1.0, 0.0, 0.0), Q=-1: True
CHECK off-cone point#1 has Q!=0 (Q=-3.75): True
CHECK off-cone S=-W=c/Q at point#1 z=(0.5, 2.0, 0.0, 0.0), Q=-3.75: True
CHECK off-cone point#2 has Q!=0 (Q=-5.0): True
CHECK off-cone S=-W=c/Q at point#2 z=(1.0, 2.0, 1.0, 1.0), Q=-5: True
CHECK off-cone point#3 has Q!=0 (Q=3.75): True
CHECK off-cone S=-W=c/Q at point#3 z=(2.0, 0.5, 0.0, 0.0), Q=3.75: True
CHECK off-cone point#4 has Q!=0 (Q=7.0): True
CHECK off-cone S=-W=c/Q at point#4 z=(3.0, 1.0, 1.0, 0.0), Q=7: True
CHECK off-cone point#5 has Q!=0 (Q=3.5): True
CHECK off-cone S=-W=c/Q at point#5 z=(-2.0, 0.0, 0.5, 0.5), Q=3.5: True
CHECK off-cone point#6 has Q!=0 (Q=-2.4365534255462): True
CHECK off-cone S=-W=c/Q at point#6 z=(0.30471707975443135, -1.0399841062404955, 0.7504511958064572, 0.9405647163912139), Q=-2.43655: True
CHECK off-cone point#7 has Q!=0 (Q=1.994514293379158): True
CHECK off-cone S=-W=c/Q at point#7 z=(-1.9510351886538364, -1.302179506862318, 0.12784040316728537, -0.3162425923435822), Q=1.99451: True
CHECK off-cone point#8 has Q!=0 (Q=-2.105702756487829): True
CHECK off-cone S=-W=c/Q at point#8 z=(-0.016801157504288795, -0.85304392757358, 0.8793979748628286, 0.7777919354289483), Q=-2.1057: True
CHECK off-cone point#9 has Q!=0 (Q=-2.2232612075272082): True
CHECK off-cone S=-W=c/Q at point#9 z=(0.06603069756121605, 1.1272412069680329, 0.4675093422520456, -0.8592924628832382), Q=-2.22326: True
CHECK off-cone point#10 has Q!=0 (Q=-1.5576462298657727): True
CHECK off-cone S=-W=c/Q at point#10 z=(0.36875078408249884, -0.9588826008289989, 0.8784503013072725, -0.049925910986252896), Q=-1.55765: True
CHECK off-cone point#11 has Q!=0 (Q=-1.9479776365819612): True
CHECK off-cone S=-W=c/Q at point#11 z=(-0.18486236354526056, -0.6809295444039414, 1.2225413386740303, -0.15452948206880215), Q=-1.94798: True
CHECK spacelike unit: S=c/Q=-c: True
CHECK timelike unit: S=c/Q=c: True

# --- 6. Numerical spatial Fourier for W vs closed form (damped) ---
CHECK Fourier vs closed W at (z0,r)=(0.5,1.0), err=3.30e-05: True
CHECK Fourier vs closed W at (z0,r)=(-0.5,1.0), err=3.30e-05: True
CHECK Fourier vs closed W at (z0,r)=(1.5,0.8), err=3.30e-05: True
CHECK Fourier vs closed W at (z0,r)=(-1.2,2.0), err=3.31e-05: True
CHECK Fourier vs closed W at (z0,r)=(0.0,1.5), err=3.30e-05: True
CHECK Fourier vs closed W at (z0,r)=(2.0,3.0), err=3.32e-05: True
CHECK Fourier vs closed W at (z0,r)=(-2.5,1.1), err=3.32e-05: True
CHECK Fourier vs closed W at (z0,r)=(0.3,0.3), err=3.30e-05: True

# --- 7. Null-cone chart s=Q: full one-variable distribution ---
CHECK on {Q=0}\{0}, z0 != 0 (so sgn(z0) and theta(+/-) are smooth there): True
# chart density future: Im=0.709646 vs pi*rho(0)=0.709704, err_im=5.84e-05, err_re=1.14e-12
CHECK chart-future Im~pi rho(0), err=5.84e-05: True
CHECK chart-future Re~PV, err=1.14e-12: True
# chart density past: Im=0.709646 vs pi*rho(0)=0.709704, err_im=5.84e-05, err_re=1.14e-12
CHECK chart-past Im~pi rho(0), err=5.84e-05: True
CHECK chart-past Re~PV, err=1.14e-12: True
CHECK dQ != 0 on the nonzero null cone (Q is a local coordinate): True
# pullback of (s-i0)^{-1} therefore defines (Q-i0)^{-1} across the nonzero null cone
# combined with §3: S/c equals that single boundary value on both sheets

# --- 7b. Variable regulator s - tau^2 - i tau a(s), a of fixed sign (lem:wpA-varreg) ---
CHECK variable regulator a = 2 + sin s (z^0>0 sheet): |pairing - (PV + i pi sgn(a) phi(0))| = 1.20e-01, 1.24e-02, 1.24e-03 (tau=1e-2,1e-3,1e-4); decreasing: True
CHECK variable regulator a = -(1.5 + 0.5 cos s) (z^0<0 sheet): |pairing - (PV + i pi sgn(a) phi(0))| = 1.14e-01, 1.17e-02, 1.17e-03 (tau=1e-2,1e-3,1e-4); decreasing: True
CHECK regulator identity alpha^2 + r^2 = -(Q - 2 i tau z^0 - tau^2): variable coefficient a = 2 z^0: True
# => near each nonzero null point (z^0 of fixed sign) the damped kernel tends to (Q - i0 sgn z^0)^{-1}

# --- 8. Cross-checks: Box, homogeneity, transposition, M.6 residue ---
CHECK homogeneity degree -2 for 1/Q (off cone; BV inherits): True
CHECK Box(c/Q)=0 off the cone as rational identity (got 0): True
CHECK transposition: S^T=S (from W-W^T=-i delta_0 and E^R-E^A=delta_0): True
CHECK propagator norm: Box S = i delta on R^4 (hence 0 on Omega): True
CHECK M.6 residue with c=1/(4 pi^2): r = -i/(8 pi^2): True
# r = -I/(8*pi**2) = -i/(8 pi^2)
CHECK counterfactual Q+i0 would flip residue sign (excluded): True

# --- 9. Fourier-family vs position: radial closed-form consistency ---
CHECK off-cone -W_damped = c/(Q-2itau t-tau^2) at (z0,r)=(0.5,1.0), err=0.00e+00: True
CHECK off-cone -W_damped = c/(Q-2itau t-tau^2) at (z0,r)=(-0.5,1.0), err=0.00e+00: True
CHECK off-cone -W_damped = c/(Q-2itau t-tau^2) at (z0,r)=(1.5,0.8), err=3.47e-18: True
CHECK off-cone -W_damped = c/(Q-2itau t-tau^2) at (z0,r)=(-1.2,2.0), err=1.73e-18: True
CHECK off-cone -W_damped = c/(Q-2itau t-tau^2) at (z0,r)=(0.0,1.5), err=0.00e+00: True
CHECK off-cone -W_damped = c/(Q-2itau t-tau^2) at (z0,r)=(2.0,3.0), err=8.67e-19: True
CHECK off-cone -W_damped = c/(Q-2itau t-tau^2) at (z0,r)=(-2.5,1.1), err=8.67e-19: True
CHECK off-cone -W_damped = c/(Q-2itau t-tau^2) at (z0,r)=(0.3,0.3), err=7.04e-14: True
CHECK near-cone off-sheet agreement at t=1.05, rel=1.073e-02: True
CHECK near-cone off-sheet agreement at t=0.95, rel=9.229e-03: True
CHECK near-cone off-sheet agreement at t=-1.05, rel=3.024e-02: True
CHECK near-cone off-sheet agreement at t=-0.95, rel=2.974e-02: True
# Fourier/position consistency off cone confirmed via damped closed forms

# --- SUMMARY ---
{
  "outcome": "Certified",
  "identity": "S = (1/(4*pi^2)) * (Q - i0)^{-1}  on D'(R^4\\{0})",
  "c": "1/(4*pi^2)",
  "boundary_side": "Q-i0",
  "M6_residue_with_c": "-i/(8*pi^2)",
  "seeds": {
    "numpy_rng": [
      42,
      7,
      99,
      100
    ],
    "script_seed_doc": 20261006
  }
}
# ALL CHECKS PASSED
\end{lstlisting}

\paragraph{Work package B (\texttt{wpB/higher\_grade\_scaling\_out.txt}).}
\begin{lstlisting}
# higher_grade_scaling.py output  (XVI-M work package B; python 3.13.5, sympy 1.14.0)
# CONDITIONAL on Input (S-bv): S = c (Q-i0)^{-1}; c^2 symbolic. Gate V and single-switch collapse unchanged.

# === Section A: complete differentiation of the XVI-L subtraction formula (all N=0..4) ===
# M_{beta alpha}(lam) = binom(beta,alpha) lam^{|beta|-|alpha|} p_{beta-alpha},  p_delta = d^delta varpi(0)
CHECK A N=0 [symbolic]: M^{(rho,mu)} unit lower-triangular (Leibniz, varpi_mu(0)=1): True
CHECK A N=0 [symbolic]: M M^{-1} = 1 (exact): True
CHECK A N=0 [symbolic]: mu d(M^{-1})/dmu = -M^{-1}(mu dM/dmu)M^{-1} (exact): True
CHECK A N=0 [symbolic]: -M^{-1}(mu dM)M^{-1} = D M^{-1} - M^{-1} D, D=diag(|beta|) (dilation grading): True
CHECK A N=0 [symbolic]: mu d/dmu P_alpha (complete formula) = (E-|alpha|) P_alpha for all |alpha|<=N: True
CHECK A N=0 [symbolic]: w_alpha = (z^alpha/alpha!) T_{N-|alpha|}[1/varpi_mu] varpi_mu (closed form): True
CHECK A N=0 [symbolic]: d^beta w_alpha(0) = delta^beta_alpha (|alpha|,|beta|<=N): True
CHECK A N=0 [random_rational_0]: M^{(rho,mu)} unit lower-triangular (Leibniz, varpi_mu(0)=1): True
CHECK A N=0 [random_rational_0]: M M^{-1} = 1 (exact): True
CHECK A N=0 [random_rational_0]: mu d(M^{-1})/dmu = -M^{-1}(mu dM/dmu)M^{-1} (exact): True
CHECK A N=0 [random_rational_0]: -M^{-1}(mu dM)M^{-1} = D M^{-1} - M^{-1} D, D=diag(|beta|) (dilation grading): True
CHECK A N=0 [random_rational_0]: mu d/dmu P_alpha (complete formula) = (E-|alpha|) P_alpha for all |alpha|<=N: True
CHECK A N=0 [random_rational_0]: w_alpha = (z^alpha/alpha!) T_{N-|alpha|}[1/varpi_mu] varpi_mu (closed form): True
CHECK A N=0 [random_rational_0]: d^beta w_alpha(0) = delta^beta_alpha (|alpha|,|beta|<=N): True
CHECK A N=0: literal d^beta(z^alpha varpi_mu/alpha!)(0) equals M_of entries (explicit varpi, first 12x12 block): True
CHECK A N=0: explicit non-radial varpi=exp(q): mu d/dmu w_alpha^{(mu)} = (E-|alpha|) w_alpha^{(mu)} as functions: True
CHECK A N=1 [symbolic]: M^{(rho,mu)} unit lower-triangular (Leibniz, varpi_mu(0)=1): True
CHECK A N=1 [symbolic]: M M^{-1} = 1 (exact): True
CHECK A N=1 [symbolic]: mu d(M^{-1})/dmu = -M^{-1}(mu dM/dmu)M^{-1} (exact): True
CHECK A N=1 [symbolic]: -M^{-1}(mu dM)M^{-1} = D M^{-1} - M^{-1} D, D=diag(|beta|) (dilation grading): True
CHECK A N=1 [symbolic]: mu d/dmu P_alpha (complete formula) = (E-|alpha|) P_alpha for all |alpha|<=N: True
CHECK A N=1 [symbolic]: w_alpha = (z^alpha/alpha!) T_{N-|alpha|}[1/varpi_mu] varpi_mu (closed form): True
CHECK A N=1 [symbolic]: d^beta w_alpha(0) = delta^beta_alpha (|alpha|,|beta|<=N): True
CHECK A N=1 [random_rational_0]: M^{(rho,mu)} unit lower-triangular (Leibniz, varpi_mu(0)=1): True
CHECK A N=1 [random_rational_0]: M M^{-1} = 1 (exact): True
CHECK A N=1 [random_rational_0]: mu d(M^{-1})/dmu = -M^{-1}(mu dM/dmu)M^{-1} (exact): True
CHECK A N=1 [random_rational_0]: -M^{-1}(mu dM)M^{-1} = D M^{-1} - M^{-1} D, D=diag(|beta|) (dilation grading): True
CHECK A N=1 [random_rational_0]: mu d/dmu P_alpha (complete formula) = (E-|alpha|) P_alpha for all |alpha|<=N: True
CHECK A N=1 [random_rational_0]: w_alpha = (z^alpha/alpha!) T_{N-|alpha|}[1/varpi_mu] varpi_mu (closed form): True
CHECK A N=1 [random_rational_0]: d^beta w_alpha(0) = delta^beta_alpha (|alpha|,|beta|<=N): True
CHECK A N=1: literal d^beta(z^alpha varpi_mu/alpha!)(0) equals M_of entries (explicit varpi, first 12x12 block): True
CHECK A N=1: explicit non-radial varpi=exp(q): mu d/dmu w_alpha^{(mu)} = (E-|alpha|) w_alpha^{(mu)} as functions: True
CHECK A N=2 [symbolic]: M^{(rho,mu)} unit lower-triangular (Leibniz, varpi_mu(0)=1): True
CHECK A N=2 [symbolic]: M M^{-1} = 1 (exact): True
CHECK A N=2 [symbolic]: mu d(M^{-1})/dmu = -M^{-1}(mu dM/dmu)M^{-1} (exact): True
CHECK A N=2 [symbolic]: -M^{-1}(mu dM)M^{-1} = D M^{-1} - M^{-1} D, D=diag(|beta|) (dilation grading): True
CHECK A N=2 [symbolic]: mu d/dmu P_alpha (complete formula) = (E-|alpha|) P_alpha for all |alpha|<=N: True
CHECK A N=2 [symbolic]: w_alpha = (z^alpha/alpha!) T_{N-|alpha|}[1/varpi_mu] varpi_mu (closed form): True
CHECK A N=2 [symbolic]: d^beta w_alpha(0) = delta^beta_alpha (|alpha|,|beta|<=N): True
CHECK A N=2 [random_rational_0]: M^{(rho,mu)} unit lower-triangular (Leibniz, varpi_mu(0)=1): True
CHECK A N=2 [random_rational_0]: M M^{-1} = 1 (exact): True
CHECK A N=2 [random_rational_0]: mu d(M^{-1})/dmu = -M^{-1}(mu dM/dmu)M^{-1} (exact): True
CHECK A N=2 [random_rational_0]: -M^{-1}(mu dM)M^{-1} = D M^{-1} - M^{-1} D, D=diag(|beta|) (dilation grading): True
CHECK A N=2 [random_rational_0]: mu d/dmu P_alpha (complete formula) = (E-|alpha|) P_alpha for all |alpha|<=N: True
CHECK A N=2 [random_rational_0]: w_alpha = (z^alpha/alpha!) T_{N-|alpha|}[1/varpi_mu] varpi_mu (closed form): True
CHECK A N=2 [random_rational_0]: d^beta w_alpha(0) = delta^beta_alpha (|alpha|,|beta|<=N): True
CHECK A N=2: literal d^beta(z^alpha varpi_mu/alpha!)(0) equals M_of entries (explicit varpi, first 12x12 block): True
CHECK A N=2: explicit non-radial varpi=exp(q): mu d/dmu w_alpha^{(mu)} = (E-|alpha|) w_alpha^{(mu)} as functions: True
CHECK A N=3 [symbolic]: M^{(rho,mu)} unit lower-triangular (Leibniz, varpi_mu(0)=1): True
CHECK A N=3 [symbolic]: M M^{-1} = 1 (exact): True
CHECK A N=3 [symbolic]: mu d(M^{-1})/dmu = -M^{-1}(mu dM/dmu)M^{-1} (exact): True
CHECK A N=3 [symbolic]: -M^{-1}(mu dM)M^{-1} = D M^{-1} - M^{-1} D, D=diag(|beta|) (dilation grading): True
CHECK A N=3 [symbolic]: mu d/dmu P_alpha (complete formula) = (E-|alpha|) P_alpha for all |alpha|<=N: True
CHECK A N=3 [symbolic]: w_alpha = (z^alpha/alpha!) T_{N-|alpha|}[1/varpi_mu] varpi_mu (closed form): True
CHECK A N=3 [symbolic]: d^beta w_alpha(0) = delta^beta_alpha (|alpha|,|beta|<=N): True
CHECK A N=3 [random_rational_0]: M^{(rho,mu)} unit lower-triangular (Leibniz, varpi_mu(0)=1): True
CHECK A N=3 [random_rational_0]: M M^{-1} = 1 (exact): True
CHECK A N=3 [random_rational_0]: mu d(M^{-1})/dmu = -M^{-1}(mu dM/dmu)M^{-1} (exact): True
CHECK A N=3 [random_rational_0]: -M^{-1}(mu dM)M^{-1} = D M^{-1} - M^{-1} D, D=diag(|beta|) (dilation grading): True
CHECK A N=3 [random_rational_0]: mu d/dmu P_alpha (complete formula) = (E-|alpha|) P_alpha for all |alpha|<=N: True
CHECK A N=3 [random_rational_0]: w_alpha = (z^alpha/alpha!) T_{N-|alpha|}[1/varpi_mu] varpi_mu (closed form): True
CHECK A N=3 [random_rational_0]: d^beta w_alpha(0) = delta^beta_alpha (|alpha|,|beta|<=N): True
CHECK A N=3 [random_rational_1]: M^{(rho,mu)} unit lower-triangular (Leibniz, varpi_mu(0)=1): True
CHECK A N=3 [random_rational_1]: M M^{-1} = 1 (exact): True
CHECK A N=3 [random_rational_1]: mu d(M^{-1})/dmu = -M^{-1}(mu dM/dmu)M^{-1} (exact): True
CHECK A N=3 [random_rational_1]: -M^{-1}(mu dM)M^{-1} = D M^{-1} - M^{-1} D, D=diag(|beta|) (dilation grading): True
CHECK A N=3 [random_rational_1]: mu d/dmu P_alpha (complete formula) = (E-|alpha|) P_alpha for all |alpha|<=N: True
CHECK A N=3 [random_rational_1]: w_alpha = (z^alpha/alpha!) T_{N-|alpha|}[1/varpi_mu] varpi_mu (closed form): True
CHECK A N=3 [random_rational_1]: d^beta w_alpha(0) = delta^beta_alpha (|alpha|,|beta|<=N): True
CHECK A N=3: literal d^beta(z^alpha varpi_mu/alpha!)(0) equals M_of entries (explicit varpi, first 12x12 block): True
CHECK A N=3: explicit non-radial varpi=exp(q): mu d/dmu w_alpha^{(mu)} = (E-|alpha|) w_alpha^{(mu)} as functions: True
CHECK A N=4 [random_rational_0]: M^{(rho,mu)} unit lower-triangular (Leibniz, varpi_mu(0)=1): True
CHECK A N=4 [random_rational_0]: M M^{-1} = 1 (exact): True
CHECK A N=4 [random_rational_0]: mu d(M^{-1})/dmu = -M^{-1}(mu dM/dmu)M^{-1} (exact): True
CHECK A N=4 [random_rational_0]: -M^{-1}(mu dM)M^{-1} = D M^{-1} - M^{-1} D, D=diag(|beta|) (dilation grading): True
CHECK A N=4 [random_rational_0]: mu d/dmu P_alpha (complete formula) = (E-|alpha|) P_alpha for all |alpha|<=N: True
CHECK A N=4 [random_rational_0]: w_alpha = (z^alpha/alpha!) T_{N-|alpha|}[1/varpi_mu] varpi_mu (closed form): True
CHECK A N=4 [random_rational_0]: d^beta w_alpha(0) = delta^beta_alpha (|alpha|,|beta|<=N): True
CHECK A N=4 [random_rational_1]: M^{(rho,mu)} unit lower-triangular (Leibniz, varpi_mu(0)=1): True
CHECK A N=4 [random_rational_1]: M M^{-1} = 1 (exact): True
CHECK A N=4 [random_rational_1]: mu d(M^{-1})/dmu = -M^{-1}(mu dM/dmu)M^{-1} (exact): True
CHECK A N=4 [random_rational_1]: -M^{-1}(mu dM)M^{-1} = D M^{-1} - M^{-1} D, D=diag(|beta|) (dilation grading): True
CHECK A N=4 [random_rational_1]: mu d/dmu P_alpha (complete formula) = (E-|alpha|) P_alpha for all |alpha|<=N: True
CHECK A N=4 [random_rational_1]: w_alpha = (z^alpha/alpha!) T_{N-|alpha|}[1/varpi_mu] varpi_mu (closed form): True
CHECK A N=4 [random_rational_1]: d^beta w_alpha(0) = delta^beta_alpha (|alpha|,|beta|<=N): True
CHECK A N=4: literal d^beta(z^alpha varpi_mu/alpha!)(0) equals M_of entries (explicit varpi, first 12x12 block): True
CHECK A N=4: explicit non-radial varpi=exp(q): mu d/dmu w_alpha^{(mu)} = (E-|alpha|) w_alpha^{(mu)} as functions: True
# Section A consequence (exact, every N, every varpi):  mu d_mu w_alpha = (E-|alpha|) w_alpha ; hence  a_mu(u,alpha) = (-1)^{|alpha|+1} <u,(E-|alpha|) w_alpha>  and  calA coefficient (-1)^{|alpha|}<R u,(N-E) w_alpha> = (-1)^{|alpha|+1}<u,(E-|alpha|)w_alpha> (moments vanish).

# === Section B: finite exact end-to-end test (Euclidean toy homogeneous kernels; NOT the graph kernels) ===
CHECK B toy N=0 u=z^() r^{-4-N-|d|}: A_mu[phi] = mu d/dmu <u,W^mu phi> equals calA[phi]=<u,W^mu((N-E)phi)> exactly: True
CHECK B toy N=0 u=z^(): A_mu[phi] = sum_gamma a_gamma (-1)^|gamma| d^gamma phi(0) with a_gamma from the nu_k formula: True
CHECK B toy N=0: top-order (|gamma|=N) coefficients mu-independent: True
CHECK B toy N=0: a_gamma(mu) = (mu/mu0)^{N-|gamma|} a_gamma(mu0) for every gamma: True
# toy N=0 u=z^(): A_mu[phi] = 2*pi**2 ; coefficients {'()': '2*pi**2'}
CHECK B toy N=1 u=z^0 r^{-4-N-|d|}: A_mu[phi] = mu d/dmu <u,W^mu phi> equals calA[phi]=<u,W^mu((N-E)phi)> exactly: True
CHECK B toy N=1 u=z^0: A_mu[phi] = sum_gamma a_gamma (-1)^|gamma| d^gamma phi(0) with a_gamma from the nu_k formula: True
CHECK B toy N=1: top-order (|gamma|=N) coefficients mu-independent: True
CHECK B toy N=1: a_gamma(mu) = (mu/mu0)^{N-|gamma|} a_gamma(mu0) for every gamma: True
# toy N=1 u=z^0: A_mu[phi] = pi**2*(-23*lam - 3024)/6048 ; coefficients {'()': '-23*pi**2*lam/12096', '0': '-pi**2/2', '1': '0', '2': '0', '3': '0'}
CHECK B toy N=1 u=z^2 r^{-4-N-|d|}: A_mu[phi] = mu d/dmu <u,W^mu phi> equals calA[phi]=<u,W^mu((N-E)phi)> exactly: True
CHECK B toy N=1 u=z^2: A_mu[phi] = sum_gamma a_gamma (-1)^|gamma| d^gamma phi(0) with a_gamma from the nu_k formula: True
CHECK B toy N=1: top-order (|gamma|=N) coefficients mu-independent: True
CHECK B toy N=1: a_gamma(mu) = (mu/mu0)^{N-|gamma|} a_gamma(mu0) for every gamma: True
# toy N=1 u=z^2: A_mu[phi] = pi**2*(140 - lam)/840 ; coefficients {'()': '-pi**2*lam/840', '0': '0', '1': '0', '2': '-pi**2/2', '3': '0'}
CHECK B toy N=2 u=z^01 r^{-4-N-|d|}: A_mu[phi] = mu d/dmu <u,W^mu phi> equals calA[phi]=<u,W^mu((N-E)phi)> exactly: True
CHECK B toy N=2 u=z^01: A_mu[phi] = sum_gamma a_gamma (-1)^|gamma| d^gamma phi(0) with a_gamma from the nu_k formula: True
CHECK B toy N=2: top-order (|gamma|=N) coefficients mu-independent: True
CHECK B toy N=2: a_gamma(mu) = (mu/mu0)^{N-|gamma|} a_gamma(mu0) for every gamma: True
# toy N=2 u=z^01: A_mu[phi] = pi**2*(lam*(5350 - 179969*lam) + 201600)/4838400 ; coefficients {'()': '-179969*pi**2*lam**2/4838400', '0': '-pi**2*lam/2016', '1': '59*pi**2*lam/96768', '2': '-pi**2*lam/13440', '3': '-pi**2*lam/32256', '00': '0', '01': 'pi**2/12', '02': '0', '03': '0', '11': '0', '12': '0', '13': '0', '22': '0', '23': '0', '33': '0'}
CHECK B toy N=2 u=z^00 r^{-4-N-|d|}: A_mu[phi] = mu d/dmu <u,W^mu phi> equals calA[phi]=<u,W^mu((N-E)phi)> exactly: True
CHECK B toy N=2 u=z^00: A_mu[phi] = sum_gamma a_gamma (-1)^|gamma| d^gamma phi(0) with a_gamma from the nu_k formula: True
CHECK B toy N=2: top-order (|gamma|=N) coefficients mu-independent: True
CHECK B toy N=2: a_gamma(mu) = (mu/mu0)^{N-|gamma|} a_gamma(mu0) for every gamma: True
# toy N=2 u=z^00: A_mu[phi] = pi**2*(lam*(17062838*lam + 575) - 3360000)/1612800 ; coefficients {'()': '8531419*pi**2*lam**2/2419200', '0': '23*pi**2*lam/32256', '1': '-pi**2*lam/2016', '2': 'pi**2*lam/6720', '3': '-pi**2*lam/16128', '00': 'pi**2/8', '01': '0', '02': '0', '03': '0', '11': 'pi**2/24', '12': '0', '13': '0', '22': 'pi**2/24', '23': '0', '33': 'pi**2/24'}

# === Section C: exact J(P,m) = int_{S^3} P (Q-i0)^{-m} d sigma (two independent routes) ===
CHECK C: int_{S^3}(Q-i0)^{-2} = -2 i pi^2 (recovers M.6 residue r = -2 i pi^2 c^2): True
CHECK C route B: int_{S^3}(Q-i0)^{-2} = -2 i pi^2: True
CHECK C numeric (Q - i eps)^{-2} check for (a,b)=(0,0): |extrap-exact|<1e-4: True
CHECK C numeric (Q - i eps)^{-4} check for (a,b)=(2,2): |extrap-exact|<1e-4: True
CHECK C numeric (Q - i eps)^{-5} check for (a,b)=(0,4): |extrap-exact|<1e-4: True
CHECK C numeric (Q - i eps)^{-6} check for (a,b)=(4,2): |extrap-exact|<1e-4: True

# === Section D: per-grade coefficients (N=0 audit, then N=1, 2, 3, 4) ===
CHECK d_N = (1,4,10,20,35): True
CHECK p_{e0} = -2 z0 (d^alpha S = c p_alpha (Q-i0)^{-1-|alpha|}): True
CHECK 185 relation vectors parsed from kernel_relation_matrix.txt: True
CHECK D N=0: reconstructed label coefficients sum_b R_{lambda b} a_b equal direct per-label evaluation (all 1 labels, all gamma): True
CHECK D N=0: every relation vector (0) annihilates the top-order, radial-specialization and general-varpi coefficient vectors: True
CHECK D N=0: B_mu = A_mu - calA_mu vanishes on top order |gamma|=N (route A vs route B, exact): True
# N=0: top-order residue slots nonzero = 1 of 1 ; B_mu = 0 (top exact; lower by Section A identity)
CHECK D N=1: reconstructed label coefficients sum_b R_{lambda b} a_b equal direct per-label evaluation (all 4 labels, all gamma): True
CHECK D N=1: every relation vector (0) annihilates the top-order, radial-specialization and general-varpi coefficient vectors: True
CHECK D N=1: B_mu = A_mu - calA_mu vanishes on top order |gamma|=N (route A vs route B, exact): True
# N=1: top-order residue slots nonzero = 4 of 16 ; B_mu = 0 (top exact; lower by Section A identity)
CHECK D N=2: reconstructed label coefficients sum_b R_{lambda b} a_b equal direct per-label evaluation (all 20 labels, all gamma): True
CHECK D N=2: every relation vector (10) annihilates the top-order, radial-specialization and general-varpi coefficient vectors: True
CHECK D N=2: B_mu = A_mu - calA_mu vanishes on top order |gamma|=N (route A vs route B, exact): True
# N=2: top-order residue slots nonzero = 22 of 100 ; B_mu = 0 (top exact; lower by Section A identity)
CHECK D N=3: reconstructed label coefficients sum_b R_{lambda b} a_b equal direct per-label evaluation (all 60 labels, all gamma): True
CHECK D N=3: every relation vector (40) annihilates the top-order, radial-specialization and general-varpi coefficient vectors: True
CHECK D N=3: B_mu = A_mu - calA_mu vanishes on top order |gamma|=N (route A vs route B, exact): True
# N=3: top-order residue slots nonzero = 68 of 400 ; B_mu = 0 (top exact; lower by Section A identity)
CHECK D N=4: reconstructed label coefficients sum_b R_{lambda b} a_b equal direct per-label evaluation (all 170 labels, all gamma): True
CHECK D N=4: every relation vector (135) annihilates the top-order, radial-specialization and general-varpi coefficient vectors: True
CHECK D N=4: B_mu = A_mu - calA_mu vanishes on top order |gamma|=N (route A vs route B, exact): True
# N=4: top-order residue slots nonzero = 194 of 1225 ; B_mu = 0 (top exact; lower by Section A identity)
CHECK D N=0 audit recovered: A_mu(S^2) = calA(S^2) = -2 i pi^2 c^2 delta: True
CHECK D test: Res(d_alpha u) = d_alpha Res(u) for every label of grade N<=3 and every alpha (340 identities): True
CHECK D test: top-order residue map is Lorentz-equivariant on all 6 so(1,3) generators (pull-back action), N=0..4: True
CHECK D test N=0: residue map = exact combination of audited equivariant maps phi_{k,k}: True
CHECK D test N=1: residue map = exact combination of audited equivariant maps phi_{k,k}: True
CHECK D test N=2: residue map = exact combination of audited equivariant maps phi_{k,k}: True
CHECK D test N=3: residue map = exact combination of audited equivariant maps phi_{k,k}: True
CHECK D test N=4: residue map = exact combination of audited equivariant maps phi_{k,k}: True
# residue map in equivariant coordinates (units i pi^2 c^2): {0: {'phi_00': '-2'}, 1: {'phi_00': '1/2'}, 2: {'phi_00': '-1/12', 'phi_11': '-1/48'}, 3: {'phi_00': '1/96', 'phi_11': '1/384'}, 4: {'phi_00': '-1/960', 'phi_11': '-1/3840', 'phi_22': '-1/61440'}}
# wrote scale_derivative_coefficients.json, dilation_defect_coefficients.json, scale_defect_difference.json
# Sections A-D done in 291.1s

# === Section E (B.4): C_scale(mu) collapse / adjoint defect / variationality / matching per reading ===
CHECK E: top-order (prescription-independent) part of C_scale sits in contact degree D=4: True
CHECK E: lower-order slots with k=N-|gamma| sit in D=4-k (k=1..4 -> D=3,2,1,0): True
CHECK E: radial specialization: odd-k lower-order coefficients vanish (parity): True
CHECK E formal: collapse and adjoint defect of the top part are homogeneous of total derivative degree 4 (no cross-degree mixing): True
# formal         top D=4: collapse nnz=14112 (per N {'0': 744, '1': 4440, '2': 10278, '3': 10040, '4': 4604}), adjoint-defect nnz=16608, self-adjoint=False (1.1s)
CHECK E chi_plateau: collapse and adjoint defect of the top part are homogeneous of total derivative degree 4 (no cross-degree mixing): True
# chi_plateau    top D=4: collapse nnz=2522 (per N {'0': 744, '1': 1728, '2': 2386, '3': 1550, '4': 514}), adjoint-defect nnz=4090, self-adjoint=False (0.3s)
CHECK E prep_plateau: collapse and adjoint defect of the top part are homogeneous of total derivative degree 4 (no cross-degree mixing): True
# prep_plateau   top D=4: collapse nnz=4580 (per N {'0': 0, '1': 0, '2': 2274, '3': 4072, '4': 4604}), adjoint-defect nnz=4090, self-adjoint=False (0.5s)
CHECK E double_plateau: collapse and adjoint defect of the top part are homogeneous of total derivative degree 4 (no cross-degree mixing): True
# double_plateau top D=4: collapse nnz=490 (per N {'0': 0, '1': 0, '2': 474, '3': 514, '4': 514}), adjoint-defect nnz=0, self-adjoint=True (0.1s)
CHECK E double plateau: the collapsed prescription-independent scale tensor is NONZERO and self-adjoint (variational): True
CHECK E formal: collapsed top scale tensor is NOT self-adjoint (non-variational); full defect retained: True
CHECK E chi_plateau: collapsed top scale tensor is NOT self-adjoint (non-variational); full defect retained: True
CHECK E prep_plateau: collapsed top scale tensor is NOT self-adjoint (non-variational); full defect retained: True
CHECK E formal: N=0 component collapse has 744 nonzero keys (reproduces M.6 Theorem thm:M6-collapse(1)): True
CHECK E double plateau: N=0 and N=1 residue components collapse to 0 (N=0 consistent with M.6); N=2,3,4 carry the tensor: True
CHECK E double plateau: top scale tensor is NOT in span(E_R2, E_RC2) (exact rank test 2 -> 3): True
CHECK E double plateau: top scale tensor is NOT a multiple of E_{R2-(188813/88456)RC2}: True
CHECK E double plateau: top scale tensor is not a multiple of E_EH (degree mismatch D=4 vs 2): True
CHECK E: exact span solution verified: True
CHECK E double plateau: top scale tensor lies in the 5-dim span of Lorentz-covariant D=4 quadratic Euler tensors: True
CHECK E: exact span solution verified: True
CHECK E: exact span solution verified: True
CHECK E double plateau: d_mu E^{mu nu} != 0 for the top scale tensor (not linearized-gauge invariant): True
CHECK E double plateau: d_mu E^{mu nu} = 0 for E_R2, E_RC2 (control): True
# double plateau top tensor = Euler tensor of -71*L1/120 - 27*L2/80 - 11*L3/10 + 19*L4/30 - 13*L5/15  (covariant basis); R2 -> {'L1=h_{mn} Box^2 h^{mn}': '0', 'L2=h Box^2 h': '1', 'L3=(d^m h_{mn}) Box (d_r h^{rn})': '0', 'L4=h Box d^m d^n h_{mn}': '-2', 'L5=(d^m d^n h_{mn})^2': '1'} ; RC2 -> {'L1=h_{mn} Box^2 h^{mn}': '1/4', 'L2=h Box^2 h': '1/4', 'L3=(d^m h_{mn}) Box (d_r h^{rn})': '1/2', 'L4=h Box d^m d^n h_{mn}': '-1/2', 'L5=(d^m d^n h_{mn})^2': '1/2'}
CHECK E formal radial k=2: collapse homogeneous of degree D=2: True
CHECK E formal radial k=4: collapse homogeneous of degree D=0: True
CHECK E formal envelope k=1: ker Coll inside SA-parameter space (dimension check): True
# formal         envelope k=1 (D=3): params=120 rank Coll=120 dim SA=0 var image=0 
CHECK E formal envelope k=2: ker Coll inside SA-parameter space (dimension check): True
# formal         envelope k=2 (D=2): params=84 rank Coll=84 dim SA=0 var image=0 
CHECK E formal envelope k=3: ker Coll inside SA-parameter space (dimension check): True
# formal         envelope k=3 (D=1): params=56 rank Coll=56 dim SA=0 var image=0 
CHECK E formal envelope k=4: ker Coll inside SA-parameter space (dimension check): True
# formal         envelope k=4 (D=0): params=35 rank Coll=35 dim SA=26 var image=26 
CHECK E chi_plateau radial k=2: collapse homogeneous of degree D=2: True
CHECK E chi_plateau radial k=4: collapse homogeneous of degree D=0: True
CHECK E chi_plateau envelope k=1: ker Coll inside SA-parameter space (dimension check): True
# chi_plateau    envelope k=1 (D=3): params=120 rank Coll=120 dim SA=0 var image=0 
CHECK E chi_plateau envelope k=2: ker Coll inside SA-parameter space (dimension check): True
# chi_plateau    envelope k=2 (D=2): params=84 rank Coll=84 dim SA=0 var image=0 
CHECK E chi_plateau envelope k=3: ker Coll inside SA-parameter space (dimension check): True
# chi_plateau    envelope k=3 (D=1): params=56 rank Coll=56 dim SA=0 var image=0 
CHECK E chi_plateau envelope k=4: ker Coll inside SA-parameter space (dimension check): True
# chi_plateau    envelope k=4 (D=0): params=35 rank Coll=35 dim SA=26 var image=26 
CHECK E prep_plateau radial k=2: collapse homogeneous of degree D=2: True
CHECK E prep_plateau radial k=4: collapse homogeneous of degree D=0: True
CHECK E prep_plateau envelope k=1: ker Coll inside SA-parameter space (dimension check): True
# prep_plateau   envelope k=1 (D=3): params=120 rank Coll=120 dim SA=0 var image=0 
CHECK E prep_plateau envelope k=2: ker Coll inside SA-parameter space (dimension check): True
# prep_plateau   envelope k=2 (D=2): params=84 rank Coll=84 dim SA=0 var image=0 
CHECK E prep_plateau envelope k=3: ker Coll inside SA-parameter space (dimension check): True
# prep_plateau   envelope k=3 (D=1): params=56 rank Coll=56 dim SA=0 var image=0 
CHECK E prep_plateau envelope k=4: ker Coll inside SA-parameter space (dimension check): True
# prep_plateau   envelope k=4 (D=0): params=35 rank Coll=35 dim SA=26 var image=26 
CHECK E double_plateau radial k=2: collapse homogeneous of degree D=2: True
CHECK E double_plateau radial k=4: collapse homogeneous of degree D=0: True
CHECK E double_plateau envelope k=1: ker Coll inside SA-parameter space (dimension check): True
# double_plateau envelope k=1 (D=3): params=120 rank Coll=120 dim SA=0 var image=0 
CHECK E double_plateau envelope k=2: ker Coll inside SA-parameter space (dimension check): True
# double_plateau envelope k=2 (D=2): params=84 rank Coll=84 dim SA=49 var image=49 targets cap image=0
CHECK E double_plateau envelope k=3: ker Coll inside SA-parameter space (dimension check): True
# double_plateau envelope k=3 (D=1): params=56 rank Coll=56 dim SA=0 var image=0 
CHECK E double_plateau envelope k=4: ker Coll inside SA-parameter space (dimension check): True
# double_plateau envelope k=4 (D=0): params=35 rank Coll=35 dim SA=26 var image=26 targets cap image=1
# wrote full_scale_collapse_certificates.json (33.2s)
# wrote matched_scale_ledger.txt
# total runtime 324.3s ; 185 checks
ALL CHECKS PASSED
exit 0
\end{lstlisting}

\paragraph{Work package C (\texttt{wpC/time\_only\_helmholtz\_out.txt}).}
\begin{lstlisting}
# time_only_helmholtz.py output  (python 3.13.5, sympy 1.14.0; exact QQ; CONDITIONAL on (S-bv) via M.2 basis)
# XVI-M work package C: universal time-only Helmholtz classification (reading time_only)
CHECK dim C_coh = 3321, graph table 3282 lines: True
# input sha256: {'graph_general_table.txt': '4cd29d3c9f691340b314aecbfe7d6aef3590b9cee21dcdb486fe9fbcdd679e49', 'kernel_basis.json': 'fcbe2876abd2a6b410813f732356f829c970c22c3a333f7e4ec34909c0e0793b', 'm4_common.py': '932b716db827efc0cbb3d9eca91c633d9eecf932c50c1988008f98795337bf1e'}
CHECK dim C_reg = 3150 (N>=3), dim K block = 171 (N<=2): True
# === C.2 specialization tests ===
CHECK S_t(D_mu F) = D_mu^(t) S_t(F) for 400 random exact polynomials x 4 directions, and S_t(D^beta F) = D^(t)beta S_t(F) for random |beta|<=4: True
CHECK negative control: a D_1^(t) that differentiates sigma (sigma'(t) for d_1) violates the intertwining: True
CHECK sympy Function route: d_mu[F(sigma(t), chi(x), h(x))] = D_mu^(t) S_t(F) on 25 random polynomials (sigma a function of x0 only): True
CHECK S_t(D_mu F) = D_mu^(t) S_t(F) on 960 actual collapse / adjoint-defect expressions (120 random contact variables x 4 directions): True
# ######## reading time_only ########
# === C.3/C.4 blockwise: route 1 vs route 2, exact kernels ===
# blockwise pass 64.9s
CHECK [time_only] ROUTE 1 = ROUTE 2 for Coll_t: all 3321 columns equal exactly (nnz 446292): True
CHECK [time_only] ROUTE 1 = ROUTE 2 for H_t: all 3321 adjoint-defect columns equal exactly (nnz 573582): True
CHECK [time_only] grading: every time-only key in block (D, parity) has total derivative count D and parity of the block (blocks have disjoint rows): True
# supplementary mod-p ranks (all rows, p=1e9+7, 998244353) agree with exact ranks in every block: True
CHECK [time_only] ker Coll_t is contained in ker H_t (each ker Coll_t vector has zero time-only adjoint defect): True
CHECK [time_only] the 26 formal SA vectors (helmholtz_certificates.txt) lie in ker H_t (specialization: ker H_formal <= ker H_t): True
CHECK [time_only] D=0: dim Coll_t(ker H_t) = 26 (mod-p rank lower bound = dim ker H_t - dim ker Coll_t upper bound): True
CHECK [time_only] D=0: regression-preserving variational image dim Coll_t(ker H_t cap C_reg) = 26: True
CHECK [time_only] D=1: dim Coll_t(ker H_t) = 0 (mod-p rank lower bound = dim ker H_t - dim ker Coll_t upper bound): True
CHECK [time_only] D=1: regression-preserving variational image dim Coll_t(ker H_t cap C_reg) = 0: True
CHECK [time_only] D=2: dim Coll_t(ker H_t) = 0 (mod-p rank lower bound = dim ker H_t - dim ker Coll_t upper bound): True
CHECK [time_only] D=2: regression-preserving variational image dim Coll_t(ker H_t cap C_reg) = 0: True
CHECK [time_only] D=3: dim Coll_t(ker H_t) = 0 (mod-p rank lower bound = dim ker H_t - dim ker Coll_t upper bound): True
CHECK [time_only] D=3: regression-preserving variational image dim Coll_t(ker H_t cap C_reg) = 0: True
CHECK [time_only] D=4: dim Coll_t(ker H_t) = 0 (mod-p rank lower bound = dim ker H_t - dim ker Coll_t upper bound): True
CHECK [time_only] D=4: regression-preserving variational image dim Coll_t(ker H_t cap C_reg) = 0: True
# === CLASSIFICATION TABLE (time_only reading; conditional on (S-bv)) ===
# D | size | kerColl | imColl | kerH | Coll(kerH) | kerColl^Creg | kerColl^K | kerH^Creg | kerH^K | Coll(kerH^Creg)
# 0 | 35 | 0 | 35 | 26 | 26 | 0 | 0 | 26 | 0 | 26
# 1 | 160 | 0 | 160 | 0 | 0 | 0 | 0 | 0 | 0 | 0
# 2 | 440 | 0 | 440 | 0 | 0 | 0 | 0 | 0 | 0 | 0
# 3 | 944 | 0 | 944 | 0 | 0 | 0 | 0 | 0 | 0 | 0
# 4 | 1742 | 0 | 1742 | 0 | 0 | 0 | 0 | 0 | 0 | 0
# totals: {'size': 3321, 'size_Creg': 3150, 'size_K': 171, 'dim_ker_Coll_t': 0, 'dim_im_Coll_t': 3321, 'dim_ker_H_t': 26, 'dim_Coll_t_of_ker_H_t': 26, 'dim_ker_Coll_t_cap_Creg': 0, 'dim_ker_Coll_t_cap_K': 0, 'dim_ker_H_t_cap_Creg': 26, 'dim_ker_H_t_cap_K': 0, 'dim_regression_preserving_variational_image': 26}
# variational-image monomial structure per D: {"0": {"n_vectors": 26, "monomial_type_counts": {"sigma^(0) chi_0 h_0": 198}, "sigma_orders_present": [0], "has_chi_derivatives": false, "max_h_order": 0}}
# === C.5 functional probes (weight w = chi_I(x) sigma_eps(t); all preparation/switch derivatives retained) ===
CHECK [time_only] weighted P1 target: total derivative degree 0, parity 0: True
CHECK [time_only] weighted P1 target (time-only Euler tensor) is self-adjoint under the time-only Gate V adjoint: True
CHECK [time_only] weighted P2 target: total derivative degree 0, parity 0: True
CHECK [time_only] weighted P2 target (time-only Euler tensor) is self-adjoint under the time-only Gate V adjoint: True
CHECK [time_only] weighted EH target: total derivative degree 2, parity 0: True
CHECK [time_only] weighted EH target (time-only Euler tensor) is self-adjoint under the time-only Gate V adjoint: True
CHECK [time_only] weighted R2 target: total derivative degree 4, parity 0: True
CHECK [time_only] weighted R2 target (time-only Euler tensor) is self-adjoint under the time-only Gate V adjoint: True
CHECK [time_only] weighted RC2 target: total derivative degree 4, parity 0: True
CHECK [time_only] weighted RC2 target (time-only Euler tensor) is self-adjoint under the time-only Gate V adjoint: True
CHECK [time_only] divergence identity: d_mu E^{mu nu} = 0 for weighted R2 with w = chi sigma(t) (time-only jets): True
CHECK [time_only] divergence identity: d_mu E^{mu nu} = 0 for weighted RC2 with w = chi sigma(t) (time-only jets): True
CHECK [time_only] divergence identity: d_mu E_EH^{mu nu} = (1) (d_mu w) G_lin^{mu nu} for weighted EH (w = chi sigma(t)): True
CHECK [time_only] ultralocal weighted P1: divergence nonzero (not gauge invariant): True
CHECK [time_only] P1+P2/C_coh: explicit preimage of direction ['1', '-89/140'] reproduces the weighted target exactly: True
CHECK [time_only] P1+P2/C_coh: route-2 collapse of the preimage equals the target: True
CHECK [time_only] P1+P2/C_coh: preimage lies in ker H_t (variational): True
CHECK [time_only] P1+P2/C_coh: EL(1/2 <h, E_C[h]>) = E_C (independent sympy time-only jet route): True
CHECK [time_only] P1+P2/C_coh: EL(1/2 <h, E_C[h]>) = E_C via sympy euler_equations with sigma(x0), chi(x), h(x) Functions: True
# probe P1+P2   C_coh D=0 block=10 : dim(span cap im Coll_t)=1 ; members={'P1': False, 'P2': False} ; directions=[['1', '-89/140']] (0.1s)
CHECK [time_only] P1+P2/C_reg: explicit preimage of direction ['1', '-89/140'] reproduces the weighted target exactly: True
CHECK [time_only] P1+P2/C_reg: route-2 collapse of the preimage equals the target: True
CHECK [time_only] P1+P2/C_reg: preimage lies in ker H_t (variational): True
CHECK [time_only] P1+P2/C_reg: EL(1/2 <h, E_C[h]>) = E_C (independent sympy time-only jet route): True
CHECK [time_only] P1+P2/C_reg: EL(1/2 <h, E_C[h]>) = E_C via sympy euler_equations with sigma(x0), chi(x), h(x) Functions: True
# probe P1+P2   C_reg D=0 block=10 : dim(span cap im Coll_t)=1 ; members={'P1': False, 'P2': False} ; directions=[['1', '-89/140']] (0.0s)
# probe EH      C_coh D=2 block=84 : dim(span cap im Coll_t)=0 ; members={'EH': False} ; directions=[] (0.1s)
# probe EH      C_reg D=2 block=80 : dim(span cap im Coll_t)=0 ; members={'EH': False} ; directions=[] (0.0s)
# probe R2+RC2  C_coh D=4 block=292 : dim(span cap im Coll_t)=0 ; members={'R2': False, 'RC2': False} ; directions=[] (1.1s)
# probe R2+RC2  C_reg D=4 block=265 : dim(span cap im Coll_t)=0 ; members={'R2': False, 'RC2': False} ; directions=[] (0.6s)
# ######## reading time_only_chi_plateau ########
# === C.3/C.4 blockwise: route 1 vs route 2, exact kernels ===
# blockwise pass 36.4s
CHECK [time_only_chi_plateau] ROUTE 1 = ROUTE 2 for Coll_t: all 3321 columns equal exactly (nnz 91568): True
CHECK [time_only_chi_plateau] ROUTE 1 = ROUTE 2 for H_t: all 3321 adjoint-defect columns equal exactly (nnz 130110): True
CHECK [time_only_chi_plateau] grading: every time-only key in block (D, parity) has total derivative count D and parity of the block (blocks have disjoint rows): True
# supplementary mod-p ranks (all rows, p=1e9+7, 998244353) agree with exact ranks in every block: True
CHECK [time_only_chi_plateau] ker Coll_t is contained in ker H_t (each ker Coll_t vector has zero time-only adjoint defect): True
CHECK [time_only_chi_plateau] the 26 formal SA vectors (helmholtz_certificates.txt) lie in ker H_t (specialization: ker H_formal <= ker H_t): True
CHECK [time_only_chi_plateau] D=0: dim Coll_t(ker H_t) = 26 (mod-p rank lower bound = dim ker H_t - dim ker Coll_t upper bound): True
CHECK [time_only_chi_plateau] D=0: regression-preserving variational image dim Coll_t(ker H_t cap C_reg) = 26: True
CHECK [time_only_chi_plateau] D=1: dim Coll_t(ker H_t) = 0 (mod-p rank lower bound = dim ker H_t - dim ker Coll_t upper bound): True
CHECK [time_only_chi_plateau] D=1: regression-preserving variational image dim Coll_t(ker H_t cap C_reg) = 0: True
CHECK [time_only_chi_plateau] D=2: dim Coll_t(ker H_t) = 156 (mod-p rank lower bound = dim ker H_t - dim ker Coll_t upper bound): True
CHECK [time_only_chi_plateau] D=2: regression-preserving variational image dim Coll_t(ker H_t cap C_reg) = 156: True
CHECK [time_only_chi_plateau] D=3: dim Coll_t(ker H_t) = 0 (mod-p rank lower bound = dim ker H_t - dim ker Coll_t upper bound): True
CHECK [time_only_chi_plateau] D=3: regression-preserving variational image dim Coll_t(ker H_t cap C_reg) = 0: True
CHECK [time_only_chi_plateau] D=4: dim Coll_t(ker H_t) = 390 (mod-p rank lower bound = dim ker H_t - dim ker Coll_t upper bound): True
CHECK [time_only_chi_plateau] D=4: regression-preserving variational image dim Coll_t(ker H_t cap C_reg) = 390: True
# === CLASSIFICATION TABLE (time_only reading; conditional on (S-bv)) ===
# D | size | kerColl | imColl | kerH | Coll(kerH) | kerColl^Creg | kerColl^K | kerH^Creg | kerH^K | Coll(kerH^Creg)
# 0 | 35 | 0 | 35 | 26 | 26 | 0 | 0 | 26 | 0 | 26
# 1 | 160 | 0 | 160 | 0 | 0 | 0 | 0 | 0 | 0 | 0
# 2 | 440 | 0 | 440 | 156 | 156 | 0 | 0 | 156 | 0 | 156
# 3 | 944 | 0 | 944 | 0 | 0 | 0 | 0 | 0 | 0 | 0
# 4 | 1742 | 0 | 1742 | 390 | 390 | 0 | 0 | 390 | 0 | 390
# totals: {'size': 3321, 'size_Creg': 3150, 'size_K': 171, 'dim_ker_Coll_t': 0, 'dim_im_Coll_t': 3321, 'dim_ker_H_t': 572, 'dim_Coll_t_of_ker_H_t': 572, 'dim_ker_Coll_t_cap_Creg': 0, 'dim_ker_Coll_t_cap_K': 0, 'dim_ker_H_t_cap_Creg': 572, 'dim_ker_H_t_cap_K': 0, 'dim_regression_preserving_variational_image': 572}
# variational-image monomial structure per D: {"0": {"n_vectors": 26, "monomial_type_counts": {"sigma^(0) chi_0 h_0": 198}, "sigma_orders_present": [0], "has_chi_derivatives": false, "max_h_order": 0}, "2": {"n_vectors": 156, "monomial_type_counts": {"sigma^(0) chi_0 h_2": 1188}, "sigma_orders_present": [0], "has_chi_derivatives": false, "max_h_order": 2}, "4": {"n_vectors": 390, "monomial_type_counts": {"sigma^(0) chi_0 h_4": 2970}, "sigma_orders_present": [0], "has_chi_derivatives": false, "max_h_order": 4}}
# === C.5 functional probes (weight w = sigma_eps(t) [chi_I == 1]; all preparation/switch derivatives retained) ===
CHECK [time_only_chi_plateau] weighted P1 target: total derivative degree 0, parity 0: True
CHECK [time_only_chi_plateau] weighted P1 target (time-only Euler tensor) is self-adjoint under the time-only Gate V adjoint: True
CHECK [time_only_chi_plateau] weighted P2 target: total derivative degree 0, parity 0: True
CHECK [time_only_chi_plateau] weighted P2 target (time-only Euler tensor) is self-adjoint under the time-only Gate V adjoint: True
CHECK [time_only_chi_plateau] weighted EH target: total derivative degree 2, parity 0: True
CHECK [time_only_chi_plateau] weighted EH target (time-only Euler tensor) is self-adjoint under the time-only Gate V adjoint: True
CHECK [time_only_chi_plateau] weighted R2 target: total derivative degree 4, parity 0: True
CHECK [time_only_chi_plateau] weighted R2 target (time-only Euler tensor) is self-adjoint under the time-only Gate V adjoint: True
CHECK [time_only_chi_plateau] weighted RC2 target: total derivative degree 4, parity 0: True
CHECK [time_only_chi_plateau] weighted RC2 target (time-only Euler tensor) is self-adjoint under the time-only Gate V adjoint: True
CHECK [time_only_chi_plateau] divergence identity: d_mu E^{mu nu} = 0 for weighted R2 with w = chi sigma(t) (time-only jets): True
CHECK [time_only_chi_plateau] divergence identity: d_mu E^{mu nu} = 0 for weighted RC2 with w = chi sigma(t) (time-only jets): True
CHECK [time_only_chi_plateau] divergence identity: d_mu E_EH^{mu nu} = (1) (d_mu w) G_lin^{mu nu} for weighted EH (w = chi sigma(t)): True
CHECK [time_only_chi_plateau] ultralocal weighted P1: divergence nonzero (not gauge invariant): True
CHECK [time_only_chi_plateau] P1+P2/C_coh: explicit preimage of direction ['1', '-89/140'] reproduces the weighted target exactly: True
CHECK [time_only_chi_plateau] P1+P2/C_coh: route-2 collapse of the preimage equals the target: True
CHECK [time_only_chi_plateau] P1+P2/C_coh: preimage lies in ker H_t (variational): True
CHECK [time_only_chi_plateau] P1+P2/C_coh: EL(1/2 <h, E_C[h]>) = E_C (independent sympy time-only jet route): True
CHECK [time_only_chi_plateau] P1+P2/C_coh: EL(1/2 <h, E_C[h]>) = E_C via sympy euler_equations with sigma(x0), chi(x), h(x) Functions: True
# probe P1+P2   C_coh D=0 block=10 : dim(span cap im Coll_t)=1 ; members={'P1': False, 'P2': False} ; directions=[['1', '-89/140']] (0.1s)
CHECK [time_only_chi_plateau] P1+P2/C_reg: explicit preimage of direction ['1', '-89/140'] reproduces the weighted target exactly: True
CHECK [time_only_chi_plateau] P1+P2/C_reg: route-2 collapse of the preimage equals the target: True
CHECK [time_only_chi_plateau] P1+P2/C_reg: preimage lies in ker H_t (variational): True
CHECK [time_only_chi_plateau] P1+P2/C_reg: EL(1/2 <h, E_C[h]>) = E_C (independent sympy time-only jet route): True
CHECK [time_only_chi_plateau] P1+P2/C_reg: EL(1/2 <h, E_C[h]>) = E_C via sympy euler_equations with sigma(x0), chi(x), h(x) Functions: True
# probe P1+P2   C_reg D=0 block=10 : dim(span cap im Coll_t)=1 ; members={'P1': False, 'P2': False} ; directions=[['1', '-89/140']] (0.0s)
# probe EH      C_coh D=2 block=84 : dim(span cap im Coll_t)=0 ; members={'EH': False} ; directions=[] (0.1s)
# probe EH      C_reg D=2 block=80 : dim(span cap im Coll_t)=0 ; members={'EH': False} ; directions=[] (0.1s)
# probe R2+RC2  C_coh D=4 block=292 : dim(span cap im Coll_t)=0 ; members={'R2': False, 'RC2': False} ; directions=[] (1.5s)
# probe R2+RC2  C_reg D=4 block=265 : dim(span cap im Coll_t)=0 ; members={'R2': False, 'RC2': False} ; directions=[] (0.3s)
# === new time-only variational directions: reading time_only_chi_plateau, D>=1 ===
CHECK [time_only_chi_plateau] D=2: every new variational direction is sigma(t) x (purely SPATIAL derivative operator on h): no sigma derivative, no time derivative of h: True
CHECK [time_only_chi_plateau] D=2: EL(1/2 <h, E_C[h]>) = E_C for all 156 new directions (independent sympy time-only jet route) (2.2s): True
CHECK [time_only_chi_plateau] D=2: EL check via sympy euler_equations with sigma(x0), h(x) Functions on 3 sample directions: True
CHECK [time_only_chi_plateau] D=2: every new direction remains SA in the double-plateau reading (specialization sigma == 1): True
CHECK [time_only_chi_plateau] D=2: no new direction is SA in the (spacetime-formal sigma) chi_plateau reading: the directions are time-only-specific: True
# D=2 new directions: dim 156 ; h-derivative multi-indices used {'23': 198, '11': 198, '22': 198, '33': 198, '13': 198, '12': 198} ; support N-counts {4: 210} ; |gamma| counts {2: 210}
CHECK [time_only_chi_plateau] D=4: every new variational direction is sigma(t) x (purely SPATIAL derivative operator on h): no sigma derivative, no time derivative of h: True
CHECK [time_only_chi_plateau] D=4: EL(1/2 <h, E_C[h]>) = E_C for all 390 new directions (independent sympy time-only jet route) (6.3s): True
CHECK [time_only_chi_plateau] D=4: EL check via sympy euler_equations with sigma(x0), h(x) Functions on 3 sample directions: True
CHECK [time_only_chi_plateau] D=4: every new direction remains SA in the double-plateau reading (specialization sigma == 1): True
CHECK [time_only_chi_plateau] D=4: no new direction is SA in the (spacetime-formal sigma) chi_plateau reading: the directions are time-only-specific: True
# D=4 new directions: dim 390 ; h-derivative multi-indices used {'1123': 198, '2223': 198, '2333': 198, '1111': 198, '1122': 198, '1133': 198, '2222': 198, '2233': 198, '3333': 198, '1113': 198, '1223': 198, '1333': 198, '1112': 198, '1222': 198, '1233': 198} ; support N-counts {4: 525} ; |gamma| counts {4: 525}
CHECK [time_only_chi_plateau] the 26 ultralocal SA vectors are supported on N=4, gamma=0: True
CHECK [time_only_chi_plateau] D=2: span{ultralocal SA (26) x spatial gamma (6)} (156 independent vectors) lies in ker H_t and has dim = dim ker H_t = 156, hence equals it: True
CHECK [time_only_chi_plateau] D=4: span{ultralocal SA (26) x spatial gamma (15)} (390 independent vectors) lies in ker H_t and has dim = dim ker H_t = 390, hence equals it: True
CHECK supplementary (floating, 60 digits): Wronskian of rho, rho', ..., rho'''' (XII-C ramp profile) nonzero at s = 1/2, 1/3, 1/7: True
# Wronskian W(rho,...,rho^(4)) at s: {'1/2': '88060251340800.0', '1/3': '2.575905490314e+15', '1/7': '4.10243726191388e+17'}
# total runtime 154.1s
ALL CHECKS PASSED
\end{lstlisting}

\paragraph{Regression (\texttt{regression\_full\_out.txt}; \texttt{XVIM\_FULL\_WP=1}).}
\begin{lstlisting}
/workspace/XVI-M-artifacts/regression_M.py:828: SyntaxWarning: invalid escape sequence '\l'
  require("\label{rem:revlog-final}" in tex and "Revision log (final verification)" in tex,
/workspace/XVI-M-artifacts/regression_M.py:841: SyntaxWarning: invalid escape sequence '\m'
  require("M-C$_{\mathrm{data}}$" in tex and "explicitly specified cutoff" in tex,
/workspace/XVI-M-artifacts/regression_M.py:846: SyntaxWarning: invalid escape sequence '\m'
  require("D=0,2,4" in tex and "D=1,3" in tex and "mathcal K_{\mathrm{SA}}" in tex,
# regression_M.py output   (python 3.13.5, sympy 1.14.0)
CHECK provenance manifest lists the frozen predecessor files (18): True
CHECK frozen file unchanged: XVI-L.tex: True
CHECK frozen file unchanged: XVI-L-artifacts/README.txt: True
CHECK frozen file unchanged: XVI-L-artifacts/counterterm_map.py: True
CHECK frozen file unchanged: XVI-L-artifacts/counterterm_map_out.txt: True
CHECK frozen file unchanged: XVI-L-artifacts/kernel_inventory.py: True
CHECK frozen file unchanged: XVI-L-artifacts/kernel_inventory_out.txt: True
CHECK frozen file unchanged: XVI-L-artifacts/regression_L.py: True
CHECK frozen file unchanged: XVI-L-artifacts/regression_L_out.txt: True
CHECK frozen file unchanged: XVI-L-artifacts/source_full.py: True
CHECK frozen file unchanged: XVI-L-artifacts/source_full_out.txt: True
CHECK frozen file unchanged: XVI-L-artifacts/source_full_table.txt: True
CHECK frozen file unchanged: XVI-L-artifacts/two_edge_coeff_table.txt: True
CHECK frozen file unchanged: XVI-L-artifacts/two_edge_full.py: True
CHECK frozen file unchanged: XVI-L-artifacts/two_edge_full_out.txt: True
CHECK frozen file unchanged: XVI-L-artifacts/two_edge_support.txt: True
CHECK frozen file unchanged: XVI-L-artifacts/xvi_common.py: True
CHECK frozen file unchanged: XVI-K-artifacts/hessian_full_out.txt: True
CHECK frozen file unchanged: XVI-K-artifacts/xvi_common.py: True
CHECK graph_general_out.txt ends with ALL CHECKS PASSED: True
CHECK graph_general_out.txt: every CHECK line is True: True
CHECK kernel_relations_out.txt ends with ALL CHECKS PASSED: True
CHECK kernel_relations_out.txt: every CHECK line is True: True
CHECK contact_space_out.txt ends with ALL CHECKS PASSED: True
CHECK contact_space_out.txt: every CHECK line is True: True
CHECK contact_collapse_out.txt ends with ALL CHECKS PASSED: True
CHECK contact_collapse_out.txt: every CHECK line is True: True
CHECK helmholtz_check_out.txt ends with ALL CHECKS PASSED: True
CHECK helmholtz_check_out.txt: every CHECK line is True: True
CHECK matching_probe_out.txt ends with ALL CHECKS PASSED: True
CHECK matching_probe_out.txt: every CHECK line is True: True
CHECK fixed_data_reading_out.txt ends with ALL CHECKS PASSED: True
CHECK fixed_data_reading_out.txt: every CHECK line is True: True
CHECK counterterm_functionals_out.txt ends with ALL CHECKS PASSED: True
CHECK counterterm_functionals_out.txt: every CHECK line is True: True
CHECK scaling_anomaly_out.txt ends with ALL CHECKS PASSED: True
CHECK scaling_anomaly_out.txt: every CHECK line is True: True
CHECK covariance_constraints_out.txt ends with ALL CHECKS PASSED: True
CHECK covariance_constraints_out.txt: every CHECK line is True: True
CHECK m7_legacy_reconstruction_out.txt ends with ALL CHECKS PASSED: True
CHECK m7_legacy_reconstruction_out.txt: every CHECK line is True: True
# relation spot-check seed=42 ; points=[['-6', '-1/2', '-1', '-1'], ['8', '9/4', '-8', '-7/2'], ['-2/5', '-9/5', '-1/2', '2']]
CHECK conventions.json records regression_M relation seed 42: True
CHECK conventions.json records contact_space seed 20261005: True
CHECK conventions.json records helmholtz_check seed 7: True
CHECK three off-cone rational points from seed 42: True
CHECK parsed 185 kernel relations + 45 wave-operator generators: True
CHECK every exported relation vanishes exactly on (d^alpha S)(d^beta S), S=1/Q, at 3 random rational points off the cone (independent route): True
CHECK dimensions consistent between M.2 and M.3 outputs: True
CHECK M.1 exact specialization check recorded: True
CHECK sign example (d0 S)(d1 S) = (1/2) S d0 d1 S with S = 1/Q: True
CHECK sign example numerators -4 z0 z1 and -8 z0 z1: True
CHECK M.4 matching: archived preimage collapses EXACTLY to the Gate V Euler tensor of int R_lin^2 (double plateau): True
CHECK M.4 matching: that collapsed tensor is self-adjoint: True
CHECK M.4 Helmholtz: formal SA vector #0 is the Gate V Euler derivative of (1/2) int <h, E_C[h]> (independent route): True
CHECK M.4 headline numbers recorded (formal Collapse injective; formal dim SA = 26): True
CHECK single-switch: every tested column matches d^gamma(chi ahatbar)=d^gamma(bar_a) (2622 columns): True
CHECK single-switch: implemented columns differ from d^gamma(ahatbar) (no-chi) on 2552 sample columns: True
CHECK single-switch: implemented columns differ from d^gamma(chi^2 ahatbar) (double-chi) on 2552 sample columns: True
CHECK single-switch: every tested column matches d(chi ahatbar), not d(chi^2 ahatbar) or d(ahatbar).: True
CHECK Certificate A variable C[N=4,{(),0222},()] located in contact variables: True
CHECK Cert A weighted symmetry: w_mn K^{mn}_ab = w_ab K^{ab}_mn holds: True
CHECK Cert A weighted symmetry: reversed w_ab K^{mn}_ab = w_mn K^{ab}_mn FAILS (as required): True
CHECK Cert A numeric: w00 K00_02 = -56/3 = w02 K02_00: True
CHECK matching_probe outs cover formal / chi-plateau / double-plateau weight readings: True
CHECK formal: dim ker Collapse = 0: True
CHECK formal: dim SA = 26: True
CHECK formal: dim variational image = 26: True
CHECK formal: variational D=3 dim = 0: True
CHECK chi_plateau: dim ker Collapse = 0: True
CHECK chi_plateau: dim SA = 26: True
CHECK chi_plateau: dim variational image = 26: True
CHECK chi_plateau: variational D=3 dim = 0: True
CHECK double_plateau: dim ker Collapse = 22: True
CHECK double_plateau: dim SA = 1347: True
CHECK double_plateau: dim variational image = 1325: True
CHECK double_plateau: variational D=3 dim = 8: True
CHECK Stage2 prep_plateau ker Coll = 21: True
CHECK Stage2 ultralocal dim 26: True
CHECK Stage2 transitions: Unresolved with domain-size bound 3321 (unsupported bound 22 removed): True
CHECK Stage2 transition upper bound = 3321 = dim C_coh: True
CHECK Stage2 criterion includes (iii) adjoint-defect map and time-only caveat: True
CHECK M.5 named theorem recorded: True
CHECK EOM/field-redef ledgers empty: True
CHECK [current status] manuscript has Stage2/M.5 labels: True
CHECK [current status] manuscript still references Input (S-bv) (now certified): True
CHECK [current status] S-bv not discharged by matching (certified analytically, thm:wpA-Sbv): True
CHECK Gate V: W((mu,mu))=1 and W((mu,nu))=2 for mu<nu: True
CHECK Gate V weights on Cert A ultralocal K: w=1 diag / w=2 off symmetry holds: True
CHECK Gate V weight w_00 = 1: True
CHECK Gate V weight w_01 = 2: True
CHECK Gate V weight w_02 = 2: True
CHECK Gate V weight w_03 = 2: True
CHECK Gate V weight w_11 = 1: True
CHECK Gate V weight w_12 = 2: True
CHECK Gate V weight w_13 = 2: True
CHECK Gate V weight w_22 = 1: True
CHECK Gate V weight w_23 = 2: True
CHECK Gate V weight w_33 = 1: True
CHECK conventions.json records eight third-order double-plateau variational directions: True
CHECK conventions.json: third-order dirs are not gravitational scalar counterterms: True
CHECK [current status] conventions.json: (S-bv) not discharged by matching: True
CHECK all three readings present in matching_probe_out: True
CHECK single switch factor: implemented collapse equals d^gamma(chi ahatbar): True
CHECK scaling_anomaly_out ends with ALL CHECKS PASSED: True
CHECK M.6 sphere integral certificate present: True
CHECK M.6 S^2 residue certificate present: True
CHECK M.6 double-plateau vanishing recorded: True
CHECK M.6 formal non-variational recorded: True
CHECK [historical record] M.6 keeps M-C unclaimed: True
CHECK [historical record] scaling_anomaly_cert.json: M_C_claimable is False: True
CHECK cert residue value: True
CHECK [historical record] archived M.6 cert records its (historical) (S-bv) input: True
CHECK independent contour residue simplifies to -1: True
CHECK independent contour: Iu0=-2 i pi: True
CHECK independent: sphere = pi * Iu0 = -2 i pi^2: True
CHECK [historical record] archived M.7 summary records its (historical) (S-bv) input: True
CHECK [historical record] M.7 keeps M-C/M-D unclaimed: True
CHECK [historical record] M.7 does not claim Target II: True
CHECK M.7 abstract equivariant dims = 1,1,3,3,6: True
CHECK M.7 rank = slots - dim sol (codim) and coker = dim codomain - rank, codomain specified: True
CHECK M.7 equivariant freedom 14 = 5 K + 9 C_reg: True
CHECK M.7 formal: (dim V, pi_K V, V cap K, pi_reg V, V cap C_reg) = (1, 0, 0, 1, 1): True
CHECK M.7 formal: dim(V cap K)+dim(V cap C_reg) <= dim V: True
CHECK M.7 formal: kernel/projection rank-nullity consistency: True
CHECK M.7 chi_plateau: (dim V, pi_K V, V cap K, pi_reg V, V cap C_reg) = (1, 0, 0, 1, 1): True
CHECK M.7 chi_plateau: dim(V cap K)+dim(V cap C_reg) <= dim V: True
CHECK M.7 chi_plateau: kernel/projection rank-nullity consistency: True
CHECK M.7 double_plateau: (dim V, pi_K V, V cap K, pi_reg V, V cap C_reg) = (12, 5, 3, 9, 7): True
CHECK M.7 double_plateau: dim(V cap K)+dim(V cap C_reg) <= dim V: True
CHECK M.7 double_plateau: kernel/projection rank-nullity consistency: True
CHECK M.7 formal: 1 exact rational V basis vectors archived: True
CHECK M.7 formal: every archived V vector re-verified: zero adjoint defect and Euler(1/2<h,E_C h>) = E_C (independent jet_euler route): True
CHECK M.7 formal: four objects recomputed from global rational vectors: True
CHECK M.7 chi_plateau: 1 exact rational V basis vectors archived: True
CHECK M.7 chi_plateau: every archived V vector re-verified: zero adjoint defect and Euler(1/2<h,E_C h>) = E_C (independent jet_euler route): True
CHECK M.7 chi_plateau: four objects recomputed from global rational vectors: True
CHECK M.7 double_plateau: 12 exact rational V basis vectors archived: True
CHECK M.7 double_plateau: every archived V vector re-verified: zero adjoint defect and Euler(1/2<h,E_C h>) = E_C (independent jet_euler route): True
CHECK M.7 double_plateau: four objects recomputed from global rational vectors: True
CHECK archived equivariant vector #1 is supported on N=1: True
CHECK M.7 independent test-function check: corrected N=1 contact map is equivariant under K01,K02,K03,J12 on d_mu delta: True
CHECK M.7 independent test-function check: the old (no-eta) identification FAILS boost equivariance: True
CHECK M.7 linear divergence identity: d_mu E^{mu nu}[R2] == 0 (double plateau): True
CHECK M.7 diagnosis: old eta^{mu mu}-weighted formula nonzero for R2: True
CHECK M.7 linear divergence identity: d_mu E^{mu nu}[RC2] == 0 (double plateau): True
CHECK M.7 diagnosis: old eta^{mu mu}-weighted formula nonzero for RC2: True
CHECK M.7 linear divergence identity: d_mu E^{mu nu}[R2-(188813/88456)RC2] == 0 (double plateau): True
CHECK M.7 diagnosis: old eta^{mu mu}-weighted formula nonzero for R2-(188813/88456)RC2: True
CHECK M.7: certified C_reg preimage collapses to a divergence-free tensor: True
CHECK M.7 summary: divergence identity recorded as holding for R2/RC2/combination in all readings: True
CHECK M.7: matched coefficient 188813/88456 unchanged: True
CHECK covariance_matrix: 14 total; rank/codim and coker columns distinguished: True
CHECK M.7 legacy: old 'dp_all=5' reproduced only with the truncating reduction: True
CHECK M.8/M.9 completion ledger stub present: True
CHECK [historical record] completion ledger recommends waiting on M.8: True
CHECK [historical record] M.7 status downgraded to 'Constructed after audit' (exactly verified items only): True
CHECK [historical record] M.6: C_anom is the N=0 component (archived scope): True
CHECK [current status] conventions.json: broad M-C Not achieved (Unresolved); M-C_data Satisfied; M-C_curv Proved obstructed: True
CHECK [historical record] conventions.json keeps the audit-2 M-C status string (superseded): True
CHECK [current status] manuscript carries the (retained) M.7 audit log remark: True
CHECK [current status] manuscript no longer claims a nonzero curvature-square divergence: True
# M.7 covariance regression checks OK (corrected identities)
# === work packages A/B/C ===
CHECK wpA archive file present: kernel_identification.tex: True
CHECK wpA archive file present: kernel_identification_checks.py: True
CHECK wpA archive file present: kernel_identification_out.txt: True
CHECK wpA archive file present: kernel_conventions.json: True
CHECK wpB archive file present: scaling_prescription_identity.tex: True
CHECK wpB archive file present: higher_grade_scaling.py: True
CHECK wpB archive file present: higher_grade_scaling_out.txt: True
CHECK wpB archive file present: _wpB_state.json: True
CHECK wpB archive file present: _wpB_sectionD.json: True
CHECK wpB archive file present: scale_derivative_coefficients.json: True
CHECK wpB archive file present: dilation_defect_coefficients.json: True
CHECK wpB archive file present: scale_defect_difference.json: True
CHECK wpB archive file present: full_scale_collapse_certificates.json: True
CHECK wpB archive file present: matched_scale_ledger.txt: True
CHECK wpC archive file present: time_only_section.tex: True
CHECK wpC archive file present: time_only_helmholtz.py: True
CHECK wpC archive file present: time_only_helmholtz_out.txt: True
CHECK wpC archive file present: time_only_summary.json: True
CHECK wpC archive file present: time_only_collapse_certificates.json: True
CHECK wpC archive file present: time_only_sa_certificates.json: True
CHECK wpC archive file present: time_only_matching_certificates.json: True
CHECK wpC archive file present: time_only_specialization_tests.txt: True
CHECK wpA script reruns with exit 0: True
CHECK wpA rerun output is byte-identical to the archived kernel_identification_out.txt: True
CHECK wpA output ends with ALL CHECKS PASSED: True
CHECK wpA: every CHECK line is True: True
CHECK wpA ledger: Certified, c=1/(4 pi^2), side Q-i0, across nonzero null cone: True
CHECK [current status] conventions.json gate_Sbv_status: Certified with c=1/(4 pi^2), side Q-i0 (agrees with wpA ledger): True
CHECK [current status] conventions.json normalization table: S row certified with c and boundary side: True
CHECK r = -2 i pi^2 c^2 = -i/(8 pi^2) with certified c: True
CHECK scale unit i pi^2 c^2 = i/(16 pi^2) with certified c: True
CHECK damped W_0 closed form 1/(4 pi^2 (a^2+r^2)) (independent sympy integral): True
CHECK a^2 + r^2 = -((z0 - i delta)^2 - r^2): W_0 = -(1/4pi^2)(Q - i0 sgn z0)^{-1}: True
CHECK sgn cancellation for sgn z0 = +1: delta(Q) coefficient = c i pi (side Q-i0): True
CHECK sgn cancellation for sgn z0 = -1: delta(Q) coefficient = c i pi (side Q-i0): True
CHECK wpB output ends with ALL CHECKS PASSED (followed by 'exit 0'): True
CHECK wpB: every CHECK line is True: True
CHECK wpB indep. rho=0 (symbolic): M M^{-1} = 1: True
CHECK wpB indep. rho=0 (symbolic): mu d(M^{-1})/dmu = D M^{-1} - M^{-1} D: True
CHECK wpB indep. rho=1 (symbolic): M M^{-1} = 1: True
CHECK wpB indep. rho=1 (symbolic): mu d(M^{-1})/dmu = D M^{-1} - M^{-1} D: True
CHECK wpB indep. rho=2 (symbolic): M M^{-1} = 1: True
CHECK wpB indep. rho=2 (symbolic): mu d(M^{-1})/dmu = D M^{-1} - M^{-1} D: True
CHECK wpB indep. rho=3 (random rational): M M^{-1} = 1: True
CHECK wpB indep. rho=3 (random rational): mu d(M^{-1})/dmu = D M^{-1} - M^{-1} D: True
CHECK wpB double plateau: rank(R2,RC2)=2 -> rank(R2,RC2,scale)=3 (scale tensor not in span(E_R2,E_Ric2)): True
CHECK conventions.json records the double-plateau scale 5-vector: True
CHECK wpB: double-plateau top scale tensor is self-adjoint/variational (certificate): True
CHECK wpB: top scale tensor not variational in formal / chi-plateau / prep-plateau readings: True
CHECK wpB: C_top has 289 entries, all in contact degree D=4: True
CHECK wpB: C_top N=0 coordinate = -2 (units i pi^2 c^2), i.e. r: True
CHECK wpB: residue map in audited equivariant coordinates as stated in the manuscript: True
CHECK wpC output ends with ALL CHECKS PASSED: True
CHECK wpC: every CHECK line is True: True
CHECK wpC: ker Coll_t = 0 for all D: True
CHECK wpC: ker H_t by D = (26,0,0,0,0): True
CHECK wpC: regression-preserving variational image = 26 ultralocals: True
CHECK wpC: chi==1 time-only ker H_t by D = (26,0,156,0,390), total 572: True
CHECK wpC: route 1 = route 2 (both readings): True
CHECK wpC time-only: only P1-(89/140)P2 is an image among P1,P2 (C_coh and C_reg): True
CHECK wpC time-only: weighted EH and R2/RC2 are not images (double-plateau matching does not survive time-only): True
CHECK wpC time-only chi==1: only P1-(89/140)P2 is an image among P1,P2 (C_coh and C_reg): True
CHECK wpC time-only chi==1: weighted EH and R2/RC2 are not images (double-plateau matching does not survive time-only): True
CHECK wpC: chi==1 time-only D=4 kernel (dim 390) supported on N=4 only (used in cor:wpB-timeonly): True
CHECK wpC independent: Wronskian W(rho,...,rho^(4))(1/2) = 88060251340800.0 (archived 8.80602513408e13); nonzero: True
CHECK XVIM_FULL_WP: wpB rerun reproduces higher_grade_scaling_out.txt modulo timing: True
CHECK XVIM_FULL_WP: wpC rerun reproduces time_only_helmholtz_out.txt modulo timing: True
# === audit 3 ===
CHECK [current status] conventions.json: separating functional basis order and ell=(1,1,-1,1,1): True
CHECK audit3: ell(v_scale) = -71/120-27/80+11/10+19/30-13/15 = -1/16 != 0 (archived conventions vector): True
CHECK audit3: ell(v_R2) = 0 and ell(v_Ric2) = 0: True
CHECK audit3: ell annihilates archived conventions vector R2: True
CHECK audit3: ell annihilates archived conventions vector RC2: True
CHECK audit3 sympy: pure gauge h_00=f(t): EL_f(L_1..L_5) = 2*(1,1,-1,1,1) f'''' (so L_j == (1,1,-1,1,1) f f'''' mod total derivatives): True
CHECK audit3 sympy: R_lin^2 and Ric_lin^2 vanish on pure gauge h_00=f(t): True
CHECK audit3 sympy: scale functional on pure gauge: EL_f = -(1/8) f'''' = 2*(-1/16) f'''' (units i/(16 pi^2)): True
CHECK audit3: basis check: (0,1,0,-2,1) in the L_1..L_5 basis gives the Gate V Euler tensor of R_lin^2: True
CHECK audit3: recollapsed C_top (double plateau, 490 entries) equals EXACTLY the Euler tensor of the displayed I_top_scale (units i pi^2 c^2): True
CHECK audit3: f'''' coefficient of E^00[f e_00] of the recollapsed C_top = -1/8 = 2*ell(v_scale): True
CHECK audit3: lem:bv(c) Beta-integral convolution (1/2pi) hat u_m * hat u_n = hat u_{m+n}, 1<=m,n<=3: True
CHECK [current status] audit3: wpA damping parameter renamed epsilon -> tau (script and fragment): True
CHECK audit3: wpA output contains the variable-regulator check (lem:wpA-varreg): True
CHECK [current status] manuscript defines \label{thm:wpA-Sbv}: True
CHECK [current status] manuscript defines \label{lem:wpA-varreg}: True
CHECK [current status] manuscript defines \label{thm:wpB-B0}: True
CHECK [current status] manuscript defines \label{thm:wpB-dp}: True
CHECK [current status] manuscript defines \label{prop:wpB-sep}: True
CHECK [current status] manuscript defines \label{rem:wpB-topfull}: True
CHECK [current status] manuscript defines \label{cor:wpB-MC}: True
CHECK [current status] manuscript defines \label{def:MCcurv}: True
CHECK [current status] manuscript defines \label{cor:wpB-timeonly}: True
CHECK [current status] manuscript defines \label{thm:to-timeonly}: True
CHECK [current status] manuscript defines \label{prop:to-actual}: True
CHECK [current status] manuscript defines \label{lem:to-indep}: True
CHECK [current status] manuscript defines \label{rem:to-joint}: True
CHECK [current status] manuscript defines \label{rem:revlog}: True
CHECK [current status] manuscript defines \label{rem:revlog3}: True
CHECK [current status] manuscript defines \label{rem:revlog-final}: True
CHECK [current status] manuscript carries 'Revision log (audit 3)' and the retained cycle log: True
CHECK [current status] manuscript states the top-versus-full scale-tensor sentence (abstract, status box, wpB box, theorem): True
CHECK [current status] manuscript no longer calls the full double-plateau scale tensor variational/Lorentz covariant: True
CHECK [current status] manuscript: M-C_curv Proved obstructed; broad M-C not achieved; graph-specific wording: True
CHECK [current status] manuscript no longer marks broad M-C Proved obstructed: True
CHECK [current status] manuscript: time-only switch-jet reduction removed; joint-family condition kept: True
CHECK [current status] manuscript: corrected S scope sentence: True
CHECK [current status] manuscript: weighted curvature-square divergence wording corrected: True
CHECK [current status] manuscript: bv(c) shortcut removed; damping parameter is tau: True
CHECK [current status] run listing: annotations are # comments: True
CHECK [current status] manuscript no longer states that S-bv is conditional: True
CHECK [current status] manuscript keeps 'Euclidean motivation is not the certification': True
CHECK [current status] manuscript states the double-plateau matching does not survive the time-only reading: True
# work package A/B/C regression checks OK
# === final verification ===
CHECK [current status] provenance pin file exists: XVI-J_local.tex -> /workspace/XVI-M-artifacts/provenance/XVI-J_local.tex: True
CHECK [current status] provenance SHA-256 match for XVI-J_local.tex: True
CHECK [current status] provenance pin file exists: XVI-J_github_fetch.tex -> /workspace/XVI-M-artifacts/provenance/XVI-J_github_fetch.tex: True
CHECK [current status] provenance SHA-256 match for XVI-J_github_fetch.tex: True
CHECK [current status] provenance pin file exists: XVI-J_corpus_eqW_DF_rational_propq.txt -> /workspace/XVI-M-artifacts/provenance/XVI-J_corpus_eqW_DF_rational_propq.txt: True
CHECK [current status] provenance SHA-256 match for XVI-J_corpus_eqW_DF_rational_propq.txt: True
CHECK [current status] provenance pin file exists: XI-E_corpus_eq_h.txt -> /workspace/XVI-M-artifacts/provenance/XI-E_corpus_eq_h.txt: True
CHECK [current status] provenance SHA-256 match for XI-E_corpus_eq_h.txt: True
CHECK [current status] provenance pin file exists: XII-C_corpus_weighted_curvature.txt -> /workspace/XVI-M-artifacts/provenance/XII-C_corpus_weighted_curvature.txt: True
CHECK [current status] provenance SHA-256 match for XII-C_corpus_weighted_curvature.txt: True
CHECK [current status] provenance pin file exists: XVI-L_corpus_thm_L6.txt -> /workspace/XVI-M-artifacts/provenance/XVI-L_corpus_thm_L6.txt: True
CHECK [current status] provenance SHA-256 match for XVI-L_corpus_thm_L6.txt: True
CHECK [current status] provenance pin file exists: XVI-K_corpus_prop_hess.txt -> /workspace/XVI-M-artifacts/provenance/XVI-K_corpus_prop_hess.txt: True
CHECK [current status] provenance SHA-256 match for XVI-K_corpus_prop_hess.txt: True
CHECK [current status] provenance pin file exists: ../XVI-L.tex -> /workspace/XVI-L.tex: True
CHECK [current status] provenance SHA-256 match for ../XVI-L.tex: True
CHECK [current status] provenance pin file exists: ../XVI-K.tex -> /workspace/XVI-K.tex: True
CHECK [current status] provenance SHA-256 match for ../XVI-K.tex: True
CHECK [current status] predecessor_pins note: XVI-J eq:rational is corroboration, not frozen XVI-L: True
CHECK [current status] XVI-J eq:rational marked corroboration in pins: True
CHECK [current status] archived corpus excerpt present: XVI-J_corpus_eqW_DF_rational_propq.txt: True
CHECK [current status] archived corpus excerpt present: XI-E_corpus_eq_h.txt: True
CHECK [current status] archived corpus excerpt present: XII-C_corpus_weighted_curvature.txt: True
CHECK [current status] archived corpus excerpt present: XVI-L_corpus_thm_L6.txt: True
CHECK [current status] archived corpus excerpt present: XVI-K_corpus_prop_hess.txt: True
CHECK [current status] manuscript carries Revision log (final verification): True
CHECK [current status] manuscript: explicitly checked kernel identification: True
CHECK [current status] manuscript body (before final revlog) has no 'resting on a new input': True
CHECK [current status] manuscript: M.2 spot-check wording: True
CHECK [current status] manuscript: prop:M7-id (EH) locally constant: True
CHECK [current status] manuscript body (before final revlog) has no 'uncontrolled': True
CHECK [current status] manuscript: broad M-C matched-scale criterion and M-C_data: True
CHECK [current status] manuscript cites XVI-L thm:L6 and the undetermined-cutoff sentence: True
CHECK [current status] manuscript: prop:to-chiplateau cases display with K_SA: True
CHECK [current status] manuscript: odd-order anti-self-adjoint reason: True
CHECK [current status] wpC summary: odd-D=1 anti-self-adjoint exact, ker H_t=0: True
CHECK [current status] wpC summary: odd-D=3 anti-self-adjoint exact, ker H_t=0: True
CHECK [current status] appendix lists XVIM_FULL_WP full rerun writing regression_full_out.txt: True
CHECK [current status] manuscript has predecessor provenance table: True
# final verification checks OK
ALL CHECKS PASSED
\end{lstlisting}

\end{document}
